% !TEX root = nc-cubic-n-fold.tex
\documentclass[11pt,a4paper]{ltjsarticle}
\usepackage{amsmath,amssymb,amsthm,mathtools,bm,mathrsfs}
\usepackage{luatexja-fontspec}
\usepackage{enumitem,booktabs,array,needspace}
\usepackage{xcolor}
\usepackage[colorlinks=true,linkcolor=blue!45!black,citecolor=blue!45!black,urlcolor=blue!45!black]{hyperref}
\DeclareMathSizes{9}{9}{6}{5}
\IfFileExists{/usr/share/fonts/opentype/noto/NotoSerifCJK-Regular.ttc}{%
  \setmainjfont[Path=/usr/share/fonts/opentype/noto/]{NotoSerifCJK-Regular.ttc}%
}{\setmainjfont{HaranoAjiMincho-Regular}}
\setlength{\textwidth}{15.6cm}
\setlength{\oddsidemargin}{0.25cm}
\setlength{\evensidemargin}{0.25cm}
\setlength{\textheight}{22.6cm}
\setlength{\topmargin}{-0.2cm}
\setlength{\parindent}{1em}
\setlength{\parskip}{1pt}
\numberwithin{equation}{section}
\newtheorem{theorem}{定理}[section]
\newtheorem{maintheorem}{主定理}
\renewcommand{\themaintheorem}{\Alph{maintheorem}}
\newtheorem{lemma}[theorem]{補題}
\newtheorem{proposition}[theorem]{命題}
\newtheorem{corollary}[theorem]{系}
\newtheorem{remark}[theorem]{注意}
\newtheorem{definition}[theorem]{定義}
\newtheorem{example}[theorem]{例}
\renewcommand{\proofname}{証明}
\newcommand{\C}{\mathbb C}
\newcommand{\CC}{\mathbb C}
\newcommand{\PP}{\mathbb P}
\newcommand{\OO}{\mathcal O}
\newcommand{\Sym}{\operatorname{Sym}}
\newcommand{\Spec}{\operatorname{Spec}}
\newcommand{\Proj}{\operatorname{Proj}}
\newcommand{\Hom}{\operatorname{Hom}}
\newcommand{\End}{\operatorname{End}}
\newcommand{\Ext}{\operatorname{Ext}}
\newcommand{\Aut}{\operatorname{Aut}}
\newcommand{\GL}{\operatorname{GL}}
\newcommand{\PGL}{\operatorname{PGL}}
\newcommand{\Gr}{\operatorname{Gr}}
\newcommand{\rank}{\operatorname{rank}}
\newcommand{\qgr}{\operatorname{qgr}}
\newcommand{\QGr}{\operatorname{QGr}}
\title{非可換三次 $n$-fold のモジュライ\\
\large 構造対象付きモジュライ・点対応・Poisson 幾何}
\author{}
\date{2026年11月20日・研究稿 v9}
\hypersetup{pdftitle={非可換三次 n-fold のモジュライ},
pdfsubject={構造対象付きモジュライ・点対応・Poisson幾何}}
\begin{document}
\maketitle
\vspace{-0.6\baselineskip}
\begin{abstract}
三つの非可換生成元と $n-1$ 個の中心一次生成元を持つ三次超曲面の量子化を扱う。
中心平坦性と普遍低次数中心条件の下で任意の Veronese から元の代数を復元し、
良い一ブロック表示を局所的に持つ $\mathcal O$-marked 圏族の滑らかな DM stack を構成する。
相対双対と全同値の降下により、次数一 shift を標識から外す。
大域表示の Brauer 障害は $n\not\equiv1\pmod3$ では消え、
$n\equiv1\pmod3$ ではどの非空開上でも位数三である。
点スキームは相対的に楕円成分と型 $(2,2,2,3)$ の定常成分へ分かれ、
非空開上で埋込み点対応から中心三次商を復元できる。
八重被覆の正規化 trace と Hessian、および任意基底上の強い affine 持上げの一意性により、
良い一ブロック型の圏 stack を分解可能な Poisson 対の非空開と同定する。
既存の被約分類の下で、この像は三次のモジュライを支配する唯一の Poisson 成分の
一般部分の開を与える。
他方、固定入力の残余復元定理の下で、trace--Hessian 三次と残余復元三次は
任意の次元の明示 Hesse 族で幾何学的一般点において異なる。
同じ周囲 Poisson tensor の Casimir を選び直しても、この相違は解消できない。
混合係数の一般点での二次障害を打ち消す代数的 Hessian--trace 補正を構成し、
その全同値への降下と、適切な稠密開上での次数九の有限 \'etale 性を示す。
補正後の active 平面切断は量子側の楕円点曲線に一致する。
固定 formality の全次数では、元の support が保存され、
残余三次が tensor を形式スカラー倍まで決定する。
さらに既知の harmonic graph weights と共変式の還元により、
全混合係数の五次 tensor、残余三次の四次項、およびスカラー補正を決定する。
有限 Hessian 補正との四次の Jacobi 類は全次元の非空開で非零であり、
特殊ファイバーの同定を保つ形式的な残余比較を妨げる。
この差を打ち消す追加の代数的 Poisson 射を構成し、全 O-同値への降下と、
三次式の六次未満・tensor の七次未満での一致を証明する。
分離型 Hesse 族では、固定入力の下で補正後三次への残余同値を有限拡大上で得る。

別に、三次超曲面の Jacobi スキームを Koszul flattening と Pl\"ucker 関係によって調べる。
$3\le n\le7$ の一般三次では Grassmann 錐に一致する。
$n=8$ の一般三次では、その頂点に長さ330、次数二に集中する平方零の冪零根基を持つことを、
次数三の厳密計算証明書で示す。この拡大は完備局所環としても分裂しない。
既存の被約分類の下で一般の全次元でも非零部分は Grassmann 錐に一致するが、
特殊な滑らかな三次には非零点で非被約となる例を構成する。
高次元の冪零部分の最高次数と極大イデアルによる消滅指数は一様に有界ではない。
最後に既知の Hochschild 計算と圏変形理論から、$n\ne4$ の通常 Artin 基底上で
残余圏の一次切断変形を幾何的変形として認識する。
全ての内在的成分の被覆、一般混合係数での全次数比較と残余同値の降下は未解決である。
\end{abstract}
\clearpage
\setcounter{tocdepth}{1}
\vspace{-0.5\baselineskip}
\begingroup\small
\tableofcontents
\endgroup

\section{序論と主結果}
\label{sec:introduction}

\subsection{目的と既存研究からの出発点}
本研究の基本目標は、一般次元の非可換三次超曲面について、構造対象を保つ
$\mathcal O$-marked moduli を構成し、point scheme を調べ、
Poisson 多様体のモジュライへの射と Kuznetsov 成分の比較を結び付けることである。
対象は Abel 圏と構造対象の対 $(\mathcal C,\mathcal O)$ とし、
各既約成分の一般点で成立する、次元に一様な構成を優先する。
次数付き表示や次数一 shift は、その内在的な回復を示すまでは補助データとして扱う。
一ブロックの表示空間の一般点での結果を、全モジュライの各成分へ自動的に広げない。
また、稠密 Zariski 開上の general な主張と、可算個の真の閉集合を除く
very general な主張を区別する。後者の点集合だけから代数的な族の射を結論しない。
三様体では、中心一次空間、二次包絡、中心三次の直線、
および Kuznetsov 残余圏を結び付ける構成がある \cite{Threefold}。
この構成を一般次元へ移すとき、次元だけを置き換えることはできない。
Serre 関手が回復する twist の次数が変わり、
残余圏の Hochschild 群には四次元と七次元で追加成分が現れ、
Poisson 方程式のスキーム構造にも八次元から新しい現象が生じる。

本稿の直接の出発点は、リポジトリ \texttt{nc-cubic} の
論文A \cite{NCCubicA} と論文B \cite{NFNCCubicB} である。
参照するソースは次の固定 commit に統一する：
\begin{center}
\texttt{eabe601c1341777742964584dde5314c0bcf64a9}.
\end{center}
論文Aは既に任意次元の Poisson 正規形と代数的中心量子化を扱い、
論文Bは既に $n\ge3$ の一般点における rotation と指定三次周期を保つ
可換残余模型を扱う。
これらを今回新しく得た一般次元化として数えない。
本稿は、これらの研究稿を明示した入力として、次数付きモジュライ、
任意次数の中心復元、点関手、Jacobi スキーム、および残余圏の形式変形を整理する。
引用した未公表稿の全証明を本稿内で再証明したという意味ではない。

関連する公刊研究との関係は次の通りである。
Etingof--Ginzburg \cite{NFMEG} は三変数 potential の PBW 変形と中心元を構成し、
Mori--Smith \cite{MoriSmith} は三生成元三二次関係の Calabi--Yau 代数を分類する。
Behrend--Noohi \cite[Sections~2.3, 3.1.2]{NFMBehrendNoohi} は
次数スカラーによる rigidification、重み付き代数の gerbe、
determinant character による中立化、および偏極付き Grothendieck 圏を扱う。
これらの一般的方法を新規とは数えず、本稿では全 $\mathcal O$-同値の復元を用いて
shift の標識を外した局所一ブロック型の圏 stack と同定し、
階数 $(n-1,3)$ に依存する、全ての非空開上の Brauer 障害を決定する。
Berger--Taillefer \cite{NFMBeTa} の potential 解消と PBW--CY 定理も入力とし、
本稿では任意個の中心差の Koszul 解消と有限射影模型によって、
非平坦・非被約な底変更にも両立する相対双対と trace を構成する。
体上の CY 性や PBW 定理自体を新規成果としては数えず、
$n$ に依存する全同型問題や torsor の降下と区別する。
Le Bruyn--Smith--Van den Bergh \cite{LBSV} の一点中心拡大における
楕円曲線と八点、および二次関係の復元に対して、本稿は任意個数の中心変数を扱い、
有限平坦な八重被覆と高次元の二十四次完全交叉を得る。
さらに Rains の中心公式 \cite[\S9, p.~1414]{NFMEG} の境界評価から
中心三次の復元を強め、相対 trace と尺度付き Hessian による
Poisson 対への射を構成する。この射と既存量子化の単純な逆写像、
および残余圏からの可換復元とは区別する。
形式量子化の中心と Poisson 中心の同型は
Etingof--Ginzburg\cite[Proposition~8.1.1, equation~(8.1.3)]{NFMEG}
にも既に現れ、その証明では unimodularity から相対形式 Poisson tensor を
Jacobian potential で表している（Section~8.1, pp.~1407--1408）。
本稿で扱うのは、有限の非零パラメータでの三次部分に対する
trace と Hessian の明示的処方、その任意底変更、全同値に関する降下と逆構成である。
Dito\cite[Section~2, equation~(3); Section~3, Theorem~2]{NFMDito} の一般 Poisson
星積の二次項と Weyl gauge、および Shoikhet\cite[Sections~1.3, 2.3.1]{NFMShoikhet}
の内向き wheel の消滅を用いて、固定入力の残余 potential の二次計算を
全混合係数へ広げる。星積や wheel 消滅自体は既知結果であり、
本稿では trace--Hessian 処方との差が全次元で非零となる開条件を調べる。
二次代数と二次 Poisson 構造の形式対応は既に
Kontsevich\cite[Section~2.3, Proposition~6]{NFMKontsevichAlg}にある。
また点曲線の Hessian 公式は Mori--Smith\cite[Theorems~1.2, 7.2]{MoriSmith}の結果である。
本稿はこれらを既知の入力として、混合係数の trace を加えた代数的補正を
O-marked moduli 上へ降下し、全混合係数での一般次数九と残余比較への適合を調べる。
さらに Lindberg--Pym\cite[Section~1.6; Theorem~6.9 and Remark~6.10]{NFMLindbergPym}
の量子トーラスの交換比 $q=e^\hbar$ を既知の入力として、
生成元固定の逆 tensor の五次項を三変数三次の共変式から決定する。
全混合項の定数評価には
Banks--Panzer--Pym\cite[Corollary~1.3; Section~6]{NFMBanksPanzerPym}の
harmonic graph weights と、Buring\cite[Sections~3.4--3.6]{NFMBuringGraphs}の
公開データを用いる。線型入力と二次 tensor に有効な73項を抽出し、
元の weight と多重度、生成元固定再帰を有理数として検証する。
既知の graph 積分値自体を新規成果とは数えず、
本稿の貢献は混合項の一様な公式、残余三次の四次比較類とその非消滅、
および全同値へ降下する追加補正への適用である。
形式対応の存在、点曲線公式、量子トーラスの交換比そのものを新規成果とは数えない。
次数計算に用いる Hesse 還元は、純 Hessian 写像に対する
Ciliberto--Ottaviani\cite[Theorem~5.7]{NFMCilibertoOttaviani}の議論にもあり、
四次の $S$ は古典的な Aronhold 不変式\cite{NFMOttavianiAronhold}である。
$I,J$ も Hessian 写像と Aronhold 四次式の方向微分で表示できる
（補題\ref{lem:nfm-classical-mixed-covariants}）。
本稿ではこれらを一般混合係数と商 stack の次数、および固定逆 tensor の係数へ適用する。
Mayanskiy \cite{Mayanskiy} の三様体における Schouten 方程式の計算に対して、
本稿は高次元での Jacobi イデアル自体を比較する。
affine 超曲面の Poisson 延長を gradient の Koszul 複体で記述する手法には
Fraenkel \cite[Proposition~4.1, Theorem~4.3]{Fraenkel} の先行研究がある。
また、一般の Lie 代数 law の Jacobi スキームに冪零元が現れることは
Carles \cite[Section~8, Lemma~8.35]{Carles} および
Carles--M\'arquez \cite[Section~3]{CarlesMarquez} に既知である。
本稿の Jacobi 係数写像は Landsberg--Ottaviani の Young flattening
\cite{LandsbergOttavianiYoung}、Gesmundo--Landsberg の Koszul flattening
\cite[Proposition~6.1]{GesmundoLandsbergPartials}に同定される。
従って普遍階数上界自体は既知であり、本稿は座標による再証明を与える。
apolar Gorenstein 代数の Betti 数を計算した Boij \cite[Remark~6]{NFMBoijBetti}
との関係も明示し、$N=9,10$ の二次階数の達成を新規とは数えない。
Jacobi 節で扱う中心的対象は、三次超曲面に付随する Hessian--Jacobi イデアルと
Grassmann 錐のスキーム論的な差、次数三の乗法が決める長さ330の平方零構造、
および特殊三次の非零点の局所構造である。
Hessian と apolar 積、Thom--Sebastiani 分解の関係については
Bricalli--Favale--Pirola \cite[Proposition~2.1, Theorem~3.4]{NFMBFFHessian}を参照し、
本稿の横断切片と Jacobi スキームの retract の主張と区別する。
Lin--Zhang \cite{NFLZSerre}、Pirozhkov \cite{NFPirozhkov}、
Huybrechts--Rennemo \cite{NFHuybrechtsRennemo} の既知の圏論的不変量と Torelli 理論は、
本稿の残余変形の入力である。
既知の公式からの帰結と、本稿に証明を与える拡張は以下で区別する。
これらの拡張について、独立な新規性の網羅的確認はなお完了していない。

\subsection{中心復元と構造対象付きモジュライ}
以下 $n\ge3$、$N=n+2$、$r=n-1$ とする。
生成元空間を $V$、中心一次空間を $U\subset V$、$\dim U=r$ と書く。
点の空間は $W=V^*$ であり、$\PP(W)$ 上の三次式は
$\Sym^3W^*=\Sym^3V$ の元である。
周囲の二次代数を $B$、中心三次商を $A=B/(w)$ と書く。

\begin{maintheorem}[中心復元と局所一ブロック型モジュライ]
可換 $\C$-代数上の次数一生成代数 $A,A'$ が、それぞれ
中心多項式環 $C=\Sym U,C'=\Sym U'$ 上平坦、
$\rank U=\rank U'\ge2$ であり、
普遍的に $Z(A)_d=\Sym^dU$ と $Z(A')_d=\Sym^dU'$ を満たすとする。
このとき $d$-Veronese 間の次数付き同型は fpqc 局所で $A\simeq A'$ へ持ち上がり、
持上げ層は $\mu_d$-torsor である。

一ブロック PBW・楕円中心条件を満たす非空開上で、
次数スカラーを除いた次数付き周囲代数の同型問題は、
有限安定化子開に制限すると滑らかな既約有限型 DM stack $\mathfrak B_n$ となる。
中心三次直線の stack $\mathfrak C_n$ からの射
$\tau_n:\mathfrak C_n\to\mathfrak B_n$ は、
階数 $s_n$ のベクトル束の加法群の torsor であり、
\[
 \dim\mathfrak B_n=b_n=\frac{n^2+7n-4}{2},\qquad
 s_n=\binom{n+1}{3},\qquad
 \dim\mathfrak C_n=b_n+s_n=\binom{n+2}{3}+3n-2.
\]
普遍相対双対と全同値比較の良い開に制限すると、
$\mathfrak C_n^{\mathrm{good}}$ は fpqc 局所的にこれらの表示を持つ
$\mathcal O$-marked Grothendieck Abel 圏族の stack
$\mathfrak M_n^{\mathrm{loc}}$ に同値である。
大域の次数付き表示を妨げる gerbe の類 $\beta$ は
$\gcd(n-1,3)\beta=0$ を満たす。
$n\equiv1\pmod3$ では、その制限は任意の非空開上で位数がちょうど三である。
\end{maintheorem}

正確な普遍中心条件と証明は定理\ref{thm:nfm-central-veronese}に、
表示空間と全次数付き同型の比較は命題\ref{prop:nfm-graded-presentation}に、
有限安定化子と torsor の証明は定理\ref{thm:nfm-ambient-dimension}、
\ref{thm:nfm-central-torsor}に置く。
相対 Serre 関手が与える $d=r$ にこの復元を適用するため、
楕円自己同型 $\sigma$ には
$\sigma^j\ne1$ $(1\le j\le\max(3,r))$ を課す。
中心スカラーの重みを考慮した束
$\mathcal Hom(\mathcal L,\Sym^3\mathcal U)$ が商 stack 上に降下する。
$\mathcal U$ や中心三次束そのものの降下を主張するものではない。

\begin{center}
\begin{tabular}{r|r|r|r|r}
$n$&可換三次の次元 $M_n$&周囲 $b_n$&中心三次 $s_n$&対 $c_n$\\\hline
3&10&13&4&17\\
4&20&20&10&30\\
5&35&28&20&48\\
6&56&37&35&72\\
7&84&47&56&103\\
8&120&58&84&142
\end{tabular}
\end{center}
ここで $M_n=\binom{n+4}{3}-(n+2)^2=\binom{n+2}{3}$ である。
$c_n-M_n=3n-2$ という数値一致から、可換モジュライへの射の存在や
そのファイバーの次元を結論してはいない。
良い相対 PBW 表示チャートでは、定理\ref{thm:nfm-relative-ambient-cy}により
真正な双加群双対 $\Theta_{B/R}\simeq B(N)[-N]$ を構成する。
定理\ref{thm:nfm-relative-line-cohomology}と\ref{thm:nfm-relative-serre}は、
任意底変更後の compact generation、相対 properness、相対 Serre 関手
$S=(1-n)[n]$ とその trace を与える。
固定表示模型間の係数作用と全 Hom の fpqc 降下を構成することにより、
定理\ref{thm:nfm-fixed-model-comparison}は
$\underline{\operatorname{Isom}}_{\mathrm{gr}}/\mathbf G_m
\simeq\operatorname{Eq}_{\mathcal O}$ を与える。
大域的な次数付き代表の障害は $\operatorname{Pic}(R)[n-1]$ の線束であり、
次数一加群がともに自由なら $\operatorname{Pic}(R)[\gcd(n-1,3)]$ に属する。
$n\equiv1\pmod3$ では、滑らかな affine 楕円曲線上で非零障害の例も構成する
（命題\ref{prop:nfm-nontrivial-three-torsion}）。
切断飽和性自体は中心正則対によって全整数次数で証明する
（補題\ref{lem:nfm-central-saturation}）。
定理\ref{thm:nfm-local-moduli}では、係数変更を圏内の係数作用として定義し、
表示 gerbe 上の重み付き次数加群の tails により対象の効果化を構成する。
表示 gerbe の $H^2$ の類と、指定同値の上の $H^1$ の線束障害は異なる。
定理\ref{thm:nfm-brauer-dense-open}は、$n\equiv1\pmod3$ では
どの非空開にも $\C(s,t)$ 上で非零 Brauer 障害を持つ対象が存在することを示す。
全ての内在的成分の一般点に局所一ブロック表示がある、という分類は主張しない。

\subsection{相対点対応と Poisson 対への射}
\begin{maintheorem}[相対点スキーム、中心復元、Poisson 対]
一ブロック模型の楕円境界、固定点なしの shift、および中心三次の正則切断条件の下で、
\[
 \operatorname{Pt}(A)=E\amalg D,\qquad
 D\subset\PP^{n+1}\text{ は型 }(2,2,2,3)\text{ の完全交叉},
\]
\[
 \dim D=n-3,\qquad\deg D=24,\qquad\omega_D=\OO_D(7-n).
\]
三つの偏微分で定まる $C$ は $\PP^{n-2}$ への次数八の有限平坦被覆である。
$n=4$ で $D$ が滑らかなら、$D$ は種数 $37$ の連結曲線である。
これらは任意底変更と両立する相対完全交叉と八重被覆を与え、
点商の関手、shift のグラフ、中心射影は
$\mathfrak M_n^{\mathrm{loc}}$ の良い開上へ降下する。
ambient は一般には Brauer--Severi 束であり、上の双対化層の式は
各ファイバーでの式である。

非空開上では、埋込み点対応から $B$ と中心直線 $[w]$、従って $A$ を復元できる。
さらに尺度込みの周囲 Poisson tensor $\Pi$ に対し、
中心三次束と古典三次 Casimir 束の間に、純中心三次を固定する同変同型
\[
 \mathfrak d_\Pi:Z(B)_3\xrightarrow{\sim}
 \mathcal R_\Pi=\C F\oplus\Sym^3U
\]
を構成できる。この同型は任意の基底変更と両立する。
従って滑らかな三次を与える非空開
$\mathfrak M_n^{\mathrm{Pois}}\subset\mathfrak M_n^{\mathrm{loc}}$ 上で
\[
 (\mathcal C,\mathcal O)\longmapsto
 \bigl(V(\mathfrak d_\Pi(w)),\Pi|_{V(\mathfrak d_\Pi(w))}\bigr)
\]
が、偏極類を持つ三次 Poisson 対への射を定める。
\end{maintheorem}

相対点関手と束の正確な降下は定理\ref{thm:relative-point-nfold}、
\ref{thm:point-omarked-descent}に、中心直線の復元は
定理\ref{thm:point-central-recovery}に置く。
Poisson 射の標的は、斉次 quadratic Poisson tensor と滑らかな三次 Casimir の
局所閉スキームの $\PGL(V)$ 商として具体的に構成する
（定理\ref{thm:point-trace-hessian}）。
定理\ref{thm:poisson-strong-lift}により affine lift は実は Poisson 対から一意に回復する。
従って追加の lift 標識なしに、Brauer--Severi 周囲空間と偏極類を持つ対を扱える。
$n\equiv1\pmod3$ で実際の相対線束 $\OO(1)$ を標識に加えるには、
上記の Brauer 障害への対処が必要である。

この射は、最大可換代数商の八重有限平坦構造が与える
正規化 trace と、境界 Jacobian 環の次数三部分への Hessian 同型を使う。
中心元の単純な可換化は古典 Casimir の制限と一致しない
（命題\ref{prop:point-casimir-restriction-obstruction}）。
逆構成は次のように強められる。

\begin{maintheorem}[Poisson 対の開部分との同型]
任意の可換 $\C$-代数上の滑らかな三次 $X_H\subset\PP(W)$ について、
各相対双ベクトルは一意な斉次二次 affine 持上げ
$\Pi=\iota_{dH}P$ を持ち、$\iota_{dH}\Pi=0$ を多項式として満たす。
この同定と Poisson 条件の同値は任意基底変更に両立する。

$P$ が非零分解可能な Poisson 対の stack を $\mathfrak P_n^{\mathrm{dec}}$ とする。
周囲代数単独の有限安定化子を要求せず、良い PBW・中心・相対双対・全同値比較の
入力と滑らかな $H_w$ を課した圏 stack $\mathfrak M_n^{\mathrm{sm}}$ には同型
\[
 \mathfrak M_n^{\mathrm{sm}}\simeq
 \mathfrak P_n^{\mathrm{good}}\subset\mathfrak P_n^{\mathrm{dec}}
\]
がある。右辺は非空開であり、両辺は次元 $\binom{n+2}{3}+3n-2$ の
滑らかな既約有限型 DM stack である。逆写像は $\Pi$ の二次関係と
$[w]=\mathfrak d_\Pi^{-1}([H])$ による。表示 gerbe の位数の区分も同じである。

既存の一般三次の被約 Poisson 分類を入力すると、一般三次の非空開上で
これは全非零 Poisson モジュライへの開埋込みとなる。
全 Poisson モジュライのうち滑らかな三次のモジュライを支配する既約成分は一つで、
本構成はその一般部分の開を同定する。
\end{maintheorem}

証明は第\ref{sec:poisson-inverse}節に置く。
これは全 intrinsic 圏モジュライの成分分類でも、Kuznetsov 成分の同値定理でもない。
後者の比較では、むしろ次の障害がある。

\begin{maintheorem}[残余復元との比較障害]
固定入力論文Bの完備 MF 移送と一般点復元を仮定する。
全ての $n\ge3$ に対し、一ブロック中心三次族で
\[
 H^{\mathrm{TH}}=F_0+hxyz,\qquad
 H^{\mathrm{res}}=F_0+
 \frac{h+h^2/3+h^5/12}{1-h^4/4}xyz,\qquad
 F_0=\frac{x^3+y^3+z^3+\sum_{i=1}^{n-1}u_i^3}{3}
\]
となるものがある。二つの滑らかな一般ファイバーは代数閉拡大後も射影同型でない。
従って非可換残余圏と $\operatorname{Ku}(V(H^{\mathrm{TH}}))$ との間には
rotation を保つ dg Morita 同値がない。
$n\not\equiv1\pmod3$ では rotation を指定しない dg Morita 同値もない。
同じ周囲 tensor の滑らかな三次 Casimir の選択だけでは、この障害を解消できない。

一般の混合 Hesse 表示では、短い中心公式から $h^2ABC/6$ を引いて
Hessian 二次 jet を持つ中心元に規格化する。指定した生成元固定の formality 模型の下で
\[
 W_{\mathrm{res}}-H^{\mathrm{TH}}
 =h^2\left(\frac{\mathscr H_F}{24}
       +\frac{\operatorname{tr}(\mathscr H_F)}{12}\right)+O(h^3),
 \qquad \mathscr H_F=\det\operatorname{Hess}_{x,y,z}F.
\]
括弧内の全 Jacobi 環の class は非空 Zariski 開で非零である。
従って特殊ファイバーを恒等に標識した形式比較は、一般の表示係数においても
既に二次 jet で成立しない。
\end{maintheorem}

厳密 Hesse 係数は固定入力の計算であり、次元に一様な比較障害への適用、
固定 tensor の全 Casimir に対する障害、混合係数上の二次 class を本稿で証明する
（第\ref{sec:nfm-residual-th-obstruction}節）。
一般混合係数における有限 jet の障害から、標識を外した全ての一般ファイバーの
圏非同値までを結論してはいない。


\subsection{残余三次に適合する補正}

\begin{maintheorem}[代数的補正と形式的な残余比較]
分解可能な Poisson 対を $(H,P)$ と表し、$K=\operatorname{supp}P$、
$U=\operatorname{ann}K$ とする。$H$ と $H|_{\PP(K)}$ が滑らかな開で、
active 臨界代数の正規化 trace を $\tau_{H,K}$ と置く。
座標不変な三次式
\[
 \mathcal D_P(H)=\langle\det(D_K^2H),P^{\otimes2}\rangle,
 \qquad H^\sharp=H+\frac{\mathcal D_P(H)}{24}
                      +\frac{\tau_{H,K}(\mathcal D_P(H))}{12}
\]
は、非空開で Poisson 対の代数的射 $\Phi:(H,P)\mapsto(H^\sharp,P)$ を定める。
この射は任意底変更と全 O-同値に両立し、始域と終域を適切な非空開へ縮めると、
表現可能な有限 \emph{\'etale} 射で階数 $9$ となる。
補正後の active 平面切断は量子側の楕円点曲線に等しく、
指定した中心正規化では $W_h\equiv H^\sharp\pmod{h^4}$ となる。

一方、滑らかな先頭三次と滑らかな active 制限を持つ一ブロック族について、
固定 formality の生成元固定模型では、形式 Poisson tensor $\pi_h$ は
元の中心一次空間を厳密に保存し、残余 potential $W_h$ に対して
\[
 \pi_h=a_h\,\iota_{dW_h}P_0,\qquad a_h=h+O(h^2)
\]
と一意に書ける。この tensor は周囲二次代数だけから決まる。
その相対化は被約有限型の普遍 PBW 係数模型とその任意底変更で成立する。

二次関係の係数が $h$ に依存せず、Hessian 正規化中心元を用いると、
$\pi_h$ と $a_h$ は奇級数、$W_h$ は偶級数である。
$g=F|_{U=0}$、$S(xyz)=1$ と規格化した四次不変式を $S$ とすると、
全混合係数について
\[
a_h=h+\frac{S(g)}{240}h^5+O(h^7),\qquad
 \pi_h=h\Pi_F+\frac{h^3}{24}\Pi_{\mathscr H_F}
                 +h^5\Pi_{Q_5}+O(h^7)
\]
を得る。また active 制限は
\[
 [h^5]\pi_h|_{U=0}=\frac{S(g)}{80}\Pi_g.
\]
従って $S(g)\ne0$ の稠密開では、固定した $h$ と生成元に関して
代数的 tensor 補正と形式逆 tensor は五次から異なる。
五次 tensor の全混合係数は、明示的な相対三次式 $I(F),J(F)$ を用いて
\[
 \pi_5(F)=\Pi_{Q_5(F)},\qquad
 Q_5(F)=\frac{4I(F)+J(F)-2S(g)F}{240}
\]
と書ける。指定した中心元
$\operatorname{sym}(F-h^2\mathscr H_F/12+h^4\Psi_\lambda(F))$ の残余三次は
\[
 W_h=F+\frac{h^2}{24}\mathscr H_F
       +h^4\left(\Psi_\lambda(F)+\frac{I(F)}{24}-\frac{J(F)}{80}\right)+O(h^6)
\]
である。全 Jacobi 環での有限補正三次との差は
\[
 [h^{-4}(W_h-H^\sharp)]_{h=0}
 =\frac{[I(F)]}{60}+\frac{[\tau_F(I(F))]}{40}
                      +\frac{S(g)[\tau_F(F)]}{80}.
\]
この類は各 $n\ge3$ の混合係数の非空開で非零であり、
$\C[h]/(h^5)$ 上で特殊ファイバーの同定を保つ三次の同型を妨げる。
固定 MF・window 入力の下では、同じ標識を保つ残余圏の同値も妨げる。

一方、式\eqref{eq:nfm-fourth-invariant-correction}の座標不変な三次式
$\mathcal E_P(H)$ による
\[
 (H,P)\longmapsto
 \left(H^\sharp+\mathcal E_P(H),
       \left(1+\frac{S_P(H)}{240}\right)P\right)
\]
は、適切な非空開で任意底変更と全 O-同値へ降下する。
指定中心族に沿い、この出力三次は $W_h$ と $h^6$ を法として、
出力 tensor は $\pi_h$ と $h^7$ を法として一致する。


分離型 $g_\lambda+G(U)$、
$g_\lambda=(x^3+y^3+z^3)/3+\lambda xyz$、$1+\lambda^3\ne0$ では、
$G$ の偏微分の共通零点が原点だけであると仮定する。
量子化尺度 $s$ に対する補正後三次は
\[
 H^{\mathrm{corr}}=g_\lambda+G+
               \frac{s^2}{24}\det\operatorname{Hess}(g_\lambda).
\]
固定入力論文Bの完備 MF・相対 window・固定標的への有限降下と
coherent 残余圏の飽和性を仮定すると、
全 $n\ge3$ において、ある有限拡大 $K'/\C(s)$ 上で
$(\mathcal K_{B/(w)},\rho)\simeq
(\operatorname{Ku}(V(H^{\mathrm{corr}})),\rho_{\mathrm{corr}})$
が指定三次周期を保って成立する。
\end{maintheorem}

代数的射と次数は定理\ref{thm:nfm-algebraic-hessian-correction}、
\ref{thm:nfm-correction-degree-nine}、形式構成は定理\ref{thm:nfm-residual-formal-poisson}、
jet と五次の相違は系\ref{cor:nfm-mixed-even-jet}、
\ref{cor:nfm-fifth-restriction-obstruction}、混合項の還元は
命題\ref{prop:nfm-mixed-fifth-reduction}、\ref{prop:nfm-mixed-fifth-explicit}、
四次 potential と比較障害は定理\ref{thm:nfm-mixed-fourth-potential}、
系\ref{cor:nfm-corrected-fourth-obstruction}、
追加の代数的射は定理\ref{thm:nfm-fourth-algebraic-correction}、
split 比較は
定理\ref{thm:nfm-split-corrected-residual}に置く。
単元倍を許した一般混合係数の全次数比較、
および残余同値の全 O-降下は未解決である。
有限 jet の標識付き障害から、一般ファイバーで標識を外した Ku の非同値は結論しない。
従来の $\Phi$ の一般次数九と点曲線の厳密な一致を、
追加補正の射にも自動的に移すことはしない。

\subsection{Jacobi スキームの次元境界}
\begin{maintheorem}[Jacobi スキームの境界と平方零構造]
滑らかな可換三次 $X_F\subset\PP^{n+1}$ の自然な Jacobi スキームについて、
$3\le n\le7$ の一般三次では
\[
 \operatorname{Pois}(X_F)\simeq\widehat{\operatorname{Gr}}(3,n+2)
\]
である。
$n\ge8$ では任意の滑らかな $F$ についてこの等式は成立せず、
頂点完備化も抽象局所環として同型でない。
論文Aの一般三次に対する被約分類を入力すると、この場合の一般の
$\operatorname{Pois}(X_F)$ は非被約であり、
冪零根基の次数二部分の次元は少なくとも
\[
 \epsilon_N=5\binom{N+1}{6}-N\binom N4+\binom N3>0\quad(N=n+2\ge10)
\]
である。$n=8$ ではさらに
\[
 J_{10}/I_F\simeq\C(-2)^{330},\qquad
 H_{R_{10}/I_F}(t)=H_{R_{10}/J_{10}}(t)+330t^2
\]
となる。冪零根基は極大イデアルで消える平方零イデアルであり、
この拡大は完備局所環としても分裂しない。
$n=8$ のこの結論は、既存の被約分類を用いない独立な計算証明による。
さらに全ての $n\ge3$ の一般三次について
\[
 \operatorname{Pois}(X_F)\setminus\{0\}
 =\widehat{\operatorname{Gr}}(3,N)\setminus\{0\},\qquad
 \operatorname{Proj}(R_N/I_F)=\operatorname{Gr}(3,N)
\]
がスキームとして成り立つ。
$n\ge8$ では冪零根基は頂点に支持される非零の有限長加群であり、
頂点の局所環は深さ零である。
一方、全ての $n\ge11$ で、非零 Poisson 点の局所環に
330次元の平方零部分空間を持つ滑らかな特殊三次が存在し、射影化しても非被約である。
\end{maintheorem}

Jacobi スキームの主張は\S\ref{sec:nfold-jacobi-scheme}で証明する。
Jacobi イデアル $I_F$ の $N\binom N4$ 個の生成二次式の間には
$\binom N3$ 個の普遍的な独立従属関係がある。
Pl\"ucker イデアル $J_N$ の二次部分の次元 $5\binom{N+1}{6}$ との比較が上の下界を与える。
$N\le10$ の最大階数には、明示的な整数小行列が法 $32003$ で非零である証明書を付す。
$n=8$ では Pl\"ucker の線形 syzygy を330次元の商へ射影し、
$39600\times39600$ 小行列を法二で検証することにより、次数三以降の差を消す
（定理\ref{thm:nfold-eightfold-square-zero}）。
標数二の三次幾何を使うのではなく、元の整数係数行列の階数を特性零へ移す。
非零分解可能 $P$ の support を $A$ とし、
$K_{F,A}=\{w\in W:T_F(A,w,W)=0\}$ と置くと、法方向の接空間は
$\Lambda^3K_{F,A}$ となる（定理\ref{thm:nfold-normal-tangent}）。
これが一般三次での非零部分の一致と頂点支持を与える。
さらに、support を固定した横断切片そのものを $\Lambda^3K_{F,A}$ 上の
polar Jacobi スキームとして記述する（定理\ref{thm:nfold-transverse-slice}）。
Thom--Sebastiani 直和にはスキームの retract があり、
これを用いて特殊三次の非零点へ冪零構造を移す。
高次数の Hilbert 関数の比較から、冪零部分の最高次数と
頂点の極大イデアルによる消滅指数が次元によらず一様に有界ではないことも示す
（命題\ref{prop:nfold-higher-nilradical}）。
非被約性の結論に用いる被約分類の依存と、方程式数の比較だけで従う結論を分ける。
また、これらの Poisson スキームを全 Abel 圏の Artin 変形空間と同一視したわけではない。

\subsection{残余圏と現在の到達範囲}
\begin{maintheorem}[残余圏の通常の Artin 形式認識]
滑らかな三次 $X$ の dg 残余圏を $\mathcal K_X=\operatorname{Ku}(X)$ とし、
$\mathcal D_X=\operatorname{Ind}(\mathcal K_X)$ とする。
$n\ge3$、$n\ne4$ なら、特殊ファイバーを同定した通常の可換局所 Artin 基底上で
\[
 \operatorname{Def}_X\simeq
 \tau_{\le1}\operatorname{CatDef}^{c}_{\mathcal D_X}.
\]
両辺の同型類は $\binom{n+2}{3}$ 次元の滑らかな形式円板となる。
$n=4$ では圏変形の形式 hull の通常 Artin 同型類関手が
二十二次元形式円板となり、幾何的変形の円板は二十次元である。
\end{maintheorem}

証明は定理\ref{thm:nf-residual-formal}に与える。
既知の sector 公式と幾何的 Kodaira--Spencer 単射性から得る接写像同型を、
実際の幾何的変形の全次数の自然変換と合わせる。
$HH^3$ の消滅に依存しないため、$HH^3=\C^2$ となる $n=7$ も含む。
圏変形の形式 hull と実際の compactly generated 変形の比較には
Lurie \cite{NFDAGX} の tame 圏に対する結果を用いる。
未切断の高次群oidや導来 Artin 基底上での同値は主張しない。

rotation $\mathsf R$ と Serre 関手 $S$ の既知の関係は
$\mathsf R^3=[2]$、$S=\mathsf R^{-(n-1)}[n]$ である。
$n\not\equiv1\pmod3$ では rotation を Serre と shift から回復できるが、
$n\equiv1\pmod3$ ではこの方法は使えない。
これは四次元の形式認識の例外とは異なる区分である。

本稿の各定理には証明または明示的な入力を記し、
計算補助の主張には再現可能な証明書を添える。
全 $\mathcal O$-marked moduli の代数性と各成分の一般対象の被覆は未完成である。
良い局所一ブロック型では、全 $\mathcal O$-同値に関して降下する
trace--Hessian 射と逆構成により、Poisson 対の開部分を同定した。
一方、この処方が選ぶ underlying cubic と残余復元の出力は一般には同一視できず、
固定した周囲 tensor の純中心補正だけでは修復できない例がある。
この障害に対して、支配的な代数的補正 $\Phi$ と、固定 formality による
全次数の形式 Poisson 対を構成した。split 族では、補正後三次と残余復元の
全次数の比較が固定入力の下で成立する。
一般混合係数における残余三次の四次項と高次補正、形式構成の代数化、
および残余同値の全 O-同値への降下を次に調べる。
各成分の一般点と必要な係数拡大上で、有限表示と同値核の降下を確認し、
稠密開への拡張を目指す。
boundary や全 Abel 形式近傍の記述は、この一般次元の構成に必要な障害を解消する
範囲で扱い、本稿に得た個別次元の証明と計算資産はその検査例として保持する。

% !TEX root = ../cubic-moduli-descent.tex
\section{中心 Veronese 錐による次数一の復元}
\label{sec:nfm-veronese}

本節では $n\geq3$, $N=n+2$, $r=n-1$ とする。
周囲の二次代数を $B$、中心三次商を $A=B/(w)$ と書く。
次数一の生成元空間を $V$、その中心部分を $U\subset V$ とする。
一般次元では相対 Serre 関手が与えるのは $r$ 次の twist である。
これから次数一を回復するためには、$r$-Veronese の中心に含まれる
Veronese 錐を用いる。
以下の代数的復元は、Serre 双対性の構成とは独立である。

\subsection{中心 Veronese の同型の持上げ}

\begin{lemma}[Veronese 錐の同型層]
\label{lem:nfm-veronese-cone}
$R$ を可換 $\mathbf C$-代数、$U,U'$ を同じ階数 $s\geq2$ の
有限局所自由 $R$-加群、$d\geq1$ とする。
$C=\operatorname{Sym}_R U$, $C'=\operatorname{Sym}_R U'$ と置く。
次数を $d$ で割った Veronese 代数について、自然な写像
\[
 \underline{\operatorname{Isom}}_R(U,U')
 \longrightarrow
 \underline{\operatorname{Isom}}_{\mathrm{gr},R}(C^{(d)},(C')^{(d)})
\]
は fpqc 層として $\mu_d$-torsor である。
\end{lemma}

\begin{proof}
基底上で局所的に $U,U'$ を自由としてよい。
Veronese 環の Proj はそれぞれ $\mathbf P^{s-1}_R$ であり、
その新しい次数一の線束は通常の $\mathcal O(d)$ である。
次数付き同型はこの線束を保つ射影空間の同型 $f$ を誘導する。
各幾何学的ファイバーでは $f^*\mathcal O(1)=\mathcal O(1)$ なので、
$M=f^*\mathcal O(1)\otimes\mathcal O(-1)$ は各ファイバーで自明である。
射影空間の $H^0(\mathcal O)=R$, $H^i(\mathcal O)=0$ $(i>0)$ と
cohomology and base change により、$M=\pi^*L$ となる基底の可逆層 $L$ がある。
この議論は有限表示のデータを noetherian 部分環へ降ろしてもよいので、
$R$ の noetherian 性を必要としない。
基底を被覆して $L$ を自明化すると、次数一切断への作用から
$f$ は線型同型 $g:U\to U'$ に持ち上がる。
ここでは projective bundle の切断の同定を用いた
\cite[Section~27.21]{NFMStacksProjective}。

与えた Veronese 同型と $\operatorname{Sym}^{(d)}g$ の差は Proj 上で恒等である。
$\mathcal O(d)$ 上の自己同型は基底の単元 $a\in R^\times$ であり、
$C^{(d)}$ は新次数一で生成されるので、この差は新次数 $j$ に $a^j$ を掛ける。
忠実平坦な有限 \'etale 被覆
$R\to R[t]/(t^d-a)$ 上で $t$ を選び、$g$ を $tg$ に置き換えれば差を消せる。
二つの持上げの比は射影空間上で恒等なので $\lambda\operatorname{id}$ であり、
$\lambda^d=1$ が必要十分である。
従って各同型の持上げ層は $\mu_d$ により単純推移的に作用され、局所的に非空である。
\end{proof}

\begin{theorem}[中心を用いた任意次数の Veronese 復元]
\label{thm:nfm-central-veronese}
$R$ を可換 $\mathbf C$-代数、$d\geq1$ とする。
$A,A'$ は非負次数付きで次数一生成の $R$-代数、$A_0=A'_0=R$ とする。
階数 $s\geq2$ の有限局所自由加群 $U,U'$ に対し、中心的な単射
\[
 C=\operatorname{Sym}_R U\longrightarrow A,
 \qquad C'=\operatorname{Sym}_R U'\longrightarrow A'
\]
があり、$A,A'$ はそれぞれ $C,C'$ 上平坦とする。
さらに任意の可換 $R$-代数 $R_1$ への変更後に
\[
 Z(A_{R_1})_d=\operatorname{Sym}^d_{R_1}U_{R_1},
 \qquad
 Z(A'_{R_1})_d=\operatorname{Sym}^d_{R_1}U'_{R_1}
\]
が成り立つと仮定する。このとき制限写像
\[
 \underline{\operatorname{Isom}}_{\mathrm{gr},R}(A,A')
 \longrightarrow
 \underline{\operatorname{Isom}}_{\mathrm{gr},R}(A^{(d)},(A')^{(d)})
\]
は $\mu_d$-torsor である。
その核の $\lambda\in\mu_d$ は、元の次数 $j$ に $\lambda^j$ を掛ける。
\end{theorem}

\begin{proof}
$U$ を自由にし、基底を $u_1,\ldots,u_s$ とする。
$Q=A^{(d)}$ と置く。
$z\in Z(Q)$ と次数一の元 $a\in A_1$ について
$u_1^{d-1}a$ は $Q$ に属し、
\[
 u_1^{d-1}[z,a]=[z,u_1^{d-1}a]=0
\]
である。$A$ の $C$-平坦性により $u_1$ は $A$ 上の非零因子なので、
$[z,a]=0$。次数一生成性から
$Z(Q)=Z(A)^{(d)}$ を得る。
従って $Z(Q)_1=\operatorname{Sym}^dU$ であり、
その生成する中心部分代数は内在的に $D=C^{(d)}$ と同定される。
また $Z(A)_1=U$ も従う。実際、$z\in Z(A)_1$ なら
$u_i^{d-1}z\in C_d$ なので、$i=1,2$ を用いて
$z\in C[u_1^{-1}]\cap C[u_2^{-1}]=C$ と書ける。
局所化した $C$ は局所化した $A$ に単射であり、Laurent 多項式の比較でこの交叉を取れる。
次数一部分を取れば $z\in U$ である。従って全ての次数付き同型は
指定した中心多項式環を保つ。

$Q\simeq Q'$ は $D\simeq D'$ を誘導する。
補題\ref{lem:nfm-veronese-cone}により、fpqc 局所でこの同型を
$C\simeq C'$ へ持ち上げる。
以下では Veronese 側の次数を $d$ 倍して元の次数に戻す。
各 $u_i$ について積写像
\[
 Q[(u_i^d)^{-1}]
 \otimes_{D[(u_i^d)^{-1}]}C[u_i^{-1}]
 \longrightarrow A[u_i^{-1}]
\]
は同型である。実際、左辺の第二因子は
\[
 C[u_i^{-1}]
 =\bigoplus_{j=0}^{d-1}u_i^jD[(u_i^d)^{-1}]
\]
と分解し、右辺の次数が $j$ modulo $d$ の部分も
$u_i^jQ[(u_i^d)^{-1}]$ と一致する。
よって $Q\simeq Q'$ と選んだ $C\simeq C'$ は、全てのこの局所化を同定する。

$C$ の二変数 $u_1,u_2$ に関する列
\[
 0\longrightarrow C\longrightarrow
 C[u_1^{-1}]\oplus C[u_2^{-1}]
 \longrightarrow C[(u_1u_2)^{-1}]
\]
は最初の二箇所で完全である。
これは多項式の Laurent 展開から任意の係数環上で直接確かめられる。
$A$ は $C$ 上平坦なので、これを tensor して
\[
 A=A[u_1^{-1}]\cap A[u_2^{-1}]
 \quad\text{in }A[(u_1u_2)^{-1}]
\]
を得る。従って局所化の同型は一意に $A\simeq A'$ を与える。
構成は局所化の重なりで一致し、$U$ の自明化を変えても一意性により降下する。

$C$ への持上げを固定すると $A$ の持上げは一意であり、
$C$ の二つの持上げの差 $u_i\mapsto\lambda u_i$ は
上の modulo $d$ 分解で次数 $j$ に $\lambda^j$ を掛ける。
よって持上げ層はちょうど $\mu_d$-torsor である。
全ての段階が係数変更と両立する。
\end{proof}

\begin{remark}
この定理は低次数中心の仮定を省略した Veronese 復元を主張しない。
また $s\geq2$ は最後の二変数交叉に使われる。
同型持上げの $\mu_d$ は次数スカラー $\mathbf G_m$ の部分群であり、
構造対象付き Abel 圏の新しい慣性ではない。
一つの同型の大域持上げ障害は $H^1_{\mathrm{fppf}}(-,\mu_d)$ にあり、
対象の有効降下に現れる $H^2$ の障害とは区別する。
\end{remark}

\subsection{一ブロック族での低次数中心}

\begin{proposition}[有限個の楕円位数条件]
\label{prop:nfm-center-bound}
$k$ を代数閉体、$\operatorname{char}k=0$、$U$ を階数 $r\geq2$ の
次数一次元からなる空間とし、$C=\operatorname{Sym}_kU$ とする。
$B$ は三つの active 生成元を持つ濾過付き二次代数で
\[
 \operatorname{gr}_{\mathrm{act}}B=S\otimes_k C,
 \qquad S/(\Omega)=H(E,\mathcal L,\sigma)
\]
とする。ここで $S$ は三次元楕円 AS 正則代数、
$\Omega\in Z(S)_3$ は正則元、$E$ は滑らかな楕円曲線である。
$G\in Z(B)_3$ の最高 active 記号が $\Omega$ であるとし、
$\sigma^j\ne\operatorname{id}$ が
$1\leq j\leq q:=\max(3,d)$ について成り立つとする。
このとき
\[
 Z(B)_1=U,\qquad Z(B)_3=kG\oplus\operatorname{Sym}^3U.
\]
任意の $g\in\operatorname{Sym}^3U$ に対して
$A=B/(G+g)$ は $C$ 上自由であり、
\[
 \operatorname{gr}_{\mathrm{act}}A=H(E,\mathcal L,\sigma)\otimes_k C,
 \qquad Z(A)_d=\operatorname{Sym}^dU.
\]
従って $d=r=n-1$ の場合に定理\ref{thm:nfm-central-veronese}を適用できる。
\end{proposition}

\begin{proof}
$H=H(E,\mathcal L,\sigma)$ は捻った Laurent 関数環
$k(E)[t,t^{-1};\sigma]$ に埋め込まれる。
次数 $j>0$ の非零中心元 $a t^j$ があれば、
次数一切断との可換性から $\sigma^j$ はその全ての比を固定する。
$\deg\mathcal L=3$ の完全線型系は $E$ を埋め込むので、
これらの比は $k(E)$ を生成する。従って $\sigma^j=\operatorname{id}$ となる。
ゆえに $Z(H)_j=0$ $(1\leq j\leq q)$ である。
$S/(\Omega)=H$ から
$Z(S)_1=Z(S)_2=0$, $Z(S)_3=k\Omega$ を得る。

$B$ の中心元の最高 active 記号も中心である。
次数一では正の active 成分が消え、$Z(B)_1=U$。
次数三では active 三次の成分は $k\Omega$ に属するので $G$ の倍を引き、
次いで active 次数二、一の成分を消せば純中心三次だけが残る。

$G+g$ の最高記号は正則元 $\Omega$ なので、
最高次数に関する帰納法により $G+g$ は左右正則で、表示した商の associated graded を得る。
$H$ の各次数の $k$-基底を斉次に持ち上げれば、
それらは $A$ の $C$-基底を与える。
最後に $a\in Z(A)_d$ の正の最高 active 次数 $j\leq d$ は
$Z(H)_j\otimes C_{d-j}=0$ に属するため不可能である。
よって $a\in C_d$ である。
\end{proof}

\begin{remark}[相対版で必要な普遍性]
\label{rem:nfm-relative-center}
有限型被約 $\mathbf C$-基底上で、上記の PBW 分裂と中心三次の
最高記号の直和因子性が係数変更と両立しているとする。
このとき $U\subset B_1$、$RG\oplus\operatorname{Sym}^3U\subset B_3$、
$\operatorname{Sym}^dU\subset A_d$ を除いた商は有限局所自由である。
各商から三つの active 生成元との交換子を取る写像は全閉点で単射なので、
最大小行列式の可逆開上で分裂単射になる。
従って中心の三つの等式は任意の、非平坦・非被約な底変更後にも成立する。
この分裂単射の確認をせずに、幾何ファイバーの中心計算だけを普遍的な等式へ拡張しない。

固定した $n$ では除外する条件は有限個である。
楕円曲線の平行移動を動かす空間で $\sigma^j=1$ は閉条件であり、
無限位数の平行移動は全てのこれらの閉集合を避ける。
従って所定の楕円パラメータ空間における共通開は非空である。
三様体の低次数中心条件だけで任意の $n$ を扱うことはできない。
\end{remark}

\section{一ブロックの次数付きモジュライと中心三次 torsor}
\label{sec:nfm-graded-stack}

\subsection{表示空間と全ての次数付き同型}

$V=X\oplus U$, $\dim X=3$, $\dim U=r$ とし、
$X$ の基底を $x,y,z$ とする。
$F\in\operatorname{Sym}^3V/\operatorname{Sym}^3U$ に対して
\[
 B_F=T(V)/\bigl([U,V],\ [y,z]-\operatorname{sym}(F_x),\
 [z,x]-\operatorname{sym}(F_y),\ [x,y]-\operatorname{sym}(F_z)\bigr)
\]
と置く。
$F$ は生成元を可換化した多項式として扱う規約であり、
純中心三次の変更は関係式を変えない。
$P\subset\operatorname{GL}(V)$ を $U$ を保つ parabolic、
$\overline P=P/\mathbf G_m$ とする。
$g\in P$ の $V/U$ 上の作用を $M$ とすると、作用は
\[
 g\cdot F=(\det M)^{-1}g(F)\pmod{\operatorname{Sym}^3U}
\]
である。中心スカラーはこの作用に自明に作用する。

一ブロックの PBW・楕円・中心三次の入力を満たす $P$-不変な非空開を
$T\subset\operatorname{Sym}^3V/\operatorname{Sym}^3U$ とする。
ここでいう入力は、命題\ref{prop:nfm-center-bound}の associated graded と中心元が
相対的に存在し、注意\ref{rem:nfm-relative-center}の有限局所自由性を持つことを含む。
この開の一ブロック量子化入力は既存の一般次元量子化構成から得られる
\cite{CubicQuantization}。
具体的には、純 active 三次を滑らかな Hesse 開に取り、
active 次数を下げる混合 potential の PBW 変形を用いる。
中心変数の個数はこの三変数 PBW 計算の係数環を変えるだけである。
Hesse 座標と中心平行移動で三つの混合二次交叉項を消した開では、
Rains の有限中心公式の係数に $u$ の一次式・二次式を代入できる
\cite[Theorem~3.3.2 and Section~9]{NFMEG}。
有限個の分母と最高 active 記号の係数を反転し、
$\sigma^j\ne1$ $(1\leq j\leq\max(3,r))$ を課す。
この開を $P$ で飽和すれば、所要の非空不変開を得る。
相対 PBW 分裂と中心核の普遍性は注意\ref{rem:nfm-relative-center}の有限行列で確認する。
この一ブロック入力自体を本節で新しく証明した結果とは数えない。

\begin{proposition}[表示空間の同型問題の忠実性]
\label{prop:nfm-graded-presentation}
$F,F'\in T$ に対して、$B_F\simeq B_{F'}$ の次数付き同型は、
$g(F)=(\det M)F'$ modulo $\operatorname{Sym}^3U$ を満たす $g\in P$ に一致する。
従って $[T/\overline P]$ は、この族の次数付き周囲代数を
次数スカラーで割った同型問題を表す。
\end{proposition}

\begin{proof}
次数一生成性から代数同型は $V$ 上の線型同型で決まる。
中心一次部分が $U$ なので、この同型は $P$ に属する。
標数零での分解 $V^{\otimes2}=\bigwedge^2V\oplus\operatorname{Sym}^2V$ に関して、
$B_F$ の関係空間は二次双ベクトル
$\pi_F:\bigwedge^2V\to\operatorname{Sym}^2V$ の graph である。
従って関係空間から $\pi_F$ は一意に復元される。
$\pi_F$ は三つの active 偏微分によって決まり、
それらが全て零となる三次式は $\operatorname{Sym}^3U$ に限る。
active 三重ベクトルは $M$ により $(\det M)^{-1}$ 倍されるので、
関係空間を移す条件はちょうど表示した式である。
逆向きは関係式への代入で従う。
この議論は有限局所自由関係束を用いれば係数環上でも同じである。
\end{proof}

\begin{lemma}[有限安定化子を持つ点]
\label{lem:nfm-finite-stabilizer}
$T$ の中で $\overline P$-安定化子が有限となる部分 $T^{\mathrm{fin}}$ は
非空な不変開である。
\end{lemma}

\begin{proof}
滑らかな純 active 三次式 $F_3(x,y,z)$ を所定の楕円開に取り、
\[
 F=F_3+xQ_1(u)+yQ_2(u)+zQ_3(u),\qquad
 Q_1=\sum_{i=1}^r u_i^2,\quad Q_2=\sum_{i=1}^r\lambda_i u_i^2
\]
とする。ここで $\lambda_i$ は互いに異なり、$Q_3$ は任意である。
ここではまず、全表示アフィン空間上の $\overline P$-安定化子を調べる。
この計算自体に PBW 条件は必要ない。

$\mathfrak p$ の元を active 部分 $M$、中心部分 $D$、
active から中心への項 $T_x,T_y,T_z\in U$ に分ける。
安定化の式の純 active 三次部分は
$M(F_3)=\operatorname{tr}(M)F_3$ である。
滑らかな平面三次の射影自己同型群は有限なので $M=aI_3$ である。
実際、その楕円曲線上の平行移動のうち次数三の埋込みを保つものは三捩れ点に限られ、
原点を保つ自己同型群も標数零では有限である。
active 次数二の成分は
\[
 T_x(F_3)_x+T_y(F_3)_y+T_z(F_3)_z=0
\]
である。滑らかな $F_3$ の三つの偏微分は線型独立なので全ての $T_i$ は零となる。
active 次数一の成分から $D(Q_i)=2aQ_i$ $(i=1,2)$ を得る。
$D-aI_r=H$ と置けば
\[
 H^{\mathsf t}+H=0,\qquad
 H^{\mathsf t}\Lambda+\Lambda H=0,
 \qquad \Lambda=\operatorname{diag}(\lambda_1,\ldots,\lambda_r).
\]
よって $H$ は $\Lambda$ と可換する反対称行列であり、固有値が相異なるので $H=0$。
残る安定化子の Lie 代数は次数スカラーのみである。
従って $\overline P$ の安定化子の Lie 代数は零となる。
標数零の有限型代数群は滑らかなので、この安定化子は有限である。

一般の点でも無限小安定化子は係数に多項式的に依存する有限行列の核である。
従って核が零となる条件は全表示空間の非空開である。
表示アフィン空間は既約で、$T$ も非空開なので、両者の交わりは非空である。
\end{proof}

\begin{theorem}[一ブロックの代数的周囲モジュライ]
\label{thm:nfm-ambient-dimension}
$\mathfrak B_n=[T^{\mathrm{fin}}/\overline P]$ は滑らかな既約有限型
Deligne--Mumford stack で、その次元は
\[
 b_n=\binom{n+4}{3}-\binom{n+1}{3}
       -\bigl((n-1)^2+3(n-1)+8\bigr)
     =\frac{n^2+7n-4}{2}
\]
である。これは次数付き代数の stack についての結論であり、
周囲の全ての無偏極 Abel 圏同値を分類するとは主張しない。
\end{theorem}

\begin{proof}
$T^{\mathrm{fin}}$ は次元
$\binom{r+5}{3}-\binom{r+2}{3}$ のアフィン空間の非空開である。
$\overline P$ は次元 $r^2+3r+8$ の滑らかな連結代数群である。
従って商 stack は滑らかで既約、有限型である。
安定化子は有限で無限小部分を持たないので対角は非分岐であり、商は DM である。
商 stack の次元公式から表示した値を得る。
\end{proof}

\subsection{中心三次の射影空間と加法 torsor}

\begin{theorem}[一般次元の中心三次 torsor]
\label{thm:nfm-central-torsor}
$T^{\mathrm{fin}}$ 上で
\[
 \mathcal E=Z(\mathcal B)_3,\qquad
 \mathcal W=\operatorname{Sym}^3\mathcal U,\qquad
 \mathcal L=\mathcal E/\mathcal W
\]
と置く。$\mathcal L$ は可逆層である。
$\mathbf P_{\mathrm{lines}}$ を一次元部分束の規約とし、
\[
 J=\mathbf P_{\mathrm{lines}}(\mathcal E)
      \setminus\mathbf P_{\mathrm{lines}}(\mathcal W)
\]
とする。この $J\to T^{\mathrm{fin}}$ は
$\mathcal{H}om(\mathcal L,\mathcal W)$ の加法群の torsor であり、
商によって
\[
 \tau_n:\mathfrak C_n=[J/\overline P]\longrightarrow\mathfrak B_n
\]
を得る。射 $\tau_n$ はスキームで表現可能、アフィン、滑らか、全射で、相対次元は
\[
 s_n=\binom{n+1}{3}
\]
である。$\mathfrak C_n$ は滑らかな既約有限型 DM stack で、
\[
 c_n=\dim\mathfrak C_n
 =\binom{n+4}{3}-(n+2)^2+3n-2.
\]
$J$ の全ての線は左右正則な中心三次元を与え、
その商は一般次元の三次 Hilbert 級数を持つ。
\end{theorem}

\begin{proof}
$0\to\mathcal W\to\mathcal E\to\mathcal L\to0$ を考える。
$J$ の $S$-点である直線 $M\subset\mathcal E_S$ は
$M\to\mathcal L_S$ が同型であるという条件を満たす。
従って $J$ はこの完全列の分裂の層を表す。
二つの分裂の差は $\mathcal{H}om(\mathcal L,\mathcal W)$ の一意な切断であり、
局所自由層の列は Zariski 局所で分裂する。
よって $J$ は所定の加法群の torsor である。

中心スカラーの重みは $\mathcal U$ 上で一、
$\mathcal E,\mathcal W,\mathcal L$ 上で三である。
従って $\mathcal{H}om(\mathcal L,\mathcal W)$ 上の重みは零で、
$\overline P$ へ降下する。一方 $J$ も直線の空間なので中心スカラーは自明に作用する。
分裂の差による作用と torsor 同型は同変であり、商へ降下する。
任意の $S\to\mathfrak B_n$ を $\overline P$-torsor 上で表示して基底変更すると、
得られる射はアフィン空間束である。
fpqc アフィン降下 \cite[Lemma~35.37.1, Tag~0245]{NFMStacksAffineDescent} により、
元の基底でもアフィンスキームで表現される。
滑らかさ、全射性、相対次元も降下する。

局所的に一つの中心元 $G$ を選べば $J$ の全ての線は $k(G+g)$
$(g\in\operatorname{Sym}^3U)$ と書ける。
命題\ref{prop:nfm-center-bound}により最高 active 記号は常に $\Omega$ であり、
正則性と商の PBW 基底が得られる。
特に次数 $j$ の Hilbert 関数は
$\binom{j+N-1}{N-1}-\binom{j+N-4}{N-1}$ である。
最後に滑らかなアフィン空間束と前定理から stack の性質と次元公式を得る。
\end{proof}

\begin{remark}[ファイバーの意味]
周囲代数とその同定を固定した幾何学的 $2$-ファイバーは
$\mathbf A^{s_n}$ の torsor である。
周囲の同型類だけを固定した同型類の集合には有限群
$\operatorname{Aut}_{\mathrm{gr}}(B)/\mathbf G_m$ の作用が残る。
また $\mathcal U,\mathcal E$ 自体を $\mathfrak B_n$ 上の通常の束として降下させたわけではない。
原点を選ばずに torsor を零切断付きベクトル束と同一視しない。
\end{remark}

\begin{proposition}[複数ブロックでの正しい相対形]
\label{prop:nfm-multiblock-torsor}
一般に $0\to W\to E\to L\to0$ が有限局所自由層の完全列で
$\operatorname{rank}L=k$ なら、
$\mathbf P_{\mathrm{lines}}(E)\setminus\mathbf P_{\mathrm{lines}}(W)$ は
$\mathbf P_{\mathrm{lines}}(L)$ 上の
$\mathcal{H}om(\mathcal O(-1),W)$ の加法 torsor である。
従って $k\geq2$ では固定周囲上の中心線の空間にも
射影的な $k-1$ 次元の方向が残る。
\end{proposition}
\begin{proof}
$\mathbf P_{\mathrm{lines}}(L)$ 上の普遍直線
$\mathcal O(-1)\subset L$ に沿って元の完全列を引き戻す。
$\mathbf P(E)\setminus\mathbf P(W)$ の点は、その引戻し列の分裂と一対一に対応する。
前定理と同じ分裂の差の議論を適用する。
\end{proof}


% !TEX root = ../nc-cubic.tex
\section{一ブロック族の普遍相対双対性}
\label{sec:nfm-relative-duality}

本節は、次数付きモジュライから構造対象付き Abel 圏へ移る際の
相対双対入力を構成する。
$N=n+2\geq5$、$r=N-3=n-1$ とする。
三様体稿\cite{Threefold}の potential 解消と導来局所化の方法を、
中心変数が $r$ 個の場合に展開する。
有限中心性、楕円平行移動の有限位数性、あるいは tails の dg smoothness は使わない。

\subsection{中心差の Koszul 解消と真正な双加群双対}

次の仮定は、既に用いた一ブロックの相対 PBW 入力を、
双加群解消の構成に必要な形で明記したものである。
$R$ は被約な有限型 $\mathbf C$-代数、
$Q=R[u_1,\ldots,u_r]$ とし、全ての $u_\alpha,x,y,z$ の次数を一とする。
$Q$ 上の cyclic potential
\[
 \Phi=xyz-yxz-\operatorname{Symm}(F),\qquad
 B=Q\langle x,y,z\rangle/(\partial_x\Phi,\partial_y\Phi,\partial_z\Phi)
\]
を考える。純中心三次は cyclic 微分で消えるので任意に除いてよい。
active 次数の最高三次部分を $\Phi_3$、その Jacobian 代数を
$S=R\langle x,y,z\rangle/(\partial\Phi_3)$ とする。
以下では次を仮定する。
\begin{enumerate}
\item 各 $S_a$ は有限射影 $R$-加群であり、全ての幾何学的ファイバーは
三次元二次 Koszul CY 代数で Hilbert 級数 $(1-t)^{-3}$ を持つ。
\item active PBW 同型 $\operatorname{gr}_{\mathrm{act}}B=S\otimes_RQ$ があり、
その濾過は $Q$-加群として分裂する。
この同型と分裂は任意の係数変更後にも PBW 性を与える。
\end{enumerate}
滑らかな純 active 三次を持つ Jacobian 族では、体上の入力は
Mori--Smith\cite[Theorem~4.2]{MoriSmith}と
Berger--Taillefer\cite[Theorems~3.1, 3.6]{NFMBeTa}による。
相対 PBW 分裂は入力論文\cite{NCCubicA}の表示チャートで構成される。
以下の結論はこれらの体上の定理を任意係数環へ直接引用したものではなく、
有限射影複体と普遍恒等式から導く。

\begin{theorem}[任意個の中心変数を持つ相対 CY 解消]
\label{thm:nfm-relative-ambient-cy}
上の仮定の下で、$B$ は有限な相対二次 Koszul 双加群解消
\[
 K_B^{-i}=B\otimes_RW_i\otimes_RB\quad(0\leq i\leq N),
 \qquad K_B\xrightarrow{\sim}B
\]
を持つ。$W_i$ は内部次数 $i$、階数 $\binom Ni$ の有限射影加群であり、
$W_0=R$、$W_N$ は可逆である。
この解消とその片側 augmentation は任意の可換 $R$-代数への係数変更後にも完全である。
また active 変数と中心変数の順序を固定すると、真正な次数付き双加群として
\begin{equation}
\label{eq:nfm-relative-ambient-cy}
 \Theta_{B/R}:=
 \mathbf R\underline{\operatorname{Hom}}_{B\otimes_RB^{\mathrm{op}}}
       (B,B\otimes_RB^{\mathrm{op}})
 \simeq B(N)[-N]
\end{equation}
を得る。この同型も任意の、非平坦・非被約な係数変更と両立する。
\end{theorem}

\begin{proof}
まず $Q$ に相対的な長さ三の potential 複体を作る。
$x_1=x,x_2=y,x_3=z$、$r_i=\partial_{x_i}\Phi$ とし、
$D=B\otimes_QB^{\mathrm{op}}$ と置く。
自由基底 $e_i,f_i,e$ の内部次数をそれぞれ $1,2,3$ とする。
\[
 K_Q:\quad 0\longrightarrow D(-3)\longrightarrow D(-2)^3
 \longrightarrow D(-1)^3\longrightarrow D\longrightarrow0
\]
の微分を
\[
 d_1(e_i)=\delta_i,\qquad
 d_2(f_i)=\sum_jH_{ji}e_j,\qquad
 d_3(e)=\sum_i\delta_i f_i
\]
で定める。ここで
\[
 \delta_i=x_i\otimes1-1\otimes x_i^{\mathrm{op}},\qquad
 H_{ji}=\sum_{r_i=a x_j b}a\otimes b^{\mathrm{op}}
\]
とし、最後の和は $r_i$ を自由 $Q$-代数の語で展開して係数付きで取る。
普遍微分の telescoping と cyclic 恒等式 $\sum_i[x_i,r_i]=0$ により
\[
 \sum_jH_{ji}\delta_j=r_i\otimes1-1\otimes r_i^{\mathrm{op}},\qquad
 \sum_i\delta_iH_{ji}=r_j\otimes1-1\otimes r_j^{\mathrm{op}}.
\]
従って $K_Q$ は複体である。
active 濾過の随伴複体は $S$ の potential 解消を $Q$ へ延長したものである。
$S$ の解消の各内部次数は有限射影 $R$-加群の有界複体であり、
増大複体は全幾何学的ファイバーで完全である。
有限射影複体を右端から局所的に分裂することにより、
この完全性は任意係数変更後にも保たれる。
active 濾過は各全次数で有限なので、元の $K_Q\to B$ も普遍的に完全である。

自己双対性は fibre ごとの CY 性から推測しない。
反自己同型 $\tau(a\otimes b^{\mathrm{op}})=b\otimes a^{\mathrm{op}}$ により
右 $D$-加群を通常の双加群へ戻す規約を取ると、
cyclic 語の二文字の位置を交換する恒等式から
\[
 H_{ji}=\tau(H_{ij}),\qquad \tau(\delta_i)=-\delta_i
\]
を得る。
双対複体の微分を $d^j(f)=(-1)^{j+1}f d$ とすれば、
\[
 1^\vee\mapsto e,\quad e_i^\vee\mapsto-f_i,\quad
 f_i^\vee\mapsto e_i,\quad e^\vee\mapsto-1
\]
は複体の同型 $K_Q^\vee\simeq K_Q(3)[-3]$ を与える。
従って $\Theta_{B/Q}\simeq B(3)[-3]$ である。
これは Berger--Taillefer の potential 解消の自己双対性を
同じ符号規約で係数環 $Q$ 上に書いた計算である。

次に $C=B\otimes_RB^{\mathrm{op}}$ と置く。
PBW 分裂により $B$ は $Q$ 上射影で、$C$ は
$Q\otimes_RQ=R[u_L,u_R]$ 上平坦である。
従って
\[
 z_\alpha=u_{\alpha,L}-u_{\alpha,R}\quad(1\leq\alpha\leq r)
\]
は $C$ 上の中心正則列で、$D=C/(z_1,\ldots,z_r)$ である。
順序を指定したこの正則列の Koszul 解消を双対化すると、
真正な $D$--$C$ 双加群として
\[
 \mathbf R\underline{\operatorname{Hom}}_C(D,C)\simeq D(r)[-r]
\]
を得る。左 $D$-作用は定義域の右掛算による前合成から生じ、
唯一の cohomology 上では通常の左掛算となる。
中心列であるため normalizing automorphism は現れない。
change of rings の随伴により
\[
 \mathbf R\underline{\operatorname{Hom}}_C(B,C)
 \simeq \mathbf R\underline{\operatorname{Hom}}_D
       (B,D(r)[-r])
 \simeq B(r+3)[-r-3]
\]
となる。全ての右作用を保持したままの随伴であり、
反転規約によって \eqref{eq:nfm-relative-ambient-cy} を得る。
$D$ は $C$-perfect、$B$ は $D$-perfect なので $B$ は $C$-perfect でもある。

最後に二次 Koszul 模型を指定する。
$V=B_1$、$R_2=\ker(V^{\otimes2}\to B_2)$ とし、
$W_i\subset V^{\otimes i}$ を全ての位置の $R_2$ の交わりで定める。
各体値では、全次数を固定した有限の active 濾過の bar spectral sequence により、
$\operatorname{gr}_{\mathrm{act}}B=S\otimes k[u]$ の Koszul 性が $B$ へ移る。
Hilbert 級数は $(1-t)^{-N}$ なので
$\dim W_i=\binom Ni$、$W_{N+1}=0$ である。
$W_i$ を定める有限行列に対し、被約な有限型基底上の一定階数の判定を
$2\leq i\leq N+1$ に適用する。
これにより $W_i$ は有限射影で、その形成は任意係数変更と可換する。
標準 Koszul 双加群複体の各全次数は有限射影複体で、全ファイバーで完全であるから、
同じ局所分裂の議論で普遍的に完全である。
片側で見れば $B$ の有界な射影解消なので分裂し、augmentation との tensor も完全となる。
先に構成した双対同型をこの有限解消で表すと、各内部次数の cone は有界な
flat 加群の完全複体であるから、任意係数変更後にも同型である。
\end{proof}

\subsection{中心三次商の全 twist と compact 対象}

以下では可換環 $R$ 上の $B$ について、定理\ref{thm:nfm-relative-ambient-cy}の
有限 Koszul 解消と双対同型が任意係数変更に両立して与えられているとする。
また $w\in Z(B)_3$、$A=B/(w)$ について、全整数 $a$ の列
\begin{equation}
\label{eq:nfm-relative-central-split}
 0\longrightarrow B_{a-3}\xrightarrow{w}B_a\longrightarrow A_a\longrightarrow0
\end{equation}
が完全で、$A_a$ は有限射影 $R$-加群であると仮定する。
従ってこの列は係数線型に分裂し、任意底変更後にも正則である。
一ブロックの中心三次 $G+g$ では、PBW 最高記号の正則性と
各次数の一定階数によりこの入力が得られる。

\begin{theorem}[普遍的線コホモロジーと相対 properness]
\label{thm:nfm-relative-line-cohomology}
任意の可換 $R$-代数 $T$ と整数 $i,j$ に対し、$a=j-i$ と置くと
\begin{equation}
\label{eq:nfm-relative-line-cohomology}
 \mathbf R\operatorname{Hom}_{\operatorname{QGr}A_T}
   (\mathcal O(i),\mathcal O(j))\simeq
 \begin{cases}
  (A_T)_a[0],&a\geq0,\\
  0,&-r<a<0,\\
  (A_T)_{-a-r}^{\vee}[-n],&a\leq-r.
 \end{cases}
\end{equation}
ここで双対は $T$ 上で取る。
この同型は任意の係数変更と両立する。
全ての $\mathcal O(i)$ は compact であり、
$\bigoplus_{i=0}^{N-1}\mathcal O(i)$ は
$D(\operatorname{QGr}A_T)$ の compact generator である。
compact 対象間の RHom は $T$-perfect で、その形成は任意の導来係数変更と可換する。
\end{theorem}

\begin{proof}
基底を $T$ に変更して $R=T$ と書く。
$M_{<q}=B/B_{\geq q}$ は有限射影係数加群であり、
$K_q=M_{<q}\otimes_BK_B$ は有限生成射影右 $B$-加群の長さ $N$ の解消である。
従って
\[
 C_B(M)=\varinjlim_q\underline{\operatorname{Hom}}_B(K_q,M)
\]
は通常の derived torsion $\mathbf R\Gamma_B(M)$ を計算する。
各項は exact な関手で filtered colimit と直和を保ち、横方向の長さが $N$ なので、
非有界複体にも有限幅の全複体として適用できる。

この模型が torsion 複体上で恒等となることも確認する。
内部次数が上に有界な加群 $L$ では、各固定内部次数について
$\underline{\operatorname{Hom}}_B(K_q,L)$ は $q\gg0$ で
$\underline{\operatorname{Hom}}_B(K_B,L)$ に安定し、その cohomology は $L$ のみである。
任意の torsion 加群は有限個の斉次元で生成される、内部次数が上に有界な
部分加群の filtered colimit である。
従って $\mathbf R\Gamma_B(L)=L$ となり、有界振幅による truncation 比較から
同じ結論が torsion cohomology を持つ非有界複体にも成立する。
torsion 部分圏の包含は完全で、その右随伴が $\Gamma_B$ である。
上で計算した導来随伴の単位が torsion 複体上で同型となるので、
$D(\operatorname{Tors}B)$ は $D(\operatorname{Gr}B)$ に完全忠実に埋め込まれ、
その像は torsion cohomology を持つ複体全体である。
また $\pi_B$ は完全なので、その右随伴 $\omega_B$ は K-injective 複体を保つ。
$\pi_B\omega_B=\operatorname{id}$ から
$\pi_B\mathbf R\omega_B=\operatorname{id}$ が従う。
従って非有界導来 Gabriel 商は
$D(\operatorname{Gr}B)/D_{\mathrm{tors}}(\operatorname{Gr}B)
\simeq D(\operatorname{QGr}B)$ であり、局所化には三角
\[
 \mathbf R\Gamma_B(M)\longrightarrow M\longrightarrow
 \mathbf R\omega_B\pi_BM\longrightarrow
\]
があり、$\mathbf R\omega_B$ は直和を保つ。
従って全ての $\mathcal O_B(i)$ は compact generators となる。
augmentation Koszul 解消を tails に送ると $N+1$ 個の twist の完全列を得る。
その両端の係数が可逆なので、任意の連続する $N$ 個の twist が全 twist を thick に生成する。

有限射影係数加群 $M$ に対し、項別の tensor--Hom 随伴は
\[
 \mathbf R\underline{\operatorname{Hom}}_B(M,B)
 \simeq \Theta_{B/R}\otimes_B^{\mathbf L}M^\vee
 \simeq M^\vee(N)[-N]
\]
を与える。追加の左作用も保持される。
$M=M_{<q}$ として colimit を取ると、真正な双加群として
\[
 \mathbf R\Gamma_B(B)\simeq B^\vee(N)[-N]
\]
となる。
局所化三角の内部次数 $a$ により、周囲の線コホモロジーは
$B_a$ $(a\geq0)$、零 $(-N<a<0)$、
$B_{-a-N}^{\vee}[-N+1]$ $(a\leq-N)$ である。

中心商の導来圏を扱う前に、通常の derived torsion の比較
\begin{equation}
\label{eq:nfm-relative-torsion-comparison}
 i_*\mathbf R\Gamma_A(M)\simeq\mathbf R\Gamma_B(i_*M)
\end{equation}
を確認する。
$J$ を injective $A$-加群とする。induction と restriction の随伴により
\[
 \operatorname{Ext}^q_B(B/B_{\geq l},i_*J)
 \simeq\operatorname{Hom}_A(\operatorname{Tor}^B_q(B/B_{\geq l},A),J)
\]
である。
$q>1$ の Tor は零である。
$q=1$ の kernel は、元の次数が $l-3,l-2,l-1$ の部分にのみ存在するので、
$l'\geq l+3$ に対する遷移は零である。
colimit を取ると $i_*J$ は $\Gamma_B$-acyclic となる。
先の有限幅模型を injective 解消へ適用すれば、非有界複体にも比較が成立する。
従って $\mathbf R\Gamma_A$ も有界振幅を持ち、直和を保ち、torsion 複体上では恒等である。
$B$ の場合と同じ導来随伴の議論により、
$D(\operatorname{Tors}A)\simeq D_{\mathrm{tors}}(\operatorname{Gr}A)$ である。
また $\pi_A$ の完全性と $\pi_A\mathbf R\omega_A=\operatorname{id}$ から
$D(\operatorname{Gr}A)/D_{\mathrm{tors}}(\operatorname{Gr}A)
\simeq D(\operatorname{QGr}A)$ を得る。
この局所化の fibre は $\mathbf R\Gamma_A$ なので、$\mathbf R\omega_A$ も直和を保つ。

中心商への induction と restriction は、この tails の非有界導来圏に随伴
$\mathbf Lj^*\dashv j_*$ を誘導する。
restriction は exact、保守的、直和保存である。
ここで $A$-torsion と忘却後の $B$-torsion は一致し、
$-\otimes_B^{\mathbf L}A$ は $w:M(-3)\to M$ の cone なので torsion 複体を保つ。
従ってこれらの随伴を導来 Gabriel 商へ降ろせる。
また
\[
 j_*\mathcal O_A(a)=\operatorname{Cone}
 \bigl(\mathcal O_B(a-3)\xrightarrow{w}\mathcal O_B(a)\bigr).
\]
周囲の線コホモロジーをこの cone に適用する。
$a\geq0$ では \eqref{eq:nfm-relative-central-split} が $A_a$ を与え、
$-r<a<0$ では両項が零である。
$a\leq-r$ では高次 cohomology 上の写像が
$B_{-a-r}^{\vee}\to B_{-a-N}^{\vee}$ となる。
これは双加群としての自然性により $w$ の双対であり、全射で kernel は
$A_{-a-r}^{\vee}$ である。
cone の次数は $N-2=n$ なので \eqref{eq:nfm-relative-line-cohomology} を得る。

$j_*$ が直和を保つので $\mathbf Lj^*\mathcal O_B(i)=\mathcal O_A(i)$ は compact である。
$j_*$ の保守性から、周囲の compact generator の像も generator となる。
compact 部分はその thick 閉包である\cite[Theorems~3.4, 3.8]{NFMKellerDg}。
線対象間の RHom は表示式により係数 perfect で、全ての係数変更と両立する。
両変数について有限 cone、shift、直和因子で閉じれば、全 compact 対象に結論が延びる。
\end{proof}

\subsection{両側局所双対から相対 Serre trace へ}

以下の双対は内部次数を反転する係数双対
$(M^\vee)_a=\mathbf R\operatorname{Hom}_R(M_{-a},R)$ である。
双加群を双対化すると左右の作用も交換する。

\begin{theorem}[中心三次商の普遍相対 Serre 双対]
\label{thm:nfm-relative-serre}
前定理の入力の下で、任意の可換 $R$-代数 $T$ に対し
$D(\operatorname{QGr}A_T)^c$ の相対 Serre 自己同値は
\[
 \mathsf S=(-r)[n]=(1-n)[n]
\]
である。すなわち任意の compact $K,L$ について、両変数に自然な同型
\[
 \mathbf R\operatorname{Hom}_T
   \bigl(\mathbf R\operatorname{Hom}(K,L),T\bigr)
 \simeq\mathbf R\operatorname{Hom}(L,K(-r)[n])
\]
がある。この同型と合成 pairing・trace は任意の係数変更と両立する。
\end{theorem}

\begin{proof}
$T=R$ と書く。前定理の証明で
\eqref{eq:nfm-relative-torsion-comparison}と非有界 Gabriel 局所化を構成した。
その $\mathbf R\Gamma_A$ は torsion cohomology の導来部分圏への右随伴である。

中心正則元の二項解消と coinduction の随伴から、有限射影係数の $A$-加群 $M$ について
\[
 \mathbf R\underline{\operatorname{Hom}}_B(M,B)
 \simeq\mathbf R\underline{\operatorname{Hom}}_A(M,A)(3)[-1].
\]
前定理で得た $B$ 上の双対を代入すると
$\mathbf R\underline{\operatorname{Hom}}_A(M,A)
 \simeq M^\vee(r)[-N+1]$ である。
$M=A/A_{\geq l}$ のとき、最初は左 $B$ と右 $A$ の作用を保持して計算するが、
唯一の cohomology 上の左 $B$-作用は全射 $B\to A$ を通る。
従ってこれは真正な $A$-双加群の同型である。
colimit と次数ごとの係数双対を取ることにより
\begin{equation}
\label{eq:nfm-relative-local-dual}
 \Gamma:=\mathbf R\Gamma_A(A)\simeq A^\vee(r)[-N+1],\qquad
 \Gamma^\vee\simeq A(-r)[N-1]
\end{equation}
を真正な双加群として得る。

ここから trace へ進むため、導来局所化を双加群のレベルで記述する。
$E=\pi A$ に左 $A$-作用を残した injective 解消 $E\to I$ を取り、
\[
 Q_a=\operatorname{Hom}^{\bullet}(I,I(a))
\]
とする。左作用付き対象の圏は、完全な monad
$M\mapsto\bigoplus_{a\geq0}A_a\otimes_RM(-a)$ の加群圏である。
その自由関手が完全なので、作用を忘れた $I(a)$ も K-injective である。
合成により $Q=\bigoplus_aQ_a$ は内部次数付き dg 代数となり、
左作用は strict な写像 $A\to Q$ を与える。
線対象族の dg Morita 表示\cite[Theorem~3.8(c)]{NFMKellerDg}により
\[
 D(\operatorname{QGr}A)\simeq D(\operatorname{Gr}Q).
\]
この表示で $\mathbf R\omega$ は $A$ への restriction、$\pi$ は extension である。
従って $Q\otimes_A^{\mathbf L}Q\simeq Q$ であり、
基礎 $A$-双加群として $Q\simeq\mathbf R\omega\pi A$ となる。

局所化三角 $\Gamma\to A\to Q\to$ を係数双対化する。
両側局所化 $\mathcal L^e(M)=Q\otimes_A^{\mathbf L}M\otimes_A^{\mathbf L}Q$ は
$Q$-双加群上では恒等である。
実際、任意の $Q$-双加群 $M$ に対し、結合則と homological epimorphism の恒等式から
\[
 Q\otimes_A^{\mathbf L}M
 \simeq(Q\otimes_A^{\mathbf L}Q)\otimes_Q^{\mathbf L}M
 \simeq M
\]
となり、右側の tensor にも同じ計算を適用する。
特に $Q^\vee$ は積を双対化した真正な $Q$-双加群なので
$\mathcal L^e(Q^\vee)\simeq Q^\vee$ である。
一方 $A^\vee$ は内部次数が非正に有界な torsion 加群なので
$\mathcal L^e(A^\vee)=0$ である。
従って \eqref{eq:nfm-relative-local-dual} を使うと、真正な $Q$-双加群の三角
\[
 Q^\vee\longrightarrow0\longrightarrow Q(-r)[N-1]
 \longrightarrow Q^\vee[1]
\]
を得る。その境界写像から
\[
 Q^\vee\simeq Q(-r)[N-2]=Q(-r)[n]
\]
が従う。必要なら係数環上 K-flat な模型を取り、$Q^e$ を導来 tensor で定める。
これにより高次 cohomology から来る両側作用も保持される。

perfect な次数付き $Q$-加群 $K,L$ に対し、評価から来る標準の dg 双対性は
\[
 \mathbf R\operatorname{Hom}_R(\mathbf R\operatorname{Hom}_Q(K,L),R)
 \simeq\mathbf R\operatorname{Hom}_Q
       (L,K\otimes_Q^{\mathbf L}Q^\vee)
\]
である。有限自由 shift では双対の定義であり、両変数について thick に延びる。
先の双加群同型を代入すれば主張を得る。
trace は $L=K$ の双対同型と恒等射での評価によって定める。
この構成は両側作用を保つ dg 写像から来るので、合成 pairing の両変数自然性も従う。

最後に係数変更を確認する。
前定理の有限射影模型は、線対象間の標準 Hom 比較が任意係数変更で同型であることを与える。
導来係数延長は torsion 加群を torsion cohomology の複体へ送るので、
dg 導来商へ降り、合成を保つ比較
$Q_R\otimes_R^{\mathbf L}T\simeq Q_T$ を与える。
局所化三角の fibre も、$B$ 上の有限幅の torsion 模型から標準写像ごと係数変更と可換する。
各内部次数の $Q$ は係数 perfect なので係数双対も可換する。
従って上の境界写像、評価、trace の全構成が任意の係数変更と両立する。
\end{proof}

\begin{remark}[供給された入力と残るモジュライ上の問題]
良い相対 PBW 表示チャートと中心三次の正則性の下で、
各次数の有限射影性、全 twist の compact 性、compact 部分の相対 properness、
真正な相対 Serre 双対、および compact Hom の任意底変更は、本節で供給された。
CY の規格化に用いた座標の順序は双対の構成のためのもので、
Abel 圏の対象に新しい標識を追加したものではない。
固定表示模型間の全同値の係数変更と fpqc 降下は、
次節の定理\ref{thm:nfm-fixed-model-comparison}で構成する。
一般の Abel 圏族がこの表示族に局所的に属することと、
表示チャートから一般の圏モジュライへの対象の効果化は別の問題として残る。
本節は tails の dg smoothness、全 Abel モジュライの代数性、
あるいは大域残余射の存在を結論しない。
\end{remark}

% !TEX root = ../nc-cubic-n-fold.tex
\section{構造対象だけを保つ Abel 同値への条件付き移行}
\label{sec:nfm-abel-conditional}

\begin{lemma}[中心正則対による全次数の切断飽和性]
\label{lem:nfm-central-saturation}
$R$ を可換環、$A=\bigoplus_{j\geq0}A_j$ を
$A_0=R$ を満たす次数一生成の $R$-代数とし、$A_1$ は有限生成とする。
階数 $s\geq2$ の有限局所自由 $R$-加群 $U$ と中心的な次数付き写像
\[
 C=\operatorname{Sym}_R U\longrightarrow A
\]
があり、$A$ は $C$-加群として平坦であると仮定する。
$\operatorname{QGr}A$ は、次数付き右 $A$-加群の圏を
\[
 \operatorname{Tors}A
 =\{M\mid \text{各斉次元 }m\in M\text{ は、ある }q
       \text{ について }mA_{\geq q}=0\}
\]
で割った Gabriel 商とし、$\mathcal O_A=\pi A$ と置く。
このとき全整数 $j$ について、左乗法による標準写像
\begin{equation}
\label{eq:nfm-central-saturation}
 A_j\xrightarrow{\ \sim\ }
 \operatorname{Hom}_{\operatorname{QGr}A}
     (\mathcal O_A,\mathcal O_A(j))
\end{equation}
は同型である。ここで $j<0$ なら $A_j=0$ とする。
特に $\operatorname{End}(\mathcal O_A)=R$ である。
結論は任意の可換 $R$-代数への係数変更後にも成立し、
標準写像は係数変更と両立する。
\end{lemma}

\begin{proof}
$A_1$ の有限生成性と次数一生成性により各 $A_q$ は有限生成である。
従って $\operatorname{Tors}A$ は部分対象・商・拡大・任意直和に閉じた
localizing 部分圏であり、表示した Gabriel 商が存在する。
拡大に閉じることは、$mA_{\geq q}$ が torsion 部分に入ったとき、
有限生成な $A_q$ の像の有限個の生成元を同時に消す上界を取り、
次数一生成性を用いれば従う。

まず、$A(j)$ に非零 torsion 部分加群は存在しない。
これは $U$ を局所自明化し、その基底元 $u$ を一つ取れば分かる。
$u$ は $C$ 上非零因子であり、$C$-平坦性によって $A$ 上も非零因子となる。
torsion 元はある $u$ の冪で零になるので、元自身が零である。
この消滅は基底の局所被覆から従う。

Gabriel 商の射の表示
\cite[Chapitre~III, \S1, p.~365]{NFMGabrielQuotient} では、
射は $A/I$ が torsion である斉次右イデアル $I\subset A$ と
標的の torsion 部分対象についての filtered colimit で与えられる。
今は標的 $A(j)$ の torsion 部分は零である。
また $A/I$ は $1$ の像で生成されるので、これが torsion なら
$A_{\geq m}\subset I$ となる $m\geq1$ が存在する。
従って
\begin{equation}
\label{eq:nfm-gabriel-colimit}
 \operatorname{Hom}_{\operatorname{QGr}A}
       (\mathcal O_A,\mathcal O_A(j))
 \simeq\varinjlim_{m\geq1}
       \operatorname{Hom}_{\operatorname{Gr}A}
       (A_{\geq m},A(j)).
\end{equation}
この議論には $R$ や $A$ の noetherian 性を必要としない。

右辺の射を、次数零の右 $A$-線型写像
$f:A_{\geq m}\to A(j)$ で表す。
$U$ が自由な基底開上で、二つの異なる基底元を $u,v$ とする。
多項式環の二変数 \v{C}ech 列を $C$-平坦な $A$ と tensor すれば
\[
 A=A[u^{-1}]\cap A[v^{-1}]
       \quad\text{in } A[(uv)^{-1}]
\]
を得る。各局所化への写像は単射である。
$\alpha=f(u^m)$、$\beta=f(v^m)$ と置くと、中心性と右線型性から
\[
 \alpha v^m=f(u^mv^m)=f(v^mu^m)=\beta u^m.
\]
よって
\[
 a=\alpha u^{-m}=\beta v^{-m}
 \in A[u^{-1}]\cap A[v^{-1}]=A.
\]
$f$ の次数から $a\in A_j$ である。
任意の斉次元 $b\in A_{\geq m}$ に対して
\[
 f(b)u^m=f(bu^m)=f(u^mb)=\alpha b=a u^m b=a b u^m
\]
となり、$u$ の非零因子性から $f(b)=ab$ となる。
従って $f$ は左乗法による写像 $A\to A(j)$ に一意に延長する。
一意性は、延長の差が $u^m$ を零に送れば差の係数も零になることによる。
局所自明化上で得られる元 $a$ はこの一意性により重なりで一致し、
$R$-加群 $A_j$ の切断として貼り合わさる。
ここでは一般の Hom に局所化が可換することを仮定していない。
既に与えた $f$ を局所化し、回復した係数 $a$ を貼り合わせている。
以上を \eqref{eq:nfm-gabriel-colimit} に適用すれば
\eqref{eq:nfm-central-saturation} が従う。

最後に、任意の $R\to R_1$ に対して
$C_{R_1}=\operatorname{Sym}_{R_1}U_{R_1}$ であり、
$A_{R_1}\simeq A\otimes_C C_{R_1}$ は $C_{R_1}$ 上平坦である。
次数一の有限生成性と非負次数性も保たれる。
従って同じ証明を $R_1$ 上で繰り返せる。
同型は左乗法という標準写像から作ったものであるから、
係数変更の比較もこの写像の可換図式によって与えられる。
\end{proof}


\begin{lemma}[一般 Veronese の Gabriel 復元]
\label{lem:nfm-gabriel}
$A$ は可換環 $R$ 上の非負次数付き、次数一生成代数で、
$A_1$ は有限生成とする。任意の $d\geq1$ に対して
\[
 \operatorname{QGr}A\simeq\operatorname{QGr}A^{(d)}
\]
である。全整数 $i$ を用いた切断関手
\[
 N\longmapsto\bigoplus_{i\in\mathbf Z}
 \operatorname{Hom}(\mathcal O_A(-di),N)
\]
はこの同値の下で Gabriel 右随伴に一致し、完全忠実である。
\end{lemma}
\begin{proof}
$Q=A^{(d)}$、$F(M)=M^{(d)}$ とし、
$G(N)=N^{\mathrm{rescale}}\otimes_{Q^{\mathrm{rescale}}}A$ とする。
rescale は次数を $d$ 倍する操作である。
$FG=1$ であり、counit $GFM\to M$ の核と余核は
$d$ の倍数の次数で零となる。
次数一生成性から、その斉次元は高々 $d-1$ 個の生成元を掛ければ
$d$ の倍数の次数へ達して零になるので、これらは tails-torsion である。

$F$ が torsion を反映することも同様である。
斉次元 $m$ に次数を $d$ の倍数へ合わせる長さ $<d$ の生成語を掛け、
有限個の生成語の像を同時に消す次数上界を取ればよい。
よって $F,G$ は torsion 同値を保ち、商上で互いに逆となる。
$F(\mathcal O_A(-di))=\mathcal O_Q(-i)$ であるから、
表示した切断関手は $\operatorname{Gr}Q\to\operatorname{QGr}Q$ の Gabriel 右随伴である。
その完全忠実性が最後の主張を与える。
\end{proof}

\begin{theorem}[相対全 $\mathcal O$-保存同値の比較]
\label{thm:nfm-full-hom-conditional}
$R$ を可換 $\mathbf C$-代数、$n\geq3$, $d=n-1$ とする。
$A,A'$ は定理\ref{thm:nfm-central-veronese}の普遍中心条件と平坦性を満たすとする。
さらに任意の可換係数変更後に次の入力を仮定する。
\begin{enumerate}
\item 各 $A_j,A'_j$ は有限局所自由である。
中心平坦性と補題\ref{lem:nfm-central-saturation}により、
全整数 $j$ で標準切断写像は同型となり、負次数の Hom は零である。
特に $\operatorname{End}(\mathcal O_A)=R$ であり、
$C_d\subset A_d$ と $A'_d$ の対応する単射により $A_d,A'_d$ は忠実である。
\item 全ての $\mathcal O_A(i)$ は $D(\operatorname{QGr}A)$ で compact であり、
compact 部分は係数環上 proper な相対 Serre 双対を持ち、
その Serre 関手は $(-d)[n]$ に同型である。$A'$ も同様である。
\item 係数作用の延長、固定模型間の Hom の fpqc 降下、
標準の Gabriel 局所化が与えられた族の係数変更と両立する。
\end{enumerate}
このとき自然な fpqc 層の同型
\[
 \underline{\operatorname{Isom}}_{\mathrm{gr},R}(A,A')/\mathbf G_m
 \simeq
 \underline{\operatorname{Eq}}_{\mathcal O,R}
 (\operatorname{QGr}A,\operatorname{QGr}A')
\]
がある。右辺は構造対象を保つ全ての完全線型同値を自然同型で割った層である。
構造対象上で恒等である自然自己同型は自明である。
\end{theorem}

\begin{proof}
完全線型同値 $F$ は複体に項別に作用し、非有界導来圏と compact 部分を同定する。
相対 Hom と係数双対を保つので、表現対象の一意性から Serre 関手と可換する。
$H=\mathsf S^{-1}[n]=(d)$ とすると、$F(\mathcal O_A)=\mathcal O_{A'}$ と
一つの可換同型 $FH\simeq H'F$ の正負の反復によって、
全ての $\mathcal O_A(di)$ とその間の Hom・合成が対応する。
補題\ref{lem:nfm-central-saturation}の切断飽和性により、次数付き同型
$\phi:A^{(d)}\simeq(A')^{(d)}$ を得る。
一つの可換同型を一貫して反復しているので、
この同型は orbit 列全体の section $\mathbf Z$-代数の作用を記述する。

補題\ref{lem:nfm-gabriel}により、その全 section 代数上の作用は
元の関手 $F$ を自然同型まで決定する。
定理\ref{thm:nfm-central-veronese}で $\phi$ を fpqc 局所的に
$A\simeq A'$ へ持ち上げれば、その構造移送と $F$ は自然同型となる。
これで層としての全射性を得る。

構造移送が標識付き恒等関手となる次数付き自己同型を考える。
各 $\mathcal O_A(di)$ 上の自然同型の成分は基底の単元であり、
隣接 Hom $A_d$ の忠実性と合成の整合性から、その Veronese 上の作用は
新次数 $j$ に $a^j$ を掛ける一つの次数スカラーとなる。
fpqc 局所で $a$ の $d$ 乗根を取り、それを元の次数スカラーで除くと、
残る作用は定理\ref{thm:nfm-central-veronese}の $\mu_d$ に属する。
これも元の次数スカラーなので、核はちょうど $\mathbf G_m$ である。

自然自己同型の各 orbit 線上の成分も、隣接 Hom の忠実性から全て等しい。
$\mathcal O_A$ 上で恒等なら全て一であり、Gabriel 切断の完全忠実性から
関手全体で恒等である。最後に仮定した Hom 降下と係数変更の整合性により、
局所的に構成した比較は元の基底へ降下する。
\end{proof}

\begin{remark}[証明の状態と適用の境界]
定理\ref{thm:nfm-relative-ambient-cy}、\ref{thm:nfm-relative-line-cohomology}、
\ref{thm:nfm-relative-serre}により、良い相対 PBW 表示チャートでは、
有限 Koszul 対角、真正な相対双加群双対、compact 性、相対 properness、
Serre trace とその非平坦係数変更が供給される。
残る入力は、一般の Abel 圏族の係数作用、全同値の fpqc 降下、
およびそれらと表示模型の係数変更との整合性である。
これらを備えたモジュライ問題に対して、$\mathfrak C_n$ は
全 $\mathcal O$-保存 Abel 問題の全部分 stack へ移る。
この条件付き結論から全 Abel モジュライの代数性・成分性や
Kuznetsov 残余射の幾何化を推論しない。
\end{remark}

% !TEX root = ../cubic-moduli-descent.tex
\section{局所一ブロック型の構造対象付き圏モジュライ}
\label{sec:nfm-local-moduli}

本節では、固定表示模型間の全同値比較を対象の効果化へ進める。
扱う族の範囲を、良い一ブロック模型を fpqc 局所的に持つ族に限定する。
この局所表示条件を満たすこと自体を、任意の非可換三次多様体に対して
証明するわけではない。
一方、この範囲では次数一 shift を標識として要求せず、
構造対象を持つ Grothendieck 圏の族を代数的 stack として分類する。

次数スカラーを除く商と、その上の次数別に重みを持つ twisted algebra は
Behrend--Noohi\cite[Section~2.3]{NFMBehrendNoohi}の一般的構成に既に現れる。
同論文の偏極付き Grothendieck 圏は自己同値もデータに含む
\cite[Section~3.1.2]{NFMBehrendNoohi}。
ここで必要な追加点は、定理\ref{thm:nfm-fixed-model-comparison}により
\emph{全ての} $\mathcal O$-保存同値を表示の同型へ戻すことと、
twisted algebra から大域の $\mathcal O$-marked 圏を実際に構成することである。
後半では、一ブロックの中心階数と active 階数から
持上げ障害の正確な位数を調べる。

\subsection{係数変更を備えた圏族}

$R$ を可換 $\mathbf C$-代数とする。
$R$-線型 Grothendieck Abel 圏 $\mathcal C$ と可換 $R$-代数 $T$ に対し、
\[
 \mathcal C_T=\operatorname{Mod}_T(\mathcal C)
\]
は $T$-作用付き対象の圏を表す。
係数延長は copower $X\mapsto T\otimes_RX$ であり、
反復係数変更には標準の結合同型を用いる。
したがって $(\mathcal C,\mathcal O)$ は、各 affine $R$-scheme 上の
圏と構造対象 $(\mathcal C_T,T\otimes_R\mathcal O)$ を定める。
この係数変更構造を持つ族だけを以下で扱い、
幾何学的ファイバーの圏を単に並べたものとは区別する。

\begin{lemma}[Grothendieck 圏内の忠実平坦降下]
\label{lem:nfm-grothendieck-fpqc}
$\mathcal C$ を $R$-線型 Grothendieck Abel 圏とする。
任意の忠実平坦 $R\to T$ に対して、
\[
 T\otimes_R-:\mathcal C\longrightarrow\mathcal C_T
\]
は完全かつ保守的であり、$\mathcal C$ は $\mathcal C_T$ の
通常の fpqc descent datum の圏に同値である。
また、二つのこのような係数圏族の間の線型完全同値、自然同型、
および指定対象上の同定は fpqc 降下を満たす。
\end{lemma}

\begin{proof}
加群 $M$ による copower $M\otimes_R X$ は、自由加群による表示を用いて
余核として構成でき、$M$ と $X$ のそれぞれについて colimit を保つ。
Lazard の定理\cite[Theorem~10.81.4, Tag~058G]{NFMStacksLazard}により、
平坦 $T$ は有限自由加群の filtered colimit である。
Grothendieck 圏の filtered colimit は完全なので、$T\otimes_R-$ は完全である。

保守性のために、任意の $X$ に対して $X\to T\otimes_RX$ が単射であることを示す。
忠実平坦性により $R\to T$ は純単射で、$Q=T/R$ は平坦である。
実際、任意の $R$-加群 $M$ に対する $M\to T\otimes_RM$ は、
$T$ を tensor した後には積写像による左逆を持つため単射であり、
$\operatorname{Tor}_1^R(Q,M)=0$ が従う。
$Q=\varinjlim Q_i$ を有限自由加群による表示とし、
$T_i=T\times_QQ_i$ と置くと
\[
 0\longrightarrow R\longrightarrow T_i\longrightarrow Q_i\longrightarrow0
\]
は分裂完全で、$\varinjlim T_i=T$ である。
これらに $-\otimes_RX$ を施し、完全な filtered colimit を取ると
$X\to T\otimes_RX$ は単射となる。
完全性と合わせて保守性を得る。

対象と射の降下は、補題\ref{lem:nfm-tails-fpqc}の等化子による証明を適用する。
必要な入力は、今証明した完全性・保守性と Abel 圏内の核だけである。
具体的に、descent datum を持つ $Y\in\mathcal C_T$ の coaction を
$\theta:UY\to T\otimes_RUY$ と書けば、
\[
 X=\ker\bigl(\theta-\eta:UY\rightrightarrows T\otimes_RUY\bigr)
\]
が降下対象である。cocycle 恒等式と $T$-作用の counit 恒等式から
$T\otimes_RX\simeq Y$ を得る。
同じ等化子は射にも適用できる。

関手と自然同型の降下には命題\ref{prop:nfm-equivalence-descent}の証明を用いる。
局所同値を降下対象ごとに適用して関手を作り、逆同値も降下させる。
得られた関手の係数変更が元の関手と一致することは、自由 $T$-作用付き対象と
標準の reflexive coequalizer
\[
 T\otimes_RT\otimes_RUY\rightrightarrows T\otimes_RUY\longrightarrow Y
\]
上で確認する。同値は colimit を保つので、この確認で全対象を扱える。
指定対象上の同定は降下対象間の射として降下する。
\end{proof}

以後 $r=n-1\ge2$, $N=n+2=r+3$ とする。
第\ref{sec:nfm-graded-stack}節の $T^{\mathrm{fin}}$ を、
必要ならさらに $P$-不変な非空開 $T^{\mathrm{good}}$ へ制限する。
この開上の一ブロック模型は、普遍中心条件、中心上の平坦性、
次数別有限局所自由性、および
定理\ref{thm:nfm-fixed-model-comparison}に必要な相対双対性を
全ての係数変更後に満たすものとする。
相対 PBW 分裂と純 active Koszul CY 入力は
第\ref{sec:nfm-relative-duality}節の仮定である。
これらは本節で追加の分類定理として証明するものではない。
中心三次の torsor をこの開へ制限して $J^{\mathrm{good}}$ と書き、
\[
 \mathfrak C_n^{\mathrm{good}}=[J^{\mathrm{good}}/\overline P]
\]
と置く。以下の議論はこの stack の任意の非空開部分にも制限できる。

\begin{definition}[局所一ブロック型の圏族]
\label{def:nfm-local-category}
$\mathfrak M_n^{\mathrm{loc}}(R)$ の対象は、
$R$-線型 Grothendieck Abel 圏と指定対象の対 $(\mathcal C,\mathcal O)$ であって、
ある忠実平坦 $R$-代数 $T$ と $j\in J^{\mathrm{good}}(T)$ に対し
\[
 (\mathcal C_T,\mathcal O_T)
 \simeq (\operatorname{QGr} A_j,\mathcal O_{A_j})
\]
となるものとする。
射は $R$-線型完全同値 $F$ と同定
$\alpha:F(\mathcal O)\xrightarrow{\sim}\mathcal O'$ の対、
二射は $\alpha$ を保つ自然同型とする。
圏族の係数変更は上で定めた $\operatorname{Mod}_T$ によるものとし、
次数一 shift、次数一生成元の束、および局所表示は対象の標識に含めない。
\end{definition}

\subsection{twisted graded module による対象の効果化}

$x\in\mathfrak C_n^{\mathrm{good}}(R)$ とし、
対応する $\overline P$-torsor を $Q\to\operatorname{Spec}R$ とする。
中心拡大
\[
 1\longrightarrow\mathbf G_m\longrightarrow P
 \longrightarrow\overline P\longrightarrow1
\]
に関する $Q$ の持上げ gerbe を $\mathcal G_x$ と書く。
これは $x$ による $[J^{\mathrm{good}}/P]$ の引戻しである。
普遍次数付き代数は $\mathcal G_x$ 上の代数
$\mathcal A=\bigoplus_{j\ge0}\mathcal A_j$ を与え、
inertia の $\lambda\in\mathbf G_m$ は $\mathcal A_j$ 上で $\lambda^j$ として作用する。
この規約を「重み $j$」と呼ぶ。

\begin{lemma}[重み付き tails の大域模型]
\label{lem:nfm-weighted-tails}
$\operatorname{Gr}^{\mathrm{wt}}_{\mathcal G_x}(\mathcal A)$ を、
各 $M_j$ が重み $j$ の準連接層であるような次数付き右
$\mathcal A$-加群 $M=\bigoplus_{j\in\mathbf Z}M_j$ の圏とする。
fpqc 局所的に通常の graded torsion となる加群全体を
$\operatorname{Tors}^{\mathrm{wt}}$ と書けば、
\[
 \mathcal C_x=
 \operatorname{Gr}^{\mathrm{wt}}_{\mathcal G_x}(\mathcal A)
 /\operatorname{Tors}^{\mathrm{wt}},\qquad
 \mathcal O_x=\pi\mathcal A
\]
は $R$-線型 Grothendieck Abel 圏とその指定対象を定める。
任意の $R$-代数 $T$ について標準同値
\[
 (\mathcal C_x)_T\simeq\mathcal C_{x_T}
\]
があり、構造対象と反復係数変更を保つ。
gerbe を中立化する fpqc 被覆上では、これは通常の
$(\operatorname{QGr}A,\mathcal O_A)$ に同値である。
\end{lemma}

\begin{proof}
代数的 stack 上の準連接層の圏は Grothendieck 圏である
\cite[Proposition~96.15.2, Tag~0781]{NFMStacksQCohGrothendieck}。
$\mathbf G_m$-gerbe 上の準連接層には重み分解があり、
各重みへの射影は完全で colimit を保つ。
この主張は gerbe を中立化した後の $\mathbf G_m$-表現の重み分解で確認でき、
fpqc 降下により大域化する。
従って各重みの圏、さらにその全整数にわたる積も Grothendieck 圏である。
ここで generator は各重みへ射影した generator の直和から得られる。

重み付き次数付き層から自由右 $\mathcal A$-加群を作る関手は
\[
 (E_i)_i\longmapsto
 \left(\bigoplus_{i\le j}E_i\otimes\mathcal A_{j-i}\right)_j
\]
である。加群の核、余核、filtered colimit は各次数で計算でき、
この自由加群関手を generator に適用すると加群圏の generator が得られる。
したがって $\operatorname{Gr}^{\mathrm{wt}}(\mathcal A)$ は Grothendieck 圏である。
次数一の有限生成性により torsion は部分対象、商、拡大、直和で閉じる。
このことは fpqc 局所で通常の次数付き加群の証明に帰着する。
従ってこれは localizing 部分圏で、表示した Gabriel 商も Grothendieck 圏である。

次に係数変更を調べる。
$\mathcal G_{x_T}\to\mathcal G_x$ は affine base change であるから、
その上の準連接層は $\mathcal G_x$ 上の $T$-作用付き準連接層に等しい。
この同一視は affine atlas 上の加群の同一視と、その fpqc 降下から従い、
inertia の重みおよび $\mathcal A$-作用を保つ。
従って
\[
 \operatorname{Gr}^{\mathrm{wt}}_{\mathcal G_{x_T}}(\mathcal A_T)
 \simeq\operatorname{Mod}_T
 \bigl(\operatorname{Gr}^{\mathrm{wt}}_{\mathcal G_x}(\mathcal A)\bigr).
\]
Gabriel 商の完全左随伴と完全忠実右随伴に $T$-作用を持たせる。
補題\ref{lem:nfm-coefficient-objects}と同様に、その核はちょうど
$\mathcal A_T$-torsion である。
これで任意の、非平坦なものも含む係数変更との整合を得る。

最後に $\mathcal G_x$ を中立化すると、各重み $j$ の準連接層は通常の
$R$-加群に一意に対応し、$\mathcal A$ は通常の次数付き代数になる。
重み付き加群・torsion・構造対象はそれぞれ通常のものへ移る。
\end{proof}

\begin{theorem}[局所一ブロック型の $\mathcal O$-marked モジュライ]
\label{thm:nfm-local-moduli}
上記の良い表示の入力の下で、
\[
 \mathfrak C_n^{\mathrm{good}}
 \xrightarrow{\ \sim\ }
 \mathfrak M_n^{\mathrm{loc}},\qquad
 x\longmapsto(\mathcal C_x,\mathcal O_x)
\]
は fpqc stack の同値である。
右辺の $2$-groupoid は通常の groupoid と同値であり、
指定対象上で恒等となる自然自己同型は全て恒等である。
特に $\mathfrak M_n^{\mathrm{loc}}$ は滑らかな既約有限型 DM stack であり、
\[
 \dim\mathfrak M_n^{\mathrm{loc}}
 =\binom{n+4}{3}-(n+2)^2+3n-2.
\]
この結論は、任意の局所模型の descent datum に対し、
大域の Grothendieck Abel 圏とその構造対象が有効に存在することを含む。
\end{theorem}

\begin{proof}
まず射を比較する。二つの表示 $A=B/(w)$ と $A'=B'/(w')$ の間の
次数付き同型は、次数二の関係空間から $B\simeq B'$ を誘導する。
実際、三次関係は次数二の関係を変えず、$B$ はその二次関係で定義される。
得られた同型は三次の核直線を互いに移すので、
命題\ref{prop:nfm-graded-presentation}により $J^{\mathrm{good}}$ の $P$-作用の射に一致する。
定理\ref{thm:nfm-fixed-model-comparison}により、次数スカラーを除いたこの射の層は
全ての $\mathcal O$-保存同値の層に一致する。
したがって主張する射は fpqc 局所で完全忠実である。
補題\ref{lem:nfm-grothendieck-fpqc}により、この完全忠実性は大域化する。

局所模型の同値の自然自己同型はスカラーであり、
その $\mathcal O$ 上の成分を恒等に指定するとスカラーは一となる
\textup{(}補題\ref{lem:nfm-natural-scalars}\textup{)}。
一般の対象でも、自然自己同型を良い被覆へ延長してこの結論を適用し、
忠実平坦降下で恒等性を確認できる。
従って各同値の groupoid は離散 groupoid に同値である。
ここで構造対象の標識を忘れた二自己同型の $\mathbf G_m$ と、
表示の次数スカラーを除いた通常の stack とを混同しない。

次に $(\mathcal C,\mathcal O)\in\mathfrak M_n^{\mathrm{loc}}(R)$ を取り、
定義の被覆上で $(\mathcal C_T,\mathcal O_T)$ を模型 $A_j$ と同定する。
重なりで得られる marked equivalence の類は、前段により
商 stack の射に一意に対応する。
元の係数圏族の結合同型は cocycle を与え、
marked natural automorphism の自明性により三重重なりでの曖昧さはない。
商 stack の対象降下から $x\in[J^{\mathrm{good}}/\overline P](R)$ を得る。
補題\ref{lem:nfm-weighted-tails}で構成した $(\mathcal C_x,\mathcal O_x)$ は
同じ局所模型と同じ同値の cocycle を持つ。
両者間の局所同値は補題\ref{lem:nfm-grothendieck-fpqc}で降下するので、
$(\mathcal C,\mathcal O)\simeq(\mathcal C_x,\mathcal O_x)$ である。

同じ議論は、最初から大域の $(\mathcal C,\mathcal O)$ が与えられていない
局所対象と marked equivalence の descent datum にも適用できる。
商 stack 上の対象を先に降下し、重み付き tails を作ることで、
大域の Grothendieck 圏を得る。
したがって対象の効果化を定義に含めて仮定したのではない。
滑らかさ、既約性、有限型性、DM 性と次元は
定理\ref{thm:nfm-central-torsor}から従う。
\end{proof}

\begin{remark}[証明した範囲]
局所的に良い一ブロック表示を持つという条件の下で、
標識は $(\mathcal C,\mathcal O)$ だけで十分である。
ただし、この条件が全ての内在的モジュライ成分の一般点で成立すること、
この stack が未構成の全モジュライの開かつ閉な成分をなすこと、
任意の Abel 圏の変形が局所表示条件を保つことは、この定理には含まれない。
これらには独立の局所表示・変形・成分の定理が必要である。
\end{remark}

\subsection{大域次数付き表示の gerbe と三捩れ障害}

\begin{proposition}[表示の持上げ gerbe]
\label{prop:nfm-presentation-gerbe}
写像
\[
 \rho_n:[J^{\mathrm{good}}/P]\longrightarrow
 \mathfrak M_n^{\mathrm{loc}}
\]
は $\mathbf G_m$-gerbe である。
その class を $\beta$ とし、$e=\gcd(r,3)=\gcd(n-1,3)$ と置くと
\[
 r\beta=3\beta=0,\qquad e\beta=0.
\]
任意の affine $S$ と $x\in\mathfrak M_n^{\mathrm{loc}}(S)$ に対し、
$\beta_x=0$ であることと、$x$ が同じ良い一ブロック型の
大域次数付き代数の tails として表示されることは同値である。
ここで次数一は有限局所自由でよく、自由な基底や中心三次直線の生成元は要求しない。

特に $n\not\equiv1\pmod3$ ならこの gerbe は全体で中立である。
$ar+3b=1$ となる整数 $a,b$ を一つ固定すると、
\[
 \chi(g)=\det(g|_U)^a\det(g|_{V/U})^b
\]
がその中立化を与える。
\end{proposition}

\begin{proof}
次数スカラーは $J^{\mathrm{good}}$ に自明に作用するため、
$P$ から $\overline P$ への商の fibre は
その torsor の $P$-持上げの gerbe である。
$[J^{\mathrm{good}}/P]$ 上には、重み一の階数 $r$ の束 $\mathcal U$ と
階数三の束 $\mathcal V/\mathcal U$ がある。
その determinant はそれぞれ重み $r$ と三の可逆層である。
gerbe 上の重み $m$ の可逆層は $m\beta=0$ を与える。
これは局所表示の三重重なりの scalar cocycle $c_{ijk}$ に対し、
可逆層の遷移写像が $c_{ijk}^{m}$ を coboundary として消すことから分かる。
従って表示した annihilator を得る。

$\beta_x=0$ は持上げ gerbe の中立性、すなわち $P$-torsor への持上げの存在と同値である。
持上げを取れば $P$-同変な普遍代数が $S$ 上へ降下し、
次数一は有限局所自由、中心一次部分は階数 $r$ の部分束になる。
三次関係は中心三次の可逆部分束で生成される。
逆にこの種の大域次数付き表示の、中心一次部分を保つ frame torsor は
$P$-torsor となり、$x$ を $[J^{\mathrm{good}}/P]$ へ持ち上げる。

$e=1$ のとき、$\chi(\lambda\operatorname{id}_V)=\lambda$ である。
写像 $p\mapsto\chi(p)^{-1}p$ は $\overline P$ の $\ker\chi$ への同型を与え、
\[
 P\simeq\mathbf G_m\times\ker\chi,\qquad
 \overline P\simeq\ker\chi
\]
と中心拡大を分裂する。
従って全ての $\overline P$-torsor はこの分裂により持ち上がる。
同じことは、gerbe 上の普遍次数 $j$ の代数束を
$\mathcal A_j\otimes\chi^{-j}$ に置き換えると重みが零になることからも分かる。
これは同じ局所代数を表すが、選んだ整数 $a,b$ に依存する中立化である。
\end{proof}

\begin{remark}[対象の障害と同値の障害]
$\beta_x\in H^2_{\mathrm{fpqc}}(S,\mathbf G_m)$ は
\emph{対象}が大域次数付き表示を持つかどうかの障害である。
一つの持上げが存在すれば、他の持上げとの差は基底の可逆層 $L$ であり、
代数は $\mathcal A_j\mapsto\mathcal A_j\otimes L^{\otimes j}$ と変わる。
一方、命題\ref{prop:nfm-lifting-line}の $L_F$ は、
\emph{既に二つの大域模型を固定した後}の指定同値の持上げ障害である。
後者の $H^1$ の現象から前者の $H^2$ の効果化を推測してはならない。
また $\beta_x=0$ は次数一が自由であること、中心三次直線が自明であること、
あるいは次数一 shift の標準的な選択を意味しない。
\end{remark}

\begin{theorem}[三捩れ障害はどの非空開でも消えない]
\label{thm:nfm-brauer-dense-open}
$n\equiv1\pmod3$、すなわち $r=3m$ とする。
任意の非空開部分 stack
$\mathfrak U\subset\mathfrak M_n^{\mathrm{loc}}$ に対し、
\[
 K=\mathbf C(s,t),\qquad x\in\mathfrak U(K)
\]
で $\beta_x\ne0$ となるものが存在する。
従って $\beta|_{\mathfrak U}$ は位数がちょうど三であり、
$\rho_n$ はどの非空開部分上でも中立化しない。
この主張は非空開全体についてのもので、特殊次元や境界点の例に依存しない。
\end{theorem}

\begin{proof}
まず $K$ 上の非自明な degree 三の中心 division algebra を固定する。
$\zeta$ を原始三乗根とし、
\[
 D=(s,t)_3
 =K\langle a,b\rangle/
 (a^3-s,\ b^3-t,\ ba-\zeta ab)
\]
と置く。$a^i b^j$ $(0\le i,j<3)$ は基底であり、中心は $K$ である。
division 性は $K'=\mathbf C(s)((t))$ へ延長して確認できる。
$L=\mathbf C(s)(a)$、$a^3=s$ とすると $[L:\mathbf C(s)]=3$ で、
$D\otimes_KK'$ は
\[
 f_0+f_1b+f_2b^2,\qquad f_i\in L((t)),\qquad
 bf=\sigma(f)b,\quad\sigma(a)=\zeta a
\]
と書ける。非零元について
\[
 v(f_0+f_1b+f_2b^2)
 =\min_i\bigl(v_t(f_i)+i/3\bigr)
\]
と定める。
三つの候補値は小数部分が異なるので最小項は一意であり、
積の最小項も一意で消えない。従って積の valuation は和に等しく、
この有限次元代数に零因子はない。
有限次元代数の非零元による乗法は単射なら全射なので、これは division algebra である。
よって $D$ 自身も division algebra であり、その Brauer class は非零である。

$U\simeq\mathbf C^3\otimes\mathbf C^m$、$V/U\simeq\mathbf C^3$ を選び、
\[
 \operatorname{GL}_3\longrightarrow P,\qquad
 g\longmapsto\operatorname{diag}(g^{\oplus m},g)
\]
を考える。中心スカラーは $V$ 全体の同じスカラーへ移るから、
$\operatorname{PGL}_3\to\overline P$ を誘導する。
この写像に沿って $P\to\overline P$ を引き戻した中心拡大は
$\operatorname{GL}_3\to\operatorname{PGL}_3$ そのものである。
$D$ に対応する $\operatorname{PGL}_3$-torsor
\cite[Tag~03R6]{NFMStacksBrauerForms}をこの写像で押し出して
$\overline P$-torsor $Q$ を得ると、その持上げ障害は $[D]$ である。
Brauer class の符号規約を逆に取っても、以下の非零性には影響しない。

$\mathfrak U$ の逆像を $J'\subset J^{\mathrm{good}}$ とする。
これは非空な $\overline P$-不変開である。
表示 potential の空間
\[
 W_F=\operatorname{Sym}^3V/\operatorname{Sym}^3U
\]
には $\det(V/U)^{-1}$ で捩じった作用を用いており、
次数スカラーの重みは零である。
従って $W_F$ は $\overline P$ の線型表現であり、
$Q\times^{\overline P}W_F$ は $K$ 上の通常の有限次元ベクトル空間となる。
$T^{\mathrm{good}}_Q$ はその非空開である。
さらに $J^{\mathrm{good}}_Q\to T^{\mathrm{good}}_Q$ は
ベクトル束の加法群による torsor であり、特に滑らかである。

非空開 $J'_Q$ の $T^{\mathrm{good}}_Q$ における像は非空開である。
$K$ は無限体なので、この像には $K$-有理点 $z$ がある。
$z$ 上の fibre $J^{\mathrm{good}}_{Q,z}$ は有限次元ベクトル空間の加法群の torsor であり、
体上では自明である。
ここでは characteristic zero のため、有限 Galois 拡大上の一点の
Galois 軌道の affine barycenter を取ることでも有理点を構成できる。
従ってこの fibre は affine 空間であり、その非空開 $J'_{Q,z}$ も $K$-有理点を持つ。
この一点は $Q\to J'$ の同変写像、すなわち
$x\in[J'/\overline P](K)=\mathfrak U(K)$ を与える。
$x$ の持上げ gerbe の障害は固定した $Q$ の障害 $[D]\ne0$ である。
命題\ref{prop:nfm-presentation-gerbe}の $3\beta=0$ と合わせて、位数はちょうど三となる。
\end{proof}

\begin{remark}[一般点と幾何学的一般点]
この非中立性は、非空な稠密開に制限するだけでは
大域次数付き表示の障害を取り除けないことを意味する。
代数閉体上の個々の幾何学的点では gerbe は中立であり、矛盾はない。
定理は商 stack の幾何学的一般点に非零の体の Brauer class が残るという意味ではなく、
\emph{全ての非空開部分 stack 上で普遍 gerbe が非中立である}という主張である。
\end{remark}

\begin{remark}[Poisson 標的の偏極に対する制約]
後に構成する射がこの表示から得られる ambient 射影空間束を保つとする。
その ambient は $\overline P\to\operatorname{PGL}(V)$ による
Brauer--Severi 束であり、符号を除いて class は $\beta$ である。
もし標的を、実際の相対可逆層 $L$ を持ち、$p_*L$ が階数 $N$ で
基底変更と両立し、$L$ の全切断がその ambient 埋込みを回復する三次族に限定すれば、
ambient は $\mathbf P((p_*L)^\vee)$ となり Brauer class は零でなければならない。
従って $n\equiv1\pmod3$ では、その種の標的への ambient と両立する射を
この stack の非空開全体に作るには、gerbe の中立化という追加データが必要である。
これは Poisson 射そのものの不存在ではない。
Brauer--Severi ambient 内の三次族、または実際の可逆層を選ばない偏極類を
許す標的は、引き続き有効な候補である。
\end{remark}

% !TEX root = ../nc-cubic-n-fold.tex
\section{一ブロック量子化の点スキームと二十四次の定常成分}
\label{sec:points}

この節では $k$ を標数零の代数閉体、$n\geq3$、$r=n-1$、
$V=\langle x,y,z,u_1,\ldots,u_r\rangle$ を生成元空間とする。
$F\in\Sym^3V$ と $h\in k^\times$ に対し、$u_\alpha$ を中心とし
\begin{equation}\label{eq:point-relations}
 [y,z]=h\operatorname{sym}(F_x),\qquad
 [z,x]=h\operatorname{sym}(F_y),\qquad
 [x,y]=h\operatorname{sym}(F_z)
\end{equation}
を課した二次代数を $B$ とする。
$g=F|_{u=0}$ は滑らかな平面三次式と仮定する。
さらに $B/(u)$ の双線型関係の零スキームが、滑らかな平面三次曲線
$E\subset\PP(\langle x,y,z\rangle^*)$ 上の固定点を持たない自己同型
$\sigma$ のグラフであると仮定する。
$w\in Z(B)_3$ の境界への像は、この $E$ の全点加群を消すものとし、
$A=B/(w)$ と置く。
これらは、適切な開にある楕円型一ブロック中心商の仮定である。
古典曲線 $V(g)$ と量子境界 $E$ は区別する。

一点の中心拡大における楕円曲線と八点、および点対応からの
二次関係の復元は Le Bruyn--Smith--Van den Bergh
\cite[Theorem~4.2.2, Proposition~4.2.5 and \S4.3]{LBSV} に先行する。
三様体稿 \cite[\S23]{Threefold} の二中心変数の議論を、ここでは
$r$ 個の中心変数へ拡張する。定常成分は $n>3$ で正次元となる。

\begin{theorem}[点関手の分解]\label{thm:point-nfold}
$w$ の tensor 代表の可換化を $w_{\mathrm{ab}}$ とし、
\begin{equation}\label{eq:point-CD}
 C=V(F_x,F_y,F_z),\qquad
 D=V(F_x,F_y,F_z,w_{\mathrm{ab}})\subset\PP(V^*)
\end{equation}
と置く。$D$ は代表の選択に依存しない。
ここで点加群は、基底 $S$ 上で次数零から生成され、各次数成分が可逆層であり、
$M_0\simeq\OO_S$ の同定を備えるものとする。同型はこの同定を保つ。
この正規化された点加群の関手は
\begin{equation}\label{eq:point-disjoint}
 \operatorname{Pt}(B)=E\amalg C,\qquad
 \operatorname{Pt}(A)=E\amalg D
\end{equation}
で表される。shift は $E$ 上で $\sigma$、定常成分上で恒等である。
$C$ は型 $(2,2,2)$、次元 $n-2$、次数 $8$ の完全交叉であり、
中心座標への射
\begin{equation}\label{eq:point-critical-cover}
 f:C\longrightarrow\PP^{r-1}_{u_1:\cdots:u_r}
\end{equation}
は次数 $8$ の有限平坦射である。
$w_{\mathrm{ab}}$ が $C$ の各既約成分上で非零ならば、
$D$ は型 $(2,2,2,3)$ の完全交叉であり、
\begin{equation}\label{eq:point-invariants}
 \dim D=n-3,\qquad \deg D=24,\qquad
 \omega_D\simeq\OO_D(7-n).
\end{equation}
\end{theorem}

\begin{proof}
二次関係の双線型化の零スキームを $\Gamma_B$ とし、その二座標を $p,q$ と書く。
中心関係は全ての一次生成元 $v_i$ に対し
\[
 u_\alpha(p)v_i(q)-v_i(p)u_\alpha(q)=0
\]
を与える。$u_\alpha(p)$ が可逆な chart では、第二因子の一座標を
可逆に取ってこの式を使えば $u_\alpha(q)$ も可逆である。
両者を一に規格化すると $p=q$ となり、残りの関係は
$F_x=F_y=F_z=0$ である。この計算は非被約な基底上でも成り立つ。

一方、$p,q$ の二座標の小行列式が可逆な chart では、中心関係の
二本から全ての $u_\alpha(p),u_\alpha(q)$ が零となる。
従ってこの部分は $\Gamma_\sigma$ に等しい。
幾何点で中心座標が全て零なら境界の対応上にあり、$\sigma$ が
固定点を持たないため、この小行列式 chart がその全体を覆う。
$C$ と平面 $V(u)$ は交わらない。交点があれば
$g_x=g_y=g_z=0$ となり、Euler の公式と $V(g)$ の滑らかさに反する。
従って二種類の chart は全体を覆い、scheme として
\[
 \Gamma_B=\Gamma_\sigma\amalg\Delta_C
\]
を得る。これは自己同型のグラフなので、二次関係を満たす無限点列は
初期点から一意に定まる。

中心三次関係を加えると、$E$ 上では仮定によって三重線型化が零である。
滑らかな $E$ 上の切断として零なので、任意の基底変更でも零である。
定常列上では三重線型化は $w_{\mathrm{ab}}(p)$ となる。
二つの tensor 代表の差の可換化は $(F_x,F_y,F_z)$ に属するので、
$D$ は代表に依存しない。これで点関手の主張を得た。

有限平坦性を環上で証明する。$S=k[u_1,\ldots,u_r]$ と置き、
\[
 T=S[x,y,z]/(F_x,F_y,F_z)
\]
を active 次数で濾過する。三つの最高記号 $g_x,g_y,g_z$ は
$k[x,y,z]$ の正則列である。従って、濾過 Koszul 複体の最高記号からの
帰納法により
\[
 \operatorname{gr}T\simeq
 S\otimes_k k[x,y,z]/(g_x,g_y,g_z)
\]
となる。具体的には、最高次数の関係を Koszul syzygy で消し、
active 次数を下げれば、初期イデアルに新しい生成元が生じないことが分かる。
右辺の Jacobian 環の Hilbert 級数は $(1+t)^3$ である。
その斉次基底を持ち上げると、$T$ は $S$ 上自由であり、全次数付き加群として
\begin{equation}\label{eq:critical-free-module}
 T\simeq S\oplus S(-1)^3\oplus S(-2)^3\oplus S(-3)
\end{equation}
を得る。この分解自体の標準性は主張しない。
$C\cap V(u)=\varnothing$ と合わせて Proj を取れば、$f$ は有限平坦で
\begin{equation}\label{eq:critical-pushforward}
 f_*\OO_C\simeq\OO\oplus\OO(-1)^3\oplus\OO(-2)^3\oplus\OO(-3)
\end{equation}
となる。

環 $T$ の次元は $r$ なので、三つの二次式の高さは三であり、
多項式環の Cohen--Macaulay 性から正則列となる。
従って $C$ の次元と次数が従う。
完全交叉は埋込み素点を持たないので、各既約成分で非零な
$w_{\mathrm{ab}}$ は非零因子である。
正則列に三次式を加えて次元と次数を得る。
随伴公式により
$\omega_D=\OO_D(2+2+2+3-(n+2))=\OO_D(7-n)$ である。
\end{proof}

\begin{corollary}[中心射影と二十四次の行列式]\label{cor:point-norm}
定理の完全交叉条件の下で、$\mathcal V=f_*\OO_C$ と置くと
\begin{equation}\label{eq:point-norm-sequence}
 0\longrightarrow\mathcal V(-3)
 \xrightarrow{\cdot w_{\mathrm{ab}}}\mathcal V
 \longrightarrow f_*\OO_D\longrightarrow0
\end{equation}
は完全である。この乗法行列の行列式は非零な
$\OO_{\PP^{r-1}}(24)$ の切断であり、$f_*\OO_D$ の
零次 Fitting イデアルを定める。
$n=4$ で $D$ が滑らかならば、$D$ は次数 $24$、種数 $37$ の連結曲線である。
\end{corollary}

\begin{proof}
中心座標が基点を持たないので $\OO_C(1)=f^*\OO(1)$ である。
$f$ は有限、$w_{\mathrm{ab}}$ は非零因子だから、中心三次切断の完全列を
押し出して表示した完全列を得る。階数 $8$ より行列式の次数は $3\cdot8=24$。
一般点で乗法は同型なので行列式は非零である。
これは scheme 論的像の被約方程式や、$D$ の射影先への双有理性を主張するものではない。
正次元の射影完全交叉では Koszul 解消と射影空間の線束コホモロジーから
$H^0(D,\OO_D)=k$ であり、従って連結である。
$n=4$ では $\omega_D=\OO_D(3)$ なので
$2g(D)-2=3\cdot24=72$ となる。
\end{proof}

\begin{remark}[滑らかな定常成分の非空性]\label{rem:point-nonempty}
三つの一般の二次式 $Q_i(u)$ に対して
\[
 F=(x^3+y^3+z^3)/3-xQ_1(u)-yQ_2(u)-zQ_3(u)+G(u)
\]
とすると、$C$ は $x^2=Q_1,y^2=Q_2,z^2=Q_3$ で与えられる。
$\PP^{r-1}$ 上の三つの分岐二次超曲面を滑らかで交わりが横断的になるよう選べば、
その二重被覆のファイバー積は滑らかである。
局所的には分岐因子の方程式を座標の一部に取り、その平方根を新座標にすればよい。
$r=2$ では分岐点を全て異なるものに取る。
純中心三次式を動かす線型系は $C$ 上基点を持たないため、
Bertini により一般の中心三次切断 $D$ は滑らかである。
これは $n=3$ では被約な二十四点を意味する。
境界の滑らかさ・固定点なしの条件と、後の Hessian 条件も、係数空間の
非空開条件として共通に課せる。ここで示すのは該当する開の非空性であり、
全ての分離型や全ての中心三次切断が滑らかだという主張ではない。
\end{remark}

\begin{theorem}[点対応からの包絡復元]\label{thm:point-envelope}
定理\ref{thm:point-nfold}の完全交叉条件を仮定し、
$R_B\subset V\otimes V$ を二次関係空間とする。
全点スキーム上の shift のグラフ
$\Gamma_A=\Gamma_\sigma\amalg\Delta_D$ について
\begin{equation}\label{eq:point-envelope}
 H^0(\mathcal I_{\Gamma_A}(1,1))=R_B
\end{equation}
が成立する。従って、埋込み点スキームと shift から二次包絡 $B$ と、
式\eqref{eq:point-relations}の尺度込みの双ベクトル $h\Pi_F$ を復元できる。
\end{theorem}

\begin{proof}
完全交叉のイデアルは飽和しており、三次生成元は二次部分に寄与しないから
$Q=I(D)_2=\langle F_x,F_y,F_z\rangle$ である。
従って $\Delta_D$ 上で消える双線型形式は
$\bigwedge^2V\oplus\operatorname{sym}(Q)$ で、次元は $\binom N2+3$ である。
$U=\langle u_1,\ldots,u_r\rangle$ を含む反対称部分は
次元 $\binom N2-3$ であり、境界上で零となる。
これで割った六次元の候補空間は
\[
 \bigwedge^2(V/U)\oplus
 \operatorname{sym}\langle g_x,g_y,g_z\rangle
 \subset(V/U)^{\otimes2}
\]
に単射する。

$L=\OO_E(1)$、$M=\sigma^*L$ とすると、乗法
$H^0(L)\otimes H^0(M)\to H^0(LM)$ は全射である。
$L\not\simeq M$ の場合は $L$ の基点のない pencil の完全列と
$H^1(ML^{-1})=0$、$L\simeq M$ の場合は平面三次に二次関係がないことから従う。
従って $\Gamma_\sigma$ の双線型イデアルは三次元であり、これは境界代数の
三つの関係で生成される。候補空間に含まれるこの三次元核を戻せば
全体の核の次元は $\binom N2$ である。
既に $R_B$ を含み、$R_B$ は反対称成分への射影が同型なので、両者は等しい。
関係を $v_i\otimes v_j-v_j\otimes v_i-\operatorname{sym}(q_{ij})$ と一意に書けば
$\sum q_{ij}\partial_{v_i}\wedge\partial_{v_j}$ が復元される。
本模型では定義から $h\Pi_F$ に等しい。
\end{proof}

\begin{proposition}[古典 Casimir 空間]\label{prop:point-casimir}
上の完全交叉条件に加えて、$g$ の Hessian に対する定数行列の条件
\begin{equation}\label{eq:point-hessian-scalar}
 M\operatorname{Hess}(g)\text{ が対称}\quad\Longrightarrow\quad M=cI_3
\end{equation}
を仮定する。このとき埋込み点スキームから
\[
 U=I(E)_1,\qquad Q=I(D)_2,
\]
および
\begin{equation}\label{eq:point-casimir}
 \mathcal R_Q=\{H\in\Sym^3V:\partial_aH\in Q\text{ for all }a\in\operatorname{ann}(U)\}
 =kF\oplus\Sym^3U
\end{equation}
を復元できる。従って $\dim\mathcal R_Q=1+\binom{n+1}{3}$ であり、
古典 Jacobian Poisson 直線 $[\Pi_F]$ が定まる。
\end{proposition}

\begin{proof}
$E$ はその境界平面を張る。
$n=3$ では正次元成分、$n>4$ では一次元成分として $E$ を識別できる。
$n=4$ では $D$ は連結な次数 $24$ の完全交叉曲線なので、
次数 $3$ の連結成分 $E$ と区別できる。
よって $U$ と $Q$ は埋込みデータから得られる。
$H\in\mathcal R_Q$ に対して
$\nabla_{x,y,z}H=M\nabla_{x,y,z}F$ となる定数行列 $M$ がある。
偏微分の可換性と $u=0$ への制限から式\eqref{eq:point-hessian-scalar}が適用でき、
$M=cI$ である。従って $H-cF$ は中心変数だけの三次式となる。
逆包含は直接従う。
$F$ を $cF+G(u)$ に替えることは Jacobian 括弧を $c$ 倍するだけなので、
中心空間と表示した三次空間は $[\Pi_F]$ を定める。

さらにこの空間は affine Poisson 括弧の三次 Casimir 全体でもある。
Casimir $H$ の active 勾配は分数体上で $a/b$ 倍の $\nabla F$ となる。
$F_x,F_y,F_z$ は正則列で共通因子を持たないので、既約表示の分母 $b$ は単元である。
斉次次数を比較して $a/b$ は定数となり、同じ結論を得る。
この最後の Casimir 計算自体には式\eqref{eq:point-hessian-scalar}を必要としない。
\end{proof}

\begin{remark}[明示的な非空条件と限界]
$g=(x^3+y^3+z^3)/3+\lambda xyz$ では
$\lambda(\lambda^3-8)\ne0$ のとき式\eqref{eq:point-hessian-scalar}が成り立つ。
実際、$M\operatorname{Hess}(g)$ の対称性は対角成分を等しくし、非対角成分には
$\lambda m_{13}=2m_{21}$、$\lambda m_{21}=2m_{32}$、$\lambda m_{32}=2m_{13}$
ともう一組の循環式を課すので、全非対角成分が零になる。
滑らかさの条件 $\lambda^3\ne-1$ と合わせても非空である。
この条件による勾配復元の原理と、本稿での任意 $r$ の点関手への適用を区別する。
また $\mathcal R_Q/\Sym^3U$ の直線を得ることは、その中の特定の三次式を
選ぶことではない。量子中心から残余圏が定める可換三次式への補正は別問題であり、
点スキームを有限係数の Poisson 零スキームと同一視してはいない。
\end{remark}

% !TEX root = ../nc-cubic-n-fold.tex
\section{相対点関手と構造対象付きモジュライへの降下}
\label{sec:relative-points}

\subsection{相対点関手と任意の基底変更}
\label{subsec:relative-points}

点スキームを族にする際には、中心三次の生成元ではなくその直線を用いる。
また、各ファイバーでの完全交叉性と、非平坦な基底変更後にも保たれる
相対完全交叉性を区別する。
以下の $S$ は良い一ブロック表示チャートの任意の基底変更とし、
$\mathcal V$ は階数 $N=n+2$ の生成元束、
$\mathcal U\subset\mathcal V$ は階数 $r=n-1$ の中心部分束、
$\mathcal X=\mathcal V/\mathcal U$ とする。
$\mathbb P(\mathcal V^*)$ は $\mathcal V$ の可逆商を表す規約とし、
その標準全射を $\mathcal V\otimes\OO\to\OO(1)$ と書く。

局所的に $\mathcal V=\langle x,y,z,u_1,\ldots,u_r\rangle$ と表示し、
式\eqref{eq:point-relations}の関係を用いる。
$g=F|_{u=0}$ は各幾何学的ファイバーで滑らかな平面三次式とする。
境界の双線型関係スキームは、相対的に滑らかな平面三次曲線
$\mathcal E\to S$ の固定点を持たない $S$-自己同型
$\sigma$ のグラフであると仮定する。
中心三次直線を $\ell\subset Z(\mathcal B)_3$ とし、
その境界への像は $\mathcal E$ 上の普遍点加群を消すものとする。
この最後の条件は切断の恒等式であり、幾何点での消滅だけを仮定するのではない。
$\mathcal A=\mathcal B/(\ell)$ と置く。

三つの偏微分が張る階数三の部分束を
$\mathcal Q\subset\operatorname{Sym}^2\mathcal V$ とする。
中心三次の可換化は、代表に依存しない写像
\[
 \ell\longrightarrow
 \bigl(\operatorname{Sym}\mathcal V/(\mathcal Q)\bigr)_3
\]
を与える。
これにより
\[
 \mathcal C=V(\mathcal Q),\qquad
 \mathcal D=V(\mathcal Q,\ell_{\mathrm{ab}})
 \subset\mathbb P(\mathcal V^*)
\]
を定める。各幾何学的ファイバーの $\mathcal C$ 上で
$\ell_{\mathrm{ab}}$ が非零因子を生成する開を、点スキームの良い開と呼ぶ。
これは以下の行列式によって検出される開条件である。

\begin{theorem}[普遍的な相対点スキーム]
\label{thm:relative-point-nfold}
以上の点スキームの良い開上では、任意の $S$-スキーム $S'$ に対して
正規化された相対点加群の関手が
\[
 \operatorname{Pt}_{S'}(\mathcal B_{S'})
   =\mathcal E_{S'}\amalg\mathcal C_{S'},\qquad
 \operatorname{Pt}_{S'}(\mathcal A_{S'})
   =\mathcal E_{S'}\amalg\mathcal D_{S'}
\]
で表される。ここで各次数成分は $S'$ 上の可逆層で、次数零を
$\OO_{S'}$ と同定し、次数零から生成されることを要求する。
shift は $\mathcal E$ 上で $\sigma$、$\mathcal C,\mathcal D$ 上で恒等である。

中心射影
\[
 f:\mathcal C\longrightarrow\mathbb P(\mathcal U^*)
\]
は次数八の有限局所自由射である。
$\mathcal C\to S$ と $\mathcal D\to S$ は平坦であり、
それぞれ型 $(2,2,2)$ と $(2,2,2,3)$ の相対完全交叉である。
これらの主張、正則埋込み、および点関手の同定は、任意の基底変更と両立する。
各ファイバーの次元と次数は、それぞれ $(n-2,8)$ と $(n-3,24)$ である。

$p_D:\mathcal D\to S$ とすると、相対双対化層は
\begin{equation}\label{eq:relative-point-dualizing}
 \omega_{\mathcal D/S}\simeq
 \OO_{\mathcal D}(7-n)\otimes
 p_D^*\bigl(\det\mathcal V\otimes
              (\det\mathcal Q)^{-1}\otimes\ell^{-1}\bigr).
\end{equation}
従って基底上の determinant 因子を自明化していない族では、
右辺を単に $\OO_{\mathcal D}(7-n)$ と書かない。
\end{theorem}

\begin{proof}
まず有限局所自由性を係数環上で確認する。
局所的に $S=\operatorname{Spec}R$ として
\[
 K=R[x,y,z]/(g_x,g_y,g_z),\qquad K_i=K_i^{\mathrm{gr}}
\]
と置く。三つの偏微分の Koszul 複体の各次数部分は有限自由複体である。
各幾何学的ファイバーでは正次数のホモロジーが零であり、
$K_i$ の階数は $i=0,1,2,3$ に対して $1,3,3,1$、それ以外では零である。
有限自由複体で可逆な小行列式を順に分離すると、
各次数の複体は局所的に可縮な複体と次数零に置いた自由加群の直和になる。
従って $K_i$ は表示した階数の有限局所自由加群であり、
この Koszul 完全性と $K_i$ の形成は任意の係数変更と両立する。
これは被約基底上の階数計算だけを非被約基底へ移す議論ではない。

$R_u=R[u_1,\ldots,u_r]$ および
\[
 T=R_u[x,y,z]/(F_x,F_y,F_z)
\]
を考える。active 次数の最高記号は $g_x,g_y,g_z$ である。
最高記号の Koszul 関係で最高次数を消す帰納法により
\[
 \operatorname{gr}_{\mathrm{act}}T\simeq R_u\otimes_R K
\]
を得る。この帰納法は全ての係数変更後にも同じである。
従って $T$ は長さ四の有限濾過を持ち、その次数付き随伴商は
$R_u(-i)\otimes_R K_i$ $(0\leq i\leq3)$ である。
各 $K_i$ が局所自由なので、局所的に斉次持上げを選べば
\begin{equation}\label{eq:relative-critical-free}
 T\simeq R_u\oplus R_u(-1)^3\oplus
                 R_u(-2)^3\oplus R_u(-3)
\end{equation}
となる。この分解の標準性は主張しない。
とくに $T$ は $R_u$ 上階数八で有限局所自由である。

滑らかな $g$ の偏微分は射影平面上に共通零点を持たないから、
$\mathcal C$ は $V(\mathcal U)$ と交わらない。
よって中心射影は全域で定義され、\eqref{eq:relative-critical-free} の Proj により
有限局所自由性を得る。
$\mathcal C$ の三つの方程式についても、上の Koszul 複体は完全であり、
その商は $S$ 上平坦である。
射影空間は $S$ 上平坦かつ有限表示なので、これは相対正則埋込みであり、
任意の基底変更後にも正則である
\cite[Tags~063P, 063U]{NFMStacksRelativeRegular}。

三次直線の像は各ファイバー上の非零因子である。
平坦かつ有限表示な $\mathcal C\to S$ 上の局所主イデアルに対する
ファイバー判定により、$\mathcal D\subset\mathcal C$ は相対有効 Cartier 因子である
\cite[Tag~062Y]{NFMStacksRelativeCartier}。
従って $\mathcal D\to S$ は平坦で、この Cartier 性は
非平坦なものを含む全ての基底変更で保たれる
\cite[Tag~056Q]{NFMStacksRelativeCartier}。
完全交叉型とファイバーの次元・次数が従う。

点関手については、定理\ref{thm:point-nfold}の二種類の chart の議論を
そのまま係数環上で用いる。
中心座標の一つが可逆なら中心関係から二つの点は等しく、
残りの式は $\mathcal Q=0$ となる。
二点の二次小行列式の一つが可逆なら、全中心座標が零となって
境界の $\sigma$ のグラフを得る。
固定点なしの仮定と $\mathcal C\cap V(\mathcal U)=\varnothing$ により
これらは全体を覆う。従って scheme として
\[
 \Gamma_{\mathcal B}=\Gamma_\sigma\amalg\Delta_{\mathcal C}
\]
である。この対応は自己同型のグラフなので、無限点列は初期点から一意に決まる。
三次直線の多重線型化は境界上で恒等的に零であり、
定常列上では $\ell_{\mathrm{ab}}$ の評価となる。
これで任意の基底上の点関手と shift の主張を得る。

最後に、射影空間の相対標準束は
$\OO(-N)\otimes\det\mathcal V$ である。
二次三本と三次直線の conormal determinant は
$\OO(-9)\otimes\det\mathcal Q\otimes\ell$ なので、
相対随伴公式から\eqref{eq:relative-point-dualizing}を得る。
\end{proof}

\begin{corollary}[相対 norm と良い点スキーム開]
\label{cor:relative-point-norm}
$p_U:\mathbb P(\mathcal U^*)\to S$、
$\mathcal H=f_*\OO_{\mathcal C}$ と置く。
点スキームの良い開上には、任意の基底変更と両立する完全列
\[
 0\longrightarrow p_U^*\ell\otimes\mathcal H(-3)
 \longrightarrow\mathcal H
 \longrightarrow f_*\OO_{\mathcal D}\longrightarrow0
\]
がある。最初の写像の行列式は
\[
 \mathcal N\in
 H^0\bigl(\mathbb P(\mathcal U^*),
       \OO(24)\otimes p_U^*\ell^{-8}\bigr)
\]
であり、$f_*\OO_{\mathcal D}$ の零次 Fitting イデアルを定める。
各ファイバーで $\mathcal N$ が零切断でないという条件は基底の開条件であり、
上の非零因子条件と同値である。
良い一ブロック表示族ではこの開は非空である。
\end{corollary}

\begin{proof}
中心射影から $\OO_{\mathcal C}(1)=f^*\OO(1)$ である。
相対 Cartier 因子の完全列を有限射で押し出せば表示した完全列を得る。
$\mathcal H$ の階数は八なので、その determinant は
$\OO(24)\otimes p_U^*\ell^{-8}$ に値を持つ。
各ファイバーでは $T$ は多項式環 $k[u]$ 上自由であり、
その上の三次乗法の単射性は行列式が非零多項式であることと同値である。
$T$ は $k[u]$ 上自由なので、三次乗法の核に射影化で消える
非零部分加群はない。従ってこれは $\mathcal C$ 上の非零因子条件とも同値である。
行列式の全係数が同時に零になる条件は閉条件なので、所定の条件は開である。

非空性のため、一つの良い周囲代数とその中心三次の持上げを固定する。
純中心三次を加える線型系は $\mathcal C$ 上基点を持たない。
有限個の既約成分のいずれかで恒等的に零となる条件はそれぞれ真の affine
線型条件であるから、その和集合を避けることができる。
これは一ブロック表示成分の非空開についての主張である。
\end{proof}

\begin{proposition}[相対点対応からの二次包絡]
\label{prop:relative-point-envelope}
$\rho=\sigma\amalg\operatorname{id}_{\mathcal D}$ とし、
$\Gamma=\Gamma_\rho\subset
\mathbb P(\mathcal V^*)\times_S\mathbb P(\mathcal V^*)$ と置く。
$p$ をこの積から $S$ への射とすると
\begin{equation}\label{eq:relative-point-envelope}
 p_*\mathcal I_\Gamma(1,1)=\mathcal R_{\mathcal B}
 \subset\mathcal V\otimes\mathcal V
\end{equation}
であり、左辺の形成は任意の基底変更と両立する。
また、埋込みイデアルの低次数部分は任意の基底変更と両立して
\begin{equation}\label{eq:relative-point-ideal-pieces}
 p_*\mathcal I_{\mathcal D}(2)=\mathcal Q,
 \qquad
 p_*\mathcal I_{\mathcal C}(3)=\mathcal Q\otimes\mathcal V,
 \qquad
 p_*\mathcal I_{\mathcal D}(3)\big/
       p_*\mathcal I_{\mathcal C}(3)\simeq\ell
\end{equation}
となる。ここで後二式の $p$ は射影空間から $S$ への射を表す。
\end{proposition}

\begin{proof}
型 $(2,2,2,3)$ の Koszul 解消に $\OO(2)$ または $\OO(3)$ を tensor する。
射影空間の相対次元は $n+1\geq4$、イデアル層のこの解消の長さは三なので、
射影空間の最高次コホモロジーは全次数零の計算に寄与しない。
中間コホモロジーも零であり、残る次数零の列から
\eqref{eq:relative-point-ideal-pieces}を得る。
二次生成元の一次係数による syzygy はないので、
$\mathcal Q\otimes\mathcal V\to\operatorname{Sym}^3\mathcal V$ は単射である。
三次生成元の tensor 代表を替えた差はその像に入るため、最後の同型は代表に依存しない。
この相対 Koszul 計算は全ての基底変更後にも同じである。
とくに $n=3$ で $\mathcal D$ が零次元でも成立するが、
$\operatorname{Sym}^3\mathcal V\to p_*\OO_{\mathcal D}(3)$ の全射性は主張しない。

従って対角成分上で消える双線型形式の束は
$\bigwedge^2\mathcal V\oplus\operatorname{sym}(\mathcal Q)$ である。
その中の $\mathcal U\wedge\mathcal V$ を割ると、境界での候補は
\[
 \bigwedge^2\mathcal X\oplus
 \operatorname{sym}\langle g_x,g_y,g_z\rangle
 \subset\mathcal X^{\otimes2}
\]
となる。境界の $L=\OO_{\mathcal E}(1)$ に対して
$L\otimes\sigma^*L$ は相対次数六であるから、その切断束は階数六で、
任意の基底変更と両立する。
乗法 $p_*L\otimes p_*\sigma^*L\to p_*(L\otimes\sigma^*L)$ は
各幾何学的ファイバーで全射であることを
定理\ref{thm:point-envelope}の pencil の議論で確認できる。
従って相対的にも分裂全射であり、三次元の核は境界の二次関係束そのものである。
これを上の候補束に引き戻すと、核は $\mathcal R_{\mathcal B}$ に等しい。
全てが部分束と分裂全射で計算されるので、この等式も任意の基底変更で保たれる。
\end{proof}

\subsection{$\mathcal O$-marked モジュライへの点データの降下}
\label{subsec:point-descent}

上の条件を満たす $P$-不変な非空開を $J^{\mathrm{pt}}\subset J^{\mathrm{good}}$ とし、
\[
 \mathfrak M_n^{\mathrm{pt}}
  =[J^{\mathrm{pt}}/\overline P]
  \subset\mathfrak M_n^{\mathrm{loc}}
\]
と置く。最後の同定には定理\ref{thm:nfm-local-moduli}を用いる。
これは良い一ブロック表示に fpqc 局所的に同値な圏族の開であり、
全ての内在的成分の一般点を覆うとは主張しない。

\begin{theorem}[点関手、shift、中心射影の $\mathcal O$-marked 降下]
\label{thm:point-omarked-descent}
$\mathfrak M_n^{\mathrm{pt}}$ 上には Brauer--Severi 束
\[
 \mathfrak P_V=[\mathbb P(\mathcal V^*)/\overline P],
 \qquad
 \mathfrak P_U=[\mathbb P(\mathcal U^*)/\overline P]
\]
と、それぞれに付随するデータ
\[
 \mathfrak Y=\mathfrak E\amalg\mathfrak D\hookrightarrow\mathfrak P_V,
 \qquad \rho:\mathfrak Y\xrightarrow{\sim}\mathfrak Y,
 \qquad f:\mathfrak C\longrightarrow\mathfrak P_U
\]
がある。$f$ は次数八で有限局所自由であり、
$\mathfrak C\hookrightarrow\mathfrak P_V$ と
$\mathfrak D\hookrightarrow\mathfrak P_V$ はそれぞれ fpqc 局所的に
型 $(2,2,2)$ と $(2,2,2,3)$ の相対完全交叉である。
これらは全ての $\mathcal O$-保存同値に関して自然であり、基底変更と両立する。

さらに $\mathfrak Y$ は次の標識を増やさない点関手を表す。
対象 $(\mathscr C,\mathcal O)$ に対し、基底 $T$ 上の点とは全射
\[
 q:\mathcal O_T\twoheadrightarrow\mathcal P
\]
であって、fpqc 局所的に良い表示
$(\mathscr C_T,\mathcal O_T)\simeq
 (\operatorname{QGr}\mathcal A_T,\pi\mathcal A_T)$ の下で、
正規化された点加群 $\mathcal A_T\twoheadrightarrow M$ の tails に同型となるものとする。
射は $q$ を保つ同型とする。
この条件は良い表示の選択によらず、当該 groupoid は自己同型を持たない。
従って $\mathfrak Y$ 上にはこの意味での普遍点商が存在する。
次数一 shift の大域的な選択を点関手の定義に加える必要はない。
\end{theorem}

\begin{proof}
次数付き同型は生成元の射影空間、全関係、点加群、およびその正規化 shift を移す。
次数スカラー $\lambda$ は射影空間上で恒等であり、点加群の次数 $j$ には
$\lambda^j$ を掛ける標準同型を与える。
次数零の同定は保たれるので、この作用は正規化点加群の同型類上で自明である。
従って幾何データは $P$ から $\overline P$ へ降下する。
定理\ref{thm:nfm-fixed-model-comparison}により、全ての
$\mathcal O$-保存同値は局所的に次数付き同型から来て、
二つの持上げは次数スカラーだけ異なる。
よってこの降下は表示同型だけでなく全ての $\mathcal O$-保存同値に関して自然である。
有限局所自由性、平坦性、および相対正則埋込みも fpqc 降下する。

点関手と tails の関係を確認する。
点加群は各次数から上へ生成されるので非零 torsion 部分加群を持たない。
二つの点加群の tails が同型なら、その同型と逆は十分大きい次数で互いに逆な
次数付き写像により表せる。これは点加群が次数零の一つの可逆層で生成され、
各打切り $M_{\geq m}$ も一つの有限局所自由次数成分から生成されることによる。
従って十分大きい次数以後の点列が等しい。
点対応 $\rho$ が自己同型なので、その初期点も等しい。
相対点関手の表示から、元の正規化点加群は同型である。
また $q$ は全射なので、$q$ を保つ自己同型は恒等だけである。
これにより点加群と上述の局所点商の同型類を同定できる。

全 $\mathcal O$-同値の局所持上げによって、局所点商という条件は
良い表示から独立である。
圏族と対象の有効 fpqc 降下により、この点商は貼り合わさる。
正規化 shift は局所的に
\[
 M\longmapsto M(1)_{\geq0}\otimes M_1^{-1}
\]
であり、初期点を $\rho$ で移す。
次数スカラーに関してこの式の次数 $j$ の重みは再び $j$ になるので、
これも点商の関手とそのグラフとして降下する。
ここでは圏全体の次数一 shift を大域に選んだわけではない。
\end{proof}

\begin{proposition}[偏極、双対化層、包絡の正確な降下形]
\label{prop:point-polarization-descent}
命題\ref{prop:nfm-presentation-gerbe}の表示 gerbe
\[
 \mathfrak G=[J^{\mathrm{pt}}/P]
   \longrightarrow\mathfrak M_n^{\mathrm{pt}}
\]
の類を $\beta$ とし、$e=\gcd(r,3)$ とする。
整数 $a,b$ を $ar+3b=e$ となるように一組固定し、
$\chi=(\det\mathcal U)^a\otimes(\det\mathcal X)^b$ と置く。
gerbe 上の線束
\begin{equation}\label{eq:point-descended-polarization}
 \OO_{\mathbb P(\mathcal V^*)}(e)\otimes p_V^*\chi^{-1},
 \qquad
 \OO_{\mathbb P(\mathcal U^*)}(e)\otimes p_U^*\chi^{-1}
\end{equation}
はそれぞれ $\mathfrak P_V,\mathfrak P_U$ 上の相対 ample 線束に降下する。
中心射影によるそれらの引戻しは $\mathfrak C$ 上で一致する。
別の整数解を用いた降下線束との差は、モジュライ基底上の線束の引戻しである。

特に ambient Brauer--Severi 束上で、
$n\not\equiv1\pmod3$ なら幾何学的ファイバーで
$\OO(1)$ となる偏極を選べる。
$n\equiv1\pmod3$ では幾何学的に $\OO(3)$ となる偏極は常に選べるが、
$\OO(1)$ の存在は $\beta$ の消滅と同値である。
素の標準 $P$-線型化を持つ $\OO(1)$ は、いずれの場合にも重み一なので
そのまま $\overline P$ へ降りるわけではない。

一方、$\omega_{\mathcal D/S}$、$f_*\OO_{\mathcal C}$、
相対 norm の切断とその Fitting 因子、および
\eqref{eq:relative-point-envelope}が定める二次包絡の
次数スカラーを除いた同型類は、追加の偏極を選ばずに降下する。
通常の生成元束 $\mathcal V$ や二次関係束 $\mathcal R_{\mathcal B}$ 自体の
降下を主張するものではない。
\end{proposition}

\begin{proof}
gerbe の中心スカラーの重みは、
$\mathcal V,\mathcal U,\mathcal X$ 上で一、
$\OO(1)$ 上で一、$\mathcal Q$ 上で二、$\ell$ 上で三である。
$\chi$ の重みは $e$ なので、\eqref{eq:point-descended-polarization} の重みは零となる。
重み零の準連接層の fpqc 降下を用いれば表示した線束を得る。
その相対 ample 性は局所的な $\OO(e)$ の ample 性から従う。
二つの $\chi$ の比は重み零の基底線束なので、選択の差についての主張も従う。
$\OO(1)$ 型線束の存在と Brauer 類の消滅の同値、および
$r\beta=3\beta=0$ は命題\ref{prop:nfm-presentation-gerbe}による。

式\eqref{eq:relative-point-dualizing}の基底因子の重みは $N-9=n-7$ であり、
$\OO(7-n)$ の重みと相殺する。
norm の値を取る線束の重みも $24-8\cdot3=0$ である。
$\mathcal C$ とその構造層には次数スカラーが自明に作用するので、
$f_*\OO_{\mathcal C}$ は通常の有限局所自由層として降下する。

包絡の関係束 $\mathcal R_{\mathcal B}$ は重み二である。
従って一般には、これをモジュライ基底上の通常の部分束と記すことはできない。
しかし埋込みグラフのイデアルそのものは降下し、各表示上での
\eqref{eq:relative-point-envelope}による復元は射影的座標変更と両立する。
よって二次包絡は次数スカラーを除いた同型問題として復元される。
反対称部分への射影を用いて関係から復元する
$\bigwedge^2\mathcal V\to\operatorname{Sym}^2\mathcal V$ は重み零であるため、
対応する尺度込みの相対 Poisson 双ベクトルも ambient Brauer--Severi 束上へ降下する。
この結論だけから中心三次に対応する古典 Casimir を選択したことにはならない。
\end{proof}

\begin{remark}[普遍点加群と普遍点商]
普遍点加群の次数 $j$ の可逆層は gerbe 上で重み $j$ を持つ。
したがって、普遍点商が $\mathcal O$-marked な圏族に降下するという主張と、
全次数成分が通常の線束として降下するという主張は異なる。
前者は上で証明したが、素の線型化を持つ次数成分は通常の線束として降下しない。
モジュライ基底上に通常の次数付き代数族を置き、その点加群として記述するには、
表示 gerbe の中立化を選ぶ必要がある。
これは点スキーム自身への引戻しで障害が消えることとは矛盾しない。
実際、点スキームへの引戻し gerbe 上では普遍点加群の $M_1$ が重み一の線束であり、
これによってその gerbe は中立化される。
そこで次数 $j$ ごとに $\mathcal A_j$ と $M_j$ の双方を
$M_1^{-j}$ で捻れば、点空間上では通常の次数付き代数とその普遍点加群を得る。
またここでの $\beta\in H^2(\mathbf G_m)$ は対象の表示障害であり、
命題\ref{prop:nfm-lifting-line}の一つの同値の持上げを妨げる
$\operatorname{Pic}$ 上の線束とは別の障害である。
\end{remark}

% !TEX root = ../nc-cubic-n-fold.tex
\section{点対応による中心復元と Poisson 三次対への射}
\label{sec:poisson-morphism}

\subsection{中心三次の評価と境界 Jacobi 線}

この節では一ブロックの良い点スキーム開をさらに縮める。
全ての主張は $n\ge3$、$r=n-1\ge2$ に一様である。
次数三中心の束を $\mathcal E$、純中心三次の部分束を
$\mathcal W=\Sym^3U$、商線束を $\mathcal L=\mathcal E/\mathcal W$ とする。
固定ファイバーでは
\[
 S=k[U],\qquad T=\Sym(V)/(F_x,F_y,F_z),\qquad
 J_g=T/(U)=k[x,y,z]/(g_x,g_y,g_z)
\]
と書く。$T$ は $B$ の最大可換\emph{代数}商である。
これは交換子の張る線形部分空間による商 $\mathrm{HH}_0(B)$ とは異なる。
代数の可換化は中心上の評価写像
\[
 \alpha:\mathcal E\longrightarrow T_3
\]
を与える。tensor 代表を可換化し、三つの偏微分のイデアルで割る表示と同じ写像である。
純中心三次は $J_g$ で零となるので、合成は
\begin{equation}\label{eq:point-boundary-jacobi-map}
 \epsilon:\mathcal L=\mathcal E/\mathcal W\longrightarrow (J_g)_3
\end{equation}
を定める。右辺は階数一である。
$C\cap V(U)=\varnothing$ であるから、この商を $C$ と境界平面との
scheme 理論的交叉への制限と解釈してはならない。
これは次数付き切断環の Artin 商であり、
$T_3=H^0(C,\mathcal O_C(3))$ を用いて
\[
 (J_g)_3=T_3/(U\cdot T_2)
\]
と求めるものである。

\begin{proposition}[境界 Jacobi 線と中心評価]\label{prop:point-jacobi-evaluation}
$\epsilon$ が同型である条件は、表示変更群 $P$ に不変な非空開条件である。
この開上で $\alpha:\mathcal E\to T_3$ は余核が局所自由な単射であり、
任意の基底変更後にも単射である。
ここで分裂は局所的な存在の主張であり、標準的な大域分裂は選ばない。
\end{proposition}

\begin{proof}
$g$ の滑らかさより三つの偏微分は正則列であり、
$J_g$ の Hilbert 級数は $(1+t)^3$ である。
従って $(J_g)_3$ は線束となり、\eqref{eq:point-boundary-jacobi-map}
が同型であることは線束間の写像の非消滅条件である。
全ての写像は代数 $B$ とその中心一次空間 $U$ から定まるので $P$-共変である。

非空性を、公刊された Etingof--Ginzburg の Rains 中心公式
\cite[\S9, p.~1414]{NFMEG} で確認する。
\[
 g=(x^3+y^3+z^3)/3+\lambda xyz,\qquad 1+\lambda^3\ne0,
 \qquad \delta=1-h\lambda/2\ne0
\]
とし、$q=(1+h\lambda/2)/\delta$、$t=h/\delta$ と置く。
境界の関係は $xy-qyx-tz^2=0$ およびその巡回形となる。
同文献の式で全ての低次係数を零とした中心元は
\begin{equation}\label{eq:point-rains-homogeneous}
 \Psi=t(q+1)\bigl[
 t(t^3+1)y^3+(q^3-t^3)yzx-q(t^3+1)zyx+t(q^3-t^3)z^3
 \bigr].
\end{equation}
$J_g$ では $x^3=y^3=z^3=-\lambda xyz$ であり、
\begin{equation}\label{eq:point-rains-socle}
 [\Psi_{\mathrm{ab}}]
 =-\frac{4h^4(1+\lambda^3)}{\delta^6}[xyz]\ne0.
\end{equation}
この計算では $[xyz]\ne0$ も必要である。
例えば $x^2y=-\lambda y^2z=\lambda^2xz^2=-\lambda^3x^2y$
から、次数三の二種類の変数を含む単項式は全て零である。
立方は上の式で $xyz$ の倍数になり、$(J_g)_3$ は一次元なので $xyz$ は非零である。
係数恒等式の直接の証明書は
\[
 \Psi_{\mathrm{ab}}-\kappa xyz
 =t^2(q+1)\bigl((t^3+1)y g_y+(q^3-t^3)z g_z\bigr),\qquad
 \kappa=-4h^4(1+\lambda^3)/\delta^6
\]
である。
良い Sklyanin 開では中心三次は一次元であるから、その任意の非零生成元にも
同じ非消滅が成立する。混合中心係数を加えても境界への像は変わらない。
従来の PBW・滑らかな楕円境界・必要な有限範囲の位数除外の各非空開と交わらせられる。
全ての有限位数を除く場合には、これとは別に very general 条件として扱う。

次に $\mathcal W\to T_3$ は単射で、局所的には直和因子である。
実際、active 次数の濾過と $J_g$ の斉次基底の持上げにより
$T\simeq S\oplus S(-1)^3\oplus S(-2)^3\oplus S(-3)$ となり、
最初の項に $\mathcal W=S_3$ が含まれる。
商 $T_3/\mathcal W$ への $\mathcal L$ の写像は、さらに $(J_g)_3$ へ写したとき
同型である。従ってこれは局所的な直和因子であり、
$\mathcal E\to T_3$ も単射、余核も局所自由である。
この議論は有限局所自由な加群と単元を用いるため、任意の基底変更後にも成立する。
\end{proof}

\begin{theorem}[点対応から中心商までの復元]\label{thm:point-central-recovery}
上の開と点スキームの完全交叉条件を仮定する。
埋込み点スキームと shift のグラフから、二次包絡 $B$ に加えて
中心三次直線 $[w]\subset Z(B)_3$、従って $A=B/(w)$ を復元できる。
復元は表示変更と両立する。
相対的には定理\ref{thm:relative-point-nfold}の良い点スキーム開に
$\epsilon$ の同型条件を加えれば、中心直線の復元も任意の基底変更と両立する。
\end{theorem}

\begin{proof}
定理\ref{thm:point-envelope}で $B$ を復元した後、$U=Z(B)_1$、
$Q=\langle F_x,F_y,F_z\rangle$、$C=V(Q)$ と $\alpha$ は全て定まる。
固定点部分は $D$ である。完全交叉イデアルの次数三部分から
\[
 I(D)_3/I(C)_3=k[w_{\mathrm{ab}}]\subset T_3
\]
を得る。命題\ref{prop:point-jacobi-evaluation}の単射 $\alpha$ による
この直線の逆像は、ただ一つの直線 $kw\subset\mathcal E$ である。
従って二次関係にこの三次関係を加えることで $A$ が決まる。

$n=3$ で $D$ が零次元でも、型 $(2,2,2,3)$ の斉次完全交叉イデアルは
飽和している。この場合に $T_{D,3}=H^0(D,\mathcal O_D(3))$ という
誤った全射性を使う必要はなく、多項式からの制限写像の核を取ればよい。
相対版では命題\ref{prop:relative-point-envelope}および
式\eqref{eq:relative-point-ideal-pieces}により、二次包絡と
$p_*\mathcal I_D(3)/(Q\cdot V)$ の直線部分束が任意の基底変更と両立して定まる。
中心評価の余核は局所自由なので、その逆像の形成も基底変更と両立する。
\end{proof}

\begin{remark}[中心変数の個数と復元]
中心三次の復元では $r\ge2$ が本質的である。
比較のため $r=1$ とし、$C$ が八つの被約点である周囲代数を固定する。
純中心三次を加えるパラメータのうち有限個を除けば、
中心三次の可換化は $C$ の全ての点で非零となるので $D=\varnothing$ である。
従って中心商の点対応は同じ $\Gamma_\sigma$ だけになり、中心直線を復元できない。
一点中心拡大の周囲二次関係を復元する
\cite[Proposition~4.2.5]{LBSV}と、本稿の $r\ge2$ における中心商の復元を区別する。
\end{remark}

\begin{proposition}[定常三次切断と古典 Casimir の相違]
\label{prop:point-casimir-restriction-obstruction}
上の開で、古典 Casimir 空間を $\mathcal R=kF\oplus\Sym^3U$ とする。
その $T_3$ への制限像は $(J_g)_3$ への写像の核に入る。
一方、$w\notin\Sym^3U$ ならば $[w_{\mathrm{ab}}]$ はこの核に入らない。
特に、
\[
 \alpha(Z(B)_3)\cap \operatorname{im}(\mathcal R\longrightarrow T_3)
 =\Sym^3U.
\]
従って $w_{\mathrm{ab}}|_C$ を古典 Casimir の制限と同一視する構成、
および純中心三次だけを加えてその同一視を実現する構成は成立しない。
$D$ が上記の Cartier 切断であるとき、古典 Casimir $H$ の切断により
$D=C\cap V(H)$ を scheme として実現することもできない。
\end{proposition}

\begin{proof}
$F$ の $J_g$ への像は $g$ の類であり、
$3g=xg_x+yg_y+zg_z$ なので零である。
純中心三次の像も零である。
一方、中心三次の像は同型 $\epsilon$ を経て零になるのが純中心三次の場合に限る。
さらに $S_3\subset T_3$ は単射なので、交叉についての等式も従う。

幾何学的な切断としての主張のために、$T_3=H^0(C,\mathcal O_C(3))$ を確認する。
$T$ は次元 $r\ge2$ の Cohen--Macaulay な斉次完全交叉環なので、
irrelevant ideal に関する局所コホモロジーの次数零と一次は消える。
切断環との標準完全列から、全ての次数で $T_j=H^0(C,\mathcal O_C(j))$ となる。
また $H^0(C,\mathcal O_C)=k$ である。
二つの非零因子切断が同じ Cartier 零スキームを定めるなら、
その比は $C$ の大域可逆関数なので $k^\times$ に属する。
従って同じ $D$ を与える古典 Casimir が存在すれば、
その制限は $w_{\mathrm{ab}}|_C$ の非零定数倍となり、既に示した障害に反する。
\end{proof}

\begin{remark}[古典制限写像そのものの一般単射性]
上の障害は古典 Casimir の制限写像の核が大きいことによるものではない。
$\mathcal R\to T_3$ 自体は非空開上で単射となる。
実際、$F=g-xu_1^2$ と取れば、Euler の式から $[F]=-(2/3)xu_1^2$ であり、
active 次数の濾過によりこれは純中心三次の像に属さない。
従って $kF\oplus S_3$ の制限は単射であり、単射性は係数空間の開条件である。
三変数制限の Hessian 剛性条件や良い点スキーム条件とも交わらせられる。
一方、$F=g+G(u)$ では $[F]=G(u)$ となるから、この単射性は全係数点の主張ではない。
この相違は補正を用いる Poisson モジュライへの射の非存在を意味しない。
次の trace と Hessian による構成は、この相違を保ったまま中心三次の両空間を同定する。
\end{remark}

\subsection{Trace と Hessian による Poisson 三次対への射}

\begin{lemma}[任意基底上の三次 Poisson 中心]\label{lem:point-relative-casimirs}
定理\ref{thm:relative-point-nfold}の良い周囲族を考える。
任意の、非被約・非平坦なものも含む係数環への基底変更後に、
三次 Poisson Casimir の加群は
\[
 Z_{\mathrm{Pois}}(\Sym(V))_3=R_0F\oplus\Sym^3U
\]
となる。ここで局所表示の係数環を $R_0$ と書いた。
従ってこれは普遍的な局所自由加群であり、その形成は任意の基底変更と両立する。
Hessian のスカラー剛性条件は、この主張には不要である。
\end{lemma}

\begin{proof}
$A_0=\Sym_{R_0}(V)$ とし、$f=(F_x,F_y,F_z)$ と置く。
相対点定理の仮定により、この三元は任意の基底変更後にも正則列である。
三次式 $H$ が Jacobian Poisson 括弧の Casimir であることは
$\nabla_{x,y,z}H\wedge f=0$ と同値である。
正則列の Koszul cochain 複体の $A_0^3$ における完全性から
$\nabla H=a f$ となる $a\in A_0$ が一意に存在する。
両辺の斉次次数を比較し、$A_0\to A_0^3$, $a\mapsto af$ の単射性を使うと
$a\in R_0$ である。
従って $H-aF$ の三つの active 偏微分は零となる。
$R_0$ は標数零の体上の代数なので、$H-aF$ は純中心三次である。
逆包含は定義から明らかである。
滑らかな三変数制限は各点で非零なので、表示した和は局所的に直和となる。
\end{proof}

尺度を含む周囲 Poisson tensor $\Pi=h\Pi_F$ を固定する。
その support は点ベクトル空間 $V^*$ の三次元部分空間
$K=\operatorname{ann}(U)$ である。
ここでは既存の二次包絡復元から得られる $\Pi$ の尺度を捨てない。
$\mathcal R=kF\oplus\Sym^3U$ を補題\ref{lem:point-relative-casimirs}の
三次 Poisson Casimir の空間とし、
$\mathcal L_{\mathrm{cl}}=\mathcal R/\Sym^3U$ と置く。
尺度込みの $\Pi$ を用いているので、命題\ref{prop:point-casimir}で
$Q$ だけから復元するための Hessian 剛性条件を追加する必要はない。

有限自由 $S$-代数 $T$ は階数八なので、正規化した trace
\begin{equation}\label{eq:point-normalized-trace}
 \operatorname{tr}:T\longrightarrow S,
 \qquad a\longmapsto\frac18\operatorname{Tr}_{T/S}(m_a)
\end{equation}
が定まる。これは次数を保ち、$S$ 上で恒等である。
任意の係数環上でも有限局所自由代数の乗法行列の trace で定義されるため、
任意の基底変更と可換する。
以下、$\operatorname{tr}(H)$ は $H$ を $T$ に送ってから trace を取る意味とする。
従って
\begin{equation}\label{eq:point-trace-splittings}
 \begin{aligned}
 \mathcal E&=\mathcal W\oplus\mathcal E^0,
 &\mathcal E^0&=\ker(\operatorname{tr}\circ\alpha),\\
 \mathcal R&=\mathcal W\oplus\mathcal R^0,
 &\mathcal R^0&=\ker(\operatorname{tr}|_{\mathcal R})
 \end{aligned}
\end{equation}
は標準的な分裂であり、後半は標準的な持上げ
$s_{\mathrm{cl}}:\mathcal L_{\mathrm{cl}}\to\mathcal R^0$ を与える。

\begin{lemma}[尺度付き Hessian 同型]\label{lem:point-hessian-line}
尺度込みの $\Pi$ から、表示に依存しない線束同型
\begin{equation}\label{eq:point-hessian-line}
 \nu_\Pi:\mathcal L_{\mathrm{cl}}\xrightarrow{\sim}(J_g)_3
\end{equation}
が定まる。局所的に非零の $\overline F\in\mathcal L_{\mathrm{cl}}$ を選び、
$\Pi=\iota_{dF}P$ となる $P\in\Lambda^3K$ を取れば、
\[
 \nu_\Pi(\overline F)=
 \left[\mathscr H(P,F)\right],\qquad
 \mathscr H(P,F)=\frac16\sum P^{ijk}P^{abc}F_{ia}F_{jb}F_{kc}.
\]
特に $\Pi=h\Pi_F$、$P=h\partial_x\wedge\partial_y\wedge\partial_z$ ならば
$\nu_\Pi(\overline F)=h^2[\det\operatorname{Hess}(g)]$ である。
\end{lemma}

\begin{proof}
この縮約はテンソルとして定義され、基底変更に不変である。
$F$ に純中心三次を加えても $K$ に沿う Hessian は変わらない。
さらに $F$ を $cF$ に替えると、同じ $\Pi$ を与える $P$ は $c^{-1}P$ となり、
縮約は $c$ 倍となる。従って表示した式は商直線上の線形写像を定める。
Hesse 座標では
\[
 [\det\operatorname{Hess}(g)]=8(1+\lambda^3)[xyz]\ne0
\]
なので同型である。滑らかな三変数三次式を Hesse 形にする変換と
縮約の共変性から一般の場合が従う。
族の場合にも、全ての幾何ファイバー上の非消滅により線束同型となる。
\end{proof}

\begin{theorem}[Trace--Hessian による中心三次の同定]
\label{thm:point-trace-hessian}
境界 Jacobi 写像 $\epsilon$ が同型となる良い一ブロック開上で、
\begin{equation}\label{eq:point-trace-hessian-map}
 \mathfrak d_\Pi:\mathcal E\xrightarrow{\sim}\mathcal R,
 \qquad
 w\longmapsto
 -12\,s_{\mathrm{cl}}\bigl(\nu_\Pi^{-1}(\epsilon(\overline w))\bigr)
 +\operatorname{tr}(\alpha(w))
\end{equation}
は線型同型であり、純中心三次上で恒等である。
これは基底変更および次数付き表示変更と両立する。
従って中心三次直線の torsor の同型
\begin{equation}\label{eq:point-central-classical-torsors}
 \mathbb P(\mathcal E)\setminus\mathbb P(\mathcal W)
 \xrightarrow{\sim}
 \mathbb P(\mathcal R)\setminus\mathbb P(\mathcal W)
\end{equation}
を定める。

$H_w=\mathfrak d_\Pi(w)$ が滑らかな射影三次超曲面を定める非空開上で、
\begin{equation}\label{eq:point-poisson-pair}
 [w]\longmapsto (X_w,\pi_w),\qquad
 X_w=V(H_w),\quad \pi_w=\Pi|_{X_w}
\end{equation}
は次数付き表示のスカラー商上の射となる。
より具体的には、$\mathcal Z_N^{\mathrm{sm}}$ を
\[
 \mathbb P(\Sym^3V)\times
 \bigl(\Sym^2V\otimes\Lambda^2V^*\bigr)
\]
内で、三次式の滑らかさおよび $[\Pi,\Pi]=0$、$\Pi(dH,-)=0$ を満たす
$([H],\Pi)$ の局所閉スキームとすれば、標的を代数的商
$[\mathcal Z_N^{\mathrm{sm}}/\mathrm{PGL}(V)]$ と取れる。
斉次 affine lift のデータを忘れることで、Brauer--Severi 周囲空間内の
滑らかな埋込み三次 Poisson 対とその偏極\emph{類}へ射を得る。
大域的な $\mathcal O(1)$ 自体の降下を仮定する必要はない。
定理\ref{thm:point-omarked-descent}と局所的な良い表示のモジュライ比較により、
この射は同じ開の $\mathcal O$-marked な族についても表示に依存しない。
\end{theorem}

\begin{proof}
\eqref{eq:point-trace-splittings} の両分裂の下で、写像は $\mathcal W$ 上の恒等と、
商直線間の同型 $-12\nu_\Pi^{-1}\epsilon$ の直和である。
従って線型同型であり、非純中心の直線を非純中心の直線へ送る。
全ての構成は最大可換代数商、有限局所自由代数の trace、
尺度付き Poisson tensor およびその三次 Casimir から定まる。
これらは表示変更に共変で、任意の基底変更と両立する。
次数付きスカラーは両三次空間で重み三を持つため、
射影化後には自明に作用する。
このことが表示のスカラー商とその fpqc forms 上での降下を与える。

$H_w$ は $\Pi$ の Casimir であるから、その斉次イデアルは Poisson イデアルである。
括弧は次数零を持つので $\operatorname{Proj}$ 上にも Poisson tensor を定める。
一般点では中心座標が非零で Euler 方向が active support の外にあるため、
得られる射影 Poisson tensor は階数二である。

滑らかな三次式の非空性も各固定周囲で確認できる。
$F+G(u)$ を $G\in\Sym^3U$ で動かす線型系は $V(U)$ の外で基点を持たない。
基点部分 $V(U,g)$ では active 勾配が非零なので、各 $F+G$ は滑らかである。
Bertini により一般の $G$ はその外でも滑らかな超曲面を与える。
同型\eqref{eq:point-central-classical-torsors}で引き戻せば、各固定周囲の
中心三次 torsor に非空開を得る。
全パラメータ空間では、既に課した良い点スキーム開や
滑らかな定常成分の非空開とも交わらせられる。
後者を全ての固定周囲で主張しているのではない。
\end{proof}

\begin{remark}[規格化・既存量子化・残余圏との比較]
式\eqref{eq:point-trace-hessian-map}の同型自体は、係数 $-12$ を
任意の固定した非零定数に替えても成立する。
この値は入力研究\cite[Proposition \texttt{prop:pg-ad-intrinsic-hessian}]{NCCubicA}の
二次 jet に合わせたものである。
すなわち、指定した Weyl 表示の二次 jet
\[
 \Omega_{F,h}=\operatorname{sym}
 \left(F-\frac{h^2}{12}\det\operatorname{Hess}_{x,y,z}(F)+O(h^4)\right)
\]
を満たす族に限れば、$\epsilon(\overline\Omega)
=-h^2[\det\operatorname{Hess}(g)]/12+O(h^4)$ なので、
$\mathfrak d_{h\Pi_F}(\Omega_{F,h})=F+O(h^2)$ となる。
ここで指定するのは、同入力命題で純中心二次補正を調整した中心元である。
同入力の短い混合 Hesse 公式には $h^2ABC/6$ という純中心項もあり、
それを引く前の式を上の $O(h^4)$ 表示と同一視しない。
有限の非零 $h$ で常に入力の $F$ と一致するとは主張しない。
またこの同変な構成の存在は、自然な逆量子化の一意性を意味しない。

点スキームの定常切断 $w_{\mathrm{ab}}|_C$ と $H_w|_C$ は異なる。
前者の境界 Jacobi 類を消すのではなく、その非零類を Hessian 同型によって
古典 Casimir 方向の係数へ読み替え、純中心成分を trace で指定している。
これにより Poisson 三次対への射は構成できるが、
$X_w$ と残余圏が回復する可換三次多様体 $X_{\mathrm{res}}$ の同型、
および $\operatorname{Ku}(A)\simeq\operatorname{Ku}(X_w)$ は別の比較問題として残る。
特定の量子化 chart の写像を全ての内在的モジュライ成分の写像に拡張したとは主張しない。
\end{remark}

% !TEX root = ../nc-cubic-n-fold.tex
\section{強い affine 持上げと Poisson 対からの逆構成}
\label{sec:poisson-inverse}

体上の三次超曲面の双ベクトル表示
\eqref{eq:nfold-bivector}は入力論文\cite[命題
\texttt{prop:pg-represent}]{NCCubicA}にある。
ここではその任意基底上の形を証明し、Poisson 射の標的から affine lift の
追加データを外す。さらに良い一ブロック型では逆構成を与え、
前節の射がどの開部分を同定するかを明示する。
この逆構成は、得られる三次の Kuznetsov 成分との同値を意味しない。

\subsection{任意基底上の強い持上げ}

$R$ を可換 $\mathbf C$-代数、$V$ を階数 $N=n+2\ge5$ の有限局所自由
$R$-加群、$W=V^*$ とする。局所的に
$H\in\Sym^3V$ が $R$ 上滑らかな三次超曲面
$X\subset\mathbb P(W)$ を定めるとする。
$E=\sum z_i\partial_{z_i}$ を Euler ベクトル場と書く。
二次係数の斉次 affine 双ベクトル
$\Pi\in\Sym^2V\otimes\Lambda^2W$ が
\[
 \iota_{dH}\Pi=0
\]
を多項式の恒等式として満たすとき、$H$ に対する\emph{強い持上げ}と呼ぶ。
これは $H$ による剰余での接性より強い条件である。

\begin{theorem}[強い affine 持上げの存在と一意性]
\label{thm:poisson-strong-lift}
任意の $\pi\in H^0(X,\Lambda^2T_{X/R})$ は一意な強い持上げを持つ。
その持上げは、一意な定数三重ベクトル $P\in\Lambda^3W$ を用いて
\[
 \Pi=\iota_{dH}P
\]
と書ける。これらの同定は任意の基底変更と両立する。
$H$ の生成元を変えたときには $H\mapsto cH$, $P\mapsto c^{-1}P$ となる。
従って三次方程式の直線束を $\mathcal L_H$ と書けば、全体として
\begin{equation}\label{eq:poisson-relative-bivectors}
 p_*\Lambda^2T_{X/R}\simeq\mathcal L_H\otimes\Lambda^3W
\end{equation}
である。特に左辺は有限局所自由で、任意の基底変更と両立する。
また $\pi$ が Poisson であることと、強い持上げ $\Pi$ が affine Poisson
双ベクトルであることは同値である。この同値も非被約基底を含めて成立する。
\end{theorem}

\begin{proof}
基底を局所化して $V$ と方程式直線を自明化する。
$\mathbb P=\mathbb P^{N-1}_R$ と書く。
Euler 列とその二次外冪の列は
\[
 0\to\OO\to W\otimes\OO(1)\to T_{\mathbb P/R}\to0,\qquad
 0\to T_{\mathbb P/R}\to\Lambda^2W\otimes\OO(2)
      \to\Lambda^2T_{\mathbb P/R}\to0
\]
である。相対次元が $N-1\ge4$ なので、射影空間の線束コホモロジーから
\[
 R^ip_*\bigl(\Lambda^2T_{\mathbb P/R}(-3)\bigr)=0
 \quad(i=0,1),\qquad R^1p_*T_{\mathbb P/R}=0
\]
を得る。これらの消滅は任意の係数環上で成立する。
従って $\mathbb P$ 上の双ベクトルの $X$ への制限は同型であり、
各切断は二次係数の affine 双ベクトル $q$ で表せる。
この affine 表示の核は、一次係数ベクトル場 $L$ に対する $E\wedge L$ である。

$q$ が $X$ 上で $\Lambda^2T_{X/R}$ に値を持つため、
$\iota_{dH}q$ の $T_{\mathbb P/R}(3)|_X$ における像は零である。
制限列と Euler 列の核を用いて
\[
 \iota_{dH}q=HL+aE
\]
と書ける。ここで $L$ は一次係数のベクトル場、$a$ は三次式である。
もう一度 $dH$ と縮約すると $H(L(H)+3a)=0$ を得る。
滑らかな相対三次の方程式 $H$ は、任意基底変更後にも多項式環上の非零因子である。
従って $a=-L(H)/3$ であり、
\[
 \Pi=q-\tfrac13E\wedge L
\]
は強い持上げである。
二つの強い持上げの差は $\Delta=E\wedge L$ と書ける。
その縮約は $3HL-L(H)E=0$ であるから、さらに $E$ と外積を取ると
$3H\Delta=0$ となる。$H$ の非零因子性から $\Delta=0$ であり、一意性を得る。

滑らかな三次の偏微分 $H_1,\ldots,H_N$ の Koszul 複体は、
各次数の有限自由複体をファイバー上の完全性から局所的に分離することで、
任意基底変更後にも完全である。
従って二次係数の $\Pi$ に対する $\iota_{dH}\Pi=0$ は
$\Pi=\iota_{dH}P$ と同値である。
係数次数が零なので $P$ は定数三重ベクトルであり、
次の Koszul 項の係数次数は $-2$ だから、その $P$ は一意である。
構成の一意性によって局所表示は貼り合わさり、全ての基底変更と座標変更に両立する。

最後に Jacobi 条件を確認する。
標準体積形式の divergence を $\delta$ と書くと、
$\Pi=\iota_{dH}P$ は $\delta\Pi=0$ を満たす。
$J=[\Pi,\Pi]$ は $\delta J=0$ および $\iota_{dH}J=0$ を満たす。
$J$ の射影化が $X$ 上で零なら、Euler 列の三次外冪と制限列により
\[
 J=HT+E\wedge B,\qquad
 T\in\Lambda^3W,\quad B\in\Sym^2V\otimes\Lambda^2W
\]
と書ける。同じ Euler コホモロジー計算がこの表示を任意基底上で与える。
$\delta(E\wedge B)=NB-E\wedge\delta B$ より
\[
 E\wedge B=-N^{-1}E\wedge\iota_{dH}T,\qquad
 0=\iota_{dH}J=(1-3/N)H\iota_{dH}T.
\]
$N\ne3$、$H$ は非零因子であり、定数三重ベクトル上の
$\iota_{dH}$ は上述の Koszul 計算で単射なので、$T=0$、従って $J=0$ となる。
逆向きは制限するだけである。
\end{proof}

\begin{corollary}[一ブロック support の回復]
\label{cor:poisson-support-recovery}
定理の $P$ が非零分解可能なら、その support
$K\subset W$ は階数三の部分束として回復される。
$U=\operatorname{ann}(K)$ は $\Pi$ の一次 Poisson 中心に等しい。
従って滑らかな三次の Poisson 対とその超平面偏極類から、
強い持上げ、中心一次部分、および一ブロック support を回復できる。
\end{corollary}
\begin{proof}
相対的な「非零分解可能」とは、非零 Pl\"ucker 錐のスキームを通るという条件とする。
これは、階数三の部分束 $K$ と
$\mathcal L_H\otimes\det K$ の非零切断を与えることに等しい。
従って support の形成は任意基底変更と両立する。
一次式 $u$ が Poisson 中心に属すれば
$\iota_{dH}(\iota_{du}P)=0$ である。
定数二重ベクトル上でも Koszul 微分は単射だから、$\iota_{du}P=0$、
すなわち $u\in\operatorname{ann}(K)$ である。逆向きは直接従う。
\end{proof}

\subsection{分解可能な Poisson 対と良い量子化の開}

$\mathcal H_N^{\mathrm{sm}}\subset\mathbb P(\Sym^3V)$ を滑らかな三次式の開、
$\mathcal L_H\subset\Sym^3V\otimes\OO_{\mathcal H_N^{\mathrm{sm}}}$ を普遍直線とする。
$\mathcal Z_N^{\mathrm{dec}}$ を、$([H],\Theta)$ で
$0\ne\Theta\in\mathcal L_H\otimes\Lambda^3W$ が分解可能なものの相対スキームとする。
具体的には、$\mathcal H_N^{\mathrm{sm}}\times\operatorname{Gr}(3,W)$ 上の
直線束 $\mathcal L_H\otimes\det\mathcal K$ の零切断を除いた空間である。
従ってこれは滑らかで既約である。
次数スカラーの重みは $3-3=0$ なので $\operatorname{PGL}(V)$ が作用する。
\[
 \mathfrak P_n^{\mathrm{dec}}
   =[\mathcal Z_N^{\mathrm{dec}}/\operatorname{PGL}(V)]
\]
と置く。前定理によりこれは、強い lift に対応する三重ベクトルが
非零分解可能な滑らかな三次 Poisson 対と超平面偏極類のモジュライである。
affine lift 自体を追加標識として保持する必要はない。
偏極類を保つ同型は、局所的に $\OO_X(1)$ を同定し全切断を取ることで、
ambient の射影同型へ一意に延長する。
ここでは $p_*\OO_X(1)=V$ および任意基底変更との整合を使う。

有限安定化子を要求する前の良い一ブロック PBW 開を
$T^{\mathrm{alg}}$ と書く。
ここには、相対 PBW・普遍中心・相対 Serre 双対・全
$\mathcal O$-同値比較に必要な入力を全て課すが、
周囲代数 $B$ 単独の安定化子が有限という条件は課さない。
これらの入力が成立することは、対象とする表示チャートごとに必要である。
中心三次 torsor を $J^{\mathrm{alg}}$ とする。
その上で境界 Jacobi 評価 $\epsilon$ の同型条件、相対点スキームの良い条件、
および $H_w=\mathfrak d_\Pi(w)$ の滑らかさを課した非空開を
$J^{\mathrm{sm}}$ と書く。
\[
 \mathfrak M_n^{\mathrm{sm}}=[J^{\mathrm{sm}}/\overline P]
\]
はこれらの良い\emph{対}を局所模型とする $\mathcal O$-marked 圏族の stack である。
必要なら旧 $T^{\mathrm{good}}$ を縮め、旧 Poisson 射の定義域
$\mathfrak M_n^{\mathrm{Pois}}$ がこの stack の有限周囲安定化子部分の
非空開になるようにする。
旧 $\mathfrak M_n^{\mathrm{loc}}$ 全体は $H_w$ の滑らかさ等を要求していないので、
その全体を $\mathfrak M_n^{\mathrm{sm}}$ の開と同一視しない。

\begin{theorem}[Poisson 対からの逆構成と開部分の同定]
\label{thm:poisson-good-open-isomorphism}
$\mathfrak P_n^{\mathrm{dec}}$ には次の条件で定まる非空開
$\mathfrak P_n^{\mathrm{good}}$ があり、trace--Hessian 射は同型
\begin{equation}\label{eq:poisson-good-open-isomorphism}
 \mathfrak M_n^{\mathrm{sm}}\xrightarrow{\sim}
 \mathfrak P_n^{\mathrm{good}}
\end{equation}
を与える。
\begin{enumerate}
\item 強い持上げ $\Pi$ から定める二次代数
\[
 B_\Pi=T(V)/
 \bigl(v\otimes v'-v'\otimes v-\operatorname{sym}\Pi(v,v')\bigr)
\]
が $T^{\mathrm{alg}}$ の良い模型であり、境界評価 $\epsilon$ が同型である。
\item 一意な中心直線
\[
 [w]=\mathfrak d_\Pi^{-1}([H])\subset Z(B_\Pi)_3
\]
が相対点スキームの良い条件を満たす。
\end{enumerate}
同型は任意の基底変更と両立し、全ての自己同型群を同定する。
$\mathfrak M_n^{\mathrm{sm}}$ は滑らかな既約有限型 DM stack で、次元は
\[
 \binom{n+4}{3}-(n+2)^2+3n-2
 =\binom{n+2}{3}+3n-2
\]
である。周囲代数 $B$ の安定化子が正次元であっても、
ここに列挙した良い入力と $H_w$ の滑らかさがあれば適用できる。
\end{theorem}

\begin{proof}
support を局所的に固定し、$K$ の体積元 $P_0$ を一つ選ぶ。
$\Theta=H\otimes tP_0$ と書けば、対応する二次代数の potential は
$F=tH$ modulo $\Sym^3U$ である。
$H\mapsto cH$ に対して $t\mapsto c^{-1}t$ なので、この $F$ は変わらない。
support を保つ frame の変更では、従来と同じ
$\det(g|_{V/U})^{-1}g(F)$ の作用となる。
従って条件(1)は $P$-不変な良い表示開の引戻しであり、開条件である。
その上では trace--Hessian 同型が任意基底変更で定義される。
条件(1)に含まれる active 楕円三次の仮定により、
$H|_{\mathbb P(K)}$ が滑らかであるため $[H]$ は純中心方向でなく、
条件(2)の中心直線は一意に定まる。条件(2)も相対 norm による開条件である。
従来の非空開 $J^{\mathrm{Pois}}$ の像を含むので、得られる開は非空である。

この逆構成から $A=B_\Pi/(w)$ を作る。
次数一生成元束は一般には表示 gerbe 上にあるが、
第\ref{sec:nfm-local-moduli}節の重み付き tails の構成は
周囲の有限安定化子を用いないので、$\mathcal O$-marked 圏族を与える。
全 $\mathcal O$-同値比較も同じ入力だけを使うため、
この圏族の同型問題は $[J^{\mathrm{sm}}/\overline P]$ に一致する。

量子側から出発すると、元の $\Pi$ は $H_w$ に対する強い持上げである。
定理\ref{thm:poisson-strong-lift}の一意性により、Poisson 対から回復する持上げは
元の $\Pi$ と等しい。従って $B$ をその二次関係から回復し、
trace--Hessian 同型の逆で元の中心直線 $[w]$ を得る。
逆に Poisson 側から作った $A$ に trace--Hessian 射を適用すると、
定義により元の $[H]$ と $\Pi|_X$ に戻る。
これらは係数変更と表示変更に両立する、相互に逆な構成である。

射も一致する。Poisson 対の偏極類を保つ同型は ambient の射影同型を与え、
強い持上げの一意性から $\Pi$ を保つ。
従って $K,U,B_\Pi$ を保ち、trace--Hessian の共変性から $[w]$ も保つ。
逆に全 $\mathcal O$-同値は局所的に次数付き同型に持ち上がるため、
Poisson 対の同型を与える。二つの線型持上げの差は次数スカラーである。
これにより幾何点の全単射だけでなく、fpqc stack の同型を得る。

DM 性を直接確認する。
滑らかな三次 $H$ の射影安定化子の Lie 代数に属する一次ベクトル場 $L$ は
$L(H)=cH$ を満たす。Euler の公式を引けば
\[
 \sum_i(L_i-cz_i/3)H_i=0
\]
となる。滑らかな三次の偏微分の正則列には一次係数の syzygy がないため、
$L=cE/3$ であり、射影 Lie 代数では零である。
従って射影安定化子は有限で、標数零では非分岐である。
Poisson 対の安定化子はその部分群なので、同じ結論を持つ。
$J^{\mathrm{sm}}$ は滑らかな既約表示空間上の加法 torsor の非空開であり、
$\overline P$ は滑らかなので、商は滑らかかつ既約である。
表示空間と群の次元を引けば主張の次元を得る。
\end{proof}

\begin{corollary}[拡大した対の stack における Brauer 障害]
\label{cor:poisson-enlarged-brauer}
$[J^{\mathrm{sm}}/P]\to\mathfrak M_n^{\mathrm{sm}}$ の表示 gerbe の類を
$\beta^{\mathrm{sm}}$ とする。
$n\not\equiv1\pmod3$ では中立であり、
$n\equiv1\pmod3$ ではどの非空開上でも位数が正確に三である。
\end{corollary}
\begin{proof}
determinant character による $r\beta=3\beta=0$ と
$\gcd(r,3)=1$ の場合の中立化には、有限安定化子は使われない。
$\gcd(r,3)=3$ の場合、旧 $\mathfrak M_n^{\mathrm{Pois}}$ は
既約な $\mathfrak M_n^{\mathrm{sm}}$ の非空開なので稠密である。
任意の非空開との交わりに定理\ref{thm:nfm-brauer-dense-open}を適用すれば、
障害が非零の体値点を得る。これで位数三が従う。
\end{proof}

\subsection{全 Poisson モジュライの一般部分との関係}

\begin{corollary}[一般三次上の開埋込みと支配的成分]
\label{cor:poisson-generic-open-immersion}
系\ref{cor:nfold-vertex-support}と注意\ref{rem:nfold-relative-punctured}が適用される
滑らかな三次式の非空不変開を $\mathcal H_N^\circ$ とする。
この開上では、非零三次 Poisson 対の全モジュライは
$\mathfrak P_n^{\mathrm{dec}}$ の制限に等しい。
従って\eqref{eq:poisson-good-open-isomorphism}の制限は、
全非零 Poisson モジュライへの非空な開埋込みである。
その像は滑らかな三次のモジュライを支配する。

滑らかな三次 Poisson 対の全モジュライには、三次のモジュライを支配する
既約成分が一つだけある。本構成はその一般部分の非空開を同定する。
特殊な三次の locus 上にのみ存在する他の既約成分を排除する主張ではない。
\end{corollary}
\begin{proof}
相対非零 Jacobi スキームと非零 Pl\"ucker 錐の等式は、引用した結果により
任意の基底変更と両立する。三次方程式の直線束の捻りと射影群の商を戻すと、
表示したモジュライの等式を得る。
$\mathcal Z_N^{\mathrm{dec}}\to\mathcal H_N^{\mathrm{sm}}$ は滑らかな全射であり、
各幾何ファイバーは非零 Grassmann 錐で既約である。
従ってその任意の非空開の像は基底の非空開となり、支配性が従う。
一般の三次上で全 Poisson スキームの台は一つの Grassmann 錐であるため、
一般ファイバーは既約である。有限型の全モジュライで基底を支配する
既約成分は一般ファイバーの既約成分に対応するので、一つだけである。
頂点に支持される非被約構造は、この台の既約性の主張を変えない。
\end{proof}

\begin{remark}[開埋込みの範囲]
分解可能な三重ベクトルの条件は、全ての滑らかな三次を動かす空間では
一般に閉条件であり、非零条件を加えても自動的に開にはならない。
従って定理\ref{thm:poisson-good-open-isomorphism}の無条件の標的は
$\mathfrak P_n^{\mathrm{dec}}$ である。
全非零 Poisson モジュライへの開埋込みには、上の一般三次の開への制限を用いる。
また、一般対象の局所一ブロック表示をこの Poisson 対のモジュライの範囲で得たことと、
未構成の全 intrinsic $\mathcal O$-marked 圏モジュライの全成分を分類することは異なる。
\end{remark}

% !TEX root = ../nc-cubic-n-fold.tex
\section{JacobiスキームとGrassmann錐の次元境界}
\label{sec:nfold-jacobi-scheme}

本節は、Poisson構造の被約な分類と、その未被約スキーム構造を区別する。
$N=n+2\ge5$、$W=\mathbb C^N$とし、$F\in\operatorname{Sym}^3W^*$が
滑らかな三次超曲面$X_F\subset\mathbb P(W)$を定めるとする。
$\widehat{G}_N\subset\Lambda^3W$を$\operatorname{Gr}(3,W)$のPl\"uckerアフィン錐とし、
$J_N\subset R_N:=\operatorname{Sym}((\Lambda^3W)^*)$をその斉次イデアルとする。
双ベクトルのスカラーは割らない。

\subsection{Hessian二次式によるスキームの表示}
座標を$z_1,\ldots,z_N$、$F_i=\partial_iF$、
$F_{ijr}=\partial_i\partial_j\partial_rF$とする。
$P\in\Lambda^3W$に対し
$q_P=\iota_{dF}P$と置く。
既存稿\cite[命題\texttt{prop:pg-represent}]{NCCubicA}の
Euler列とgradientのKoszul複体による表示は、
\begin{equation}\label{eq:nfold-bivector}
 \Lambda^3W\xrightarrow{\sim}H^0(X_F,\Lambda^2T_{X_F}),
 \qquad P\longmapsto q_P|_{X_F}
\end{equation}
を与える。$F$を$aF$へ取り替えると$P$は$a^{-1}P$へ取り替える。
以下のイデアルはこの選択に依存せず同型となる。

$H\in\operatorname{Sym}^2W^*$について、
\begin{align}\label{eq:nfold-hessian-B}
 B_H(P)(\alpha,\beta,\gamma,\eta)
 &=H(P(\alpha,\beta),P(\gamma,\eta))\notag\\
 &\quad+H(P(\beta,\gamma),P(\alpha,\eta))
       +H(P(\gamma,\alpha),P(\beta,\eta))
\end{align}
と定める。ここで$P(\alpha,\beta)\in W$は
$\eta(P(\alpha,\beta))=P(\alpha,\beta,\eta)$で定める。
$B_H(P)$は四つの共変ベクトルに関して交代的であり、$\Lambda^4W$の元である。
$H_r=(F_{ijr})_{ij}$と置き、$I_F\subset R_N$を
\begin{equation}\label{eq:nfold-jacobi-quadrics}
 B_{H_r}(P)(dz_a,dz_b,dz_c,dz_d),
 \qquad 1\le r\le N,\quad a<b<c<d
\end{equation}
で生成される二次斉次イデアルとする。

\begin{lemma}\label{lem:nfold-jacobi-scheme}
Poisson構造の自然なJacobi零スキームは
\begin{equation}\label{eq:nfold-poisson-scheme}
 \operatorname{Pois}(X_F)=\operatorname{Spec}(R_N/I_F)
\end{equation}
である。この同定は任意の可換$\mathbb C$-代数、特に非被約なArtin係数への
基底変更と両立する。また
\begin{equation}\label{eq:nfold-quadratic-inclusion}
 I_F\subset J_N,\qquad
 \dim_{\mathbb C}(I_F)_2\le N\binom N4
\end{equation}
が成り立つ。
\end{lemma}
\begin{proof}
既存稿のJacobi延長の議論を、係数環を含めて確認する。
$E=\sum z_i\partial_i$と標準体積形式のdivergence $\delta$を用いると、
$q_P$は$\delta q_P=0$および$\iota_{dF}q_P=0$を満たす。
$J=[q_P,q_P]$の射影化が$X_F$で零なら、Euler列と制限列から
$J=FT+E\wedge B$、$T\in\Lambda^3W$、
$B\in\operatorname{Sym}^2W^*\otimes\Lambda^2W$と書ける。
$\delta J=0$と
$\delta(E\wedge B)=NB-E\wedge\delta B$により
$E\wedge B=-N^{-1}E\wedge\iota_{dF}T$。
さらに$\iota_{dF}J=0$から
$(1-3/N)F\iota_{dF}T=0$を得る。
$N\ne3$、$F$は零因子でなく、$F_i$は定数係数で独立なので$T=0$、$J=0$となる。
この議論は\cite[命題\texttt{prop:pg-jacobi-extension}]{NCCubicA}と同じものである。

アフィンJacobiatorの$ijk$成分は、非零定数倍と符号を除いて
\begin{equation}\label{eq:nfold-J-linearized}
 \sum_a F_a L_a^{ijk},\qquad
 L_a^{ijk}=
 \sum_{r,b}\bigl(P^{ira}P^{jkb}+P^{jra}P^{kib}+P^{kra}P^{ijb}\bigr)F_{br}
\end{equation}
である。$F_1,\ldots,F_N$は二次式の正則列なので、
一次式を係数とするsyzygyは存在しない。
従って$J=0$と全ての$L_a^{ijk}=0$は同値であり、後者は
\eqref{eq:nfold-jacobi-quadrics}に一致する。
ここで用いた写像
$\bigoplus_a\operatorname{Sym}^1W^*\to\operatorname{Sym}^3W^*$、
$(\ell_a)\mapsto\sum\ell_aF_a$は、有限次元$\mathbb C$-ベクトル空間の
分裂単射である。Euler列から得た有限次元の同型も分裂し、
$F$による乗法も各次数で分裂単射である。
従って任意の係数環へテンソルしても、上の同値は保たれる。
これは複素点の集合の等式だけでなく、零スキームの等式を示す。

$P$が分解可能なら、そのsupportの三次元部分空間に沿う
Jacobian括弧となるため、任意の$F$についてJacobi恒等式を満たす。
従って全ての生成二次式は$\widehat G_N$上で消え、$I_F\subset J_N$となる。
最後の上界は\eqref{eq:nfold-jacobi-quadrics}の本数である。
\end{proof}

\begin{lemma}\label{lem:nfold-plucker-count}
Grassmann錐の二次イデアルの次元は
\begin{equation}\label{eq:nfold-plucker-count}
 q_N:=\dim(J_N)_2
 =5\binom{N+1}{6}=N\binom N5-\binom N6
\end{equation}
である。
\end{lemma}
\begin{proof}
Pl\"ucker環の標準単項式基底を用いる
\cite[\S2.4, Theorem~2.3]{BCMStandardMonomial}。
三元部分集合を増加列として書き、成分ごとの順序を用いると、
次数二の標準単項式は$p_Ap_B$、$A\le B$である。
従って$(J_N)_2$の次元は、辞書式に$A\le_{\mathrm{lex}}B$と並べた
三元部分集合対で、成分ごとには比較できないものの個数に等しい。
二つ以上の添字を共有する対は比較可能である。
和集合の大きさが五の場合、固定した五元集合上には比較不能な対が五個ある。
添字を$1,\ldots,5$へ移した表示は
\[
 (125,134),\ (125,234),\ (135,234),\ (145,234),\ (145,235).
\]
和集合の大きさが六の場合も五個であり、
\[
 (126,345),\ (136,245),\ (145,236),\ (146,235),\ (156,234)
\]
である。従って$q_N=5\binom N5+5\binom N6=5\binom{N+1}{6}$を得る。
\end{proof}

\begin{lemma}[Jacobi 二次式の普遍的な一次従属]
\label{lem:nfold-universal-row-relations}
$N\ge5$とし、$T\in\operatorname{Sym}^3W^*$を任意の対称三次テンソルとする。
式\eqref{eq:nfold-jacobi-quadrics}の係数写像を
\[
 \Phi_T:W\otimes\Lambda^4W^*\longrightarrow (J_N)_2,
 \qquad e_r\otimes\alpha\longmapsto B_{H_r}(P)(\alpha)
\]
と書く。ただし$\alpha$に関する評価は線形に延長する。
座標に依存しない分裂単射
\[
 \iota:\Lambda^3W^*\longrightarrow W\otimes\Lambda^4W^*,
 \qquad\beta\longmapsto\sum_{r=1}^N e_r\otimes(dz_r\wedge\beta)
\]
について$\Phi_T\iota=0$であり、従って
\[
 \dim(I_F)_2\le N\binom N4-\binom N3.
\]
さらに、全ての対称三次テンソルを同時に考えた共通核は
\[
 \bigcap_{T\in\operatorname{Sym}^3W^*}\ker\Phi_T=\operatorname{im}\iota
\]
に等しい。この共通核の等式だけから、固定した一般の$T$で核が
$\operatorname{im}\iota$に等しいとは結論しない。
\end{lemma}
\begin{proof}
任意の$a<b<c$について、Jacobi二次式の間に
\[
 \sum_{r=1}^N B_{H_r}(P)(dz_r,dz_a,dz_b,dz_c)=0
\]
がある。実際、左辺を展開すると
\[
 \sum_{r,i,j}T_{ijr}
 \bigl(p_{rai}p_{bcj}+p_{abi}p_{rcj}+p_{bri}p_{acj}\bigr)
\]
である。第一項・第三項では$(r,i)$、第二項では$(r,j)$を交換すると、
$T$は不変で該当する$p$は符号を変えるので、各和は零である。
これが$\Phi_T\iota=0$を示す。
縮約写像
\[
 \kappa:W\otimes\Lambda^4W^*\longrightarrow\Lambda^3W^*,
 \qquad e_r\otimes\alpha\longmapsto\iota_{e_r}\alpha
\]
に対し$\kappa\iota=(N-3)\operatorname{id}$だから、
$\iota$は分裂単射であり、$\binom N3$個の独立な従属関係を得る。

共通核を調べる。$u=\sum e_r\otimes\alpha_r$をその元とし、
任意の$\ell\in W^*$に対して$T=\ell^3$を代入する。
$H_r=\ell(e_r)\ell^2$であるから、
$\alpha(\ell):=\sum\ell(e_r)\alpha_r$は
全ての$P$に対して$B_{\ell^2}(P)$を消す。
$\ell\ne0$なら$\iota_\ell P$は$\Lambda^2\ker\ell$の任意の元を走り、
$B_{\ell^2}(P)$はその外積二乗の非零定数倍である。
二重ベクトルの外積二乗は$\Lambda^4\ker\ell$を張るので、
$\alpha(\ell)$の$\ker\ell$への制限は零、すなわち
$\ell\wedge\alpha(\ell)=0$である。

最後の条件から$\alpha(\ell)=\ell\wedge\beta$となる
一定の$\beta\in\Lambda^3W^*$が存在することを座標で確認する。
$\ell=dz_i$を代入すると$\alpha_i=dz_i\wedge\beta_i$と書ける。
$\ell=dz_i+dz_j$を代入すると
$dz_i\wedge dz_j\wedge(\beta_j-\beta_i)=0$である。
任意の三元添字集合$K$について、$i\notin K$の場合の
$\beta_i$の$dz_K$係数は$i$によらない。この共通係数を$\beta$の
$dz_K$係数とすれば、全ての$i$について
$\alpha_i=dz_i\wedge\beta$を得る。
従って$u=\iota(\beta)$である。
\end{proof}


\subsection{一般三次における被約化と未被約構造}
ここで既存稿\cite[系\texttt{cor:pg-generic}]{NCCubicA}を明示的な入力とする：
一般の滑らかな三次式$F$について、全ての非零Poisson構造は一般階数二であり、
\eqref{eq:nfold-bivector}の下で分解可能な三重ベクトルに対応する。
これは既存稿の全次元正規形と高階数層の次元評価の帰結である。
したがって、ある非空Zariski開集合上で
\begin{equation}\label{eq:nfold-reduced-cone}
 \sqrt{I_F}=J_N
\end{equation}
となる。以下では、この被約化の主張からは決まらないスキーム構造を調べる。

\begin{theorem}[八次元以上のJacobiスキームの非被約性]
\label{thm:nfold-high-dimensional-nilpotents}
$n\ge8$、すなわち$N\ge10$とする。
任意の滑らかな三次式$F$に対し$I_F\ne J_N$であり、
頂点での完備局所環も$\widehat G_N$の頂点完備化と同型ではない。
さらに\eqref{eq:nfold-reduced-cone}を満たす一般の$F$について、
$\operatorname{Pois}(X_F)$は非被約であり、その被約化は$\widehat G_N$である。
冪零根基の次数二部分は少なくとも
\begin{equation}\label{eq:nfold-nilpotent-lower-bound}
 \epsilon_N:=5\binom{N+1}{6}-N\binom N4+\binom N3
 =\frac{(N+1)N(N-1)(N-2)(N-4)(N-9)}{144}>0
\end{equation}
次元を持つ。
特に$n=8$では$\epsilon_{10}=330$である。
\end{theorem}
\begin{proof}
補題\ref{lem:nfold-universal-row-relations}と\ref{lem:nfold-plucker-count}から
\[
 \dim (J_N)_2-\dim (I_F)_2\ge\epsilon_N>0
\]
となり、$I_F\subsetneq J_N$である。
\eqref{eq:nfold-reduced-cone}の下では$J_N/I_F$が冪零根基であり、
その二次部分は少なくとも$\epsilon_N$次元である。

頂点の極大イデアルを$\mathfrak m$とする。
$I_F$も$J_N$も二次斉次式で生成され、一次部分は零なので、
対応する局所環の$\mathfrak m/\mathfrak m^2$はともに次元$\binom N3$であるが、
$\mathfrak m^2/\mathfrak m^3$の次元には少なくとも$\epsilon_N$の差がある。
この数は局所環の抽象同型および完備化で不変だから、完備局所環は同型でない。
冪零根基の次数二の元が局所化で消えないことも、定数項が非零の多項式を掛けて
$I_F$に入るなら、その二次成分を比較して元の二次式が$I_F$に入ることから分かる。
\end{proof}

\begin{remark}[被約な分類と形式変形]
この定理は、一般の高次元三次でも非零Poisson構造の複素点が
Grassmann錐をなすことと両立する。
相違は非被約係数を許したJacobiスキームに現れる。
ここから全Abel圏変形の形式環との同一視を結論してはいない。
その比較には、相対formality、非可換模型の平坦性、および降下の独立した証明が必要である。
少なくとも、被約なPoisson分類だけから全Artin変形を
Grassmann錐の完備化に置き換えることはできない。
\end{remark}

\subsection{有限次元の最大階数を与える計算証明書}
\begin{theorem}[八次元以下の一般Jacobi行列の階数]
\label{thm:nfold-low-dimensional-generic}
$3\le n\le8$、すなわち$5\le N\le10$とする。
滑らかな三次式のパラメータ空間には非空Zariski開集合が存在し、
その上で
\[
 \dim(I_F)_2=min\left\{5\binom{N+1}{6},\,
                  N\binom N4-\binom N3\right\}
\]
が成り立つ。従って$3\le n\le7$の一般の$X_F$のPoissonスキームは、被約化を取る前から
$\operatorname{Gr}(3,N)$のアフィン錐である。
この範囲では、局所的に三次式の代表を選んだ相対Poissonスキームも
この開集合上でGrassmann錐との積になる。
$n=8$では一般に$\dim(I_F)_2=1980$、
$\dim(J_{10}/I_F)_2=330$である。
また$N=9,10$では一般の$F$について
$\ker\Phi_{T_F}=\operatorname{im}\iota$である。
$n=3$で全ての滑らかな三次に対する等式は既知であり
\cite[Theorem~1, \S2.2]{Mayanskiy}、ここで新たに用いるのは$n=4,5,6,7$の
一般点のイデアル一致と$n=8$の不足次元の厳密な判定である。
\end{theorem}
\begin{proof}
補題\ref{lem:nfold-jacobi-scheme}から$(I_F)_2\subset(J_N)_2$である。
補題\ref{lem:nfold-universal-row-relations}も合わせると、
表示した最小値が階数の上界である。
従ってJacobi二次式の係数行列がこの上界を達成する三次テンソルが一つあれば十分である。
この階数条件は三次式の係数に関する非空Zariski開条件であり、
滑らかな三次式の非空開集合と交わる。
$N\le9$ではPl\"uckerイデアルが二次式で生成されることから、
二次部分の一致はイデアル全体の一致を与える。
$N=10$では$\dim(J_{10})_2=2310$から$1980$を引いて$330$を得る。
$N=9,10$の核の主張は、その次元が$\binom N3$に等しいことと
普遍的な分裂単射$\iota$から従う。

再現コード\path{computations/verify_jacobi_rank.py}と
証明書\path{computations/jacobi_rank_certificate.json}は、各$N$に対して
対称整数テンソル$T_{ijr}$を明示している。
対応する三次式は
$F_T=\frac16\sum_{i,j,r}T_{ijr}z_i z_j z_r$であり、Hessianの$z_r$係数は$T_{ijr}$である。
行は$(r,a<b<c<d)$、列は辞書式に並べたPl\"ucker変数の非順序二次単項式で添字付ける。
行列の各行は
\begin{equation}\label{eq:nfold-certificate-matrix}
 \sum_{i,j}T_{ijr}
 \bigl(p_{abi}p_{cdj}+p_{bci}p_{adj}+p_{cai}p_{bdj}\bigr)
\end{equation}
の係数であり、添字の重複を持つ$p$は零、その他は交代符号で増加順に直す。
証明書は小行列の行番号・列番号と、法$32003$での非零行列式を含む：
\begin{center}
\begin{tabular}{c|r|r|r}
$N$ & Jacobi二次式の本数 & 小行列の次数 & 行列式$\bmod 32003$\\
\hline
5 & 25 & 5 & 15512\\
6 & 90 & 35 & 14374\\
7 & 245 & 140 & 7973\\
8 & 560 & 420 & 508\\
9 & 1134 & 1050 & 13073\\
10 & 2100 & 1980 & 16095
\end{tabular}
\end{center}
法素数で非零の整数行列式は特性零でも非零である。
ここでは有限体での零を特性零の恒等式と見なす推論を使っていない。
証明書の生成と、指定された小行列の行列式の再計算は別の手順で実行する。

最後に$N\le9$で、ある小行列式が可逆な係数開集合では、
Jacobi二次式から$(J_N)_2$の固定基底を係数環上で解き出せる。
従って相対イデアルも$J_N$の係数拡大に等しく、相対積表示を得る。
\end{proof}

\subsection{非零分解可能点の接空間と頂点に集中する冪零構造}

非被約性が非零のPoisson構造の近傍にも残るかどうかは、二次方程式の
総数だけからは分からない。以下では滑らかな三次に対する接空間を計算し、
被約化がGrassmann錐となる場合には、余分なスキーム構造が頂点だけに
支持されることを示す。

$T_F\in\operatorname{Sym}^3W^*$を三階微分テンソル
$T_F(v,w,u)=\partial_v\partial_w\partial_uF$とする。
非零分解可能な$P\in\Lambda^3W$の三次元supportを$A$と書き、
\begin{equation}\label{eq:nfold-transverse-kernel}
 K_{F,A}:=\{w\in W\mid T_F(a,w,u)=0
                 \text{ for all }a\in A,\ u\in W\}
\end{equation}
と置く。これは$W\to A^*\otimes W^*$、
$w\mapsto T_F(-,w,-)|_{A\times W}$の核であり、$P$の尺度に依存しない。

\begin{theorem}[非零分解可能点の法方向]\label{thm:nfold-normal-tangent}
$F$が滑らかな三次超曲面を定めるとする。任意の非零分解可能な$P$について
$A\cap K_{F,A}=0$であり、自然な同型
\begin{equation}\label{eq:nfold-normal-tangent}
 \frac{T_P\operatorname{Pois}(X_F)}{T_P\widehat G_N}
       \simeq\Lambda^3K_{F,A}
\end{equation}
が成り立つ。右辺から左辺への写像は
$\Lambda^3K_{F,A}\subset\Lambda^3W$を接ベクトルとみなすことによる。
従って
\[
 \dim T_P\operatorname{Pois}(X_F)
 =3(N-3)+1+\binom{\dim K_{F,A}}3.
\]
さらに、$P$のあるZariski近傍で$I_F=J_N$となることと
$\dim K_{F,A}\le2$は同値である。
\end{theorem}
\begin{proof}
$0\ne w\in A\cap K_{F,A}$なら$T_F(w,w,u)=0$が全ての$u\in W$について
成り立つ。三次式の同次性から$dF(w)(u)=T_F(w,w,u)/2$であるので、
$[w]\in\mathbb P(W)$は$X_F$の特異点となる。これは滑らかさに反する。
従って$K:=K_{F,A}$を含む補空間$B$を選んで$W=A\oplus B$とできる。
$P=e_1\wedge e_2\wedge e_3$となる$A$の基底を選ぶ。

Grassmann錐の$P$での接空間は
$\Lambda^3A\oplus(\Lambda^2A\otimes B)$である。
$I_F\subset J_N$だから、この接空間は全てのJacobi方程式の微分で消える。
従って法方向の代表は一意に
\[
 Q=\sum_{a=1}^3e_a\wedge\beta_a+\gamma,
 \qquad\beta_a\in\Lambda^2B,\quad\gamma\in\Lambda^3B
\]
と書ける。$B$の基底を$e_4,\ldots,e_N$とし、$b<c$を$B$の添字とする。
$\{a,a',a''\}=\{1,2,3\}$について、Jacobi二次式の
$(a',a'',b,c)$成分を$P$で微分すると、非零の符号を除いて
\begin{equation}\label{eq:nfold-normal-equations}
 T_F(e_a,q_{bc},e_r)=0\quad(1\le r\le N),\qquad
 q_{bc}:=\sum_{j=1}^3\beta_j^{bc}e_j
             +\sum_{j=4}^N\gamma^{bcj}e_j
\end{equation}
を得る。実際、\eqref{eq:nfold-certificate-matrix}の積のうち、
$P$の成分が非零となるのは二つの$A$添字を持つ因子だけである。
従って\eqref{eq:nfold-normal-equations}は$q_{bc}\in K$を意味する。
$K\subset B$なので、まず全ての$\beta_j^{bc}$が零となる。
残る条件は$\gamma$の全ての二重縮約が$K$に属することであり、
これは$\gamma\in\Lambda^3K$と同値である。
最後の同値は$B=K\oplus B'$に沿う外積の分解で、$B'$を含む成分を
二重縮約すれば$B'$方向のベクトルが得られることから従う。

逆に$Q=\gamma\in\Lambda^3K$なら、上の$(2A,2B)$成分の微分は零である。
残りの可能な$(3A,1B)$成分の微分は
$\sum_{a=1}^3\sum_{j\in B}T_F(e_a,e_j,e_r)\beta_a^{bj}$であり、
$\beta=0$なので零である。$A$添字が一つ以下の成分の微分は、
$P$の非零成分を含めないので初めから零である。
これで\eqref{eq:nfold-normal-tangent}を得る。
$\Lambda^3K$の接空間商への写像は補空間$B$を選ばずに定まるので、
同型は自然である。

$\dim K\le2$なら、Jacobi零スキームとGrassmann錐の接空間は一致する。
$P\ne0$ではGrassmann錐は滑らかである。
周囲の正則局所環を$\mathcal R$、その極大イデアルを$\mathfrak m$とし、
$I=I_F\mathcal R\subset J=J_N\mathcal R$と置く。
滑らかな埋込みなので$J/\mathfrak mJ\to\mathfrak m/\mathfrak m^2$は単射である。
接空間の一致から$I$の像が$J/\mathfrak mJ$を生成し、
$J=I+\mathfrak mJ$を得る。有限生成加群$J/I$にNakayamaの補題を適用して
$J=I$となる。有限生成性により、この等式は$P$のあるZariski近傍で成り立つ。
逆に局所イデアルが等しければ接空間も等しいため、
$\Lambda^3K=0$、すなわち$\dim K\le2$である。
\end{proof}

\begin{lemma}[非零層の横断性を与える係数開集合]
\label{lem:nfold-generic-transverse}
$N\ge6$とする。
滑らかな三次式のうち、ある三次元部分空間$A\subset W$に対して
$\dim K_{F,A}\ge3$となるものは、三次式の係数空間内で
余次元少なくとも$3N$の閉部分集合の中に含まれる。
従って、滑らかな三次式の非空Zariski開集合上では、
全ての非零分解可能な$P$の近傍で$I_F=J_N$となる。
この局所的な結論には、一般三次の被約Poisson分類を用いない。
$N=5$では同じ局所的な結論が全ての滑らかな三次に対して成り立つ。
\end{lemma}
\begin{proof}
$N\ge6$とする。交わらない三次元部分空間の対$(A,L)$の空間は
$\operatorname{Gr}(3,W)^2$の開部分であり、次元は$6(N-3)$である。
一つの対を固定して$W=A\oplus L\oplus C$と分解する。
条件$T_F(A,L,W)=0$は、$F$の単項式のうち、
$A$と$L$の座標を少なくとも一つずつ含むものの係数が零であることを意味する。
これらは型$(2A,1L)$、$(1A,2L)$、$(1A,1L,1C)$に分かれ、
独立な線形条件の数は
\[
 \binom42\,3+3\binom42+3\cdot3\,(N-6)=9N-18
\]
である。従って、この条件を満たす$(F,A,L)$のincidenceの次元は
\[
 \dim\operatorname{Sym}^3W^*-(9N-18)+(6N-18)
 =\dim\operatorname{Sym}^3W^*-3N
\]
である。その係数空間への像の閉包の余次元も少なくとも$3N$となる。

$F$が滑らかで$\dim K_{F,A}\ge3$なら、
$K_{F,A}$内に三次元部分空間$L$を選べる。
定理\ref{thm:nfold-normal-tangent}により$A\cap L=0$であり、
$T_F(A,L,W)=0$だから、$F$は上のincidenceの像に属する。
像の閉包の補集合を滑らかな三次式の開集合と交わらせれば、
求める非空開集合を得る。そこでの局所イデアルの一致は
定理\ref{thm:nfold-normal-tangent}による。
$N=5$では$A\cap K_{F,A}=0$だけで$\dim K_{F,A}\le2$となる。
\end{proof}

\begin{remark}
このincidenceの議論は、Grassmann錐上の全非零点での局所構造を制御する。
Grassmann錐から離れたPoisson構造の不在は、この議論だけからは結論しない。
次の系では、その点に既存の被約分類を用いる。
\end{remark}

\begin{corollary}[冪零構造の支持]\label{cor:nfold-vertex-support}
$F$は滑らかで、さらに$\sqrt{I_F}=J_N$と仮定する。
このとき全ての非零分解可能な$P$について$\dim K_{F,A}\le2$である。
従って
\begin{equation}\label{eq:nfold-punctured-equality}
 \operatorname{Pois}(X_F)\setminus\{0\}
       =\widehat G_N\setminus\{0\}
\end{equation}
はスキームとしての等式であり、
\begin{equation}\label{eq:nfold-projective-equality}
 \operatorname{Proj}(R_N/I_F)\simeq\operatorname{Gr}(3,W).
\end{equation}
冪零根基$\mathcal N_F=J_N/I_F$は有限長加群である。
特に、ある整数$s$について
\begin{equation}\label{eq:nfold-nilradical-annihilator}
 \mathfrak m_0^s\mathcal N_F=0,
 \qquad \mathfrak m_0=(p_{ijk})\subset R_N
\end{equation}
となり、$I_F$の斉次極大イデアルに関する飽和は$J_N$に等しい。
これらの結論は全ての$n\ge3$の一般の滑らかな三次に適用できる。
$n\ge8$では定理\ref{thm:nfold-high-dimensional-nilpotents}により
$\mathcal N_F\ne0$なので、その支持は正確に頂点である。
\end{corollary}
\begin{proof}
ある$P$について$\dim K_{F,A}\ge3$と仮定し、
$K_{F,A}$内の三次元部分空間をsupportとする非零分解可能な
$Q\in\Lambda^3K_{F,A}$を選ぶ。
$P$と$Q$はそれぞれ全てのJacobi二次式を満たす。
その混合項は全て$T_F(A,K_{F,A},W)=0$により零となるので、
$P+tQ$も任意の$t\in\mathbb C$についてJacobi方程式を満たす。
$A\cap K_{F,A}=0$であり、$t\ne0$なら$P+tQ$のsupportは六次元である。
従って$P+tQ$は分解可能でなく、$\sqrt{I_F}=J_N$に反する。

定理\ref{thm:nfold-normal-tangent}により、$\widehat G_N$の全ての
非零点の近傍で$I_F=J_N$となる。
被約化の仮定により他の点は存在しないので、
\eqref{eq:nfold-punctured-equality}と\eqref{eq:nfold-projective-equality}を得る。
$\mathcal N_F$は有限生成加群で支持が頂点に含まれるため、有限長であり、
ある$s$について\eqref{eq:nfold-nilradical-annihilator}が成り立つ。
$I_F\subset J_N$、$J_N$の飽和性、およびこの包含から
$(I_F:\mathfrak m_0^\infty)=J_N$が従う。
一般の三次への適用には、既存稿の被約分類
\eqref{eq:nfold-reduced-cone}を用いる。
\end{proof}

\begin{remark}[相対的な非零層]\label{rem:nfold-relative-punctured}
三次式の代表を選び、\eqref{eq:nfold-reduced-cone}が成立する非空係数開集合を
$U$とする。上の結論は各ファイバーだけでなく、
$U$上の相対Jacobiスキームの非零部分にも適用できる。
実際、$J_N$は固定されたGrassmannイデアルの係数拡大であり、
$R_N/J_N$もそのイデアルも係数環上平坦である。
従って相対商加群$J_N/I$のファイバーは$J_N/I_F$である。
各非零閉点でこのファイバーが零であることから、Nakayamaの補題により
相対商加群自体が零となる。従って相対非零Poissonスキームは
$U\times(\widehat G_N\setminus\{0\})$である。
この主張は任意の係数変更で保たれる。
三次式の代表の変更に伴う直線束の捻りを戻せば、
一般三次上の非零Poisson対の全スキームは滑らかであり、
その被約化だけを取り出す必要はない。
\end{remark}

\begin{corollary}[頂点の局所コホモロジー]\label{cor:nfold-embedded-vertex}
系\ref{cor:nfold-vertex-support}の仮定の下で
\[
 H^0_{\mathfrak m_0}(R_N/I_F)=J_N/I_F
\]
が成り立つ。特に$n\ge8$の一般の滑らかな三次について、
$\mathfrak m_0$は$R_N/I_F$の埋め込まれた随伴素イデアルであり、
\[
 \operatorname{Ass}_{R_N}(R_N/I_F)=\{J_N,\mathfrak m_0\},\qquad
 \operatorname{depth}\bigl((R_N/I_F)_{\mathfrak m_0}\bigr)=0.
\]
従ってJacobiスキームは頂点でCohen--Macaulayではない。
\end{corollary}
\begin{proof}
$J_N/I_F$は有限長なので全て$\mathfrak m_0$-torsionである。
一方、$R_N/J_N$は正の次元を持つ整域であり、
$H^0_{\mathfrak m_0}(R_N/J_N)=0$である。
完全列
$0\to J_N/I_F\to R_N/I_F\to R_N/J_N\to0$から最初の等式を得る。
$n\ge8$では$J_N/I_F\ne0$なので、有限長加群の非零socleに属する元は
ちょうど$\mathfrak m_0$で消される。
従って$\mathfrak m_0$は随伴素イデアルで、頂点の深さは零である。
唯一の極小素イデアルは$J_N$であり、
$\dim(R_N/J_N)=3(N-3)+1>0$だから、これは埋め込まれた素イデアルである。
上の完全列から随伴素イデアルは$J_N$と$\mathfrak m_0$以外には存在しない。
同じ次元の正値性からCohen--Macaulay性も成り立たない。
\end{proof}

\subsection{高次数の下界と冪零部分の次数幅}

頂点に支持されることは、冪零部分が次数二だけに集中することを意味しない。
この点はGrassmann環のHilbert関数と生成二次式の数から定量化できる。
$m_N=\binom N3$、$a_N=N\binom N4-\binom N3$と置き、$d\ge2$について
\begin{align}
 h_N(d)&:=\dim(R_N/J_N)_d
 =\frac{2\prod_{i=1}^3\prod_{j=1}^d(N+j-i)}{d!(d+1)!(d+2)!},
 \label{eq:nfold-grassmann-hilbert}\\
 L_{N,d}&:=\binom{m_N+d-1}{d}-h_N(d)
            -a_N\binom{m_N+d-3}{d-2}
 \label{eq:nfold-higher-defect}
\end{align}
と定める。

\begin{proposition}[高次数の差と次数幅の非一様性]
\label{prop:nfold-higher-nilradical}
任意の滑らかな$F$と$d\ge2$について
\[
 \dim(J_N/I_F)_d\ge L_{N,d}
\]
である。特に$N\ge13$では$L_{N,3}>0$であり、$N=13$では
\[
 \dim(J_{13}/I_F)_3\ge179179.
\]
さらに、各固定した$d\ge2$について、十分大きな$N$では$L_{N,d}>0$となる。
従って、一般三次の有限長冪零根基の非零成分が現れる最高次数、
およびそのsocleが現れる最高次数は、$N$によらない定数では抑えられない。
また$\mathfrak m_0^s(J_N/I_F)=0$となる最小の$s$も一様には抑えられない。
この主張は冪零根基の冪零指数が非有界であることを意味しない。
\end{proposition}
\begin{proof}
Grassmann環の標準単項式基底\cite[\S2.4, Theorem~2.3]{BCMStandardMonomial}は、
次数$d$では形$(d,d,d)$の半標準Young表に対応する。
Weylの次元公式をこの長方形に適用すると、その分母は
\[
 \prod_{j=1}^d j(j+1)(j+2)
 =\frac{d!(d+1)!(d+2)!}{2}
\]
であり、分子は式\eqref{eq:nfold-grassmann-hilbert}の積になる。
従って
\[
 \dim(J_N)_d=\binom{m_N+d-1}{d}-h_N(d).
\]
$I_F$は二次式で生成されるので、乗法
$(R_N)_{d-2}\otimes(I_F)_2\twoheadrightarrow(I_F)_d$は全射である。
補題\ref{lem:nfold-universal-row-relations}により
$\dim(I_F)_2\le a_N$であり、これらを引けば下界を得る。

$d=3$では展開により
\[
 L_{N,3}=
 \frac{N(N-4)(N-2)(N-1)(N+1)}{25920}
 \bigl(17N^4-258N^3+631N^2-270N-360\bigr).
\]
括弧内を$N=t+13$で書くと
$17t^4+626t^3+7807t^2+34726t+21480$となる。
従って$N\ge13$で正であり、$N=13$を代入すると$179179$を得る。

一般の固定した$d\ge2$については
\[
 L_{N,d}=
 \left(\frac1{6^dd!}-\frac2{d!(d+1)!(d+2)!}\right)N^{3d}
       +O(N^{3d-1}).
\]
実際、$m_N=N^3/6+O(N^2)$であり、最後の生成式数による項は
$O(N^5N^{3d-6})=O(N^{3d-1})$である。
先頭係数は$(d+1)!(d+2)!>2\cdot6^d$により正である。
最後の不等式は$d=2$で成立し、$d$を一つ増やすと左辺の倍率が
$(d+2)(d+3)>6$となるため全ての$d\ge2$で成立する。
よって十分大きな$N$で$(J_N/I_F)_d\ne0$である。

一般の三次では$J_N/I_F$は系\ref{cor:nfold-vertex-support}の有限長冪零根基である。
その最高非零次数の成分は全ての正次数変数で消えるのでsocleに含まれる。
従って各$d$に対して最高socle次数が$d$以上となる次元が存在する。
さらに$J_N$が二次式で生成されることから、$J_N/I_F$も次数二で生成される
$R_N$-加群である。従って$(J_N/I_F)_d\ne0$なら
$\mathfrak m_0^{d-2}(J_N/I_F)\ne0$であり、極大イデアルによる消滅指数も非有界である。
ここでは座標環の正次数元による作用を調べており、
冪零根基自身の積の長さを決定してはいない。
\end{proof}

\subsection{非零階数二層と特殊な高階数構造}
$M_n=\binom{n+4}{3}-(n+2)^2$を滑らかな三次$n$-foldのモジュライの次元とする。
三次式と非零分解可能な$P$の対のモジュライは
\begin{equation}\label{eq:nfold-rank-two-dimension}
 M_n+\dim\widehat G_N=M_n+3n-2
\end{equation}
次元の滑らかなDeligne--Mumford stackであり、三次モジュライへの忘却射は
相対次元$3n-2$の滑らかな全射である。
実際、三次式の係数開上では、これは
$\mathcal O(-1)$で捻ったGrassmann錐の頂点を除いた束である。
その相対次元は$3(N-3)+1=3n-2$であり、
$\operatorname{PGL}(W)$の商を取ると上式を得る。
滑らかな三次の射影自己同型群は有限なので商はDeligne--Mumfordである。
一般三次の係数開上では注意\ref{rem:nfold-relative-punctured}により、
これは非零Poisson対の全スキームを与え、被約化を取る必要はない。
この次元はPoisson対の底の次元であり、固定した量子化パラメータで得られる
非可換Abel圏の同型類の像の次元を意味しない。
また$n\ge8$では、一般のJacobiスキームの非被約性は零Poisson構造に集中する。

全ての滑らかな三次に対するGrassmann表示は、$n\ge4$では複素点の集合としても偽である。
例えば
\begin{equation}\label{eq:nfold-fermat-rank-four}
 F=z_1^3+\cdots+z_N^3,\qquad
 P=e_1\wedge e_2\wedge e_3+e_4\wedge e_5\wedge e_6
\end{equation}
は、二つの独立なJacobian括弧の和としてPoisson構造を与え、一般階数は四である。
$P$のsupportは六次元なので分解可能でない。
$N=6$ではEuler方向が像に入る条件が
$z_1^3+z_2^3+z_3^3=z_4^3+z_5^3+z_6^3=0$であり、$X_F$の真の閉集合をなす。
$N>6$では一般点のEuler方向が残りの座標成分を持つので、やはり射影化で階数は落ちない。
従って、一般の三次に関する定理と、特殊な三次の全Poisson構造の分類を区別する必要がある。

\section{残余圏の形式変形と次元に依存する算術}
\label{sec:nf-residual}

本節では $k=\mathbb C$、$n\geq3$、$N=n+2$ とし、
$X=V(F)\subset\mathbb P^{N-1}$ を滑らかな三次超曲面とする。
残余圏と rotation の規約は
\[
 \mathcal K_X=\operatorname{Ku}(X)
   =\langle\mathcal O_X,\ldots,\mathcal O_X(n-2)\rangle^\perp,
 \qquad
 \mathsf R_X=L_{\mathcal O_X}\circ(-\otimes\mathcal O_X(1))
\]
である。dg enhancement と冪等完備化を伴う圏を用い、
圏の変形を考える際には
$\mathcal D_X=\operatorname{Ind}(\mathcal K_X)$ を用いる。
ここで扱うのは、通常の可換局所 Artin $k$-代数を基底とする変形である。
全ての変形には特殊ファイバーと $X$、または $\mathcal D_X$ との同定を付ける。

\subsection{Hochschild 計算と幾何的 Kodaira--Spencer 写像}

次数付き Jacobian 環を
\[
 J_F=k[x_1,\ldots,x_N]/(F_{x_1},\ldots,F_{x_N})
\]
と置く。負の次数成分は零と解釈する。
次の計算は matrix factorization の既知の sector 公式の三次の場合であり、
本稿の新しい Hochschild 計算とは数えない。

\begin{proposition}[三次超曲面の Hochschild 次数]
\label{prop:nf-hh}
ベクトル空間として
\begin{equation}
\label{eq:nf-hh}
 HH^q(\mathcal K_X)\simeq
 \begin{cases}
   (J_F)_{3q/2},&q\text{ が偶数},\\
   0,&q\text{ が奇数}
 \end{cases}
 \ \oplus\
 \begin{cases}
   k^2,&3\mid N,\ q=N/3,\\
   0,&\text{それ以外}
 \end{cases}
\end{equation}
である。特に $HH^q=0$ for $q<0$、$HH^0=k$、$HH^1=0$ であり、
\[
 \dim HH^2(\mathcal K_X)=
 \begin{cases}
  22,&n=4,\\
  \binom{n+2}{3},&n\neq4.
 \end{cases}
\]
幾何的変形から誘導される実際の写像
\begin{equation}
\label{eq:nf-geometric-ks}
 \eta_X:H^1(X,T_X)\longrightarrow HH^2(\mathcal K_X)
\end{equation}
は単射である。$n\neq4$ では同型、$n=4$ では余核は二次元である。
\end{proposition}

\begin{proof}
滑らかさと Euler の公式により、$F$ の affine cone の特異点は原点だけである。
graded matrix factorization による dg 比較と
\cite[Theorem~5.39 の一般式とその証明]{NFBFKKernels}、
\cite[Theorem~3.5]{NFLZSerre} を使う。
後者の表示式の parity は、cohomological degree $q$ から conormal rank を引いたものと読む。
前者の一般式では Koszul コホモロジーの偶奇が区別され、
孤立特異点の場合には次数零の Jacobian 成分だけが残る。
恒等 sector の $q$ 次成分は、$q$ が偶数なら $(J_F)_{3q/2}$、
奇数なら零である。
非恒等な $\mu_3$ の二つの元は、それぞれ原点だけを固定する。
その conormal rank は $N$、determinant の重みは $-N$ なので、
不変部分が現れる条件は
\[
 q\equiv N\pmod2,\qquad
 N+\frac32(q-N)=0
\]
となる。これは $3\mid N$ かつ $q=N/3$ と同値である。
この場合には各元から一次元が寄与し、
\eqref{eq:nf-hh} を得る。

偏微分は $N$ 個の二次式からなる正則列であるから
\[
 \operatorname{Hilb}(J_F,t)=(1+t)^N,\qquad
 \dim(J_F)_3=\binom N3.
\]
Euler 列と正規列により
$H^1(X,T_X)\simeq(J_F)_3$ である。
写像 \eqref{eq:nf-geometric-ks} の単射性には、
単なるベクトル空間の同型でなく、
\cite[Lemma~5.2]{NFPirozhkov} の幾何的写像を用いる。
同補題の入力は、直交部分が例外列で生成されることと
$X$ の Hodge diamond が対角的でないことである。
前者は上の半直交分解、後者は同論文 Lemma~5.1 により $n\geq3$ の三次の場合に成立する。
次元を比較すれば結論を得る。
四次元の場合の余核二次元は
\cite[Remark~3.8 and the proof of Proposition~3.5]{NFHuybrechtsRennemo}
にも明記されている。
\end{proof}

\begin{lemma}[三次超曲面の通常の変形空間]
\label{lem:nf-geometric-smooth}
特殊ファイバーを固定した $X$ の通常の変形群oidは、
次元 $\binom N3$ の滑らかな形式円板に同値である。
各 Artin 基底上の自己同型で特殊ファイバー上恒等なものは恒等だけである。
\end{lemma}

\begin{proof}
超曲面の制限列から、$0<i<n$ と任意の整数 $a$ について
$H^i(X,\mathcal O_X(a))=0$ である。
また $X$ は Fano なので $H^n(X,\mathcal O_X)=0$ である。
Euler 列と
\[
 0\longrightarrow T_X\longrightarrow
 T_{\mathbb P^{N-1}}|_X\longrightarrow\mathcal O_X(3)
 \longrightarrow0
\]
から $H^2(X,T_X)=0$ を得る。
線型ベクトル場が $X$ に接していれば
$\sum_i\ell_iF_{x_i}=cF$ であり、Euler の公式により
\[
 \sum_i\left(\ell_i-\frac c3x_i\right)F_{x_i}=0
\]
となる。正則な二次式列の Koszul 複体には線型係数の一次 syzygy がない。
従ってこのベクトル場は scalar であり、射影空間上では零である。
よって $H^0(X,T_X)=0$ である。
標準的な小拡大の変形・障害理論から、変形関手は滑らかで、
形式自己同型群は自明となる。
接空間は前命題の $(J_F)_3$ であり、次元は $\binom N3$ である。

この変形を埋込み付き三次として扱うこともできる。
$H^1(X,\mathcal O_X)=H^2(X,\mathcal O_X)=0$ によって
$\mathcal O_X(1)$ は同型類として小拡大へ一意に延長し、
$H^1(X,\mathcal O_X(1))=0$ によってその切断を持ち上げられる。
従って局所 Artin 基底上で射影空間への埋込みが延長し、
方程式は再び三次の一式となる。
\end{proof}

\subsection{滑らかな幾何的族からの形式認識}

次の補題は、接空間の次元一致だけから全形式変形を結論するものではない。
全ての幾何的変形を送る実際の自然変換と、
その出発点の形式的滑らかさを同時に使う。

\begin{lemma}[滑らかな源と接写像同型]
\label{lem:nf-smooth-source}
$k$ を標数零の体、$\mathfrak g$ を dg Lie 代数とし、
$H^0(\mathfrak g)=0$、$V=H^1(\mathfrak g)$ と
$H^2(\mathfrak g)$ は有限次元とする。
その形式 moduli problem を通常の局所 Artin $k$-代数へ制限し、
一次切断した群oidを $\tau_{\leq1}\mathcal F_{\mathfrak g}$ と書く。
$d=\dim V$ とし、自然変換
\[
 \phi:\operatorname{Spf}k[[t_1,\ldots,t_d]]
       \longrightarrow\tau_{\leq1}\mathcal F_{\mathfrak g}
\]
の接写像が同型なら、$\phi$ は同値である。
$H^2(\mathfrak g)=0$ は仮定しない。
\end{lemma}

\begin{proof}
標数零での形式 moduli problem と dg Lie 代数の比較
\cite[Section~2]{NFDAGX} を用い、
$H=H^\bullet(\mathfrak g)$ 上へ極小 $L_\infty$ 構造
$\ell_j$ を移す。通常の Artin 環 $A$ に対する
Maurer--Cartan 元は $x\in V\otimes\mathfrak m_A$ であり、
方程式は
\begin{equation}
\label{eq:nf-mc}
 Q(x)=\sum_{j\geq2}\frac1{j!}\ell_j(x,\ldots,x)=0
       \quad\text{in }H^2(\mathfrak g)\otimes\mathfrak m_A
\end{equation}
である。$A$ ごとにこの和は有限である。
$H^0(\mathfrak g)=0$ なので gauge による異なる $x$ の同一視はなく、
一次切断の自己同型群も自明である。
従ってこの群oidは、
$k[[V^\vee]]$ を $Q$ の成分で割った完備局所環による離散群oidで表される。

$V$ の座標を選ぶと、$\phi$ は
$x=f(t)\in V\otimes(t_1,\ldots,t_d)k[[t_1,\ldots,t_d]]$
で表される。これは全 Artin 切断で両立する自然変換から得られる。
接写像の仮定により $f$ の線型項は可逆である。
形式逆写像定理から $f$ は形式座標変換である。
また $\phi$ が全ての次数で定義されるので、
$Q(f(t))=0$ は形式冪級数の恒等式である。
$f^{-1}$ を代入すれば $Q(x)=0$ が恒等的に成り立つ。
従って上の商環は $k[[V^\vee]]$ 自身であり、
$\phi$ は形式座標変換による同値である。
この議論は、高次の $L_\infty$ 演算全部が零であるとは主張しない。
\end{proof}

以下の $\operatorname{CatDef}^{c}_{\mathcal D_X}(A)$ は、
特殊ファイバーを $\mathcal D_X$ と同定した
compactly generated な $A$-線型圏の変形空間である。
$\tau_{\leq1}$ を取ると、射は圏同値を自然同型で割ったものとなる。
自然変換の scalar 自己同型まで保持する二群oidとの同一視ではない。

\begin{theorem}[四次元を除く残余圏の通常の Artin 形式認識]
\label{thm:nf-residual-formal}
$n\geq3$、$n\neq4$ とする。
滑らかな三次超曲面 $X$ について、自然な残余操作は
通常の局所 Artin $k$-代数上の群oidの同値
\begin{equation}
\label{eq:nf-formal-recognition}
 \operatorname{Def}_X
 \ \xrightarrow{\ \sim\ }\
 \tau_{\leq1}\operatorname{CatDef}^{c}_{\mathcal D_X}
\end{equation}
を与える。
両辺の特殊ファイバーは固定している。
同型類関手は次元 $\binom{n+2}{3}$ の形式円板である。
これは全 Abel 圏の変形を中心三次対で表示する定理ではない。
\end{theorem}

\begin{proof}
$X_A$ の相対半直交分解から得られる
$\operatorname{Ind}\operatorname{Ku}(X_A)$ は、
実際の compactly generated な圏の変形を与える。
相対射影・mutation の核は有限個の線束から作られ、
通常の Artin 係数変更と両立する。
線束の延長の自明化に含まれる基底の単元は、
射影核を作る $E^\vee\boxtimes E$ では相殺される。
従って上の自然変換はその選択に依存しない。
これで \eqref{eq:nf-formal-recognition} の自然変換を得る。

その微分は \eqref{eq:nf-geometric-ks} の写像である。
実際、幾何的な一次変形の Hochschild 類は
$H^1(T_X)\to HH^2(X)$ で与えられ、
相対射影核による変形の制限は $HH^2(X)\to HH^2(\mathcal K_X)$ に対応する。
両者の合成が、上で構成した相対半直交成分の一次変形を記述する。
これは幾何的変形・Hochschild 類・admissible 射影の自然性であり、
接空間の抽象的な次元一致を代用するものではない。

$\mathcal D_X$ は proper な compact 部分圏で生成されるので
tamely compactly generated である。
\cite[Theorems~5.1.9, 5.3.16, Proposition~5.3.21]{NFDAGX}
および \cite[Sections~3.3 and~4.3]{NFBKP} により、
形式 hull への自然変換
\[
 \theta:\operatorname{CatDef}^{c}_{\mathcal D_X}
       \longrightarrow\operatorname{CatDef}^{\wedge}_{\mathcal D_X}
\]
は $(-1)$-truncated である。
従って $\pi_0$ 上単射であり、
既に存在する対象における全ての高次 homotopy group 上同型である。
通常の可換 Artin 環へ制限した hull を支配する dg Lie 代数は
\[
 H^i(\mathfrak g)=HH^{i+1}(\mathcal K_X)
\]
を満たす。
命題 \ref{prop:nf-hh} により $H^0(\mathfrak g)=0$、
幾何的接写像は同型である。
補題 \ref{lem:nf-geometric-smooth} と
補題 \ref{lem:nf-smooth-source} により
\[
 \operatorname{Def}_X
 \xrightarrow{\sim}
 \tau_{\leq1}\operatorname{CatDef}^{\wedge}_{\mathcal D_X}.
\]
この全ての対象は幾何的変形から来るので、
$\theta$ の $\pi_0$ 上の像は hull の全体である。
$(-1)$-truncated 性による単射性と
$\pi_1$ の比較を合わせると主張を得る。
\end{proof}

\begin{remark}[七次元と高次自己同型]
\label{rem:nf-seven}
$n=7$ では $N=9$ であり、
$HH^3(\mathcal K_X)=k^2$ が実際に残る。
従って $HH^3=0$ を前提にする証明はこの次元に適用できない。
上の証明は滑らかな幾何的変形からの全次数の自然変換を使うため、
この場合も成り立つ。
また、一般に $HH^0=k$ に対応する高次 scalar は消えない。
例えば $A=k[\epsilon]/(\epsilon^2)$ では
$\pi_2$ に $k\epsilon$ が現れる。
定理は $\tau_{\leq1}$ と通常の基底上の主張であり、
未切断の高次群oidや導来 Artin 基底について幾何変形との同値を主張しない。
\end{remark}

\begin{proposition}[四次元の形式的障害]
\label{prop:nf-fourfold-obstruction}
$n=4$ のとき、残余圏の形式 hull の同型類関手は
二十二次元の滑らかな形式円板である。
幾何的変形の二十次元形式円板はその中の滑らかな余次元二の形式部分空間となる。
従って定理 \ref{thm:nf-residual-formal} の hull 版は四次元では成立しない。
\end{proposition}

\begin{proof}
\eqref{eq:nf-hh} より $HH^\bullet(\mathcal K_X)$ は次数 $0,2,4$ にのみ存在し、
その次元は $1,22,1$ である。
従って $\mathfrak g$ の極小模型は次数 $-1,1,3$ にのみ存在する。
$j$ 項演算の次数は $2-j$ であり、
奇数次数の $j$ 個の入力に対する出力は偶数次数になる。
極小模型には偶数次数成分がないので全ての演算は零である。
通常の Artin 基底上の $\pi_0$ は二十二次元形式円板となる。
幾何的写像の接写像は命題 \ref{prop:nf-hh} により単射である。
二十個の標的座標を選んで接写像を可逆にし、形式逆写像を使うと、
写像は残り二座標の形式グラフで表される。
標的の形式座標を変更すれば、このグラフは最後の二座標が零の部分空間となる。
\end{proof}

この余次元二は、指定した非可換中心三次族の残余圏が
非幾何的になることを示すものではない。
その族の形式写像が上の幾何的部分空間に入ることは、
中心性や matrix factorization の追加構造によって起こり得る。
また、ここで二十二次元と計算したのは明示した形式 hull である。

\subsection{rotation、Serre 安定性と三を法とする場合分け}

\begin{proposition}[rotation を回復する算術]
\label{prop:nf-rotation}
$S_X$ を $\mathcal K_X$ の Serre 自己同値とすると、
\begin{equation}
\label{eq:nf-serre-rotation}
 \mathsf R_X^3\simeq[2],\qquad
 S_X\simeq\mathsf R_X^{-(n-1)}[n]
          \simeq\mathsf R_X^{-N}[N]
\end{equation}
である。従って
\begin{equation}
\label{eq:nf-rotation-cases}
 \begin{array}{c|c|c}
 n&S_X&\mathsf R_X\\ \hline
 3q&\mathsf R_X[q]&S_X[-q]\\
 3q+1&[q+1]&\text{Serre と整数 shift の積では表せない}\\
 3q+2&\mathsf R_X^{-1}[q+2]&S_X^{-1}[q+2]
 \end{array}
\end{equation}
となる。中段の否定は、他の圏論的構造を使った rotation の復元まで否定するものではない。
\end{proposition}

\begin{proof}
\eqref{eq:nf-serre-rotation} は graded matrix factorization の周期性と
Serre 公式である
\cite[Section~3.3, Corollary~3.18, Lemma~3.19]{NFKuznetsovFCY},
\cite[proof of Proposition~3.7]{NFLZSerre}。
$n$ を三で割った余りに従って指数を整理すると表を得る。
中段では Serre と整数 shift が生成する群は整数 shift だけからなる。
もし $\mathsf R_X\simeq[a]$ なら $\mathsf R_X^3\simeq[2]$ から $[3a]\simeq[2]$ となる。
非零 proper 圏において shift は無限位数である。
実際、非零対象が非零 shift と同型なら、その自己 Ext は無限に多くの次数で非零となり、
proper 性に反する。従って $3a=2$ が必要となり、矛盾する。
\end{proof}

rotation の算術的回復と、それを用いる categorical Torelli は
\cite[Proposition~3.7, Theorems~4.1 and~5.1]{NFLZSerre}
の既知結果である。本稿で必要なのは、これを残余射の構成に使う際、
三次元で自動的だったどの条件を一般次元で記憶すべきかを明示することである。

\begin{proposition}[rotation に安定な像の生成]
\label{prop:nf-rotation-generation}
$\mathcal T\subset\mathcal K_X$ を非零の thick 部分圏とする。
$\mathsf R_X$ とその逆で安定なら $\mathcal T=\mathcal K_X$ である。
従って $n\not\equiv1\pmod3$ なら、
$S_X$ とその逆に安定な非零 thick 部分圏も全体に等しい。
特に、冪等完備な非零三角圏から $\mathcal K_X$ への完全忠実関手の像が
該当する自己同値で安定なら、その関手は同値である。
\end{proposition}

\begin{proof}
graded Orlov 比較で $\mathcal K_X$ を
$k[x_1,\ldots,x_N]/(F)$ の graded 特異圏と同定する。
任意の非零対象 $M$ に対し
$M\oplus\mathsf R_XM\oplus\mathsf R_X^2M$ は全体を生成する
\cite[Theorem~5.15]{NFBFKOrlov}。
これは次数三・孤立特異点の場合への既知定理の直接の適用である。
従って $\mathcal T$ は生成対象を含み、全体に等しい。
Serre に関する主張は \eqref{eq:nf-rotation-cases} に従う。
完全忠実関手については、冪等完備性により像が thick となるので同じ議論を使う。
像の Serre 安定性そのものはこの命題の仮定であり、
完全忠実性だけからは結論しない。
\end{proof}

\begin{remark}[四次元での Serre 安定性の不足]
\label{rem:nf-serre-insufficient}
滑らかな Pfaffian 三次四様体 $X$ には、
K3 曲面 $Y$ と同値 $\mathcal K_X\simeq D^b(Y)$ がある
\cite[Theorem~3.1]{NFKuznetsovCubicFourfolds}。
$y\in Y$ に支持された複体全体は、非零で真の thick 部分圏である。
これは $S=[2]$ に安定である。
従って四次元では「Serre に安定な非零 thick 部分圏は全体」という命題自体が偽である。
rotation 安定性を仮定する前命題は、この次元でも成立する。
\end{remark}

\subsection{既存の一般点復元と新しいモジュライ問題の境界}

研究稿 \cite{NFNCCubicB} の主定理は既に任意の $n\geq3$ を対象としている。
その入力は、$R=\mathbb C[h]_{(h)}$ 上の次数別有限自由な中心三次対
$(B_R,w_R)$ で、
特殊ファイバーが多項式環と滑らかな三次式になることである。
同稿の定理を入力すれば、一般点の全 coherent 残余圏に、
rotation と中心元からの指定三次周期を保つ可換模型が存在する。
方程式は $\mathbb C(h)$ 上、同値核はある有限拡大上に得られ、
幾何学的な射影同型類は一意である。
これは本研究で初めて得る一般次元化ではなく、
同稿の未公表の一般点復元定理を用いる結論として扱う。

以上から、一般次元のモジュライ研究では次を分ける必要がある。
第一に、線型偏極を記憶した中心三次対の族では rotation が与えられる。
第二に、構造対象 $\mathcal O$ だけを保つ全 Abel 同値へ移るには、
例外列と残余構成の独立性・相対降下を別に証明する必要がある。
第三に、可換模型の存在、模型の一意性、全 Artin 形式認識は異なる入力である。
$n\neq4$ では定理 \ref{thm:nf-residual-formal} が形式認識を与えるが、
$n\equiv1\pmod3$ では rotation の Serre による回復は依然として使えない。
$n=4$ では一般の圏変形の形式認識も成立せず、
指定された中心三次族の追加構造を使う必要がある。

定理\ref{thm:point-trace-hessian}は、局所一ブロック型の良い開上で
別の構成による Poisson 三次 $X_w$ を与える。
残余復元の三次 $X_{\mathrm{res}}$ と $X_w$ の一致、および
この同じ三次に対する残余圏同値の存在と降下は、なお比較すべき問題である。
形式変形の次元一致や零ファイバーでの一致だけでは、この比較は従わない。
本節から全 $\mathcal O$-marked Abel モジュライの成分性、
残余圏による大域射の相対次元・滑らかさ・全射性を結論しない。
そのためには、残余の相対構成、rotation と三次周期、同値核の有限表示と降下が必要である。

% !TEX root = ../nc-cubic-n-fold.tex
\section{Trace--Hessian 三次と残余復元三次の比較障害}
\label{sec:nfm-residual-th-obstruction}

本節では、Trace--Hessian 写像が与える三次と、
rotation を備えた残余圏から復元される三次を区別する。
固定入力論文Bの完備移送・一般点復元定理を用いる箇所を明示する。
同入力の全証明の独立再検証を行ったという意味ではない。
特に、可換特殊ファイバーとの一致や同じ一次変形類だけでは、
二つの復元を同一視できないことを、任意の次元の一ブロック族で示す。

\subsection{残余復元が選ぶデータ}

本稿の定理\ref{cft:thm:complete}は、次数別有限自由な完備中心三次対
$(B,\Omega)$ に対し、特殊ファイバー $F$ を持つ形式三次式 $W$ と、
内部 shift および指定三次周期を保つ graded MF 同値を構成する。
同入力の補題\texttt{cft:lem:model}, \texttt{cft:lem:center}で、
formality と星積模型への厳密同型 $\phi$ を選んだ後の $W$ は
\begin{equation}
\label{eq:nfm-formality-central-inverse}
 W=C_\pi^{-1}\phi(\Omega),\qquad
 C_\pi(f)=\sum_{j\ge0}\frac1{j!}
 \mathcal U_{j+1}(f,\pi,\ldots,\pi)
\end{equation}
として定まる。従って $W\bmod h=F$ ではあるが、$W=F$ は結論ではない。
この写像は有限次数内の形式級数であり、
量子関係式から直接回復する尺度付き二次 Poisson tensor と、
formality が構成する formal Poisson tensor とも区別する。

三次周期は $\deg t=(2,-3)$ の周期変数を係数として固定することにより保存される。
したがって、potential の単元倍を行ったときには周期の規格化も追跡する必要がある。
一方、rotation の次数一・二の積から復元するのは勾配空間であり、
指定周期だけから量子中心と古典三次の Trace--Hessian 同定が強制されるわけではない。

一般点での低次数不変量は
\[
 V_\rho=H^0\operatorname{RHom}_{\mathcal K^e}(\Delta,P_\rho),\qquad
 U_\rho=H^0\operatorname{RHom}_{\mathcal K^e}
 (\Delta,P_\rho\otimes_{\mathcal K}^{\mathbf L}P_\rho),
 \qquad
 Q_\rho=\ker\bigl(\operatorname{Sym}^2V_\rho\to U_\rho\bigr)
\]
である。固定入力の定理\texttt{gcd:thm:field-model}は、
\emph{滑らかな可換模型の幾何学的存在を得た後に}この勾配空間を積分し、
方程式の定義体を元の関数体へ降ろす。
有限体拡大が不要なのは方程式であり、同値核とその周期比較には有限拡大が残る。
この降下から、方程式が $h=0$ で正則であることや剰余が $F$ であることを
一般には結論できない。

\begin{proposition}[既知の可換模型がある場合の比較判定]
\label{prop:nfm-residual-low-product-test}
$k$ を標数零の体とし、$(\mathcal K,\rho)$ は smooth proper な dg pair で、
ある代数閉拡大体上で滑らかな三次超曲面の rotation 付き Kuznetsov pair に
同値であると仮定する。
滑らかな三次式 $H\in\operatorname{Sym}^3W$ に対して、次は同値である。
\begin{enumerate}
\item 代数閉包上で $(\mathcal K,\rho)\simeq(\operatorname{Ku}(V(H)),\rho_H)$ である。
\item 代数閉包上の線型同型 $V_\rho\simeq W$ が存在し、
$Q_\rho$ を $\langle\partial_iH\rangle$ に移す。
\end{enumerate}
成立する場合、同値はある有限代数拡大上に取れる。
元の幾何的模型との同値が指定三次周期を保っていれば、
必要な有限拡大と方程式の規格化の後、その周期も保存できる。
\end{proposition}

\begin{proof}
同値から低 rotation 積を比較する向きは、固定入力の
命題\texttt{prop:ej-geometric-low-degree}である。
逆に既存の可換模型 $Y$ を取ると、第二条件の下で $Y$ と $V(H)$ の勾配空間が一致する。
滑らかな三次式の勾配空間による射影的一意性
\textup{(}同入力の命題\texttt{prop:ej-projective-uniqueness}\textup{)}により、
両者は代数閉包上で射影同型となる。
既存の残余同値と合成すればよい。
座標変換によって方程式が $cH$ に移る場合は、$c$ の三乗根で
全座標を一様に調整して potential の厳密同型にする。
これにより内部 shift と指定 MF 周期を保てる。
座標変換と同値核は、固定入力の有限降下補題\texttt{rog:lem:descent}により
有限拡大上へ降下する。
なお幾何的模型の存在の仮定を除き、低積の次元や可積分性だけから
圏同値を結論したわけではない。
\end{proof}

\subsection{次元に一様な split Hesse 族}

$r=n-1\ge2$、$R=\mathbf C[h]_{(h)}$、$K=\mathbf C(h)$ とする。
$D_R$ を
\[
 [y,z]=h x^2+\frac{h^2}{2}(yz+zy),\qquad
 [z,x]=h y^2+\frac{h^2}{2}(zx+xz),\qquad
 [x,y]=h z^2+\frac{h^2}{2}(xy+yx)
\]
で定義し、$B_R=D_R[u_1,\ldots,u_r]$ と置く。
中心変数と混同しない係数を $a_h,b_h$ と書いて
\[
 \begin{split}
 S&=x^3+y^3+z^3,\qquad P=xyz,\qquad
 G=\frac13\sum_{i=1}^r u_i^3,\qquad F_0=S/3+G,\\
 a_h&=1+\frac{h^4}{2}+\frac{h^5}{6}-\frac{h^8}{48},\qquad
 b_h=\frac{(4-h^4)(12h+h^5-8h^2)}{48},\\
 w_h&=\frac{a_h}{3}S+b_h\operatorname{sym}(xyz)+G
 \end{split}
\]
とする。
固定入力論文Bの命題\texttt{gre:prop:input}の Sklyanin 中心公式と
次数別自由性の証明は active な三変数で完結している。
従って中心変数を任意個付加しても、$B_R$ は次数別有限自由で、
$w_h$ は中心三次元、剰余 $F_0$ は滑らかな三次式である。

\begin{proposition}[Trace--Hessian 出力の厳密計算]
\label{prop:nfm-th-hesse-exact}
上の族の一般点において、Trace--Hessian 写像の $-12$ 規格化は
\[
 H_h^{\mathrm{TH}}=\mathfrak d_\Pi(w_h)=F_0+hxyz
\]
を厳密に与える。
従ってこの射の当該 trait に沿う出力は $h=0$ まで正則に延長する。
これは Trace--Hessian 構成そのものを、中心次元が跳ぶ可換特殊点全体へ
延長したという主張ではない。
\end{proposition}

\begin{proof}
尺度付き二次 tensor は $\Pi=h\Pi_{g_h}$、$g_h=S/3+hP$ である。
最大可換代数商は
\[
 T=K[u_1,\ldots,u_r]\otimes_K
 K[x,y,z]/(x^2+hyz,\ y^2+hzx,\ z^2+hxy)
\]
となる。active な正次数元による乗法は active 次数を真に上げるので、
その trace は零である。従って $\operatorname{tr}(w_{h,\mathrm{ab}})=G$、
$s_{\mathrm{cl}}(\overline g_h)=g_h$ である。
Jacobi 線上で $x^3=y^3=z^3=-hP$ だから
\[
 \epsilon(\overline w_h)=(b_h-ha_h)[P]
 =-\frac23h^2(1+h^3)[P].
\]
一方 Hessian 同型は
\[
 \nu_\Pi(\overline g_h)=8h^2(1+h^3)[P]
\]
を与える。よって
$-12\nu_\Pi^{-1}\epsilon(\overline w_h)=\overline g_h$ であり、
trace の純中心成分 $G$ を加えると主張を得る。
\end{proof}

\begin{proposition}[指定周期付き残余復元の Hesse 係数]
\label{prop:nfm-residual-hesse-input}
固定入力論文Bの完備移送定理とその明示 Hesse 計算を用いると、
上の族の残余圏は有限拡大後に
\[
 H_h^{\mathrm{res}}=F_0+\lambda_{\mathrm{res}}(h)xyz,
 \qquad
 \lambda_{\mathrm{res}}(h)
 =\frac{h+h^2/3+h^5/12}{1-h^4/4}
\]
の三次超曲面の Kuznetsov 成分と、rotation および指定三次周期を保って同値になる。
特に
\[
 H_h^{\mathrm{res}}-H_h^{\mathrm{TH}}
 =\frac{h^2}{3}xyz+O(h^3).
\]
\end{proposition}

\begin{proof}
入力の命題\texttt{gre:prop:hesse-normalform}は、formality と中心逆像を
active な三変数だけで構成し、中心多項式環上へ自明に延長する。
この証明では中心変数が二個であることを用いていない。
正規化も active 座標だけを $A(h)^{-1/3}$ 倍するもので、全ての $u_i$、
底の $h$、周期変数 $t$ を固定する。
従って任意の $r$ で同じ active Hesse 係数を得る。

同入力の定理\texttt{gre:thm:exact-potential}は、この active MF 圏を
指定した中心商の相対楕円曲線の $\operatorname{Perf}$ と比較する。
分数体上の楕円曲線の導来同値から $j$ 不変量の一致を得て、
一次項が $h$ である分枝を選ぶことにより、表示した有理式を決定する。
二次係数は、同入力の命題\texttt{gre:prop:jet}でも、
生成元を固定する harmonic-angle formality と中心写像から独立に計算されている。
この二次計算では Weyl gauge の eye 項と中心写像の inner-wheel 項を区別しており、
中心元の可換化だけを用いているのではない。

任意の $N=r+3\ge5$ について、固定三対象への制限・Laurent 局所化・生成・有限降下は
同入力の一般次元定理を適用できる。
period を固定した完備比較と上の exact potential 座標変更を用いるため、
同値は $w_h$ と表示した規格化 $H_h^{\mathrm{res}}$ の三次周期を保つ。
分母を払った方程式を使う場合は、この周期の規格化を別途保持する。
\end{proof}

\subsection{有限 jet と幾何学的一般点の障害}

\begin{lemma}[同じ低次 jet を持つ三次族の最初の差]
\label{lem:nfm-first-difference-jet}
$F$ を滑らかな三次式とし、$H,W\in k[[h]][x_1,\ldots,x_N]_3$ が
$H\bmod h=W\bmod h=F$ および $H\equiv W\pmod{h^m}$ を満たすとする。
既に固定した低次の標識の下で、$h^{m+1}$ を法とする両族の
射影座標同型が存在するための障害は
\[
 \left[\frac{W-H}{h^m}\bmod h\right]\in\operatorname{Jac}(F)_3
\]
である。これが非零なら、剰余を恒等とする形式射影同型は存在しない。
\end{lemma}

\begin{proof}
方程式の単元倍は、その形式三乗根による scalar 座標変換に吸収する。
滑らかな三次超曲面には無限小射影自己同型がない。
従って低次で同じ方程式を保つ、剰余が恒等の座標変換は、
射影 scalar を除いて $h^m$ を法として恒等でなければならない。
$I+h^mM$ による最初の変化は $\sum_{i,j}M_{ij}x_jF_{x_i}$ であり、
これは Jacobian イデアルの次数三部分全体を動く。
Euler の公式によって $F$ の scalar 倍も同じ部分に含まれる。
これで一段の存在の必要十分条件と、非零の場合の障害を得る。
\end{proof}

\begin{theorem}[Trace--Hessian と残余復元の比較障害]
\label{thm:nfm-th-residual-obstruction}
上の split Hesse 族について、次が成立する。
\begin{enumerate}
\item 二つの出力の $h^2$ の差の class は
$\tfrac13[xyz]\ne0$ in $\operatorname{Jac}(F_0)_3$ である。
従って可換特殊ファイバーの標識と底のパラメータ $h$ を保つ形式射影同型は存在しない。
\item 二つの一般ファイバーは、共通の代数閉拡大体上でも射影同型でない。
\item 固定入力の残余復元定理を用いると、非可換側の残余圏と
$\operatorname{Ku}(V(H_h^{\mathrm{TH}}))$ の間に rotation を保つ dg Morita 同値は存在しない。
さらに $n\not\equiv1\pmod3$ では、rotation を指定しない dg Morita 同値も存在しない。
\end{enumerate}
\end{theorem}

\begin{proof}
$\operatorname{Jac}(F_0)=k[x,y,z,u_1,\ldots,u_r]/(x^2,y^2,z^2,u_1^2,\ldots,u_r^2)$
なので $xyz$ の class は非零である。第一項は前補題から従う。

一般点では $g_\lambda=S/3+\lambda xyz$ の symmetrizer は scalar だけである。
実際、Hessian の対称条件は $\lambda(\lambda^3-8)\ne0$ のとき
対角成分を等しくし、非対角成分を零にする。
具体的には各 Hessian 係数行列に対する上三角三成分を並べた
九元一次方程式の係数行列に、行列式 $\lambda^2(\lambda^3-8)^2$ の
八次小行列がある。恒等行列は核に入るので、核はちょうど scalar である。
小行列の指定と厳密検算を付属コードに収録する。
$\lambda=h$ と $\lambda=\lambda_{\mathrm{res}}(h)$ はともにこの条件を満たす。
この場合の分解の一意性は、一般的な分解定理なしに確認できる。
$H(x)=T_H(x,x,x)$ に対し
\[
 E(H)=\{L:T_H(Lu,v,w)=T_H(u,Lv,w)\ \text{for all }u,v,w\}
\]
と置く。この空間は線型座標変換によって共役される。
滑らかな三次の $T_H$ は第一引数について非退化である。
実際、$T_H(a,-,-)=0$ なる $a\ne0$ があれば、その方向の座標に方程式が依存せず、
$[a]$ が特異点となる。
この非退化性の下で $E(H)$ は可換な単位的部分代数になる。
$L,M\in E(H)$ に対し、三つの引数間で作用を移すと
\[
 \begin{aligned}
 T_H(LMu,v,w)&=T_H(u,MLv,w),\\
 T_H(LMu,v,w)&=T_H(Mu,v,Lw)
             =T_H(u,Mv,Lw)=T_H(u,LMv,w).
 \end{aligned}
\]
非退化性から $LM=ML$、さらに同じ等式から $LM\in E(H)$ を得る。
Hessian を $\operatorname{Hess}(H)=\sum_i x_iH_i$ と書けば、条件は
$L^{\mathsf t}H_i=H_iL$ である。
中心変数 $u_i$ の係数行列は $2E_{ii}$ なので、各中心行の非対角成分は零である。
active な係数行列との条件から、中心列の active 成分も零となる。
実際、それらは active な全 Hessian 係数行列の共通核に入るが、
その核は滑らかさにより零である。
active ブロックには前段の scalar 条件が適用できるから
\[
 E(H)=K I_3\oplus\bigoplus_{i=1}^r K E_{u_i u_i}
 \simeq K^{r+1}.
\]
この代数の原始冪等元のうち階数三のものは active ブロックへの射影ただ一つであり、
残りは全て階数一である。この計算は任意の体拡大後にも成立する。
従って全体の射影同型があれば、唯一の階数三の像を保ち、
active 平面三次の同型を与える。

Hesse 平面三次の $j$ 不変量は、
\cite[Section~2, equations~(4)--(5)]{NFMArtebaniDolgachev} の
Weierstrass 形に $(t_0,t_1)=(1,3\lambda)$ を入れると
\[
 j(\lambda)=\frac{27\lambda^3(8-\lambda^3)^3}{(1+\lambda^3)^3}
\]
であり、厳密有理式の展開は
\[
 j(\lambda_{\mathrm{res}}(h))-j(h)
 =13824h^4+4608h^5+512h^6+O(h^7)\ne0
\]
となる。従って任意の体拡大後も active 三次は同型でなく、第二項を得る。

rotation を保つ Kuznetsov pair の同値があれば、低 rotation 積から
二つの滑らかな三次の射影同型を得るので矛盾する。
$n\not\equiv1\pmod3$ では命題\ref{prop:nf-rotation}により rotation は
Serre 関手と cohomological shift の積から回復できる。
任意の dg Morita 同値はそれらを保つため、rotation の標識を外しても同じ矛盾を得る。
$n\equiv1\pmod3$ で rotation を外した同値の不存在は、この議論からは結論しない。
\end{proof}

\begin{corollary}[Casimir の選択だけでは比較を修復できない]
\label{cor:nfm-casimir-selection-obstruction}
上の split 族で、尺度付き周囲 tensor $\Pi=h\Pi_{g_h}$ を固定する。
その三次 Casimir $H$ が滑らかな超曲面を定めるなら、
$V(H)$ は $V(H_h^{\mathrm{res}})$ と幾何学的に射影同型でない。
従って Trace--Hessian 写像の非零の数値規格化や、
純中心三次成分の選択を替えるだけでは、この例の比較を修復できない。
\end{corollary}

\begin{proof}
三次 Casimir は $H=cg_h+G'(u)$ と書ける。
滑らかさから $c\ne0$ であり、$G'$ の勾配の共通零点も原点だけである。
二つの部分の対称三重線形形式はそれぞれ第一引数について非退化なので、
前証明と同じ mixed-block の条件から
$E(H)=E(cg_h)\times E(G')$ となる。
第一因子は体 $K$ なので、$G'$ のさらなる分解の有無にかかわらず、
active 部分への射影 $(1,0)$ は階数三の原始冪等元である。
もし $V(H)\simeq V(H_h^{\mathrm{res}})$ なら、
この冪等元は後者の唯一の階数三の原始冪等元へ移る。
従って $g_h$ と $g_{\lambda_{\mathrm{res}}}$ の平面三次は同型となり、
前定理の $j$ 不一致に反する。
\end{proof}

\begin{remark}[一般性と適用範囲]
この族の ambient は中心変数への連続自己同型を持つため、
周囲代数だけに有限安定化子を要求した $T^{\mathrm{fin}}$ には属さない。
従って上の定理だけから、その特定の開上で一般対象が比較不能であるとは結論しない。
一方、量子中心で切った pair の有限安定化子条件は ambient の条件とは異なる。
Trace--Hessian 出力が滑らかなら、pair の安定化子はその滑らかな三次の
射影自己同型群へ単射するため有限である。
全同値比較と相対双対入力を保ち、ambient の有限性を要求しない
定理\ref{thm:poisson-good-open-isomorphism}の pair モジュライを用いればよい。

split 族でも、$T$ 上の $w_{h,\mathrm{ab}}$ の乗法は
$G$ 倍と active な次数を上げる冪零作用の和である。
従ってその determinant は $G^8\ne0$ で、$w_{h,\mathrm{ab}}$ は非零因子となる。
よって定常点スキームの完全交叉条件は満たす。
定常点スキーム自体の滑らかさは主張せず、上の比較にも必要ない。
PBW 性は $D_R$ の次数別自由性と中心多項式変数の付加から得られる。
点曲線の平行移動切断は
$p=(1-h^2/2:-1-h^2/2:-h)$ で、$h=0$ では原点 $e=(1:-1:0)$ に特殊化するが、
形式近傍では $p\ne e$ である。
標数零の楕円形式群の乗算 $[m]$ は線型項 $m$ を持つ可逆な形式写像なので、
任意の整数 $m>0$ について $[m]p\ne e$ である。
従って一般点の平行移動は無限位数であり、必要な有限個の位数除外を満たす。
Sklyanin の中心入力と既存の相対双対定理を適用する範囲で、
中心条件・相対双対性・$\epsilon$ の非消滅を確認できる。

混合 potential の一般点における、特殊ファイバーの標識を保つ形式比較は
次小節の有限 jet で判定できる。
一般ファイバー上で標識を外した圏同値の不存在は、
その有限 jet の計算だけからは結論しない。
\end{remark}

\subsection{混合係数の一般点における二次 jet 障害}

以下の計算では入力論文Aの Hesse 中心公式と、入力論文Bの
harmonic-angle formality による完備 MF 移送を使う。
後者の存在定理を任意の係数環へ無条件に拡張するのではなく、
各複素点における形式模型の二次 jet を同じ有限式で計算する。
その有限式自体は係数空間上の正則切断となる。

$U=\langle u_1,\ldots,u_r\rangle$、$r=n-1\ge2$ とし、
次の全混合 Hesse 係数空間を考える。
\[
 F=g_\lambda+Axy+Byz+Czx+Dx+Ey+Hz+G,\qquad
 g_\lambda=(x^3+y^3+z^3)/3+\lambda xyz,
\]
ただし $A,B,C\in U$、$D,E,H\in\Sym^2U$、$G\in\Sym^3U$ である。
$F$ と $g_\lambda$ がともに滑らかな三次を定める開を $\mathcal S_n$ とする。
$\mathscr H_F=\det\operatorname{Hess}_{x,y,z}(F)$ と書き、
\[
 T_F=\operatorname{Sym}U[x,y,z]/(F_x,F_y,F_z),\qquad
 \tau_F=\tfrac18\operatorname{Tr}_{T_F/\operatorname{Sym}U}
\]
と置く。$\tau_F$ は商への写像との合成として多項式にも適用する。
これは全混合係数と両立する $\operatorname{Sym}U$-線型写像である。
量子関係は $[a,b]=h\operatorname{sym}\Pi_F(a,b)$ とし、
入力の短い中心元から $h^2ABC/6$ を引いた規格化を
\[
 \widetilde\Omega_{F,h}
 =\operatorname{sym}\left(F-\frac{h^2}{12}\mathscr H_F
                              +h^4\Psi_\lambda(F)\right)
\]
とする。中心性・次数別 PBW 性は固定入力の該当開で用いる。

\begin{proposition}[残余復元と trace 像の普遍二次 jet]
\label{prop:nfm-relative-residual-second-jet}
Kontsevich の harmonic-angle formality と、次数一生成元を固定する
厳密同型 $\phi$ を選ぶ。
入力の中心逆像 $W^{\mathrm{res}}=C_\pi^{-1}\phi(\widetilde\Omega_{F,h})$ は
\begin{equation}\label{eq:nfm-residual-universal-jet}
 W^{\mathrm{res}}=F+\frac{h^2}{24}\mathscr H_F+O(h^3)
\end{equation}
を満たす。同じ中心元の trace--Hessian 像は
\begin{equation}\label{eq:nfm-th-universal-jet}
 H^{\mathrm{TH}}
 =F-\frac{h^2}{12}\tau_F(\mathscr H_F)
      +h^4\tau_F(\Psi_\lambda(F))
\end{equation}
である。この式は $h\ne0$ で定義した射の出力が当該 trait に沿って正則に延びたもので、
可換特殊点での最大可換商を $T_F$ と同一視するものではない。
従って差の最初の候補は
\begin{equation}\label{eq:nfm-relative-comparison-class}
 \Delta_2(F)=
 \left[\frac1{24}\mathscr H_F+
             \frac1{12}\tau_F(\mathscr H_F)\right]
 \in (J_F)_3
\end{equation}
となる。この類は $\mathcal S_n$ 上の Jacobi 束の正則切断であり、
その形成は任意の係数変更と両立する。
各 $n\ge3$ について非消滅開 $\mathcal S_n^{\mathrm{diff}}$ は非空である。
\end{proposition}

\begin{proof}
Weyl 対称化を $\sigma$ と書く。
$h\mapsto-h$ と語順の反転は反対代数を同定するので、
対称化で戻した積の一次項は $\frac12\Pi_F$ である。
入力論文Bの補題\texttt{cft:lem:model}の逐次構成を
$h^2$ を法とする恒等写像から始め、次数一成分を線型座標変更で調整すると、
$\phi\sigma=1+O(h^2)$ を保って全生成元を固定できる。
ここで $P_0^{ij}=\Pi_F(z_i,z_j)$ とし、添字は全 $N=n+2$ 変数を走る。
Dito の一般 Poisson tensor に対する二次積
\cite[Section~2, equation~(3)]{NFMDito}は、$\pi=P_0/2$ という規約で
\[
 \begin{aligned}
 B_2(P_0;f,g)={}&\frac18P_0^{ij}P_0^{kl}f_{ik}g_{jl}\\
 &+\frac1{12}P_0^{ij}(\partial_iP_0^{kl})
                  (f_{jl}g_k+f_kg_{jl})
   -\frac1{24}\Theta^{ik}f_i g_k
 \end{aligned}
\]
である。線型形式 $X=a_i z_i$ の冪と $X$ に適用すると、第一項は二階微分、
第二項は反対称性により消え、二-cycle 項だけが残る。
入力の命題\texttt{gre:prop:jet}と同じ polarization により
\[
 \phi\sigma=1+h^2\rho_2+O(h^3),\qquad
 \rho_2=-\frac1{48}\Theta^{ik}\partial_i\partial_k,
 \qquad
 \Theta^{ik}=(\partial_lP_0^{ij})(\partial_jP_0^{kl})
\]
を得る。
実際、$X=a_i z_i$ の星積冪の二次項では二-cycle 項だけが残り、
$-r(r-1)X^{r-2}\Theta^{ik}a_i a_k/48$ となる。
二次積は対称なので、量子生成元の関係との比較から
formal Poisson tensor の $h^2$ 係数も零である。

$P_0$ の非零成分は active な三添字だけを持つ。
$M=(F_{ab})_{a,b=x,y,z}$ と置けば
$P_0^{ij}=\varepsilon^{ija}F_a$ であり、縮約恒等式
\[
 \Theta^{ik}
 =\varepsilon^{ija}\varepsilon^{klb}M_{al}M_{bj}
 =-2\operatorname{adj}(M)_{ik}
\]
を得る。その他の成分は零である。
従って、$M\operatorname{adj}(M)=(\det M)I_3$ より
\[
 \rho_2F=\tfrac1{24}\operatorname{tr}(\operatorname{adj}(M)M)
        =\tfrac18\det M=\tfrac18\mathscr H_F.
\]
これは混合係数を含む対称三次式についての多項式恒等式である。

中心写像について、$\mathcal U_2(f,q)=0$ は二本の辺の反対称性による。
$\mathcal U_3(f,P_0,P_0)$ に寄与し得るグラフは inner 二輪だけであり、
その harmonic-angle 重みは零である
\cite[Section~1.3, Theorem]{NFMShoikhet}。
ここではグラフの重みの消滅だけを使い、Poisson tensor の線形性を仮定しない。
逆向き wheel の非零重みとは区別する\cite[Section~2.3.1]{NFMShoikhet}。
よって $C_\pi=1+O(h^3)$ であり、
\[
 C_\pi^{-1}\phi(\widetilde\Omega_{F,h})
  =F+h^2(-\mathscr H_F/12+\rho_2F)+O(h^3)
\]
から\eqref{eq:nfm-residual-universal-jet}を得る。
生成元を固定すること、eye 項を含む星積と同じ formality の中心写像を使うことが
この係数の規約を定めている。

trace 側では、入力の有限 Hesse 公式により
\[
 \Psi_\lambda(F)|_{U=0}
       =\lambda(8-\lambda^3)g_\lambda/48
\]
である。右辺は境界 Jacobi 環で零なので、
$\epsilon(\widetilde\Omega_{F,h})=-h^2[\det\operatorname{Hess}(g_\lambda)]/12$
は厳密等式である。
$\nu_{h\Pi_F}$ が $h^2$ を含むことによる高次項の繰り下がりは、
この規格化では生じない。
従って trace--Hessian 式の非純中心方向の係数は厳密に一であり、
\eqref{eq:nfm-th-universal-jet}を得る。
一般の $F-h^2\mathscr H_F/12+O(h^4)$ という jet だけから
この結論を出したわけではない。

有限局所自由な臨界代数の trace は係数空間上正則である。
滑らかな $F$ の勾配 Koszul 複体は次数別に普遍的に完全であり、
$(J_F)_3$ は階数 $\binom{n+2}{3}$ の局所自由束となる。
従って\eqref{eq:nfm-relative-comparison-class}は任意の係数変更と両立する正則切断である。
Fermat 点
\[
 F_*=(x^3+y^3+z^3+u_1^3+\cdots+u_r^3)/3
\]
では $\mathscr H_{F_*}=8xyz$、$\tau_{F_*}(\mathscr H_{F_*})=0$ であり、
\[
 \Delta_2(F_*)=\tfrac13[xyz]\ne0,
 \qquad J_{F_*}=k[x,y,z,u]/(x^2,y^2,z^2,u_1^2,\ldots,u_r^2).
\]
これで非消滅開の非空性を得る。
\end{proof}

\begin{corollary}[一般混合係数での標識付き形式比較の障害]
\label{cor:nfm-generic-marked-comparison-obstruction}
$F\in\mathcal S_n^{\mathrm{diff}}(\C)$ を、前命題の中心量子化・次数別 PBW・
指定 formality 模型の入力が成立する開の複素点とし、特殊ファイバーを $V(F)$ と同定する。
同じ完備 MF 入力が適用できる係数体についても、同じ主張が成立する。
二つの三次族 $V(W^{\mathrm{res}})$ と $V(H^{\mathrm{TH}})$ は、
既に $\C[h]/(h^3)$ 上で、この同定を保つ同型を持たない。
固定入力の完備 MF 移送と相対 window 同定を適用する場合、
量子中心商の残余圏と $\operatorname{Ku}(V(H^{\mathrm{TH}}))$ の間にも、
同じ $h$ を固定し特殊ファイバー上恒等となる
$\C[[h]]$-線型 dg Morita 同値は存在しない。
この第二の主張は $n=4$ を含む全 $n\ge3$ に対する、入力付きの有限 jet 障害である。
\end{corollary}

\begin{proof}
両族は $h^2$ を法として $V(F)$ の定数族であり、
小拡大 $\C[h]/(h^3)\to\C[h]/(h^2)$ 上の持上げの差は
$h^2\Delta_2(F)$ である。
$H^0(V(F),T_{V(F)})=0$ なので、特殊ファイバー上恒等な同型は
$h^2$ を法としても恒等である。
そのような同型があれば差は $(J_F)_3$ で零となり、仮定に反する。

圏については、命題\ref{prop:nf-hh}の実際の幾何的 Kodaira--Spencer 写像
\[
 \eta_F:(J_F)_3\hookrightarrow HH^2(\operatorname{Ku}(V(F)))
\]
が全ての次元で単射であり、$HH^1=0$ である。
従って恒等特殊ファイバーを持つ圏同値の一次の曖昧さはなく、
二つの幾何的残余変形の持上げの差は非零の
$h^2\eta_F(\Delta_2(F))$ である。
入力の完備 MF 同値は剰余で恒等であり、各 Artin 切断と両立するので、
量子残余変形にも同じ結論が移る。
ここで全 Artin 形式認識の全射性は用いず、
四次元でも成立する幾何的接写像の単射性だけを使う。
\end{proof}

\begin{remark}[一般性と依存範囲]
全混合 Hesse 係数空間は既約であり、Hesse 座標と中心平行移動は
一般の一ブロック表示を覆う。
従って上の非消滅開は、既存の良い一ブロック入力が成立する非空開とも交わる。
ただしここで代数的に相対化したのは二次 jet とその Jacobi 類である。
全次数の残余射を係数空間上に構成したこと、任意の非被約係数環上の
完備 MF 移送を新たに証明したこと、あるいは一般ファイバー上の
標識を外した Kuznetsov 圏が非同値であることは主張しない。
特に、有限 jet 障害と先の split 族における有理式・$j$ 不変量による
幾何学的一般点の障害を区別する。
\end{remark}

% !TEX root = ../nc-cubic-n-fold.tex
\section{残余比較のための Poisson 補正}
\label{sec:nfm-poisson-correction}

比較障害の後に残る問題は、残余圏が復元する三次と同じ三次を選ぶ
Poisson 対を、一般次元の族として構成することである。
まず全 O-同値に降下する代数的補正を与え、次に固定 formality による
全次数の形式構成と比較する。最後に分離型の族では、固定入力の下で
代数的な補正後三次への全次数の残余比較を示す。
三つの構成の適用範囲を区別し、一般混合係数での全次数の一致を仮定しない。

\subsection{Hessian と trace による代数的な補正}
\label{sec:nfm-algebraic-correction}

有限 jet の障害は、周囲 Poisson tensor を変えることで取り除ける。
ここではその二次 jet を持つ、座標に依存しない代数的な処方を構成する。
この処方の全次数が残余復元と一致するとは主張しない。

$R$ を可換 $\mathbf C$-代数、$V$ を階数 $N=n+2$ の有限局所自由加群、
$W=V^*$ とする。分解可能 Poisson 対を局所的に $(H,P)$ と表し、
$K=\operatorname{supp}(P)\subset W$、$U=\operatorname{ann}(K)\subset V$ と置く。
$H$ の定める三次超曲面と、その $\mathbb P(K)$ への制限がともに滑らかとする。
$S=\operatorname{Sym}U$ および
\[
 T_{H,K}=\operatorname{Sym}V/(\partial_kH:k\in K),\qquad
 \tau_{H,K}=\tfrac18\operatorname{Tr}_{T_{H,K}/S}
\]
と置く。ここでも trace には商への写像を合成する。
active な勾配の最高次部分が三変数の正則列をなすので、
$T_{H,K}$ は $S$ 上階数八の有限局所自由代数である。
この性質と trace は任意の基底変更と両立する。

$K$ 方向の二階微分は
$D_K^2H\in\operatorname{Sym}^2K^*\otimes V$ であり、その行列式は
$(\det K^*)^{\otimes2}\otimes\operatorname{Sym}^3V$ に属する。
従って次の三次式は $K$ の座標と体積形式の選択に依存しない。
\begin{equation}\label{eq:nfm-intrinsic-hessian-correction}
 \mathcal D_P(H)=\langle\det(D_K^2H),P\otimes P\rangle,
 \qquad
 \mathcal C_P(H)=\frac1{24}\mathcal D_P(H)
       +\frac1{12}\tau_{H,K}(\mathcal D_P(H)).
\end{equation}
例えば $P=t\,\partial_x\wedge\partial_y\wedge\partial_z$ なら
$\mathcal D_P(H)=t^2\det\operatorname{Hess}_{x,y,z}(H)$ である。
trace 項は $\operatorname{Sym}^3U$ に属する。

\begin{theorem}[良い Poisson 対の代数的 Hessian 補正]
\label{thm:nfm-algebraic-hessian-correction}
active な平面三次も滑らかな $\mathfrak P_n^{\mathrm{dec}}$ の開を
$\mathfrak P_n^{\mathrm{ell}}$ と書く。
次の処方は、出力の三次と active な平面三次がともに滑らかな非空開上で、
代数的な射
\[
 \Phi:\mathfrak P_n^{\mathrm{ell}}\dashrightarrow
             \mathfrak P_n^{\mathrm{ell}},\qquad
 (H,P)\longmapsto(H^\sharp,P),\quad H^\sharp=H+\mathcal C_P(H)
\]
を定める。これは任意の基底変更、方程式の生成元変更、
全ての射影座標同型と両立する。
補正後の強い lift は厳密に
\begin{equation}\label{eq:nfm-corrected-algebraic-tensor}
 \Pi^\sharp=\iota_{dH^\sharp}P
 =\Pi+\frac1{24}\iota_{d\mathcal D_P(H)}P
\end{equation}
であり、同じ $K$ と $U$ を持つ Poisson tensor である。

この有理射は支配的で、非空開への制限は \emph{\'etale} である。
さらに定理\ref{thm:nfm-correction-degree-nine}で、始域と終域を適切に縮めると
有限 \emph{\'etale} で階数 $9$ となることを示す。
従って定理\ref{thm:poisson-good-open-isomorphism}を介し、
$\mathfrak M_n^{\mathrm{sm}}$ の非空開から Poisson 対のモジュライへの射を与える。
必要なら定義域をさらに縮め、出力も $\mathfrak P_n^{\mathrm{good}}$ に入るようにできる。
ただしこのことは、元の非可換圏とその出力の Kuznetsov 成分の
全次数の同値を意味しない。
\end{theorem}

\begin{proof}
三次方程式を $H\mapsto cH$ とすると、強い lift の表示は
$P\mapsto c^{-1}P$ となる。
Hessian 行列式は $c^3$ 倍、$P^{\otimes2}$ は $c^{-2}$ 倍になるから
\[
 \mathcal D_{c^{-1}P}(cH)=c\mathcal D_P(H).
\]
active な勾配イデアルは変わらず、trace も変わらない。
従って $\mathcal C_{c^{-1}P}(cH)=c\mathcal C_P(H)$ であり、
出力 $H^\sharp\otimes P$ は方程式の生成元に依存しない。
Hessian の縮約と有限局所自由代数の trace は線型座標変換に共変である。
特に active 変数に中心一次式を加える frame 変更も含み、
$V\twoheadrightarrow K^*$ の分裂を選ぶ必要はない。
次数スカラーの重みも相殺するので、射影群の商 stack へ降下する。
有限局所自由性と trace の基底変更により、非平坦・非被約の場合も含めて
任意の基底変更と両立する。

局所的に $K$ の座標を選ぶと、$\Pi^\sharp$ は三変数の Jacobian 括弧で、
残りの変数は中心である。従って Jacobi 恒等式は多項式恒等式として成立し、
$H^\sharp$ はその Casimir である。
trace 項の微分は $K$ 方向で零なので、式
\eqref{eq:nfm-corrected-algebraic-tensor}を得る。
出力が滑らかなら、系\ref{cor:poisson-support-recovery}により一次中心は
ちょうど $U$ である。

定義域の非空性と支配性を一様に確認する。
滑らかな三次と active な平面三次の開
\[
 B\subset\mathbb P(\operatorname{Sym}^3V)\times\operatorname{Gr}(3,W)
\]
上で、$\mathcal L_H\otimes\det\mathcal K$ の全空間 $\overline Z$ を取り、
零切断も許す。ただし零切断では $K$ を忘れない。
式\eqref{eq:nfm-intrinsic-hessian-correction}は $P=0$ でも定義でき、
補正は線束方向の二次以上である。
したがって $\Phi$ は零切断を含む開で定義され、零切断を各点ごとに固定し、
そこでの微分は恒等である。
$\overline Z$ は滑らかなので、$\Phi$ は零切断の近傍で \emph{\'etale} となる。
この近傍の零切断を除いた部分は非空で、像は開かつ稠密である。
よって元の非零部分上の有理射は支配的であり、非空開で \emph{\'etale} である。
同じ群による同変商に移すと、この性質は stack 上にも降りる。

$\mathfrak P_n^{\mathrm{good}}$ は既約な
$\mathfrak P_n^{\mathrm{ell}}$ の非空開である。
支配性により、その逆像と定義域中の
$\mathfrak P_n^{\mathrm{good}}$ の交わりも非空開となる。
そこに既存の O-marked moduli との同型を適用すればよい。
\end{proof}

\begin{corollary}[flag を残した零切断の形式近傍]
\label{cor:nfm-correction-formal-automorphism}
上の $\overline Z$ の零切断に沿う形式完備化上で、$\Phi$ は
恒等を剰余に持つ形式自己同型である。
同じことは射影群の商にも成立する。
\end{corollary}
\begin{proof}
零切断を固定する \emph{\'etale} 射に対し、零切断上の恒等写像の
持上げは各冪零厚みで一意である。
これらを順次持ち上げると、任意の有限厚み上で相互に逆な写像を得る。
その逆極限が形式逆写像を与える。一意性により群作用とも両立する。
これは明示的な補正 $\Phi$ 自身の形式可逆性であり、
formality の中心写像との全次数の一致ではない。
\end{proof}

\begin{proposition}[楕円点曲線との厳密な整合]
\label{prop:nfm-correction-point-cubic}
良い一ブロック型の ambient 二次代数 $B_\Pi$ を考える。
$g=H|_{\mathbb P(K)}$ とすると、補正後の制限は
\[
 H^\sharp|_{\mathbb P(K)}
 =g+\frac1{24}\mathcal D_P(g)
\]
であり、これは $B_\Pi$ の active な楕円点曲線の方程式である。
従って補正射は、回復した support 平面の中で、
量子側の楕円点曲線と同じ平面三次を選ぶ。
この主張は点曲線の shift を保持した対応全体や、
定常点成分の同定を含まない。
\end{proposition}

三生成元の二次 Calabi--Yau 代数の点曲線を Hessian によって記述することは、
Mori--Smith\cite[Theorems~1.2, 7.2]{MoriSmith}の既知の結果である。
以下の証明は本稿の係数規約を確認するもので、その点曲線公式自体の新規性を主張しない。
本節ではこれを混合係数の trace と組み合わせ、
全次元の Poisson 対および O-marked moduli 上の代数的補正に用いる。

\begin{proof}
trace 項は純中心三次なので、$\mathbb P(K)$ への制限で零になる。
局所座標で $P=t\,\partial_x\wedge\partial_y\wedge\partial_z$ と書き、
$p=(x,y,z)^{\mathsf t}$ とする。二次関係の双線型化は
\[
 \left(C(p)-\frac t2\operatorname{Hess}(g)(p)\right)q=0,
 \qquad C(p)q=p\times q
\]
である。対称行列 $M$ に対する三次恒等式
\[
 \det(C(p)+sM)=s\,p^{\mathsf t}Mp+s^3\det M
\]
と $p^{\mathsf t}\operatorname{Hess}(g)(p)p=6g(p)$ より、
その行列式は
\[
 -3t\left(g+\frac{t^2}{24}\det\operatorname{Hess}(g)\right)
\]
である。非零の $P$ を持つ相対チャートでは $t$ は単元に取れる。
行列式恒等式は、$\operatorname{adj}C(p)=pp^{\mathsf t}$ と
転置による $s$ の偶数次係数の消滅からも直接分かる。
この三次が滑らかな開では、係数行列の階数は曲線上の全ての幾何点で二である。
階数が一以下なら全ての余因子が消え、行列式の全偏微分も消えるので、
滑らかさに反する。さらに $Mq=0$ に $\operatorname{adj}(M)$ を掛けると
$\det(M)q=0$ となり、射影座標の局所的な可逆成分から incidence は
行列式曲線をスキームとして経由する。
可逆な二次小行列を使えば、任意基底上で核は直線部分束となる。
従って核は任意基底変更と両立する直線部分束であり、
双線型関係の incidence はこの曲線からの射のグラフとなる。
左右の変数を交換すると $t$ の符号が変わるだけなので、
右射影の三次行列式も同じ曲線を与え、その核も直線部分束である。
従って両射影は同型で、incidence は曲線の自己同型のグラフである。
これにより三次が楕円点曲線であることも任意基底上で確認できる。
$t=0$ で得られる三次はこの族の極限であり、
可換多項式環の点スキームそのものと同一視しない。
\end{proof}

\begin{corollary}[残余復元の二次 jet を実現する代数的処方]
\label{cor:nfm-corrected-residual-jet}
命題\ref{prop:nfm-relative-residual-second-jet}の規約と入力を使う。
Hessian 正規化中心元の trace--Hessian 出力に $\Phi$ を適用すると、
\[
 H^\sharp
 =F+\frac{h^2}{24}\mathscr H_F+h^4\tau_F(\Psi_\lambda(F)),
 \qquad H^\sharp\equiv W^{\mathrm{res}}\pmod{h^4}.
\]
その強い lift は
\[
 \Pi^\sharp=h\Pi_F+\frac{h^3}{24}\Pi_{\mathscr H_F}
\]
である。従って前節の二次 Jacobi 障害を、この代数的な補正により除去できる。
\end{corollary}
\begin{proof}
trace--Hessian 出力と $F$ の差は純中心三次なので、
active な勾配、Hessian、および $\tau_F$ は変わらない。
その強い lift の表示は $P=hP_0$ である。
従って $\mathcal D_{hP_0}(H^{\mathrm{TH}})=h^2\mathscr H_F$ となり、
二つの trace 項が相殺して表示した式を得る。
tensor の式は純中心項の縮約が零であることから従う。
残余二次 jet と以下の系\ref{cor:nfm-mixed-even-jet}の偶奇性により、
合同の剰余は $h^4$ まで改善する。
この計算は $h\ne0$ の対に対する代数的処方を当該 trait に沿って延長したもので、
可換特殊点全体の O-marked moduli に非零 Poisson tensor を付与するものではない。
\end{proof}

\begin{remark}[残る比較問題]
$\Phi$ は特定の formality の全次数の逆写像として定義したものではない。
ここで示したのは、任意基底の代数的構成、全 O-同値への降下、
楕円点曲線との一致、および指定中心元に対する残余三次の二次 jet との一致である。
実際の formal Poisson tensor は scalar formal unit による規格化も持ち得る。
従ってその tensor との比較には、potential と support の比較に加えて尺度を追跡する。
全次数での $\operatorname{Ku}(V(H^\sharp))$ との同値、
またはこの処方の一意性は、二次 jet の一致からは従わない。
\end{remark}

\subsection{一般混合係数における補正射の次数}
\label{sec:nfm-correction-degree}

Hessian 補正の零切断に沿う形式逆写像は、大域の有理逆写像ではない。
この相違は分離型の族に限らず、任意次元の一般混合係数で定量化できる。
以下の次数は、同じ射影座標群による商 stack の間の表現可能射の次数である。
粗モジュライの関数体の次数を、慣性群の比較なしに言い換えてはいない。
Hesse pencil と同変性による還元は、射影的な純 Hessian 写像の次数三を扱う
Ciliberto--Ottaviani\cite[Section~5, Theorem~5.7]{NFMCilibertoOttaviani}にも現れる。
本稿で数えるのは $f\mapsto f+\det\operatorname{Hess}(f)/24$ とその混合拡張であり、
純 Hessian 写像の二回合成の次数九とも異なる。

\begin{theorem}[補正射の一般次数]
\label{thm:nfm-correction-degree-nine}
任意の $n\geq3$ に対し、代数的 Hessian--trace 補正
\[
 \Phi:\mathfrak P_n^{\mathrm{ell}}\dashrightarrow
       \mathfrak P_n^{\mathrm{ell}}
\]
の一般次数は $9$ である。
より正確には、始域と終域に非空開部分 stack を取ると、$\Phi$ は
表現可能な有限 \emph{\'etale} 射で、階数は $9$ となる。
両方の開を $\mathfrak P_n^{\mathrm{good}}$ に含めてよい。
従って既存の $\mathfrak M_n^{\mathrm{sm}}\simeq
\mathfrak P_n^{\mathrm{good}}$ を介しても、同じ次数の射を得る。
特に、この補正射自身は有理同型ではない。
\end{theorem}

\begin{proof}
まず support $K$ を固定し、その体積三重ベクトル $P_0$ を選ぶ。
局所表示 $P=sP_0$ に対し、方程式の生成元変更によって
$(H,sP_0)$ を $(F,P_0)$、$F=sH$ と同一視する。
この表示における補正は
\begin{equation}\label{eq:nfm-normalized-correction-map}
 F\longmapsto F+\frac1{24}\Delta_F
                  +\frac1{12}\tau_F(\Delta_F),
 \qquad \Delta_F=\det\operatorname{Hess}_K(F).
\end{equation}
trace は $F$ の純中心部分に依存しない。

\smallskip\noindent
\textbf{active な三次の次数。}
三次元の体積座標 $(x,y,z)$ における多項式写像
\[
 \phi_3:\operatorname{Sym}^3\mathbf C^3{}^*
   \longrightarrow\operatorname{Sym}^3\mathbf C^3{}^*,
 \qquad f\longmapsto f+\tfrac1{24}\det\operatorname{Hess}(f)
\]
を考える。$g_{p,q}=p(x^3+y^3+z^3)/3+qxyz$ とおくと
\[
 \phi_3(g_{p,q})=g_{\psi_1(p,q),\psi_2(p,q)},\qquad
 \psi(p,q)=\left(p(1-q^2/4),\ q+p^3/3+q^3/12\right).
\]
この $\psi$ は、射影平面の射
\begin{equation}\label{eq:nfm-projective-correction-chart}
 [p:q:t]\longmapsto
 [3p(4t^2-q^2):12qt^2+4p^3+q^3:12t^3]
\end{equation}
へ拡張する。$t\ne0$ では最後の成分が零でなく、$t=0$ では
$pq^2=4p^3+q^3=0$ の非零解はないので基点はない。
この射は $\mathcal O(1)$ を $\mathcal O(3)$ に引き戻す。
従って有限かつ次数 $3^2=9$ である。
無限遠直線の逆像は同じ無限遠直線なので、$\psi$ 自身も有限であり、
滑らかな二次元の始域と終域の間の有限射として平坦、階数 $9$ となる。

ここで
\[
 \vartheta:\operatorname{SL}_3\times\mathbb A^2
    \longrightarrow\operatorname{Sym}^3\mathbf C^3{}^*,
 \qquad (A,p,q)\longmapsto A\cdot g_{p,q}
\]
は支配的かつ一般に有限である。
実際、$(1,1,0)$ における微分は、$\mathfrak{sl}_3$ の非対角成分から
六つの $x_i^2x_j$、対角成分から二つの $x_i^3$ の差、$p,q$ から
$x^3+y^3+z^3$ と $xyz$ を与える。これらは十次元の三次式空間の基底である。
Hessian は $\operatorname{SL}_3$ 同変なので
\[
 \vartheta\circ(1\times\psi)=\phi_3\circ\vartheta.
\]
一般次数の乗法性により
\[
 (\deg\vartheta)\,9=(\deg\phi_3)(\deg\vartheta),
 \qquad \deg\phi_3=9.
\]
この論証は Hesse frame の有限同定の個数を仮定していない。

\smallskip\noindent
\textbf{混合係数の復元。}
補空間を一時的に選び、active 次数で
\[
 F=f_3+f_2+f_1+f_0,\qquad
 f_i\in\operatorname{Sym}^iK^*\otimes
                    \operatorname{Sym}^{3-i}U
\]
と書く。$\operatorname{Hess}_K(F)$ は $f_3+f_2$ だけに依存する。
従って $\Delta_F=\delta_3+\delta_2+\delta_1+\delta_0$ と分けると
\[
 \delta_3=\det\operatorname{Hess}(f_3),\qquad
 \delta_2=\operatorname{tr}\bigl(
   \operatorname{adj}(\operatorname{Hess}(f_3))
                       \operatorname{Hess}_K(f_2)\bigr),
\]
また $\delta_1,\delta_0$ は $f_3,f_2$ だけから決まる。
式\eqref{eq:nfm-normalized-correction-map}は次の三角形の形になる。
\begin{align*}
 f'_3&=\phi_3(f_3),\\
 f'_2&=(L_{f_3}\otimes1_U)f_2,\\
 f'_1&=f_1+\delta_1/24,\\
 f'_0&=f_0+\delta_0/24+\tau_F(\Delta_F)/12,
\end{align*}
ここで
\[
 L_f(b)=b+\tfrac1{24}\operatorname{tr}\bigl(
   \operatorname{adj}(\operatorname{Hess}(f))\operatorname{Hess}(b)\bigr)
 \quad (b\in\operatorname{Sym}^2K^*).
\]
$L_f$ は一般の $f$ で可逆である。
例えば $f=g_{p,q}$ の場合、順序 $(x^2,yz)$、$(y^2,zx)$、$(z^2,xy)$
では同じ二次正方行列の三つの直和になり、各行列は
\[
 \begin{pmatrix}
  1-q^2/12&-pq/6\\ p^2/3&1+q^2/12
 \end{pmatrix}.
 \qquad
 \det L_{g_{p,q}}
   =\left(1+\frac{8p^3q-q^4}{144}\right)^3.
\]
特に $g_{1,0}$ で行列式は $1$ である。
一般の $f'_3$ には九つの $f_3$ があり、それらすべてで $L_{f_3}$ を
可逆にするよう終域の開を縮めてよい。
各 $f_3$ を選ぶと、上式から $f_2$、次に $f_1$、最後に $f_0$ が一意に決まる。
最後の trace は $f_0$ に依存しないため、そこで新たな方程式は発生しない。
よって任意の $\dim U=n-1$ に対して一般次数は同じく $9$ である。

\smallskip\noindent
\textbf{開の選択と商。}
ここまでの計算は有限型既約な係数空間の一般点の計算である。
一般に有限な写像では、始域の真の閉集合は終域を支配できない。
従って full cubic の特異点条件や good 開の補集合を除いても、
終域を縮めれば九つの一般分枝はすべて残る。
active 三次の滑らかさにより、trace を定める次数八の代数は有限局所自由で、
trace は正則である。全三次の滑らかさ、補正後の滑らかさ、
および既存の good 条件もそれぞれ非空開である。
一般に有限な写像では、取り除いた真の閉集合の像の閉包は終域を覆わない。
したがって有限性の開を取り、分岐像およびこれらの除外集合の像を取り除くと、
有限 \emph{\'etale} 階数 $9$ の射が得られる。
階数は一般の混合係数から求めており、特殊な分離型部分空間の次数を
そのまま全空間へ移したのではない。

表示 $(H,P)$ の方程式スカラーを商にした係数空間は、
$\operatorname{Gr}(3,W)$ 上の
$\operatorname{Sym}^3V\otimes\det\mathcal K$ の非零部分の開と見られる。
実際、その点は $H\otimes P$ である。
$\operatorname{PGL}(W)$ は support $K$ に推移的に作用し、
$K$ を固定して $P_0$ を選ぶと上の $F\otimes P_0$ 表示を得る。
$K$ の安定化群の作用は、体積の変化を $F$ の倍率へ吸収した作用であり、
始域と終域で同じである。従って support を固定した計算の次数 $9$ は、
support と体積 frame を動かした係数空間でも変わらない。
構成は同じ射影座標群 $G$ に関して同変であり、
有限 \emph{\'etale} 性を持つ終域の開を $G$ 不変に選べる。
一般に $G$ 同変な $X\to Y$ に対して
\[
 [X/G]\mathbin{\times}_{[Y/G]}Y\simeq X
\]
である。従って $[X/G]\to[Y/G]$ は表現可能であり、
有限 \emph{\'etale} 性と階数 $9$ はこの忠実平坦被覆から降下する。
以上により主張を得る。
\end{proof}

\begin{remark}[Hesse pencil と楕円曲線の乗法写像]
\label{rem:nfm-degree-not-multiplication}
active な入力と出力は同じ Hesse pencil に属するが、一般には同じ楕円曲線ではない。
例えば $q=0$、$p\ne0$ の入力曲線は Fermat 曲線である一方、出力の
Hesse 係数は $q'/p'=p^2/3$ であり、その $j$ 不変量は非定数である。
従って上の次数 $9$ を、入力の楕円曲線上の $[3]$ と同一視できない。
Hesse pencil とその $j$ 関数については
\cite[Section~2, Lemma~2.1, equations~(4)--(5)]{NFMArtebaniDolgachev}を参照。
Sklyanin 中心三次を別の superpotential と見たときに現れる
$(E,p)\mapsto(E,[-2]p)$\cite[Section~2]{NFMDLL}とも異なる。
後者は同じ $E$ を保つ。
今回の次数定理は、形式逆量子化の代数化や残余圏の全次数比較を主張していない。
\end{remark}

\subsection{残余復元と両立する形式 Poisson 対}
\label{sec:nfm-corrected-formal-poisson}

前節の障害は、量子代数の二次関係から直読する tensor と、
formality の逆対応が与える tensor を区別することで解消される。
本節では harmonic-angle formality を一つ固定し、形式パラメータ $h$ と
次数一生成元を保つ模型を用いる。
次数付き二次代数と二次 Poisson tensor の形式対応そのものは
Kontsevich\cite[Section~2.3, Proposition~6]{NFMKontsevichAlg}の既知結果である。
ここで明らかにするのは、一ブロックの場合の中心一次空間の保存、
残余三次による全次数の tensor の決定、および前節の障害を打ち消す補正である。
構成は形式的なものであり、非零の複素数 $h$ への収束・代数化とは区別する。

\subsubsection{次数一を固定する模型と有限係数の再帰}

$R$ を複素代数、$V_R$ を階数 $N$ の有限自由加群とし、
$S_R=\operatorname{Sym}_R V_R$ と置く。
次数一生成元を $z_1,\ldots,z_N$ とする。
次数別に $h$ 進完備な二次代数
\[
 B_h=T_R(V_R)[[h]]\Big/
 \left(z_i z_j-z_j z_i-h\operatorname{sym}(q_{ij})\right),
 \qquad q_{ij}\in\operatorname{Sym}^2_R V_R[[h]],\quad q_{ji}=-q_{ij},
\]
を考える。完備化は各内部次数で行い、最後にその直和を取る。
PBW 条件とは、この各内部次数が $S_R[[h]]$ と同じ有限局所自由加群になることをいう。
以下で存在を用いるのは、固定入力論文Bの
補題\texttt{cft:lem:model}が適用できる完備模型、
または後述の被約有限型の普遍係数模型からの基底変更である。

\begin{lemma}[生成元を固定する正規化の一意性]
\label{lem:nfm-normalized-star-model}
固定した affine-equivariant な Kontsevich 星積\cite[Section~6.3]{NFMKontsevich}について、
$h$ を固定し、全ての次数一生成元を固定する厳密同型
\[
 \phi:B_h\xrightarrow{\sim}(S_R[[h]],*_{\pi_h})
\]
が存在するなら、二次係数の formal Poisson tensor $\pi_h$ と $\phi$ は一意である。
もとの formality 同型が次数を保ち剰余で恒等なら、
形式的な線型座標変換によりこの正規化を行える。
従って正規化された対応は、係数に依存する可逆な線型座標変換にも共変である。
\end{lemma}

\begin{proof}
もとの同型の次数一部分は、剰余で恒等な可逆行列である。
その逆による線型座標変換を標的へ合成する。
グラフ写像の線型共変性により、標的は変換した tensor の同じ Kontsevich 星積であり、
全生成元が固定される。

二つの正規化された tensor が初めて $h^m$ で異なり、
$\pi'_h-\pi_h=h^mD\pmod{h^{m+1}}$ とする。
二頂点以上のグラフにこの差を入れると、他の tensor 入力が $h$ の倍数なので
$h^{m+1}$ の倍数になる。従って星積の差は
\[
 f*_{\pi'_h}g-f*_{\pi_h}g
 =\frac{h^m}{2}D(df,dg)\pmod{h^{m+1}}.
\]
同じ二次関係を両標的で評価すると、交換子の差は
$h^mD(dz_i,dz_j)$ である。
関係式の右辺には既に $h$ が掛かっているため、その差は $h^{m+1}$ に入る。
よって全ての $i,j$ について $D(dz_i,dz_j)=0$ であり、$D=0$ となる。
これは任意の最初の相違を排除する。
$B_h$ は次数一で生成されるので、tensor が等しければ $\phi$ も等しい。
座標変換した模型も同じ生成元固定条件を満たすため、最後の共変性は一意性から従う。
\end{proof}

この議論は構成的である。$\pi_{<m}=\sum_{j<m}h^j\pi_j$ が得られたとき、
\begin{equation}
\label{eq:nfm-normalized-tensor-recursion}
 \pi_m(dz_i,dz_j)=
 -[h^m]\left(
 z_i*_{\pi_{<m}}z_j-z_j*_{\pi_{<m}}z_i
 -h\operatorname{sym}_{*_{\pi_{<m}}}(q_{ij})\right)
\end{equation}
で次の係数が定まる。右辺の星積は、この段階ではグラフによる双微分式として評価する。
最終的な Maurer--Cartan 恒等式は、formality による存在と一意性から確認する。
各固定次数で使うグラフは有限個であり、右辺は二次多項式である。
従って $\pi_m$ は、$q_{ij}$ の有限個の係数と固定したグラフ重みの有限式である。

\begin{proposition}[普遍係数上での正則性と基底変更]
\label{prop:nfm-normalized-model-basechange}
$T$ を被約な有限型複素スキームとし、その上に上記の次数別 PBW 二次代数族と
形式中心三次 $\Omega_h$ があるとする。
各複素閉点で固定入力の星積模型が存在すると仮定する。
このとき正規化された $\pi_h,\phi$ および
\[
 W_h=C_{\pi_h}^{-1}\phi(\Omega_h)
\]
は、各 $h$ 係数が $T$ 上正則な形式族として一意に構成される。
その形成は $T$ からの任意の基底変更と両立する。
基底変更先は非被約でも非平坦でもよい。
\end{proposition}

\begin{proof}
局所枠で\eqref{eq:nfm-normalized-tensor-recursion}を使う。
各係数は有限式なので正則であり、前補題の共変性により局所枠の間で貼り合わさる。
各閉点で再帰は formality の正規化された解と一致する。
従って $[\pi_h,\pi_h]=0$ の各係数は、被約有限型複素スキームの全閉点で零となり、
$T$ 上で零である。これで得られた tensor の星積は結合的である。
再帰により二次関係が成立するので、生成元を固定する代数準同型 $\phi$ が存在する。
各内部次数でこれは剰余が恒等な有限自由加群間の行列であり、
$h$ 進幾何級数によって逆を持つ。

$C_{\pi_h}$ は次数三の有限局所自由加群上で $1+O(h)$ である。
従ってその逆も係数wise正則であり、$W_h$ が正則に得られる。
中心性と formality の接写像の chain 恒等式から
$[\pi_h,W_h]=0$ となる。これは固定入力の補題\texttt{cft:lem:center}と同じ
$h$ 進の HKR 単射性の議論である。
以上は全て有限係数の恒等式と $h$ 進の行列逆であるため、
任意の基底変更後も成立する。
ここで非被約基底について用いたのは、既に構成した恒等式の基底変更であり、
その閉点だけを調べて恒等式を証明したのではない。
\end{proof}

\subsubsection{中心一次空間と一ブロック support の保存}

$U_R=\langle u_1,\ldots,u_r\rangle\subset V_R$ を中心一次空間とし、
$V_R=\langle x,y,z\rangle\oplus U_R$、$P=\partial_x\wedge\partial_y\wedge\partial_z$ とする。
この分解は証明用の局所座標であり、結論はその選択に依存しない。

\begin{lemma}[線型中心の厳密な保存]
\label{lem:nfm-formal-linear-center}
上の生成元固定模型で $u\in U_R$ が $B_h$ の中心に属するなら、
$\iota_{du}\pi_h=0$ である。
さらに任意の $f\in\operatorname{Sym}U_R[[h]]$ について
\[
 C_{\pi_h}(f)=f,
 \qquad f*_{\pi_h}a=fa=a*_{\pi_h}f
\]
が成立する。従って補正後の tensor も、元の中心変数方向を微分しない。
\end{lemma}

\begin{proof}
$h$ の次数について帰納法を使う。
$h^m$ 未満の全 tensor 係数が $U_R$ を消すと仮定する。
これらの係数のみを使う正の次数のグラフで、境界入力に線型関数 $u$ を置く。
$u$ へ辺が入れば、その辺の添字は $U_R$ 方向であり、
出発頂点の tensor のその成分は零である。
辺が一本も入らない場合のグラフの総寄与は、単位の正規化によって零である。
従って $u* a-a*u$ の $h^m$ 係数に寄与するのは
HKR 項 $\pi_m(du,da)$ だけである。
$\phi$ は $u$ を固定し、$u$ は中心元なのでこの交換子は零である。
$a=z_i$ として全成分の消滅を得る。帰納法により $\pi_h$ 全体が相対的となる。

これ以降、$f(U)$ への全 active 微分は零である。
星積の正次数グラフも同じ理由で零となる。
中心写像のグラフで関数頂点 $f(U)$ に辺が入れば零であり、
一本も入らなければその関数頂点は孤立する。
孤立頂点を忘れると積分形式は必要な次元より低い配置空間からの引戻しとなるため、
正次数の重みは零である。これは $C_{\pi_h}(1)=1$ のグラフによる証明でもある。
よって $C_{\pi_h}(f)=f$ を得る。
\end{proof}

\begin{lemma}[三次 Casimir による相対二次 tensor の決定]
\label{lem:nfm-relative-casimir-tensor}
$W_h\in R[[h]][x,y,z,u_1,\ldots,u_r]_3$ とし、
その active 偏微分 $W_x,W_y,W_z$ の Koszul 複体が普遍的に完全であるとする。
二次係数の相対 bivector $Q$ が $\iota_{dW_h}Q=0$ を満たすなら、
一意な $a\in R[[h]]$ によって
\[
 Q=a\,\iota_{dW_h}P
\]
と書ける。$a$ は座標変数に依存しない。
\end{lemma}

\begin{proof}
$Q=Q^{yz}\partial_y\wedge\partial_z+
Q^{zx}\partial_z\wedge\partial_x+Q^{xy}\partial_x\wedge\partial_y$ とし、
$\alpha=Q^{yz}dx+Q^{zx}dy+Q^{xy}dz$ と置く。
Casimir 条件は $d_KW_h\wedge\alpha=0$ である。
三つの偏微分による wedge Koszul 複体の一次完全性から
$\alpha=a\,d_KW_h$ を得る。
$\alpha$ と $d_KW_h$ の係数次数はいずれも二なので、
次数零成分を取れば $a$ は $R[[h]]$ に属する。
次数零の Koszul 単射性により一意である。
全ての座標を内部次数一と数えているため、$a$ に中心変数の依存も残らない。
\end{proof}

active 制限 $g=F|_{U=0}$ が滑らかな平面三次を定める場合、
必要な Koszul 完全性は、active 勾配による八次有限局所自由代数とその次数別分裂完全性から従う。
$h$ 進変形については各固定内部次数で持上げ、完備性により極限を取る。
この議論は、その普遍係数模型からの任意の基底変更と両立する。

\begin{theorem}[残余三次に適合する形式 Poisson 対]
\label{thm:nfm-residual-formal-poisson}
良い一ブロックの PBW 中心三次族に対し、上の固定 formality による
生成元固定模型を取る。量子関係の先頭を
$[a,b]=h\operatorname{sym}\Pi_F(a,b)$、$\Pi_F=\iota_{dF}P$ とし、
$\Omega_h\bmod h=F$ とする。
$V(F)$ と active な平面三次がともに基底上滑らかなら、一意な形式スカラー $a_h=h+O(h^2)$ が存在して
\begin{equation}
\label{eq:nfm-residual-formal-factor}
 \pi_h=a_h\,\iota_{dW_h}P,
 \qquad W_h=C_{\pi_h}^{-1}\phi(\Omega_h)
\end{equation}
となる。従って $\pi_h/h$ の線は $W_h$ と元の rank 三 support によって決まり、
$h$ を逆にした後の $(V(W_h),\pi_h|_{V(W_h)})$ は非零の一ブロック Poisson 対である。

固定入力の完備 MF 移送と相対 window 比較が適用できる場合には、同じ $W_h$ に対し
\[
 \mathcal K_{B_h/(\Omega_h)}\simeq\operatorname{Ku}(V(W_h))
\]
を得る。従ってこの形式 Poisson 対では、残余比較の対象となる三次超曲面が一致する。
この圏同値は入力の結論を適用したものであり、本節で完備 MF 定理の全証明を
独立に再検証したという意味ではない。
\end{theorem}

\begin{proof}
中心逆像の構成から $[\pi_h,W_h]=0$ であり、前補題から $\pi_h$ は相対 bivector である。
Koszul 補題を適用して\eqref{eq:nfm-residual-formal-factor}を得る。
$\pi_h\equiv h\Pi_F\pmod{h^2}$ と $W_h\bmod h=F$ を比較すると
$a_h\equiv h\pmod{h^2}$ である。
従って $a_h/h$ は単元であり、$h$ の逆元を許すと tensor は非零となる。
Jacobi 恒等式は formality 模型の性質から従うが、
相対三変数の Jacobian tensor の表示からも直接従う。
$W_h$ は滑らかな $F$ の形式変形なので、対応する各 Artin 切断とその形式的極限は滑らかな射影三次の変形である。
最後は $W_h$ の定義が固定入力の完備 MF 移送で用いる中心逆像そのものであることによる。
\end{proof}

$\pi_h/h$ の次数一 Poisson 中心は厳密に $U_R[[h]]$ である。
実際、余分な active 線型中心があれば、その定数 active 勾配が
$d_KW_h$ と wedge 零となる。
同じ Koszul 完全性は、それが係数次数 $-2$ の倍数であることを要求するため零である。
従って元の中心一次空間と rank 三 support は、形式変形の途中でも保存される。

\begin{corollary}[形式中心と Casimir 空間の線型同定]
\label{cor:nfm-formal-center-casimir}
上の仮定の下で、補正された $\pi_h$ は $B_h$ だけから定まり、
選んだ中心三次には依存しない。中心逆像は線型同型
\[
 C_{\pi_h}^{-1}\phi:
 Z(B_h)_3\xrightarrow{\sim}
 \operatorname{Cas}(\pi_h)_3
 =R[[h]]W_h\oplus\operatorname{Sym}^3 U_R[[h]]
\]
を与え、純中心三次の部分加群上で恒等である。
特に $\Omega_h$ を $b\Omega_h+G(U)$ に替えると、
$W_h$ は $bW_h+G(U)$ に変わり、$\pi_h$ は変わらない。
ここでは $h$ 進形式中心を扱っており、$h=0$ での通常の中心や、
有限 Artin 切断で新たに現れる $h$-torsion の中心元と同一視しない。
\end{corollary}

\begin{proof}
$\pi_h,\phi$ の構成は二次関係だけを使う。
formality の接写像の chain 恒等式と、その vector 場上での最低次の単射性により、
$C_{\pi_h}$ は Casimir と星積中心の間の双射である。
次数を保つので次数三に制限でき、純中心三次上の恒等性は
補題\ref{lem:nfm-formal-linear-center}による。
$h$ は $R[[h]]$ 上非零因子、$a_h/h$ は単元なので、
三次 $H$ が Casimir である条件は
$d_KH\wedge d_KW_h=0$ である。
Koszul 完全性と次数により $d_KH=c\,d_KW_h$、$c\in R[[h]]$ を得る。
標数零であるため $H-cW_h$ は active 変数に依存せず、純中心三次になる。
和が直和であることも active 勾配の Koszul 単射性から従う。
\end{proof}

\subsubsection{三次 tensor のスカラー消滅}

命題\ref{prop:nfm-relative-residual-second-jet}と同じ中心元の正規化を使い、
$F$ の係数を $h$ に関して定数とする。
$\mathscr H_F=\det\operatorname{Hess}_{x,y,z}(F)$ と書く。
二次の星積の対称性と既知の残余二次 jet を
\eqref{eq:nfm-residual-formal-factor}へ代入すると、
$a_h=h(1+c_2h^2+O(h^3))$ かつ
\begin{equation}
\label{eq:nfm-corrected-poisson-third}
 \pi_h=h\Pi_F+h^3\left(\frac1{24}\Pi_{\mathscr H_F}
                  +c_2\Pi_F\right)+O(h^4)
\end{equation}
となる。Casimir 条件に加えて普遍係数の多項式性と不変式を使うことで、
この $c_2$ を決定できる。

\begin{lemma}[三変数三次形式の二次不変式]
\label{lem:nfm-no-quadratic-cubic-invariant}
三変数三次形式の係数に関する斉次二次の
$\operatorname{SL}_3$ 不変多項式は零である。
\end{lemma}

\begin{proof}
$g=\sum_{i+j+k=3}a_{ijk}x^iy^jz^k$ と書く。
行列式一の対角行列に対する重みを比較すると、不変な斉次二次式は
\[
 Q=Aa_{210}a_{012}+Ba_{201}a_{021}
   +Ca_{120}a_{102}+Da_{111}^2
\]
という形に限られる。実際、二つの係数の添字の和は
$(2,2,2)$ でなければならない。
せん断 $x\mapsto x+ty$ と $x\mapsto x+tz$ に対する
$t=0$ での微分を、それぞれ $L_y,L_z$ とする。
\[
\begin{aligned}
 L_yQ={}&3Aa_{300}a_{012}
       +(A+2C)a_{210}a_{102}+(B+4D)a_{201}a_{111},\\
 L_zQ={}&3Ba_{300}a_{021}
       +(B+2C)a_{201}a_{120}+(A+4D)a_{210}a_{111}.
\end{aligned}
\]
これらはともに零である。まず $A=B=0$ を得て、続いて
$C=D=0$ を得る。
\end{proof}

\begin{proposition}[三次の補正スカラーの消滅]
\label{prop:nfm-third-tensor-scalar-zero}
命題\ref{prop:nfm-relative-residual-second-jet}の正規化と
定理\ref{thm:nfm-residual-formal-poisson}の仮定を用いる。
$F$ の係数は $h$ に依存せず、二次関係を厳密に
$[z_i,z_j]=h\operatorname{sym}(\Pi_F(z_i,z_j))$ とする。
harmonic-angle formality による生成元固定模型では
\begin{equation}
\label{eq:nfm-third-tensor-no-scalar}
 \pi_h=h\Pi_F+\frac{h^3}{24}\Pi_{\mathscr H_F}+O(h^4),
 \qquad a_h=h+O(h^4).
\end{equation}
特に\eqref{eq:nfm-corrected-poisson-third}の $c_2$ は零である。
この恒等式は普遍係数模型からの任意の基底変更と両立する。
\end{proposition}

\begin{proof}
まず\eqref{eq:nfm-normalized-tensor-recursion}を普遍係数で評価する。
関係の係数が $h$ に依存しない場合、$\pi_m$ はその係数に関する
斉次次数 $m$ の多項式である。
これは帰納法による。星積の $h^m$ 係数の各グラフでは、
tensor 入力の $h$ 次数の和が $m$ である。
右辺 $h\operatorname{sym}_*(\Pi_F(z_i,z_j))$ の寄与では、
tensor 入力の次数の和が $m-1$ であり、さらに関係の係数が一つ入る。
従ってどちらの寄与も係数次数 $m$ を持つ。
固定次数でグラフは有限個であり、再帰に係数による除算はない。
またこの多項式構成は線型座標変更に共変である。

中心変数を含まない普遍三次形式
\[
 g=\sum_{|\alpha|=3}a_\alpha x^\alpha,
 \qquad A=\C[a_\alpha:|\alpha|=3],
 \qquad P=\partial_x\wedge\partial_y\wedge\partial_z
\]
を考え、
\[
 D=\pi_3(g)-\frac1{24}\Pi_{\det\operatorname{Hess}(g)}
\]
と置く。$D$ の係数は $A$ の斉次三次式である。
滑らかな三次形式の稠密開では、既存の Casimir の分解と二次 jet により
$D=c\Pi_g$ である。この主張は、必要なら
$g+\sum_j u_j^3$ という滑らかな分離型模型を用いて適用できる。

$v$ を $\Pi_g$ の二次単項式係数を並べたベクトル、$d$ を
$D$ の同じ係数ベクトルとする。
$d_i v_j-d_j v_i$ は稠密開で零なので $A$ 内で恒等的に零である。
従って $\operatorname{Frac}(A)$ 上で $d=cv$ となる。
ここで $c=p/q$ を互いに素な分子・分母で書くと、$cv_i\in A$ より
$q$ は全ての $v_i$ を割る。
一方、$v$ の非零成分は普遍係数 $a_\alpha$ の非零整数倍であり、
全ての $a_\alpha$ が現れる。従ってそれらの最大公約元は $1$ である。
$A$ は一意分解整域なので $q$ は単元、すなわち $c\in A$ を得る。
この分母の除去は、判別式の補集合上の正則性だけからは従わない点に注意する。

斉次性から $c$ は係数次数二である。
$P$ は $\operatorname{SL}_3$ に不変であり、再帰と Hessian は共変なので、
$d=cv$ におけるスカラーの一意性から $c$ は
$\operatorname{SL}_3$ 不変である。
補題\ref{lem:nfm-no-quadratic-cubic-invariant}により $c=0$ となる。

次に中心変数と任意の混合係数を許す。
補題\ref{lem:nfm-formal-linear-center}により全ての tensor 入力は
active 方向だけを持つ。従ってグラフと再帰の全微分は active 方向であり、
$U=0$ への制限と可換である。
$g=F|_{U=0}$ とすれば
\[
 \pi_3(F)|_{U=0}=\pi_3(g),\qquad
 \mathscr H_F|_{U=0}=\det\operatorname{Hess}(g).
\]
既存の分解
$\pi_3(F)-\Pi_{\mathscr H_F}/24=c_2(F)\Pi_F$ のスカラーは
中心変数を含む全座標に依存しない。
制限を取ると $c_2(F)\Pi_g=0$ となり、active 勾配の
次数零 Koszul 単射性から $c_2(F)=0$ を得る。
以上の普遍恒等式を基底変更すれば、非被約基底についても同じ結論を得る。
最後に $W_h=F+h^2\mathscr H_F/24+O(h^3)$ と
$\pi_h=a_h\iota_{dW_h}P$ を比較すれば $a_h=h+O(h^4)$ となる。
\end{proof}

\begin{remark}[規格化と適用範囲]
この証明が用いる中心元の情報は二次 jet だけである。
高次の中心元を Hesse 座標で記述した際の選択は、
$D=\pi_3-\Pi_{\mathscr H}/24$ の普遍多項式性と共変性に影響しない。
一方、一般の formality への変更、$h$ の再パラメータ化、
または $h$ に依存する二次関係を追加した場合には、
同じ係数式をそのまま主張しない。

偶奇性を用いた剰余の改善は、以下の系\ref{cor:nfm-mixed-even-jet}で証明する。
\end{remark}

\subsubsection{反対積による全次数の偶奇性}

ここでは harmonic-angle の Kontsevich formality を固定する。
一般の angle form や、任意に gauge 変換した formality に同じ偶奇性を
無条件に主張するものではない。
\begin{lemma}[harmonic-angle グラフの反転]
\label{lem:nfm-harmonic-parity}
同じ bivector $\pi$ を $k$ 個入れた星積の項を $B_k(\pi;f,g)$ と書く。
このとき
\[
 B_k(\pi;g,f)=(-1)^kB_k(\pi;f,g),\qquad
 f*_{-\pi}g=g*_{\pi}f.
\]
さらに関数上の接 formality 写像について
\[
 C_{-\pi}(f)=C_{\pi}(f)
\]
が成立する。最後の等式は各奇数 $k$ に対して
$\mathcal U_{k+1}(f,\pi,\ldots,\pi)=0$ であることから従う。
\end{lemma}

\begin{proof}
Kontsevich の配置空間と向きの規約を用いる
\cite[Sections~2.2 and~6.2]{NFMKontsevich}。
上半平面の反射 $r(z)=-\overline z$ は harmonic-angle 一形式を
$r^*d\varphi=-d\varphi$ に変える。
星積では境界の二点を $0,1$ に固定した表示を使い、
$z\mapsto1-\overline z$ により両境界点のラベルを交換する。
内部点は $k$ 個なので配置空間の向きは $(-1)^k$ 倍される。
辺は $2k$ 本であり角形式の積には符号が生じない。
従って両入力を交換したグラフの重みは $(-1)^k$ 倍になり、第一式を得る。
多重線型性から第二式も従う。

中心写像の $k$ 次項では、内部点は bivector の $k$ 頂点と関数の一頂点の
合計 $k+1$ 個、境界点は零個、辺は $2k$ 本である。
配置空間 $C_{k+1,0}$ の次元は $2k$ である。
商を取る前の配置空間の向きは反射により $(-1)^{k+1}$ 倍される一方、
実平行移動と正の拡大の二次元群の向きは $-1$ 倍される。
従って商の向きは $(-1)^k$ 倍される。
ここでは境界ラベルの交換はなく、同じグラフに戻る。
角形式は $2k$ 個なのでその積は不変であり、
重み $w_\Gamma$ は $w_\Gamma=(-1)^k w_\Gamma$ を満たす。
標数零では奇数 $k$ の重みは零である。
これは二次 tensor の線型性・二次性を仮定しないグラフの対称性であり、
偶数次数の中心グラフ全てが消えるという主張ではない。
\end{proof}

\begin{proposition}[生成元固定模型の偶奇性]
\label{prop:nfm-normalized-parity}
補題\ref{lem:nfm-normalized-star-model}の存在仮定の下で、
二次関係の $q_{ij}$ は $h$ に依存しないとする。
すると正規化した tensor は
\[
 \pi_{-h}=-\pi_h
\]
を満たす。Weyl 対称化を $\sigma_h$ とすれば、線型写像
$T_h=\phi_h\sigma_h$ も $T_{-h}=T_h$ を満たす。
特に $\Omega_h=\sigma_h(F_h)$、$F_{-h}=F_h$ という規格化の形式中心元に対し、
\[
 W_h=C_{\pi_h}^{-1}\phi_h(\Omega_h)
 \quad\text{は}\quad W_{-h}=W_h
\]
を満たす。
\end{proposition}

\begin{proof}
生成元を固定して語順を逆にする写像は
$B_{-h}\simeq B_h^{\mathrm{op}}$ を与える。
前補題により $*_{-\pi_h}$ は $*_{\pi_h}$ の反対積である。
従って $B_{-h}$ の生成元固定模型を二通りに書き、
補題\ref{lem:nfm-normalized-star-model}の一意性を適用すれば
$\pi_{-h}=-\pi_h$ となる。

Weyl 対称化された語の総和は語順を逆にしても不変である。
このため同じ反対積の比較で $T_{-h}=T_h$ を得る。
また $C_{\pi_{-h}}=C_{-\pi_h}=C_{\pi_h}$ なので、
$F_h$ が偶数級数なら $C_{\pi_h}^{-1}T_h(F_h)$ も偶数級数になる。
全ての等式は各係数の有限グラフ恒等式であり、
命題\ref{prop:nfm-normalized-model-basechange}で得た普遍族からの
任意の基底変更後も成立する。
\end{proof}

\begin{corollary}[全混合係数における jet の改善]
\label{cor:nfm-mixed-even-jet}
命題\ref{prop:nfm-relative-residual-second-jet}の全混合 Hesse 族で、
同じ規格化中心元 $\widetilde\Omega_{F,h}$ を用いる。
このとき
\[
 W_h=F+\frac{h^2}{24}\mathscr H_F+h^4W_4(F)+O(h^6),
 \qquad
 W_h\equiv H^\sharp\pmod{h^4}.
\]
また形式スカラー $a_h$ も奇数級数である。
命題\ref{prop:nfm-third-tensor-scalar-zero}の $c_2=0$ と合わせると
\[
 \pi_h=h\Pi_F+\frac{h^3}{24}\Pi_{\mathscr H_F}+O(h^5)
       \equiv\Pi^\sharp\pmod{h^5}.
\]
これは全 $n\ge3$、全混合係数についての合同である。
次に決定すべき potential の係数は $W_4$ であり、
$h^3$ の混合補正は存在しない。
\end{corollary}

\begin{proof}
指定中心元は
$\sigma_h(F-h^2\mathscr H_F/12+h^4\Psi_\lambda(F))$
と書かれ、括弧内は $h$ の偶数多項式である。
前命題と既知の二次 jet を合わせる。
代数的補正は厳密に
$H^\sharp=F+h^2\mathscr H_F/24+h^4\tau_F(\Psi_\lambda(F))$
であるため、第一の合同が従う。
$\pi_h=a_h\Pi_{W_h}$ に $h\mapsto-h$ を施すと
$(a_{-h}+a_h)\Pi_{W_h}=0$ となる。
相対勾配 Koszul 複体の完全性からスカラーの作用は単射であり、
$a_{-h}=-a_h$ を得る。
最後は $c_2=0$ と tensor の奇数性による。
ここでは $W_4=\tau_F(\Psi_\lambda(F))$ とも、
有限式 $H^\sharp$ が全次数の中心逆像であるとも結論していない。
\end{proof}

\begin{remark}[中心次数の filtration]
\label{rem:nfm-central-filtration}
相対グラフの全ての微分方向は $x,y,z$ であり、中心変数を微分しない。
従って $I=(U)$ による filtration を保ち、$U=0$ への制限は
生成元固定再帰、Weyl 対称化の像、中心写像およびその逆と可換である。
さらに各グラフの中心次数は、各頂点に選んだ係数の中心次数の和である。
従って tensor の係数を $I^d/I^{d+1}$ で求めるために、
元の関係の中心次数 $d$ を超える係数は不要である。
これは $\pi_h$ の境界部分を純三変数の計算から求める根拠となるが、
その情報のみから正の中心次数を持つ全混合係数は決まらない。
中心三次自身の選択に正の中心次数の係数が含まれる場合には、
もちろんその指定係数も入力に含める。
\end{remark}

\subsubsection{五次の尺度補正と代数的処方の適用範囲}

\begin{lemma}[log-canonical 模型の厳密な逆 tensor]
\label{lem:nfm-log-canonical-inverse}
$g=xyz$ とし、二次関係を
$[x,y]=h(xy+yx)/2$ およびその巡回置換とする。
harmonic-angle formality の生成元固定模型では
\[
 \pi_h(g)=\log\left(\frac{1+h/2}{1-h/2}\right)\Pi_g
       =\left(h+\frac{h^3}{12}+\frac{h^5}{80}+O(h^7)\right)\Pi_g.
\]
\end{lemma}

\begin{proof}
Kontsevich の元の星積による
$\sigma=\sum_{i<j}\lambda_{ij}x_ix_j\partial_i\wedge\partial_j$
の量子化では、原座標に関する交換比が
$q_{ij}(t\sigma)=\exp(t\lambda_{ij})$ となる
\cite[Section~1.6; Theorem~6.9 and Remark~6.10, pp.~59--60]{NFMLindbergPym}。
ここで交換比は、トーラス同変性による
$x_i\star x_j=w_{ij}x_ix_j$ の係数の比 $w_{ij}/w_{ji}$ である。
従って抽象的な量子トーラスの同型による生成元変更は不要である。

$q_h=(1+h/2)/(1-h/2)$、$\ell(h)=\log q_h$ と置き、
この恒等式を $t\mapsto\ell(h)$ で形式的に基底変更する。
$\ell(h)\Pi_{xyz}$ の星積は、原生成元で
$x\star y=q_h y\star x$ とその巡回置換を満たす。
これは指定した $h$ を保った二次関係と同値である。
単元係数の $q$-交換関係には順序単項式の PBW 基底があり、
星積側も同じ単項式の単元倍を基底に持つ。
従って生成元を固定する準同型は、各内部次数で剰余が恒等な
$\C[[h]]$-線型同型である。
補題\ref{lem:nfm-normalized-star-model}の一意性が所望の式を与える。
ここで $\ell(h)$ は標的 tensor の係数であり、
比較の際に指定されたパラメータ $h$ 自体を取り替えたのではない。
\end{proof}

\begin{proposition}[純 active 族の五次 tensor]
\label{prop:nfm-active-fifth-covariant}
普遍三変数三次形式 $g$ に対する生成元固定 tensor を
$\pi_h(g)=\sum_{m\ge1}h^m\pi_m(g)$ とする。
$x_1=x,x_2=y,x_3=z$、$\varepsilon^{123}=1$ と規約を取り、
$t_{ijk}=\partial_i\partial_j\partial_k g/6$ とし、
\[
 S(g)=54\,t_{abc}t_{def}t_{ghi}t_{jkl}
             \varepsilon^{adg}\varepsilon^{bej}
             \varepsilon^{chk}\varepsilon^{fil}
\]
と定める。添字は $1,2,3$ を走り、反復添字について和を取る。
この $S$ は古典的な四次 Aronhold 不変式であり、
四つの volume の縮約による表示も既知である
\cite[Section~3, after Remark~3.3]{NFMOttavianiAronhold}。
ここでは $S(xyz)=1$ と規格化する。
harmonic-angle formality の生成元固定模型では
\[
 \pi_5(g)=\frac1{80}S(g)\Pi_g
\]
となる。
\end{proposition}

\begin{proof}
$\pi_m(g)$ と volume によって双対な一次形式を $\alpha_m(g)$ とする。
滑らかな $g$ に純中心 Fermat 三次を加えた分離型を用いると、
Casimir の分解と $U=0$ への制限により
$\sum_m h^m\alpha_m=a_h\,dW_h$ であり、$a_h$ は座標に依存しない。
従って $d\alpha_m=0$ である。
これは普遍係数に関する多項式恒等式なので、全ての $g$ で成立する。
$E=x\partial_x+y\partial_y+z\partial_z$ に対し
\[
 G_m(g)=\frac13\iota_E\alpha_m(g)
\]
と置くと、Cartan の公式と内部次数により $dG_m=\alpha_m$ となる。
従って $G_m$ は係数次数 $m$、変数次数三の
$\operatorname{SL}_3$ 共変式である。

この型の共変式の次元は
$\operatorname{Sym}^m(\operatorname{Sym}^3\C^3)$ における
最高ウェイト $(m+2,m-1,m-1)$ の重複度に等しい。
重みの個数を
\[
 p_m(a,b,c)=[t^mX^aY^bZ^c]
     \prod_{i+j+k=3}(1-tX^iY^jZ^k)^{-1}
\]
と置く。$\delta=(2,1,0)$ とすれば、Weyl の交代和より
最高ウェイト $\lambda$ の重複度は
\[
 \sum_{\sigma\in S_3}\operatorname{sgn}(\sigma)
             p_m(\lambda+\delta-\sigma\delta)
\]
である。$m=5$、$\lambda=(7,4,4)$ では、この六項は
\[
 41-35-30+23+21-19=1
\]
となる。
一方、上に定義した $S$ は明示的な volume 縮約により不変であり、
$S(xyz)=1$ なので非零である。
従って $g\mapsto S(g)g$ は当該一次元空間を張り、
$G_5(g)=\kappa S(g)g$ を得る。
補題\ref{lem:nfm-log-canonical-inverse}と $S(xyz)=1$ により
$\kappa=1/80$ と決まる。
\end{proof}

\begin{corollary}[混合係数の制限と代数的 tensor 補正の限界]
\label{cor:nfm-fifth-restriction-obstruction}
任意の混合三次 $F$ に対し $g=F|_{U=0}$ と置くと、
\[
 \pi_5(F)|_{U=0}=\frac1{80}S(g)\Pi_g
\]
となる。従って先の Hessian 規格化の族に対する代数的補正
\[
 \pi_h^\sharp=h\Pi_F+\frac{h^3}{24}\Pi_{\mathscr H_F}
\]
は、$S(g)\ne0$ なる稠密開では、固定した $h$ と生成元の下で
実際の形式 tensor $\pi_h$ と $h^5$ の係数から異なる。
この開は全ての $n\ge3$ で非空である。
\end{corollary}

\begin{proof}
線型中心の保存により全グラフの微分は active 方向であるため、
$U=0$ への制限と再帰は可換である。
従って純 active の公式を制限に適用できる。
代数的 tensor 補正の $h^5$ 係数は零であるが、
上の制限は $S(g)\ne0$ で零でない。
五次より低い次数の一致は系\ref{cor:nfm-mixed-even-jet}による。
実際、$g=(x^3+y^3+z^3)/3+xyz$ は滑らかで $S(g)=-7$ であり、
$F=g+\sum_j u_j^3$ は全次元で滑らかな分離型を与える。
普遍係数空間の既約性から非消滅条件は稠密開を定める。
\end{proof}

\begin{remark}[規格化の検算と混合係数]
上の四つの $\varepsilon$ の縮約は $6^4$ 個の有理数項で厳密に展開できる。
得られる $S$ は通常の単項式係数に関する整数係数二十五項であり、
\[
 S\left(\frac{x^3+y^3+z^3}{3}+\lambda xyz\right)
       =\lambda(\lambda^3-8)
\]
を満たす。
共変式の同じ重複度を $m=1,\ldots,8$ で計算すると
$1,0,1,0,1,0,2,0$ となる。
これらは上の有限生成関数による整数計数であり、
Kontsevich グラフ重みや残余同値を計算したものではない。

この結果は純 active 係数についての還元である。
中心変数を含む一般の混合三次について、
$\pi_5(F)=S(F|_{U=0})\Pi_F/80$ と結論することはできない。
active 方向の局所 jet による別の共変式が混合項に寄与する可能性は残る。
また系\ref{cor:nfm-fifth-restriction-obstruction}は固定された tensor の不一致であり、
tensor のスカラーを忘れた比較、残余三次の四次 jet、
または Kuznetsov 成分の非同値までは結論しない。
\end{remark}


\subsubsection{五次の混合項と残余四次項}

次に、五次 tensor の全混合係数を二つの普遍定数へ還元する。
まず共変式の空間を決定し、続いて既知の graph weight を使って二定数を評価する。
active 微分について
\[
 T_{ijk}=F_{ijk},\qquad M_{ij}=F_{ij},\qquad q_i=F_i
\]
と置く。$T$ は座標 $x,y,z,U$ に依存せず、$F|_{U=0}=g$ だけから決まる。
全添字は $1,2,3$ を走り、$\varepsilon^{123}=1$ とする。
次の完全縮約を定義する。
\begin{align*}
 S(T)&=\frac1{24}
 T_{adg}T_{bej}T_{chk}T_{fil}
 \varepsilon^{abc}\varepsilon^{def}\varepsilon^{ghi}\varepsilon^{jkl},\\
 I(F)&=\frac1{24}
 T_{adg}T_{bej}M_{cf}M_{hk}M_{il}
 \varepsilon^{abc}\varepsilon^{def}\varepsilon^{ghi}\varepsilon^{jkl},\\
 J(F)&=\frac1{24}
 T_{adg}T_{bej}T_{chk}M_{fi}q_l
 \varepsilon^{abc}\varepsilon^{def}\varepsilon^{ghi}\varepsilon^{jkl}.
\end{align*}
$S$ は $S(xyz)=1$ の規格化の四次 Aronhold 不変式であり、
$S(g_\lambda)=\lambda(\lambda^3-8)$ である。
$I(F),J(F)$ は全変数に関する三次式である。

\begin{proposition}[五次の混合項の二定数への還元]
\label{prop:nfm-mixed-fifth-reduction}
固定 harmonic-angle formality に依存する二つの普遍定数
$\beta,\gamma\in\mathbf C$ があり、全 $n\ge3$ の全混合係数について
\[
 \pi_5(F)=\Pi_{Q_5(F)},\qquad
 Q_5(F)=\frac1{80}S(g)F
       +\beta\bigl(I(F)-S(g)F\bigr)
       +\gamma\bigl(J(F)-S(g)F\bigr).
\]
ここで $\pi_h=h\Pi_F+h^3\Pi_{\mathscr H_F}/24+h^5\pi_5+O(h^7)$ である。
係数 $1/80$ は別途得た純三変数の五次 tensor の計算を用いる。
この入力を使わない場合も、右辺は三つの普遍定数
$\alpha S(g)F+\beta I(F)+\gamma J(F)$ の勾配へ還元できる。

\end{proposition}

\begin{proof}
生成元固定再帰の $\pi_5$ は、関係 $\Pi_F$ の係数に関する
斉次次数五の微分多項式であり、active 方向だけを微分する。
再帰の Weyl 対称化も局所 jet で書けることを確認しておく。
任意の二次式 $r$ に対して
\[
 \operatorname{sym}_*(r)=r+\frac12\sum_{a,b}
      (\partial_a\partial_b r)(z_a*z_b-z_az_b)
\]
である。$r=\Pi_F(z_i,z_j)$ なら、その active 二階微分は $T$ の成分である。
片方の添字が中心なら括弧内は零となる。
active 添字の場合、正次数の単位正規化グラフは各線型入力をちょうど一回微分するので、
括弧内も下位 tensor の jet だけで書け、裸の座標は残らない。
従って帰納的に再帰の全項が以下の局所 jet 型に入る。
各係数をある点で評価すると、使える jet は $q,M,T$ だけである。
$F$ 自身は二次関係には現れない。
$q,M,T$ の全座標次数はそれぞれ $2,1,0$ であり、
$\pi_5$ の係数は全座標に関する次数二の多項式である。
従って現れ得る型は $qT^4$ と $M^2T^3$ の二つに限られる。
固定体積形式により $\pi_5$ を一形式 $\alpha_5$ と同一視する。
$\operatorname{SL}_3$ 共変性と determinant weight を合わせると、
これら二型の共変式空間の次元はそれぞれ $1,2$ である。

この小さい表現論の計算は、次の有限係数計算で直接確認できる。
$H_{d,k}(X,Y,Z)$ を
\[
 H_{d,k}=[t^k]\prod_{a+b+c=d}(1-tX^aY^bZ^c)^{-1}
\]
と置く。$\lambda$ の重複度は
$\rho=(2,1,0)$ として次の式で与えられる。
\[
 \sum_{w\in\mathfrak S_3}\operatorname{sgn}(w)
 [X^{\lambda_1+\rho_1-(w\rho)_1}
  Y^{\lambda_2+\rho_2-(w\rho)_2}
  Z^{\lambda_3+\rho_3-(w\rho)_3}]H.
\]
二型について $\lambda=(5,4,4)$ を使うと
\[
\begin{array}{c|c|c}
H&\text{六つの符号付き係数}&\text{和}\\\hline
H_{1,1}H_{3,4}&67,-58,-55,41,42,-36&1\\
H_{2,2}H_{3,3}&158,-136,-126,94,92,-80&2
\end{array}
\]
となる。

三つの一形式 $S(g)d_KF,d_KI(F),d_KJ(F)$ はこの三次元空間に属し、
一次独立である。具体的には $T$ を $g_1$ の三階微分とし、
ある点で
\[
 q=(1,0,0),\qquad
 M=\begin{pmatrix}1&1&0\\1&0&0\\0&0&0\end{pmatrix}
\]
とすると、三つの一形式を列に持つ行列は
\[
 \begin{pmatrix}
 -7&-2/3&-10/3\\
 0&-2/3&-1\\
 0&5/3&0
 \end{pmatrix},\qquad \det=-35/3\ne0
\]
となる。ここで $q,M,T$ は独立な Taylor 係数として指定できる。
実際、一つの中心変数 $u_1$ を用いて
$F=T_{ijk}x_ix_jx_k/6+u_1M_{ij}x_ix_j/2+u_1^2q_ix_i+cu_1^3$
とし、$(x,u_1)=(0,1)$ で評価すればよい。
したがってこの独立性は全ての $n\ge3$ に有効である。
したがって三つの勾配は共変式空間全体を張る。
このため $\alpha_5=d_K(\alpha SF+\beta I+\gamma J)$ と書ける。
全 graph の有効添字は三つの active 方向に限られるため、
三定数は中心変数の個数に依存しない。

純三変数の場合、$I(g),J(g)$ はいずれも
$S(g)g$ に等しい。これも小さい共変式計算で確認できる。
$H_{3,5}$ に $\lambda=(7,4,4)$ を適用した六項は
$41,-35,-30,23,21,-19$ で、和は一である。
従って次数五・三次出力の共変式は $S(g)g$ の一次元空間である。
$g=g_1$、$(x,y,z)=(1,0,0)$ で
$S=-7$、$g=1/3$、$I=J=-7/3$ を計算すれば比例定数は一となる。
純 active 計算 $\pi_5(g)=S(g)\Pi_g/80$ から
$\alpha+\beta+\gamma=1/80$ が従い、表示された形を得る。
以上の係数計算と縮約の評価は、保存した検証コードでも有理数として確認できる。
局所 jet の多項式恒等式なので、普遍模型からの任意の基底変更とも両立する。
\end{proof}

\begin{remark}[分離型で消える二方向]
\label{rem:nfm-mixed-fifth-potential}
$I-SF$ と $J-SF$ は $U=0$ に制限すると零となる。
より一般に $F=g(x+a(U))+G(U)$ という中心平行移動された分離型では、
両者は純中心三次 $-S(g)G(U)$ に等しく、その active 勾配は零である。
従ってこの分離型の計算だけでは二つの普遍定数を区別できない。
以下では二つの混合 Taylor jet で定数を決定し、
さらに中心逆像の局所性から四次 potential とその幾何的比較類を求める。
\end{remark}

% !TEX root = ../nc-cubic.tex

\begin{proposition}[混合五次 tensor の普遍定数]
\label{prop:nfm-mixed-fifth-explicit}
命題\ref{prop:nfm-mixed-fifth-reduction}の harmonic-angle 規格化において
\[
 \beta=\frac1{60},\qquad \gamma=\frac1{240}.
\]
従って全ての $n\ge3$ の混合三次 $F$ に対し
\[
 \boxed{\quad\pi_5(F)=\frac1{240}\Pi_{\,4I(F)+J(F)-2S(g)F}\quad}
\]
となる。この式は普遍係数の多項式恒等式であり、既存の普遍 PBW 模型からの
任意の基底変更と両立する。
\end{proposition}

\begin{proof}
先の二定数への還元により、二つの独立な局所 jet で正規化再帰を評価すれば足りる。
一つの中心変数 $u$ を用いた斉次三次
\[
 F_A=xyz+\frac12uz^2,\qquad F_B=xyz+u^2z
\]
を採用する。$u=1$ とし、$p=(x,y,z)=(0,0,1)$ で評価する。
いずれも $g=xyz$、$S(g)=1$ であり、完全縮約の定義から
\[
\begin{array}{c|cc}
 &I(F)-S(g)F&J(F)-S(g)F\\\hline
 F_A&-2uz^2/3&-uz^2/6\\
 F_B&-u^2z&-2u^2z/3
\end{array}
\]
を得る。従って求める二つの成分は
\[
 \pi_5(F_A)^{xy}(p)=\frac1{80}-\frac{4\beta+\gamma}{3},\qquad
 \pi_5(F_B)^{xy}(p)=\frac1{80}-\beta-\frac23\gamma.
\]
ここでは特異な二つの三次を独立 jet の評価に用いている。
一般対象の主張を特異な例から直接推定するのではなく、
既に証明した普遍多項式の二定数表示に代入している。

$P=\Pi_F$、$R=\Pi_{\det\operatorname{Hess}_K(F)}/24$ とする。
$B_k(P^{\otimes k};f,g)$ を、先頭項が $fg+\frac12 P(df,dg)$ となる
Kontsevich 星積の $k$ 頂点部分とし、$B_k[P^{k-1},R]$ は
各頂点への $R$ の挿入の和を表す。
$q^{ij}=P(dx_i,dx_j)$ とすると、生成元固定再帰の五次部分は
\begin{align*}
 \pi_5^{ij}
 ={}&-2\left(B_5[P^5;x_i,x_j]+B_3[P^2,R;x_i,x_j]\right)\\
 &+\frac12\sum_{a,b=1}^3 \partial_a\partial_b q^{ij}
 \left(B_4[P^4;x_a,x_b]+B_2[P,R;x_a,x_b]\right).
\end{align*}
偶奇性により奇数頂点部分は交換子に二倍で現れ、
右辺の最後の項は二次式の Weyl 対称化から生じる。

この有限式の評価には、Banks--Panzer--Pym\cite[Corollary~1.3; Section~6]{NFMBanksPanzerPym}の harmonic weights を
Buring\cite[Sections~3.4--3.6]{NFMBuringGraphs}がまとめた \texttt{star5.txt} と、同じ著者資料の
\texttt{kontsevich\_2\_k.txt},
\texttt{weights\_kontsevich\_2\_k.txt} ($1\le k\le5$) を用いる。
規約は Buring の学位論文 Section~3.4.1--3.4.2, pp.~74--76 にある。
前者の $k$ 次係数は $fg+hP(df,dg)$ の規約であり、
本稿の規約では $2^k$ で割る。
頂点置換と outgoing edge の交換をまとめた multiplicity は
$m(\Gamma)=2^k k!/|\operatorname{Aut}\Gamma|$ なので、
元の weight $w(\Gamma)$ に対する本稿の係数は
$w(\Gamma)/|\operatorname{Aut}\Gamma|$ となる。
この二種類のデータの一致も、有理数演算と有限置換の列挙で確認できる。

線型入力では二つの sink の入次数がともに一でなければ項が消え、
二次 tensor では aerial vertex の入次数が三以上なら項が消える。
これらの条件で残る非零項数は $k=1,2,3,4,5$ の順に
$1,1,1,19,51$ である。
上の式の三項の評価は次の表になる。
最後の列 $W$ は $\frac12\sum_{a,b}\partial_{ab}q^{xy}(B_4+B_2)$ を表す。
\[
\begin{array}{c|ccc|c}
 F&B_5^{xy}&B_3[P^2,R]^{xy}&W&-2(B_5+B_3)+W\\\hline
 F_A&1/480&0&-1/144&-1/90\\
 F_B&1/720&0&-1/240&-1/144
\end{array}
\]
全ての項は保存した73項の有限 graph certificate と\texttt{verify\_mixed\_fifth.py}で有理数として評価される。
この計算は既知の graph weight を入力とし、その積分自体を再計算してはいない。
元の graph--weight 対との符号・自己同型群を含む照合、
公刊の二変数 log-canonical 星積の五次までの係数との照合も行う。
従って
\[
 4\beta+\gamma=\frac{17}{240},\qquad
 \beta+\frac23\gamma=\frac7{360},
\]
であり、これを解くと $\beta=1/60$、$\gamma=1/240$ を得る。
\end{proof}


\subsubsection{中心逆像の四次項と一つの幾何的障害}

以下では指定した Hessian 正規化中心元を
\[
 \widetilde\Omega_h=\operatorname{sym}
 \left(F-\frac{h^2}{12}\mathscr H_F+h^4\Psi_\lambda(F)\right)
\]
とする。$S,I,J$ の規約は命題\ref{prop:nfm-mixed-fifth-reduction}と同じである。
この中心元の規格化は四次のスカラーを決めるために必要であり、
任意の単元倍に対して同じスカラー公式を主張するものではない。

\begin{lemma}[四次中心逆像の局所共変式]
\label{lem:nfm-fourth-center-locality}
固定した harmonic-angle 模型に依存する二定数 $d,e\in\C$ があって、
全混合係数に対して
\[
 W_4(F)=\Psi_\lambda(F)+d I(F)+e J(F)
\]
となる。二定数は $n\ge3$ に依存しない。
\end{lemma}

\begin{proof}
$R_{\pi_h}=\phi\circ\operatorname{sym}$ と置く。
これは可換多項式を同じ生成元の星積で完全対称化する写像である。
Kontsevich グラフと完全対称化は active な affine 座標変更に共変であり、
$R_{\pi_h}(1)=C_{\pi_h}(1)=1$ である。
したがって正次数の係数をある点で評価すると、
入力関数を少なくとも一回微分した局所 jet の多項式で表せる。
具体的には、その点を原点へ移して入力多項式を Taylor 展開すればよい。
定数項への作用は単位の保存で零となり、
残る有限個の単項式の星積は有限個のグラフで計算される。
グラフの微分添字は全て active 方向であるので、
純中心関数への正次数の作用も零である。

$h^4$ の係数には、$R$ と $C^{-1}$ の四次項のほか、
$\pi_3=\Pi_{\mathscr H_F}/24$ を二次項へ入れた項と、
$-h^2\mathscr H_F/12$ への二次作用が含まれる。
いずれも $F$ の係数次数は五、全座標次数は三である。
中心元の $h^4\Psi_\lambda(F)$ 自身への正次数の作用は
$h^4$ の係数には寄与しない。
従って $W_4-\Psi_\lambda(F)$ は、
$q_i=F_i$, $M_{ij}=F_{ij}$, $T_{ijk}=F_{ijk}$
だけの局所多項式であり、微分されない $F$ は現れない。
次数から可能な型は $qMT^3$ と $M^3T^2$ に限られる。

この二型の $\operatorname{SL}_3$ スカラー共変式の空間はいずれも一次元である。
命題\ref{prop:nfm-mixed-fifth-reduction}の Weyl 係数公式を
最高重み $(4,4,4)$ に適用すると、六項はそれぞれ
\[
 \begin{aligned}
 H_{1,1}H_{2,1}H_{3,3}&:
       &&156,-136,-136,100,109,-92,\\
 H_{2,3}H_{3,2}&:
       &&122,-105,-105,78,84,-73
 \end{aligned}
\]
で、和はともに一である。
非零の $J,I$ がそれぞれを張るので表示を得る。
有効な微分添字は三つの active 方向に限られ、二定数は $n$ に依存しない。
\end{proof}

\begin{lemma}[Hessian 写像と Aronhold 不変式による表示]
\label{lem:nfm-classical-mixed-covariants}
$\mathscr H(F)=\det\operatorname{Hess}_K(F)$ とし、
$T=D_K^3F$, $M=D_K^2F$, $q=d_KF$ と置く。
$S(T)$ は上の規格化の Aronhold 四次式である。
このとき
\[
 \begin{aligned}
 I(F)
 &=\frac1{12}\operatorname{tr}
       \bigl(\operatorname{adj}(M)\operatorname{Hess}_K(\mathscr H(F))\bigr)
   =\frac1{12}D\mathscr H_F[\mathscr H(F)],\\
 J(F)&=\frac14(dS)_T[M\odot q],\qquad
 (M\odot q)_{ijk}
       =\frac{M_{ij}q_k+M_{ik}q_j+M_{jk}q_i}{3}.
 \end{aligned}
\]
第二式の微分は対称三階 tensor $T$ の空間上で取り、
$M\odot q$ の対称化には表示した平均を用いる。
従って $I,J$ は Hessian 写像の微分と Aronhold 四次式の微分によって
計算でき、独立した新しい不変式を導入する必要はない。
\end{lemma}

\begin{proof}
第一式の trace は $M^3T^2$ 型、第二式の右辺は $qMT^3$ 型の
スカラー共変式である。
それぞれの共変式空間は前の計数により一次元なので、
$I,J$ との比例定数を一つの非零評価で決められる。
$F=xyz$ と $(x,y,z)=(1,1,1)$ を取る。
このとき $\mathscr H(F)=2F$ であり、
\[
 \operatorname{tr}\bigl(\operatorname{adj}(M)
           \operatorname{Hess}(\mathscr H(F))\bigr)
       =2\operatorname{tr}(\operatorname{adj}(M)M)=12.
\]
一方 $I=1$ である。
また $T_{123}=1$ と添字対称性だけが非零で、
$q=(1,1,1)$、$M_{ii}=0$, $M_{ij}=1$ ($i\ne j$) となる。
$S(T)=1$、$(M\odot q)_{123}=1$ であり、
四次式を微分すると $(dS)_T[M\odot q]=4$、$J=1$ となる。
これで両比例定数が定まる。
Hessian 写像の方向微分の表示は、行列式の微分公式による。
\end{proof}

体積を局所的に固定しない表示では、同じ公式で
$\mathscr H$ を $\mathscr H_P$、$S$ を $S_P$ に置き換え、
方向微分では $P$ を固定する。
これは各式の四つの $P$ の重みと一致する。

\begin{lemma}[Hesse 中心元の四次恒等式]
\label{lem:nfm-fourth-center-symbol}
\[
 F=g_\lambda+Axy+Byz+Czx+Dx+Ey+Hz+G
\]
とし、$A,B,C$ は中心一次式、$D,E,H$ は中心二次式、$G$ は中心三次式とする。
$S=\lambda(\lambda^3-8)$ と書く。
指定中心元の四次 Weyl symbol は恒等的に
\[
 \Psi_\lambda(F)=-\frac{S}{80}F-\frac1{40}I(F)
                         +\frac1{60}J(F)+\mathcal C_4(U),
\]
ただし
\[
 \mathcal C_4(U)=
  \frac{\lambda}{60}(A^3+B^3+C^3)
  +\frac{\lambda^2}{40}ABC
  +\frac{2-\lambda^3}{180}(AH+BD+CE)
  +\frac{S}{80}G
\]
である。
\end{lemma}

\begin{proof}
固定入力論文A\cite{NCCubicA}の
系\texttt{cor:pg-hesse-hessian-form}、式\texttt{eq:pg-hesse-hessian-closed}
にある有限中心公式を $h^4$ について展開する。
$Q=Axy+Byz+Czx$, $Z_2=Az^2+Bx^2+Cy^2$,
$L=Dx+Ey+Hz$, $R_2=B^2x+C^2y+A^2z$,
$R_{11}=ACx+ABy+BCz$ と置けば、その係数は
\[
 \Psi_\lambda(F)=
 -\frac S{48}g_\lambda-\frac{\lambda^4}{48}Q
 +\frac{2-\lambda^3}{36}Z_2-\frac S{144}L
 -\frac{\lambda^2}{24}R_2-\frac{\lambda^3+4}{72}R_{11}.
\]
$T_{111}=T_{222}=T_{333}=2$, $T_{123}=\lambda$
および添字対称性を、$I,J$ の完全縮約へ代入して整理すると上の恒等式を得る。
これは全係数を独立記号とした有限多項式恒等式であり、
保存した検証器は左辺と右辺の全係数差が零であることを確認する。
ここで検算したのは固定入力の中心公式の Weyl symbol であり、
非可換中心性そのものをこの計算で証明したという意味ではない。
\end{proof}



\begin{theorem}[混合残余 potential の普遍四次係数]
\label{thm:nfm-mixed-fourth-potential}
五次 tensor の係数 $\beta=1/60$, $\gamma=1/240$ を用いると、
指定中心元に対して全 $n\ge3$ で
\[
 \boxed{\quad W_4=\Psi_\lambda(F)+\frac1{24}I(F)-\frac1{80}J(F)\quad}
\]
であり、同値に
\[
 W_4=-\frac S{80}F+\frac1{60}I+\frac1{240}J+\mathcal C_4(U).
\]
さらに Casimir 分解のスカラーは
\[
 a_h=h\left(1+\frac{S}{240}h^4\right)+O(h^7)
\]
である。従って先の表示における純中心項 $G_4$ は
$\mathcal C_4(U)$ に等しい。
\end{theorem}

\begin{proof}
局所共変式の補題によって $W_4=\Psi+dI+eJ$ と書く。
まず $U=0$ へ制限する。
$I(g)=J(g)=Sg$, $\Psi(g)=-Sg/48$ および
$\pi_5(g)=S\Pi_g/80$ を用いると
\[
 a_4=S(1/30-d-e)
\]
となる。$a_4$ は座標変数に依存しないので、この式は混合係数にもそのまま有効である。
前補題の恒等式を微分し、
$\pi_5=a_4\Pi_F+\Pi_{W_4}$ と比較する。
$S\,d_KF$, $d_KI$, $d_KJ$ は既に確認した独立な三つの共変式なので
\[
 d-\frac1{40}=\beta=\frac1{60},\qquad
 e+\frac1{60}=\gamma=\frac1{240}
\]
を得る。従って $d=1/24$, $e=-1/80$, $a_4=S/240$ である。
最後の純中心項の同定も、前補題の恒等式から従う。
\end{proof}

\begin{corollary}[四次の幾何的比較類]
\label{cor:nfm-fourth-jacobi-class}
$\tau_F$ を従来の八次 active Jacobi 代数の正規化 trace とする。
全変数の Jacobi 環で
\[
 \boxed{\quad
 [W_4-\tau_F(\Psi_\lambda(F))]
 =\frac1{60}[I(F)]+\frac1{40}[\tau_F(I(F))]
             +\frac S{80}[\tau_F(F)] .\quad}
\]
特に四次の幾何的比較には $\gamma$ と $a_4$ を別途評価する必要はなく、
純中心項の別のグラフ積分も不要である。
\end{corollary}

\begin{proof}
$J$ の縮約は $q_l=F_l$ を一つ含むので、
$J=\sum_l B_l F_l$ と書ける。各 $B_l$ は全座標の一次式である。
従って $[J]=0$ かつ $\tau_F(J)=0$。
また Euler の公式から全 Jacobi 環では $[F]=0$ である。
中心 symbol の恒等式に $\tau_F$ を適用して
\[
 \mathcal C_4-\tau_F(\Psi)
       =\frac S{80}\tau_F(F)+\frac1{40}\tau_F(I)
\]
を得る。これを四次 potential の公式へ代入すればよい。
\end{proof}

\begin{proposition}[Fermat active 切断上の混合公式と一般次元の非消滅]
\label{prop:nfm-fourth-uniform-witness}
$\lambda=0$ の上記全混合族では
\[
 [W_4-\tau_F(\Psi_0(F))]
 =\frac19[ACx+ABy+BCz+AH+BD+CE].
\]
従って各 $n\ge3$ の滑らかな全混合係数空間で、
この四次比較類が非零となる Zariski 開は非空である。
\end{proposition}

\begin{proof}
$R_{11}=ACx+ABy+BCz$, $K=AH+BD+CE$ と書く。
縮約式を $\lambda=0$ で展開すると
\[
 I=-\frac83(Az^2+Bx^2+Cy^2)+\frac43R_{11},\qquad
 J=-\frac23(AF_z+BF_x+CF_y).
\]
$F_x=x^2+Ay+Cz+D$ 等を使えば
\[
 I+\frac83(AF_z+BF_x+CF_y)
        =\frac{20}{3}R_{11}+\frac83K.
\]
八つの平方自由単項式を基底とする active Jacobi 代数では
$\tau_F(x)=\tau_F(y)=\tau_F(z)=0$ なので $\tau_F(I)=8K/3$ である。
前系を適用すると所望の $[R_{11}+K]/9$ を得る。

非消滅には、全次元に共通して
\[
 F_*=\frac13\left(x^3+y^3+z^3+u^3+v^3+\sum_j w_j^3\right)
                    +u xy+v yz
\]
を使える。$n=3$ では $w_j$ はなく、それ以外では $n-3$ 個を加える。
この三次は滑らかである。実際、全偏微分が零なら $w_j=0$。
$y=0$ なら残りも全て零なので、非零解があれば $y=1$ と規格化できる。
そのとき $x^2+u=u^2+x=0$ から $xu\in\{0,1\}$、
$z^2+v=v^2+z=0$ から $zv\in\{0,1\}$ となるが、
$F_y=1+xu+zv$ は零にならない。

三次多項式の空間上の線型汎関数 $\ell$ を
\[
 \ell(vx^2)=\ell(uz^2)=1,\qquad \ell(uvy)=-1
\]
とし、その他の三次単項式では零と定める。
これは全ての $z_i\partial_{z_j}F_*$ と全ての純中心三次を消す。
非自明な確認は
$vF_x=vx^2+uvy$ と $uF_z=uz^2+uvy$ の二つだけである。
従って全 Jacobi 環上の汎関数へ降り、
\[
 \ell\left(W_4-\tau_F(\Psi_0(F_*))\right)
      =\ell(uvy)/9=-1/9\ne0
\]
を得る。実際の比較類は滑らかな係数開上の Jacobi 束の正則切断なので、
非消滅は Zariski 開である。
ここでは純中心三次の像でさらに商を取った束の一定階数性を仮定していない。
\end{proof}

\begin{corollary}[Hessian 補正後にも残る標識付き比較の障害]
\label{cor:nfm-corrected-fourth-obstruction}
各 $n\ge3$ について、前命題の四次類が非零となる開の複素点で、
指定した中心量子化・PBW・formality の入力を満たすものを取る。
残余三次 $V(W_h)$ と有限補正三次 $V(H^\sharp)$ の間には、
同じ $h$ と特殊ファイバーの同定を保つ
$\C[h]/(h^5)$ 上の同型は存在しない。
固定入力の完備 MF 移送と相対 window 比較を適用する場合には、
量子残余圏と $\operatorname{Ku}(V(H^\sharp))$ の間にも、
同じ標識を保つ $\C[[h]]$-線型 dg Morita 同値は存在しない。
後者は $n=4$ を含む全 $n\ge3$ で成立する入力付きの主張である。
\end{corollary}

\begin{proof}
二つの三次は $h^4$ を法として同じである。
小拡大 $\C[h]/(h^5)\to\C[h]/(h^4)$ 上の持上げの差は
$h^4[W_4-\tau_F(\Psi)]$ である。
滑らかな三次の $H^0(T)=0$ により、特殊ファイバー上で恒等な同型は、
逐次の小拡大を通じて $h^4$ を法として恒等となる。
従ってそのような同型が存在すれば四次の差が Jacobi 環で零となり、矛盾する。

圏についても、命題\ref{prop:nf-hh}の $HH^1=0$ と
実際の幾何的 Kodaira--Spencer 写像
$(J_F)_3\hookrightarrow HH^2(\operatorname{Ku}(V(F)))$ の単射性を使う。
共通の $h^4$ 切断の持上げの差は、非零の四次類の像である。
特殊ファイバー上恒等な同値の曖昧さは $HH^1=0$ により逐次消えるため、
この差を同値で消せない。最後に、剰余で恒等かつ Artin 切断と両立する
固定入力の MF・window 同値を用いる。
四次元の幾何的変形の全射性は必要ない。
\end{proof}

この結論は、一般ファイバー上で標識を忘れた三次超曲面または
Kuznetsov 成分の非同値を意味しない。
また、Hessian 補正 $\Phi$ の全 O-降下と次数九の定理はそのまま有効であり、
ここで否定したのは $\Phi$ の三次と残余三次の全次数での標識付き一致である。

\subsubsection{四次まで適合する代数的 Poisson 射}

四次の差も座標不変に打ち消せる。
$P\in\det K$ を使い、$S_P(H),I_P(H),J_P(H)$ を
先の完全縮約の四つの $\varepsilon$ を $P$ に置き換えて定義する。
局所体積枠で $P=sP_0$ なら、それぞれ
$s^4S(H|_K),s^4I(H),s^4J(H)$ である。
$\tau=\tau_{H,K}$ と書き、
\begin{align}
 \mathcal E_P(H)
  &=\frac{I_P(H)}{60}+\frac{J_P(H)}{240}
     -\frac{S_P(H)}{80}\bigl(H-\tau(H)\bigr)
     +\frac{\tau(I_P(H))}{40},\label{eq:nfm-fourth-invariant-correction}\\
 H^{[4]}&=H^\sharp+\mathcal E_P(H),&
 P^{[4]}&=\left(1+\frac{S_P(H)}{240}\right)P
 \label{eq:nfm-fourth-corrected-pair}
\end{align}
と置く。

\begin{theorem}[全 O-同値へ降下する四次補正]
\label{thm:nfm-fourth-algebraic-correction}
$H,H|_K,H^{[4]},H^{[4]}|_K$ が滑らかで
$1+S_P(H)/240$ が可逆な非空開上で、
\[
 \Phi^{[4]}:(H,P)\longmapsto(H^{[4]},P^{[4]})
\]
は分解可能な三次 Poisson 対の代数的射を定める。
この射は任意の基底変更・方程式の生成元の変更・全射影同型と両立し、
既存の良い O-marked moduli の対応する開へ降下する。

指定中心元の trace--Hessian 像
$H=H^{\mathrm{TH}}_{F,h}$ と $P=hP_0$ に沿って、
固定 harmonic-angle 模型では全混合係数・全 $n\ge3$ に対して
\[
 H^{[4]}\equiv W_h\pmod{h^6},\qquad
 \iota_{dH^{[4]}}P^{[4]}\equiv\pi_h\pmod{h^7}
\]
が成立する。この有限式の射と、その全次数での残余比較は別の主張であり、
後者の収束・代数化・同値核の降下はここでは結論しない。
\end{theorem}

\begin{proof}
四つの縮約は $K$ の枠によらず、$\tau$ は任意底変更と両立する。
$(H,P)$ を同じ Poisson 対を表す $(cH,c^{-1}P)$ に替えると、
$S_P(H)$ は不変、$I_P(H),J_P(H)$ は $c$ 倍になる。
従って $\mathcal E_P(H)$ と $H^{[4]}$ は $c$ 倍、
$P^{[4]}$ は $c^{-1}$ 倍となる。全射影同型への共変性も同じ縮約から従う。
また、$H$ に純中心三次 $G(U)$ を加えても、
$S,I,J$ と active 臨界代数は変わらず、
$H-\tau(H)$ も変わらない。従って $\mathcal E_P$ はこの変更に不変である。

$P^{[4]}$ は $P$ の座標に依存しない単元倍なので、
出力は同じ rank 三 support を持つ Jacobian Poisson tensor を定める。
滑らかさと単元条件は開条件であり、$P=sP_0$ の $s=0$ では
$H^{[4]}=H$、単元は一となる。滑らかな三次と active 制限を出発点に取れば、
この形式 trait の一般点で必要な開条件を満たす。
既存の good 一ブロック開との交わりも非空である。
表示上の射を全同値へ降下する議論は
定理\ref{thm:nfm-algebraic-hessian-correction}と同じである。

jet の比較では $H^{\mathrm{TH}}_{F,h}-F$ が純中心三次であることを使う。
そのため
\[
 \mathcal E_{hP_0}(H^{\mathrm{TH}}_{F,h})
 =h^4\left(\frac I{60}+\frac J{240}
          -\frac S{80}(F-\tau_F(F))+\frac{\tau_F(I)}{40}\right)
 =h^4\bigl(W_4-\tau_F(\Psi)\bigr).
\]
$H^\sharp=F+h^2\mathscr H_F/24+h^4\tau_F(\Psi)$ と偶奇性から
最初の合同を得る。
さらに $P^{[4]}=h(1+Sh^4/240)P_0$ と
定理\ref{thm:nfm-mixed-fourth-potential}の $a_h$ の式を
$\pi_h=a_h\iota_{dW_h}P_0$ へ代入すれば二つ目の合同を得る。
\end{proof}

\begin{remark}[有限補正の二つの役割]
従来の $\Phi$ は量子側の楕円点曲線との厳密な一致と一般次数九を持つ。
新しい $\Phi^{[4]}$ についてここで証明したのは、全混合係数でのより高い残余 jet の一致である。
実際、純 active 三次では $\mathcal E_{hP_0}(g)=h^4S(g)g/120$ となるため、
新しい有限式に従来の点曲線との厳密な一致をそのまま移してはいない。
新しい射の一般次数も上の議論では決定していない。
\end{remark}


有限式
\[
 W_h^{[2]}=F+\frac{h^2}{24}\mathscr H_F,\qquad
 \pi_h^{[2]}=h\iota_{dW_h^{[2]}}P
\]
は全ての $h$ に対して Jacobi 恒等式と Casimir 条件を満たす。
しかし、固定 formality の実際の tensor との差は一般に五次から現れる。
四次 potential と新たな有限補正は上で決定したが、
三次の単元倍を許した全次数比較はなお別の課題である。

補正の座標共変性は内在的に書ける。
$K$ を rank 三の active 方向、$P\in\det K$ として
\[
 \mathscr H_{P}(F)
 =\bigl(\det\operatorname{Hess}_{K}F\bigr)(P\otimes P)
\]
と定める。局所的に $P=e_1\wedge e_2\wedge e_3$ と書けば
$\mathscr H_P(F)=\det(e_i e_jF)$ である。
従ってこれは補空間やその基底の選択に依存せず、
全線型座標変換に共変である。また
\[
 \mathscr H_{bP}(aF)=a^3b^2\mathscr H_P(F).
\]
特に $(F,P)\mapsto(aF,a^{-1}P)$ という方程式直線の局所自明化変更では、
$F+h^2\mathscr H_P(F)/24$ は $a$ 倍され、対応する tensor は変わらない。
従って有限 jet の補正は、方程式の代表や volume の選択から独立に定義される。

\begin{remark}[形式的な降下と残る代数化]
固定 formality と $h$ を保つ次数付き中心三次対の線型同型は、
正規化の一意性により形式 Poisson 対の同型を誘導する。
中心元を単元 $b$ 倍すると $W_h$ は $b$ 倍され、$\pi_h$ は変わらない。
従って構成は中心直線についても自然であり、
形式的な表示 groupoid 上では座標変更と cocycle に両立する。
この主張には formal Poisson tensor を単に候補として選ぶ曖昧さはない。

一方、ここで固定した formality 自体を変えた場合の比較、
形式級数の収束・代数化、$h$ の標識を忘れる操作、
および全 intrinsic O-marked 圏族への降下は別問題である。
Kontsevich\cite[Section~2.3, Conjecture~1 and Questions~1]{NFMKontsevichAlg}の
形式対応の収束・周期に関する問題とも区別する。
本節から有限 $h$ の一般点における代数的モジュライ射を宣言しない。
既存の全 O-同値比較を用いる場合も、形式変形とその標識を保つ同値が
表示 groupoid の同値へ移るという適用条件を確認してから降下する。
\end{remark}

\subsection{点曲線が決める split 族の厳密補正}

ここでは量子化の尺度 $s$ と Hesse 係数 $\lambda$ を独立に扱う。
点曲線の Hessian 表示は Mori--Smith\cite[Theorems~1.2, 7.2]{MoriSmith}
に既にある結果であり、新規性を主張する対象ではない。
本小節の用途は、その Hessian 表示を残余復元と結び付け、
補正後の Poisson 対の尺度と定義体を明記することである。

$\lambda\in\C$ を $1+\lambda^3\ne0$ なる定数とし、
$R=\C[s]_{(s)}$、$\widehat R=\C[[s]]$、$K=\C(s)$ と置く。
$S=x^3+y^3+z^3$、$P=xyz$、$g_\lambda=S/3+\lambda P$ と書く。
$D_R$ は
\[
 \begin{aligned}
 [y,z]&=s\left(x^2+\frac\lambda2(yz+zy)\right),\\
 [z,x]&=s\left(y^2+\frac\lambda2(zx+xz)\right),\\
 [x,y]&=s\left(z^2+\frac\lambda2(xy+yx)\right).
 \end{aligned}
\]
で定まる代数とする。$r=n-1\ge2$ とし、中心変数を加えて
$B_R=D_R[u_1,\ldots,u_r]$ と置く。
$G\in\C[u_1,\ldots,u_r]_3$ の偏微分の共通零点は原点だけと仮定する。
特に $G=\sum u_i^3/3$ を全ての $n$ に一様に使える。
\[
 \begin{aligned}
 a_{\lambda,s}&=1+\frac{s^2\lambda^2}{2}
                       +\frac{s^4\lambda}{6}-\frac{s^4\lambda^4}{48},\\
 b_{\lambda,s}&=\lambda-\frac{s^2\lambda^3}{6}-\frac{2s^2}{3}
                       +\frac{s^4\lambda^2}{6}-\frac{s^4\lambda^5}{48},\\
 \Omega_D&=a_{\lambda,s}S/3+b_{\lambda,s}\operatorname{sym}(P),
 \qquad w=\Omega_D+G.
 \end{aligned}
\]
これらは既存の Hesse 中心公式の二パラメータ表示である。
$r$ 個の中心変数を付加しても、各次数の有限自由性と中心性は保存される。

\begin{proposition}[点曲線の厳密な Hessian 補正]
\label{prop:nfm-hesse-point-correction}
$\mathscr H_\lambda=\det\operatorname{Hess}_{x,y,z}(g_\lambda)$ と置き、
\[
 \begin{aligned}
 A_{\lambda,s}&=1-\frac{s^2\lambda^2}{4},\\
 B_{\lambda,s}&=\lambda+\frac{s^2(4+\lambda^3)}{12},\\
 J_{\lambda,s}&=g_\lambda+\frac{s^2}{24}\mathscr H_\lambda
       =A_{\lambda,s}S/3+B_{\lambda,s}P,\\
 \nu(\lambda,s)&=\frac{B_{\lambda,s}}{A_{\lambda,s}}
       =\frac{\lambda+s^2(4+\lambda^3)/12}{1-s^2\lambda^2/4}
 \end{aligned}
\]
とする。一般点の active な点曲線は $V(J_{\lambda,s})$ である。
この曲線の $R$ 上の滑らかな閉包は $V(S+3\nu(\lambda,s)P)$ であり、
$s=0$ のファイバーは $V(g_\lambda)$ である。
ただし $s=0$ の可換代数自体の点スキームは $\mathbb P^2$ であり、
この閉包と混同しない。

Trace--Hessian 出力は $H^{\mathrm{TH}}=g_\lambda+G$ であり、
補正後の対
\[
 H^{\mathrm{corr}}=J_{\lambda,s}+G,
 \qquad
 \Pi^{\mathrm{corr}}
 =s\,\iota_{dH^{\mathrm{corr}}}
       (\partial_x\wedge\partial_y\wedge\partial_z)
\]
は $K$ 上の滑らかな三次超曲面と非零 Poisson 構造を定める。
その active な平面切断は、もとの代数の点曲線と埋込みも含めて一致する。
\end{proposition}

\begin{proof}
Sklyanin 係数は
$a=1-s\lambda/2$、$b=-1-s\lambda/2$、$c=-s$ である。
De Laet--Le Bruyn\cite[Section~2, preprint p.~3]{NFMDLL}の中心元
（本稿の基礎節の式\texttt{gre:eq:known-center}も参照）を
$C=xyz+yzx+zxy$、$D=yxz+xzy+zyx$ により
\[
 a(b^3-c^3)C+b(c^3-a^3)D+c(a^3-b^3)S
\]
と書き、関係式 $C-D=sS+(s\lambda/2)(C+D)$ を代入すると
$-9s\Omega_D$ を得る。
一般点は滑らかな楕円型の Sklyanin 代数であり、
剰余は三変数多項式環である。
両者の Hilbert 関数が一致するので、DVR 上の構造定理により全次数が自由である。
一般点の中心性は次数四の torsionfree 性により $R$ 上へ戻る。
この次数別自由性と多項式環への特殊化により、固定入力の
$\texttt{gcd:prop:automatic-as}$ の仮定も満たされる。
従って一般点の $B$ は両側 noetherian な Koszul AS 正則代数であり、
大域次元と Gorenstein parameter は $N=n+2$ である。
また $w\bmod s=g_\lambda+G\ne0$ なので $w$ は正則で、中心商も次数別に自由である。

点曲線の行列式は
\[
 -\frac{a^3+b^3+c^3}{3abc}
 =\frac{\lambda+s^2(4+\lambda^3)/12}{1-s^2\lambda^2/4}
\]
を与える。直接計算でも
\[
 \mathscr H_\lambda=(8+2\lambda^3)P-2\lambda^2S
\]
であり、表示した $J_{\lambda,s}$ を得る。
$A_{\lambda,s}$ と $1+\nu(\lambda,s)^3$ は $R$ の単元なので曲線は滑らかである。
点自己同型を与える切断は $(a:b:c)$ であり、$s=0$ では $(1:-1:0)$ に特殊化する。

active な Jacobi 環で $x^3=y^3=z^3=-\lambda P$ であるから
\[
 b_{\lambda,s}-\lambda a_{\lambda,s}
 =-\frac23s^2(1+\lambda^3),\qquad
 \nu_{s\Pi_{g_\lambda}}(\overline g_\lambda)
 =8s^2(1+\lambda^3)[P].
\]
正次数の active 元の乗法 trace は零であり、純中心成分は $G$ となる。
従って $-12$ 規格化の Trace--Hessian 出力は $g_\lambda+G$ である。
$H^{\mathrm{corr}}\bmod s=g_\lambda+G$ は滑らかなので、
$R$ 上の族と一般ファイバーは滑らかである。
Jacobian 形の三変数 Poisson tensor は全ての $u_i$ を Casimir とし、
$H^{\mathrm{corr}}$ も Casimir であるから、表示した対を得る。
\end{proof}

\begin{theorem}[split 族での残余圏と補正後三次の比較]
\label{thm:nfm-split-corrected-residual}
上の族について、固定入力論文Bの完備 MF 移送、相対 window、
固定対象の局所化および有限降下、coherent 残余圏の飽和性の入力を用いる。
特に最後の入力は $\texttt{crs:prop:residual-smooth}$ と
$\texttt{crs:thm:coherent}$ である。
このとき、完備残余復元の Hesse 係数は全次数で $\nu(\lambda,s)$ に等しい。
さらに、ある有限代数拡大 $K'/K$ 上で
\[
 (\mathcal K_{B/(w)},\rho)
 \simeq
 (\operatorname{Ku}(V(H^{\mathrm{corr}})),\rho_{\mathrm{corr}})
\]
という dg Morita 同値が存在する。
同値は $w$ と表示した $H^{\mathrm{corr}}$ に由来する指定三次周期を保存する。
これは任意の $n\ge3$ と任意の定数 $\lambda\in\C$、$1+\lambda^3\ne0$ に成立する。
\end{theorem}

\begin{proof}
固定入力の一パラメータ式へ単に記号を代入するのではなく、
その構成を本族について確認する。
まず三変数の $D_{\widehat R}$ と $\Omega_D$ に完備移送を適用する。
$g_\lambda$ は滑らかで、上で次数別自由性と中心性を確認したので、その仮定を満たす。
巡回置換と $(x,y,z)\mapsto(x,\omega y,\omega^2z)$ が生成する有限 Heisenberg 群は、
全ての $s$ に対して関係式と $\Omega_D$ を保存する。
ここでも harmonic-angle formality と、次数一生成元を固定する模型を使う。
補題\ref{lem:nfm-normalized-star-model}の一意性と線型共変性により、
この模型は自動的に Heisenberg 群と同変である。
従って固定入力の $\texttt{gre:prop:hesse-normalform}$ と同じ不変三次空間に入り、
active potential は
\[
 W_{\mathrm{act}}=a(s)S/3+b(s)P,
 \qquad a(0)=1,\quad b(0)=\lambda
\]
となる。中心多項式変数への自明な延長により、全 potential は
$W=W_{\mathrm{act}}+G$ である。
active 座標だけを $a(s)^{-1/3}$ 倍して得る Hesse 係数を
$\mu(s)=b(s)/a(s)$ と書く。
普遍二次 jet 計算により
\[
 a(s)=1-\frac{s^2\lambda^2}{4}+O(s^3),\qquad
 b(s)=\lambda+\frac{s^2(4+\lambda^3)}{12}+O(s^3),
\]
従って
\[
 \mu(s)=\lambda+\frac{1+\lambda^3}{3}s^2+O(s^3).
\]

$D_R/(\Omega_D)$ の相対 THCR 表示も同じ証明で得られる。
すなわち、滑らかな点曲線と平行移動切断を $R$ 上で取り、
その $\mathcal O(1)$ の twisted tensor 冪の切断環への写像を次数一から構成する。
各正次数の切断加群は有限自由であり、一般点で二次関係と $\Omega_D$ が消えるので、
torsionfree 性によりこの写像は $R$ 上で $D_R/(\Omega_D)$ を経由する。
また $\Omega_D\bmod s=g_\lambda$ は非零因子であるから、
次数別 Nakayama と分裂単射によって商の各次数も有限自由となる。点曲線と平行移動切断の剰余は
$(V(g_\lambda),(1:-1:0))$ であり、剰余の商写像は通常の斉次座標環の同型である。
各次数で同じ階数の有限自由加群の間の写像なので Nakayama により同型となる。
相対 ampleness と三変数の指数零の MF 比較は固定入力の
$\texttt{gre:prop:elliptic-quotient}$、$\texttt{gre:lem:elliptic-mf}$ の証明通りである。
ここでは Fermat という条件を使わず、$V(g_\lambda)$ の滑らかさだけを使う。
従って $\C((s))$ 上で $E_\mu$ と $E_\nu$ の導来同値を得る。
両者は $(1:-1:0)$ を原点に持つ楕円曲線なので、
Antieau--Krashen--Ward\cite[Theorem~2.2]{NFMAKW}により同型となり、
$j(\mu)=j(\nu)$ である。

この等式から分枝を一意に決める。
$j(z)$ は $z=\lambda$ で正則な非定数有理関数なので、ある整数 $e\ge1$ と
形式的可逆変数 $\psi(z)=z+O(z^2)$、非零定数 $c$ により
\[
 j(\lambda+z)-j(\lambda)=c\,\psi(z)^e
\]
と書ける。$\mu-\lambda$ と $\nu-\lambda$ はともに
$(1+\lambda^3)s^2/3+O(s^3)$ で始まり、その係数は非零である。
$j$ の等式から二つの $\psi$ の値の比は $e$ 乗根となり、
その剰余は一だから比は一である。
$\psi$ の形式逆により $\mu=\nu$ を得る。
この証明は $j'(\lambda)=0$ の場合も含む。

従って全 potential は $a(s)g_\nu+G$ である。
active 座標だけを $(A_{\lambda,s}/a(s))^{1/3}$ 倍すれば、
厳密に $H^{\mathrm{corr}}=A_{\lambda,s}g_\nu+G$ へ移る。
この根は剰余が一であり、底の $s$、中心変数、周期変数 $t$ を固定する。
従って内部 shift と指定三次周期を保つ完備 MF 比較が得られる。

$H^{\mathrm{corr}}$ 自体は $K$ 上の多項式である。
固定入力の一般次元の固定対象への制限、局所化、生成性および有限降下を
この既に指定した代数的標的へ適用すると、ある有限 $K'/K$ 上で
まず線束生成残余部分について同値を得る。
ここで用いる局所化と有限降下は $\texttt{rog:lem:localization}$ と
$\texttt{rog:lem:descent}$ である。
全 coherent 残余圏への移行には、さらに上記の残余 smoothness と
coherent 飽和性の入力を用いる。これにより表示した $\mathcal K$ 全体の同値となる。
完備化上の形式座標変換そのものが $K'$ 上に代数化されることは使わない。
また、特殊点を含む $R$ 上の代数的相対圏同値はここでは結論しない。
\end{proof}

\begin{remark}[尺度と既存の一パラメータ例]
$A=A_{\lambda,s}$ と書く。
$A$ の三乗根を加え、active 座標に $x,y,z\mapsto A^{-1/3}(x,y,z)$ を施すと、
補正後の対は
\[
 \left(g_\nu+G,\ sA\,\Pi_{g_\nu}\right)
\]
へ移る。三次式を規格化するだけで tensor を $s\Pi_{g_\nu}$ と書くのは正しくない。
三重ベクトルは逆 determinant によって $A$ 倍されるためである。
補正前の表示 $(H^{\mathrm{corr}},\Pi^{\mathrm{corr}})$ は既に $K$ 上にあり、
三乗根は規格化した表示との同型のためにだけ必要となる。

前節の $\lambda=s=h$ の族では
$A=1-h^4/4$、$\nu=\lambda_{\mathrm{res}}(h)$ となる。
従って補正後の代数的三次は、前節で独立に決定した残余三次と
有限三乗根拡大後に厳密な座標変換で一致する。
この場合の残余比較は、前節の固定入力の exact formula からも直接従う。

上の $\Pi^{\mathrm{corr}}$ は、点曲線と Hessian 共変式によって選んだ
代数的 Poisson tensor である。
固定した Kontsevich 比較の formal Poisson tensor の全次数と同一だとは証明していない。
後者の尺度や gauge の全高次比較、一般混合係数での残余比較、
全 O-同値に対する残余同値の降下は別の課題として残る。
\end{remark}

\par\begingroup\small
\begin{thebibliography}{99}

\bibitem{NCCubicA}
\emph{滑らかな三次超曲面のPoisson構造：正規形・退化軌跡・代数的量子化},
\texttt{nc-cubic} 論文A、未公表研究稿 A-v61、2026年11月14日。
\href{https://github.com/uedakazushi/nc-cubic/blob/eabe601c1341777742964584dde5314c0bcf64a9/papers/paper_a.tex}{固定 commit の主原稿}。
支持節は \texttt{papers/support/papers/sections/poisson\_geometry.tex}
の命題 \texttt{prop:pg-represent}, \texttt{prop:pg-jacobi-extension}、
定理 \texttt{thm:pg-normal}、系 \texttt{cor:pg-generic}、および
\texttt{general\_dimension\_universal\_quantization.tex}
の定理 \texttt{thm:pg-all-dimensions}、命題 \texttt{prop:pg-ad-relative-chart}。

\bibitem{NFNCCubicB}
\emph{中心三次非可換変形の一般点における可換復元},
\texttt{nc-cubic} 論文B、未公表研究稿 B-v59、2026年11月14日。
\href{https://github.com/uedakazushi/nc-cubic/blob/eabe601c1341777742964584dde5314c0bcf64a9/papers/paper_b.tex}{固定 commit の主原稿}。
\texttt{papers/generic\_recovery/} の
\texttt{main\_result.tex}, \texttt{generic\_descent.tex},
\texttt{coherent\_saturation.tex}, \texttt{field\_model.tex} を参照。

\bibitem{Threefold}
\emph{非可換三次三様体のAbel圏モジュライ：二次包絡とKuznetsov残余射},
未公表研究稿 v22、2026年11月20日。
\href{https://github.com/uedakazushi/nc-cubic-threefold-moduli/blob/6c35d2c84671c2538e61836a6bf3a511392eacce/papers/nc-cubic-threefold-moduli.tex}{固定 commit の主原稿}。
点スキームは \texttt{papers/sections/point\_scheme\_poisson.tex}、
全 Hom の比較は \texttt{jacobian\_abel\_comparison.tex} を参照。

\bibitem{NFMEG}
P.~Etingof and V.~Ginzburg,
\emph{Noncommutative del Pezzo surfaces and Calabi--Yau algebras},
J. Eur. Math. Soc. \textbf{12} (2010), 1371--1416.
\url{https://doi.org/10.4171/JEMS/235}.
\url{https://ems.press/content/serial-article-files/31746},
Theorem~3.3.2, Section~9.

\bibitem{MoriSmith}
I.~Mori and S.~P.~Smith,
\emph{The classification of 3-Calabi--Yau algebras with 3 generators and 3 quadratic relations},
Math. Z. \textbf{287} (2017), 215--241.
\url{https://doi.org/10.1007/s00209-016-1824-5}.
\url{https://arxiv.org/abs/1502.07403v2}, Theorem~4.2.

\bibitem{LBSV}
L.~Le Bruyn, S.~P.~Smith and M.~Van den Bergh,
\emph{Central extensions of three dimensional Artin--Schelter regular algebras},
Math. Z. \textbf{222} (1996), 171--212.
\url{https://doi.org/10.1007/PL00004532}.
\url{https://sites.math.washington.edu/~smith/Research/Central-extns.pdf},
Theorem~4.2.2, Proposition~4.2.5, Section~4.3.

\bibitem{Mayanskiy}
E.~Mayanskiy,
\emph{Poisson cohomology of two Fano threefolds},
J. Algebra \textbf{424} (2015), 21--45.
\url{https://doi.org/10.1016/j.jalgebra.2014.08.049}.
\url{https://arxiv.org/abs/1303.5916v1}, Theorem~1, Section~2.2.
Theorem~1 は Loray--Pereira--Touzet の分類に帰属されている。
スキーム構造には Section~2.2 の Schouten 計算を用いる。

\bibitem{BCMStandardMonomial}
N.~C.~Bonala, O.~Clarke and F.~Mohammadi,
\emph{Standard monomial theory and toric degenerations of Richardson varieties in the Grassmannian},
\url{https://arxiv.org/abs/2103.16208v1}, 2021,
Section~2.4, Theorem~2.3.

\bibitem{NFMStacksProjective}
The Stacks Project Authors,
\emph{The Stacks Project}, Section~27.21, Projective bundles, Tag~01OA.
\url{https://stacks.math.columbia.edu/tag/01OA}.

\bibitem{NFMStacksAffineDescent}
The Stacks Project Authors,
\emph{The Stacks Project}, Lemma~35.37.1, Tag~0245.
\url{https://stacks.math.columbia.edu/tag/0245}.

\bibitem{NFLZSerre}
X.~Lin and S.~Zhang,
\emph{Serre algebra, matrix factorization and categorical Torelli theorem for hypersurfaces},
Math. Ann. \textbf{391} (2025), 163--177.
\url{https://doi.org/10.1007/s00208-024-02915-8}.
定理番号は \url{https://arxiv.org/abs/2310.09927v2}
の Theorem~3.5, Proposition~3.7, Theorems~4.1 and~5.1 による。

\bibitem{NFBFKKernels}
M.~Ballard, D.~Favero and L.~Katzarkov,
\emph{A category of kernels for equivariant factorizations and its implications for Hodge theory},
Publ. Math. Inst. Hautes \'Etudes Sci. \textbf{120} (2014), 1--111.
\url{https://doi.org/10.1007/s10240-013-0059-9}.
\url{https://numdam.org/item/10.1007/s10240-013-0059-9.pdf},
Theorem~5.39 の一般式とその証明、pp.~82--84。

\bibitem{NFPirozhkov}
D.~Pirozhkov,
\emph{A categorical Torelli theorem for hypersurfaces},
Bull. London Math. Soc. \textbf{56} (2024), no.~10, 3075--3089.
\url{https://doi.org/10.1112/blms.13117}.
本文の Lemmas~5.1 and~5.2、Theorem~5.3 は
\url{https://arxiv.org/abs/2208.13604v1} による。

\bibitem{NFHuybrechtsRennemo}
D.~Huybrechts and J.~V.~Rennemo,
\emph{Hochschild cohomology versus the Jacobian ring and the Torelli theorem for cubic fourfolds},
Algebraic Geometry \textbf{6} (2019), no.~1, 76--99.
\url{https://doi.org/10.14231/AG-2019-005}.
\url{https://content.algebraicgeometry.nl/2019-1/2019-1-005.pdf},
Proposition~3.5 and Remark~3.8.

\bibitem{NFDAGX}
J.~Lurie,
\emph{Derived Algebraic Geometry X: Formal Moduli Problems}, November 5, 2011.
\url{https://www.math.ias.edu/~lurie/papers/DAG-X.pdf},
Section~2, Theorems~5.1.9 and~5.3.16, Proposition~5.3.21.

\bibitem{NFBKP}
A.~Blanc, L.~Katzarkov and P.~Pandit,
\emph{Generators in formal deformations of categories},
Compos. Math. \textbf{154} (2018), 2055--2089.
\url{https://doi.org/10.1112/S0010437X18007303}.
\url{https://arxiv.org/abs/1705.00655v1}, Sections~3.3 and~4.3.

\bibitem{NFKuznetsovFCY}
A.~Kuznetsov,
\emph{Calabi--Yau and fractional Calabi--Yau categories},
J. reine angew. Math. \textbf{753} (2019), 239--267.
\url{https://doi.org/10.1515/crelle-2017-0004}.
\url{https://arxiv.org/abs/1509.07657v2},
Section~3.3, Corollaries~3.18 and~4.1, Lemma~3.19.

\bibitem{NFBFKOrlov}
M.~Ballard, D.~Favero and L.~Katzarkov,
\emph{Orlov spectra: bounds and gaps},
Invent. Math. \textbf{189} (2012), 359--430.
\url{https://doi.org/10.1007/s00222-011-0367-y}.
定理番号は \url{https://arxiv.org/abs/1012.0864v3} の Theorem~5.15 による。

\bibitem{NFKuznetsovCubicFourfolds}
A.~Kuznetsov,
\emph{Derived categories of cubic fourfolds},
in \emph{Cohomological and Geometric Approaches to Rationality Problems},
Progress in Mathematics \textbf{282}, Birkh\"auser, 2010, 219--243.
\url{https://arxiv.org/abs/0808.3351}, Theorem~3.1.
著者稿 \url{https://cims.nyu.edu/~tschinke/books/rationality/kuznetsov.pdf}.

\bibitem{NFMGabrielQuotient}
P.~Gabriel, \emph{Des cat\'egories ab\'eliennes},
Bull. Soc. Math. France \textbf{90} (1962), 323--448,
Chapitre~III, \S1, p.~365.
\url{https://doi.org/10.24033/bsmf.1583}.
\url{https://www.numdam.org/article/BSMF_1962__90__323_0.pdf}.

\bibitem{Fraenkel}
A.~McMillan Fraenkel,
\emph{Extensions of Poisson structures on singular hypersurfaces},
\url{https://arxiv.org/abs/1310.6083v2}, 2013,
Proposition~4.1, Theorem~4.3.

\bibitem{Carles}
R.~Carles,
\emph{D\'eformations dans les sch\'emas d\'efinis par les identit\'es de Jacobi},
Ann. Math. Blaise Pascal \textbf{3} (1996), no.~2, 33--62.
\url{https://doi.org/10.5802/ambp.67}.
\url{https://www.numdam.org/item/AMBP_1996__3_2_33_0.pdf},
Section~8, Lemma~8.35.

\bibitem{CarlesMarquez}
R.~Carles and M.~C.~M\'arquez,
\emph{Different methods for the study of obstructions in the schemes of Jacobi},
Ann. Inst. Fourier \textbf{61} (2011), no.~2, 453--490.
\url{https://doi.org/10.5802/aif.2620}.
\url{https://www.numdam.org/item/AIF_2011__61_2_453_0.pdf},
Section~3, Theorem~3.2, Propositions~3.9 and~3.10.

\end{thebibliography}

\endgroup
\end{document}
